\documentclass[a4paper,12pt]{article}
\usepackage{tikz} % Основной пакет для node и path
\usetikzlibrary{arrows, positioning, shapes, calc} % Полезные дополнения



%---------------------------------

% Кодировка и язык

\usepackage[T2A]{fontenc}

\usepackage[utf8]{inputenc}

\usepackage[russian]{babel}
\usepackage[unicode, colorlinks=true, linkcolor=blue, urlcolor=blue]{hyperref}



%---------------------------------

% Математика и графика

\usepackage{amsmath, amssymb, amsthm}

\usepackage{graphicx}

\usepackage[most]{tcolorbox}
\usepackage[dvipsnames]{xcolor}

\newtheorem{proposition}{Утверждение} % или Proposition

% Создаем команду \term: первый аргумент - текст, который станет синим и жирным
\newcommand{\term}[1]{\textbf{\textcolor{RoyalBlue}{#1}}}

\newcommand{\myexample}[2]{
	\medskip
	\noindent\textbf{\textcolor{ForestGreen}{Пример (#1):}} \textit{#2}
	\medskip
}



%---------------------------------

% Геометрия страницы

\usepackage{geometry}

\geometry{left=2.5cm, right=2.5cm, top=2.5cm, bottom=2.5cm}



%---------------------------------

% Списки и заголовки

\usepackage{enumitem}

\usepackage{titlesec}

\usepackage{tocloft} % Для настройки содержания



%---------------------------------

% Колонтитулы

\usepackage{fancyhdr}

\usepackage{lastpage}



%---------------------------------

% Настройка содержания

\renewcommand{\contentsname}{Содержание}

\renewcommand{\cfttoctitlefont}{\hfill\Large\bfseries}

\renewcommand{\cftaftertoctitle}{\hfill}

\setlength{\cftbeforetoctitleskip}{0pt}

\setlength{\cftaftertoctitleskip}{1em}

\setlength{\cftsecindent}{0em}

\setlength{\cftsubsecindent}{1.5em}

\setlength{\cftbeforesecskip}{0.5em}

\setlength{\cftbeforesubsecskip}{0.3em}



%---------------------------------

% Заголовки разделов

\usepackage{titlesec}
% Формат для \section: шрифт, размер, цвет и т.д.
\titleformat{\section}{\normalfont\Large\bfseries\color{RoyalBlue}}{\thesection}{1em}{}
\titleformat{\subsection}{\normalfont\large\bfseries\color{RoyalBlue}}{\thesubsection}{1em}{}


%---------------------------------

% Стиль для доказательств, замечаний и следствий

\newcommand{\myproof}{\textbf{Доказательство.}\ }

\newcommand{\myremark}{\textbf{Замечание.}\ }

\newcommand{\mycorollary}{\textbf{Следствие.}\ }



%---------------------------------

% Теоремные окружения

\theoremstyle{definition}

\newtheorem{definition}{Определение}[section]



\theoremstyle{plain}

\newtheorem{theorem}{Теорема}[section]

\newtheorem{lemma}[theorem]{Лемма}

\newtheorem{corollary}[theorem]{Следствие}



\theoremstyle{remark}

\newtheorem*{remark}{Замечание}

\newtheorem*{example}{Пример}



%---------------------------------

% Списки

\setlist[itemize]{leftmargin=1.2cm, itemsep=0.3em}

\setlist[enumerate]{leftmargin=1.2cm, itemsep=0.3em}



%---------------------------------

% Колонтитулы

\pagestyle{fancy}

\fancyhf{}

\renewcommand{\headrulewidth}{0pt}

\fancyfoot[L]{\footnotesize Конспект по алгебре}

\fancyfoot[R]{\footnotesize \thepage\ / \pageref{LastPage}}



%---------------------------------

% Типографика

\sloppy



%-------------------------------------------------

\begin{document}
	
	
	
	% Титульная страница
	
	\begin{titlepage}
		
		\centering
		
		{\LARGE\bfseries Подготовка к экзамену по алгебре\par}
		
		\vspace{0.5cm}
		
		{\Large ПИ СПБГУ\par}
		
		\vspace{0.3cm}
		
		{\Large Семестр 1\par}
		
		
		
		\vspace{2cm}
		
		
		
		
		
		{\large Январь 2026\par}
		
		
		
		\vfill
		
		
		
		{\footnotesize составлен по учебнику Шмидта и др.\par}
		
	\end{titlepage}
	
	
	
	\newpage
	
	\tableofcontents
	
	
	
	% ==================== РАЗДЕЛ 1: МНОЖЕСТВА ====================
	\newpage
	\section{Понятие множества. Действия над множествами (1 вопрос)}
	
	\subsection{Понятие множества. Элемент множества }
	
	Первое из принципиально неопределяемых нами понятий — это понятие \textit{множества}, или, говоря точнее, понятие множества и его элемента.
	
	\medskip
	
	\smallskip
	
	\noindent \textbf{Определение.} \textit{Множеством} называется коллекция объектов произвольной природы. 
	Для обозначения множеств обычно используются прописные латинские ($A, B, C, \dots$) или греческие ($\Omega, \Sigma, \dots$) буквы
	
	\medskip
	
	\noindent \textbf{Определение.} \textit{Элементами} множества называются объекты, составляющие данное множество
	
	\medskip
	\begin{tcolorbox}[colback=gray!5, colframe=gray!50, title=Пример множества]
		Пусть задано множество:
		\[ \Phi = \{1, \lambda, \text{element}, \{IV, \text{white}\}\} \]
		Данное множество состоит из четырех элементов: $1, \lambda$, слово «element» и подмножество $\{IV, \text{white}\}$. Обратите внимание, что последний элемент сам является множеством
	\end{tcolorbox}
	
	\subsection{Принадлежность, подмножество и надмножество}
	
	\noindent \textbf{Определение 1 (Принадлежность).} 
	Если $x$ является объектом, составляющим множество $A$, то говорят, что $x$ \textit{принадлежит} $A$ (или $A$ \textit{содержит} $x$), и пишут $x \in A$. В противном случае пишут $x \notin A$
	
	\smallskip
	\begin{tcolorbox}[colback=yellow!5, colframe=orange!50, size=small]
		\textbf{Важно:} Символ $\in$ связывает \textbf{элемент} и \textbf{множество}
	\end{tcolorbox}
	
	\medskip
	
	\noindent \textbf{Определение 2 (Включение).} 
	Пусть $A$ и $B$ — два множества. Говорят, что $A$ является \textit{подмножеством} множества $B$ (обозначается $A \subset B$), если каждый элемент множества $A$ одновременно является элементом множества $B$:
	\[
	A \subset B \iff \forall x \colon (x \in A \implies x \in B).
	\]
	
	\smallskip
	В этом случае также говорят:
	\begin{itemize}
		\item Множество $B$ является \textit{надмножеством} множества $A$ (пишут $B \supset A$)
		\item Множество $A$ \textit{содержится} в $B$, или $B$ \textit{содержит} $A$
	\end{itemize}
	
	
	
	\medskip
	
	\myremark
	Будьте осторожны с термином «содержит». Фраза «$B$ содержит $A$» может означать:
	\begin{enumerate}
		\item $A$ как \textbf{подмножество} ($A \subset B$);
		\item $A$ как \textbf{элемент} ($A \in B$)
	\end{enumerate}
	Обычно смысл ясен из контекста, но в строгих доказательствах лучше использовать термины «включение» и «принадлежность»
	
	\medskip
	
	\noindent \textbf{Свойства включения:}
	\begin{itemize}
		\item $\varnothing \subset A$ (пустое множество — подмножество любого множества);
		\item $A \subset A$ (рефлексивность);
		\item Если $A \subset B$ и $B \subset C$, то $A \subset C$ (транзитивность)
	\end{itemize}
	
	\subsubsection*{Обозначения принадлежности}
	
	Для краткости записи используются специальные логические символы:
	\begin{itemize}
		\item «$\in$» — элемент \textit{принадлежит} множеству (содержится в нем);
		\item «$\notin$» — элемент \textit{не принадлежит} множеству
	\end{itemize}
	
	\medskip
	\noindent \textbf{Пример использования:}
	\begin{itemize}
		\item $1 \in \Phi$ --- единица является элементом множества $\Phi$;
		\item $IV \notin \Phi$ --- римская цифра $IV$ не является элементом $\Phi$ (она является лишь частью одного из элементов, но не самостоятельным элементом верхнего уровня)
	\end{itemize}
	
	\bigskip
	\begin{tcolorbox}[colback=blue!5, colframe=blue!40!black]
		\textbf{Важный нюанс:} Запись $A \subset B$ допускает равенство $A = B$. Если нужно подчеркнуть, что $A$ не равно $B$, используют символ \textit{строгого} включения $A \subsetneq B$
	\end{tcolorbox}
	
	\subsection{Основные операции над множествами}
	
	
	
	\begin{center}
		\begin{tikzpicture}[fill opacity=0.4, scale=0.8]
			% --- ОБЪЕДИНЕНИЕ ---
			\begin{scope}[shift={(0,0)}]
				\fill[blue] (0,0) circle (1);
				\fill[blue] (1.2,0) circle (1);
				\draw (0,0) circle (1) node[left, opacity=1] {$A$};
				\draw (1.2,0) circle (1) node[right, opacity=1] {$B$};
				\node at (0.6, -1.8) [opacity=1] {$A \cup B$};
			\end{scope}
			
			% --- ПЕРЕСЕЧЕНИЕ ---
			\begin{scope}[shift={(4,0)}]
				\draw (0,0) circle (1) node[left, opacity=1] {$A$};
				\draw (1.2,0) circle (1) node[right, opacity=1] {$B$};
				\begin{scope}
					\clip (0,0) circle (1);
					\fill[red] (1.2,0) circle (1);
				\end{scope}
				\node at (0.6, -1.8) [opacity=1] {$A \cap B$};
			\end{scope}
			
			% --- РАЗНОСТЬ ---
			\begin{scope}[shift={(8,0)}]
				\begin{scope}
					\clip (-1.5,-1.5) rectangle (0.6,1.5); % Обрезаем до середины
					\fill[green!70!black] (0,0) circle (1);
					\fill[white, fill opacity=1] (1.2,0) circle (1);
				\end{scope}
				\draw (0,0) circle (1) node[left, opacity=1] {$A$};
				\draw (1.2,0) circle (1) node[right, opacity=1] {$B$};
				\node at (0.6, -1.8) [opacity=1] {$A \setminus B$};
			\end{scope}
			
			% --- СИММЕТРИЧЕСКАЯ РАЗНОСТЬ ---
			\begin{scope}[shift={(12,0)}]
				\fill[orange] (0,0) circle (1);
				\fill[orange] (1.2,0) circle (1);
				\begin{scope}
					\clip (0,0) circle (1);
					\fill[white, fill opacity=1] (1.2,0) circle (1); % "Вырезаем" пересечение белым
				\end{scope}
				\draw (0,0) circle (1) node[left, opacity=1] {$A$};
				\draw (1.2,0) circle (1) node[right, opacity=1] {$B$};
				\node at (0.6, -1.8) [opacity=1] {$A \triangle B$};
			\end{scope}
		\end{tikzpicture}
	\end{center}
	
	
	Пусть $(A, B)$ — пара множеств в некотором универсальном множестве $U$
	
	\begin{itemize}
		\item \textbf{Объединение ($A \cup B$):} Элементы, принадлежащие хотя бы одному из множеств
		\[ A \cup B = \{x \in U \colon x \in A \lor x \in B\} \]
		
		\item \textbf{Пересечение ($A \cap B$):} Элементы, принадлежащие обоим множествам одновременно
		\[ A \cap B = \{x \in U \colon x \in A \land x \in B\} \]
		
		\item \textbf{Разность ($A \setminus B$):} Элементы, входящие в $A$, но не входящие в $B$
		\[ A \setminus B = \{x \in U \colon x \in A \land x \notin B\} \]
		
		\item \textbf{Симметрическая разность ($A \triangle B$):} Элементы, входящие ровно в одно из множеств
		\[ A \triangle B = (A \setminus B) \cup (B \setminus A) \]
	\end{itemize}
	
	\subsection{Законы де Моргана}
	
	Законы де Моргана связывают операции объединения, пересечения и дополнения ($\overline{A} = U \setminus A$)
	
	\bigskip
	\begin{tcolorbox}[colback=red!5, colframe=red!70!black, title=Теорема: Законы де Моргана]
		Для любых подмножеств $A, B \subseteq U$:
		\begin{enumerate}
			\item \textbf{Отрицание объединения:} $\overline{A \cup B} = \overline{A} \cap \overline{B}$
			\item \textbf{Отрицание пересечения:} $\overline{A \cap B} = \overline{A} \cup \overline{B}$
		\end{enumerate}
		\medskip
		\textbf{Обобщение для семейств множеств $\{A_i\}_{i \in I}$:}
		\[ \overline{\bigcup_{i \in I} A_i} = \bigcap_{i \in I} \overline{A_i}, \quad \overline{\bigcap_{i \in I} A_i} = \bigcup_{i \in I} \overline{A_i} \]
	\end{tcolorbox}
	
	
	\subsection{Пустое множество}
	
	\noindent \textbf{Определение.} Множество, не содержащее ни одного элемента, называется \textit{пустым множеством} и обозначается символом $\varnothing$ (или $\emptyset$)
	
	\bigskip
	Пустое множество единственно и обладает следующими фундаментальными свойствами:
	
	\begin{enumerate}
		\item \textbf{Универсальное подмножество:} Пустое множество является подмножеством любого множества $A$:
		\[ \forall A \colon \varnothing \subset A \]
		\item \textbf{Результат операций:} 
		\begin{itemize}
			\item Пересечение множества с пустым всегда пусто: $A \cap \varnothing = \varnothing$
			\item Объединение множества с пустым не меняет само множество: $A \cup \varnothing = A$
			\item Разность множества и самого себя есть пустое множество: $A \setminus A = \varnothing$
		\end{itemize}
	\end{enumerate}
	
	\begin{tcolorbox}[colback=yellow!5, colframe=orange!60!black]
		\textbf{Важное различие:} Не путайте пустое множество $\varnothing$ и множество, содержащее пустое множество как элемент $\{\varnothing\}$
		\begin{itemize}
			\item $\varnothing$ — это «пустая коробка» (0 элементов)
			\item $\{\varnothing\}$ — это «коробка, внутри которой лежит пустая коробка» (1 элемент)
		\end{itemize}
	\end{tcolorbox}
	
	\medskip
	\myremark
	Мнемоническое правило: при вынесении «отрицания» за скобку знак объединения «переворачивается» в знак пересечения, и наоборот
	
	
	
	
	
	
	
	
	
	% ==================== РАЗДЕЛ 2: МОЩНОСТЬ МНОЖЕСТВ ====================
	\newpage
	\section{Мощность множества. Декартово произведение множеств. Счётные множества. Несчётные множества. Счётность множества рациональных чисел (2 вопрос)}
	
	\subsection{Равномощность}
	Пусть $A, B$ — пара множеств. Говорят, что $A$ \textit{равномощно} $B$ (или $A$ и $B$ имеют одинаковую мощность, или мощности $A$ и $B$ равны и т.п.), если существует биекция $A \to B$ или $A = B = \varnothing$. При этом мы будем писать $|A| = |B|$, в противном случае $|A| \neq |B|$
	
	\medskip
	
	Только что введённое понятие задаёт в произвольном множестве множеств отношение — отношение равномощности. Ясно, что это отношение рефлексивно, симметрично и транзитивно
	
	\begin{tcolorbox}[colback=blue!5, colframe=blue!40!black, title=Пример равномощности]
		Множество чётных натуральных чисел $\{2, 4, 6, \dots\}$ равномощно множеству всех натуральных чисел $\mathbb{N}$. Биекция задаётся формулой $f(n) = 2n$. Это удивительно: часть множества оказалась равномощна целому!
	\end{tcolorbox}
	
	\medskip
	
	\subsection{Сравнение мощностей}
	Если существует инъекция $A$ в $B$ или $A = \varnothing$, то говорят, что мощность $A$ не больше мощности $B$, и пишут $|A| \leq |B|$. Если при этом $|A| \neq |B|$, то говорят, что мощность $A$ меньше мощности $B$, и пишут $|A| < |B|$. В частности, если $B \neq \varnothing$, то $|\varnothing| < |B|$
	
	
	\medskip
	\subsection{Конечность, бесконечность, счётность}
	Мы будем называть множество $A$ \textit{конечным} и писать $|A| < \infty$, если это множество равномощно начальному отрезку натурального ряда, отличному от всего этого ряда (в частности, пустое множество конечно). Если этот отрезок есть $\{1, 2, \dots, n\}$, то говорят, что $n$ есть число (или количество) элементов множества $A$, и пишут $|A| = n$ (для пустого множества $|\varnothing| = 0$). \textit{Бесконечным} называют множество, не являющееся конечным
	
	\medskip
	Множество, равномощное множеству натуральных чисел, называется \textit{счётным}
	
	\begin{tcolorbox}[colback=yellow!5, colframe=orange!60!black, title=Важное замечание о бесконечности]
		Бесконечные множества обладают парадоксальным свойством: они могут быть равномощны своей собственной части. Именно это свойство часто принимают за определение бесконечности (в смысле Дедекинда). Конечные множества таким свойством не обладают
	\end{tcolorbox}
	
	\myremark
	Счётные множества — это «самые маленькие» бесконечные множества. Любое бесконечное множество содержит счётное подмножество
	
	\subsection{Декартово произведение множеств}
	
	\noindent \textbf{Определение 4 (Декартово произведение).}
	Пусть $(A, B)$ — пара множеств. \textit{Декартовым произведением} этой пары называется множество $A \times B$ всех упорядоченных пар $(a, b)$, где $a \in A$, $b \in B$:
	\[
	A \times B = \{(a, b) \mid a \in A, b \in B\}
	\]
	
	\begin{tcolorbox}[colback=gray!5, colframe=gray!50, title=Примеры декартова произведения]
		\textbf{Пример 1 (числовой):} 
		Если $A = \{1, 2, 3\}$, а $B = \{x, y\}$, то
		\[
		A \times B = \{(1, x), (1, y), (2, x), (2, y), (3, x), (3, y)\}
		\]
		Это множество можно представить в виде таблицы $3 \times 2$ или как точки на плоскости с целочисленными координатами
		
		
		\medskip
		\textbf{Пример 2 (геометрический):}
		$\mathbb{R} \times \mathbb{R} = \mathbb{R}^2$ — это всем известная декартова плоскость. Каждая точка плоскости задаётся парой координат $(x, y)$
		
		\medskip
		\textbf{Пример 3 (строки):}
		Если $A = \{\text{Иван}, \text{Пётр}\}$, $B = \{\text{любит}, \text{ненавидит}\}$, то $A \times B$ содержит пары вроде (Иван, любит), (Пётр, ненавидит) и т.д.
	\end{tcolorbox}
	
	\myremark
	Важно понимать различие: $A \times B \neq B \times A$ (порядок важен), но мощности их равны. Декартово произведение можно обобщить на любое конечное число множеств: $A_1 \times A_2 \times \dots \times A_n$
	
	\subsection{Свойства счётных множеств}
	
	\begin{proposition}[1]
		Объединение конечного и счётного множества счётно
	\end{proposition}
	\begin{proof}
		Пусть $A = \{a_1, a_2, \dots, a_m\}$, $B = \{b_1, b_2, \dots, b_n, \dots\}$. Предположим, что $A \cap B = \varnothing$. В противном случае в объединении $A \cup B$ вместо $A$ нужно будет рассмотреть множество $A \setminus B$, которое будет конечным. Итак, $A \cup B = \{a_1, a_2, \dots, a_m, b_1, b_2, \dots, b_n, \dots\}$
		
		\medskip
		Чтобы установить счётность, $A \cup B$ достаточно представить в виде его последовательности $c_1, c_2, \dots, c_n, \dots$. Действительно, в этом случае каждому элементу множества $A \cup B$ будет сопоставлено одно и только одно натуральное число — его номер в этой последовательности. Тем самым будет установлено, что $A \cup B$ равномощно $\mathbb{N}$. То есть,
		\[
		c_1 = a_1, c_2 = a_2, \dots, c_m = a_m, c_{m+1} = b_1, c_{m+2} = b_2, \dots
		\]
	\end{proof}
	
	\begin{tcolorbox}[colback=green!5, colframe=green!50!black, title=Пояснение к доказательству]
		Мы просто выписываем сначала все элементы конечного множества (их конечное число), а потом начинаем выписывать элементы счётного. Поскольку после конечного числа шагов мы всё ещё имеем бесконечную последовательность, все элементы счётного множества получат свои номера (со сдвигом на $m$)
	\end{tcolorbox}
	
	\begin{proposition}[2]
		Объединение конечного числа счётных множеств счётно
	\end{proposition}
	\begin{proof}
		Рассмотрим счётные множества $A$ и $B$. Предположим, что $A \cap B = \varnothing$.
		
		Если это не так, то воспользуемся тем, что $A \cup B = (A \setminus B) \cup B$. Множество $A \setminus B$ будет либо конечным, либо счётным. В первом случае $A \cup B$ будет счётным по предыдущему предложению, во втором — докажем ниже.
		
		\medskip
		Пусть $A = \{a_1, a_2, \dots, a_n, \dots\}$, $B = \{b_1, b_2, \dots, b_n, \dots\}$. Тогда
		\[
		A \cup B = \{a_1, b_1, a_2, b_2, \dots, a_n, b_n, \dots\}.
		\]
		Для доказательства счётности множества $A \cup B$ достаточно представить его в виде последовательности $c_1, c_2, \dots, c_n, \dots$. Сделаем это так:
		\[
		c_1 = a_1, c_2 = b_1, c_3 = a_2, c_4 = b_2, c_5 = a_3, c_6 = b_3, \dots, c_{2n-1} = a_n, c_{2n} = b_n, \dots
		\]
		Взаимно однозначное соответствие между элементами объединения и множеством натуральных чисел установлено. Доказательство счётности $A_1 \cup A_2 \cup A_3 \cup \dots \cup A_{n-1} \cup A_n$ проводится индукцией по числу множеств
	\end{proof}
	
	\begin{tcolorbox}[colback=yellow!5, colframe=orange!60!black, title=Пример к предложению 2]
		Пусть $A = \{1, 3, 5, 7, \dots\}$ (нечётные числа), $B = \{2, 4, 6, 8, \dots\}$ (чётные числа). Тогда $A \cup B = \mathbb{N}$, и наше построение даёт последовательность:
		\[
		1, 2, 3, 4, 5, 6, \dots
		\]
		что и является натуральным рядом
	\end{tcolorbox}
	
	
	
	\bigskip
	\begin{tcolorbox}[colback=red!5, colframe=red!70!black, title=Теорема 1]
		Объединение счётного числа счётных множеств счётно
	\end{tcolorbox}
	
	\medskip
	\begin{proof}
		Рассмотрим счётное число счётных множеств $A_1, A_2, A_3, \dots, A_n, \dots$. Будем предполагать, что они попарно не пересекаются; в противном случае вместо них можно рассмотреть не более чем счётные (т.е. конечные или счётные) множества
		$A_1, A_2 \setminus A_1, A_3 \setminus (A_1 \cup A_2), \dots$, объединение которых совпадает с объединением исходных множеств $A_1 \cup A_2 \cup \dots \cup A_n \cup \dots$. Итак, пусть
		\begin{align*}
			A_1 &= \{a_{11}, a_{12}, a_{13}, a_{14}, \dots, a_{1n}, \dots\}, \\
			A_2 &= \{a_{21}, a_{22}, a_{23}, a_{24}, \dots, a_{2n}, \dots\}, \\
			A_3 &= \{a_{31}, a_{32}, a_{33}, a_{34}, \dots, a_{3n}, \dots\}, \\
			A_4 &= \{a_{41}, a_{42}, a_{43}, a_{44}, \dots, a_{4n}, \dots\}.
		\end{align*}
		
		Установим порядок пересчёта элементов таблицы. Первый номер дадим такому элементу, у которого сумма индексов равна 2 (это $a_{11}$).
		Следующие два номера сопоставим тем элементам, у которых сумма индексов равна 3, причём договоримся, что номера таким элементам будем присваивать в порядке убывания первых индексов; таким образом, $2 \leftrightarrow a_{21}$, $3 \leftrightarrow a_{12}$.
		Номера 4-6 получат элементы с суммой индексов равной 4, которые пересчитываются далее в порядке убывания первых индексов, и так далее
		
		\medskip
		Иначе говоря, мы пересчитываем элементы таблицы, двигаясь последовательно по её диагоналям, начиная с левого верхнего угла. Очевидно, что при этом каждому элементу таблицы сопоставляется единственное натуральное число, а каждому натуральному числу — единственный элемент таблицы. Правило пересчёта установлено. Это и доказывает счётность объединения $A_1, A_2, A_3, \dots, A_n, \dots$
	\end{proof}
	
	\begin{tcolorbox}[colback=blue!5, colframe=blue!40!black, title=Наглядная иллюстрация диагонального метода]
		Расположим элементы в бесконечную таблицу:
		\[
		\begin{array}{cccccc}
			a_{11} & \rightarrow & a_{12} &  & a_{13} & \cdots \\
			& \swarrow &  & \swarrow &  & \\
			a_{21} &  & a_{22} &  & a_{23} & \cdots \\
			\downarrow & \nearrow &  & \swarrow &  & \\
			a_{31} &  & a_{32} &  & a_{33} & \cdots \\
			&  &  &  &  & \ddots
		\end{array}
		\]
		Двигаясь по диагоналям: (1,1); (2,1), (1,2); (3,1), (2,2), (1,3); (4,1), (3,2), (2,3), (1,4); \dots мы никогда не пропустим ни одного элемента!
	\end{tcolorbox}
	
	\myremark
	Теорема 1 — одна из важнейших в теории множеств. Она показывает, что счётное объединение счётных множеств остаётся счётным. Без неё мы бы не могли доказать счётность рациональных чисел
	
	\bigskip
	\begin{tcolorbox}[colback=red!5, colframe=red!70!black, title=Теорема 2]
		Множество рациональных чисел $\mathbb{Q}$ счётно
	\end{tcolorbox}
	
	\begin{proof}
		Докажем сначала счётность $\mathbb{Q}^+$ (множество положительных рациональных чисел).
		Представим это множество как объединение счётного числа счётных множеств. Пусть
		\begin{align*}
			A_1 &= \left\{ \frac{1}{1}; \frac{1}{2}; \frac{1}{3}; \frac{1}{4}; \cdots; \frac{1}{n}; \cdots \right\}, \\
			A_2 &= \left\{ \frac{2}{1}; \frac{2}{2}; \frac{2}{3}; \frac{2}{4}; \cdots; \frac{2}{n}; \cdots \right\}, \\
			A_3 &= \left\{ \frac{3}{1}; \frac{3}{2}; \frac{3}{3}; \frac{3}{4}; \cdots; \frac{3}{n}; \cdots \right\}, \\
			\cdots &, \\
			A_m &= \left\{ \frac{m}{1}; \frac{m}{2}; \frac{m}{3}; \frac{m}{4}; \cdots; \frac{m}{n}; \cdots \right\}, \\
			\cdots &
		\end{align*}
		Очевидно, что $\mathbb{Q}^+ = A_1 \cup A_2 \cup A_3 \cup \dots \cup A_m \cup \dots$. Тогда $\mathbb{Q}^+$ счётно как объединение счётного числа счётных множеств (теорема 1)
		
		\medskip
		Ясно, что $\mathbb{Q}^-$ равномощно $\mathbb{Q}^+$, что значит, что $\mathbb{Q}^-$ также счётно. Тогда $\mathbb{Q}$ счётно как объединение двух счётных и одного конечного множества $\{0\}$ (предложения 1 и 2)
	\end{proof}
	
	\begin{tcolorbox}[colback=green!5, colframe=green!50!black, title=Пояснение к доказательству теоремы 2]
		Почему мы не можем просто сказать, что $\mathbb{Q}$ — это пары $(p, q)$, и их счётно? Потому что одно и то же рациональное число представляется разными парами: $\frac{1}{2} = \frac{2}{4} = \frac{3}{6} = \dots$. Наше доказательство обходит эту проблему, рассматривая все дроби (включая сократимые), а потом объединяя их. Избыточность не мешает — если мы можем занумеровать множество с повторениями, то исходное множество (без повторов) будет не более чем счётным (т.е. конечным или счётным)
		
		
	\end{tcolorbox}
	
	\begin{tcolorbox}[colback=green!5, colframe=green!50!black, title=Визуализация метода (Змейка Кантора)]
	Чтобы занумеровать все элементы $\mathbb{Q}^+$, мы записываем их в таблицу (где строка — числитель, столбец — знаменатель) и идем «зигзагом» по диагоналям:
	
	\begin{center}
	\begin{tikzpicture}[scale=1.2]
	% Рисуем сетку дробей
	\foreach \x in {1,2,3,4}
	\foreach \y in {1,2,3,4}
	\node (n\x\y) at (\y,- \x) {$\frac{\x}{\y}$};
	
	% Многоточия
	\foreach \x in {1,2,3,4} \node at (5, -\x) {$\dots$};
	\foreach \y in {1,2,3,4} \node at (\y, -5) {$\vdots$};
	
	% Стрелки нумерации (Змейка)
	\draw[->, red, thick] (n11) -- (n12);
	\draw[->, red, thick] (n12) -- (n21);
	\draw[->, red, thick] (n21) -- (n31);
	\draw[->, red, thick] (n31) -- (n22);
	\draw[->, red, thick] (n22) -- (n13);
	\draw[->, red, thick] (n13) -- (n14);
	\draw[->, red, thick] (n14) -- (n23);
	\draw[->, red, thick] (n23) -- (n32);
	\draw[->, red, thick] (n32) -- (n41);
	
	\node[below right, red] at (n41) {$\dots$};
	\end{tikzpicture}
	\end{center}
	
	\small \textbf{Алгоритм:} Мы выписываем числа в один ряд: $1/1, 1/2, 2/1, 3/1, 2/2, 1/3 \dots$. Если число уже встречалось (например, $2/2 = 1/1$), мы его просто пропускаем. Таким образом, каждое рациональное число получит свой номер
	\end{tcolorbox}
	
	\myremark
	Интуитивно кажется, что рациональных чисел должно быть больше, чем натуральных — ведь между любыми двумя натуральными числами лежит бесконечно много рациональных! Однако канторовская теория мощностей показывает, что их количества (мощности) совпадают. Это первый удивительный факт на пути к осознанию природы бесконечности
	
	
	\subsection{Несчётные множества и мощность континуума}
		
	\bigskip
	\begin{tcolorbox}[colback=red!5, colframe=red!70!black, title=Теорема 3]
		Множество точек промежутка $(0; 1)$ не является счётным
	\end{tcolorbox}
	
	\begin{proof}
		Пусть $A = \{x \in \mathbb{R} \mid 0 < x < 1\}$. Предположим, что $A$ счётно, а значит, все его элементы можно пронумеровать, т.е. представить в виде последовательности: $\{a_1, a_2, a_3, \dots, a_n, \dots\}$.
		Каждый член последовательности — это число из промежутка $(0; 1)$ и поэтому его можно представить в виде десятичной дроби $0, \alpha_1 \alpha_2 \alpha_3 \dots \alpha_n \dots$, где $\alpha_1, \alpha_2, \alpha_3, \dots, \alpha_n, \dots$ — цифры от 0 до 9
		
		\medskip
		Условимся конечные десятичные дроби записывать как периодические с периодом 0. Например, число 0,27 запишем так: 0,27000\dots. Кроме того, договоримся не использовать запись чисел в виде бесконечной дроби с периодом 9, чтобы избежать неоднозначной записи. Например, будем писать 0,23000\dots вместо $0,22999\dots = 0,2(9)$. Тогда для любого числа из промежутка его запись в виде $0, \alpha_1 \alpha_2 \alpha_3 \dots \alpha_n \dots$ будет единственной
		
		\medskip
		Итак, пусть
		\begin{align*}
			a_1 &= 0, \mathbf{\alpha_{1,1}} \alpha_{1,2} \alpha_{1,3} \alpha_{1,4} \dots \alpha_{1,n} \dots \\
			a_2 &= 0, \alpha_{2,1} \mathbf{\alpha_{2,2}} \alpha_{2,3} \alpha_{2,4} \dots \alpha_{2,n} \dots \\
			a_3 &= 0, \alpha_{3,1} \alpha_{3,2} \mathbf{\alpha_{3,3}} \alpha_{3,4} \dots \alpha_{3,n} \dots \\
			\dots & \\
			a_n &= 0, \alpha_{n,1} \alpha_{n,2} \dots \alpha_{n,n-1} \mathbf{\alpha_{n,n}} \dots \\
			\dots & 
		\end{align*}
		У цифры $\alpha_{i, j}$ первый индекс $i$ — это номер элемента $a_i$ в последовательности, а второй индекс $j$ — номер знака после запятой
		
		\medskip
		Укажем такое число $b$, которое с одной стороны принадлежит промежутку $(0; 1)$, а с другой — не совпадает ни с одним членом последовательности $\{a_1, a_2, a_3, \dots, a_n, \dots\}$.
		Пусть $b = 0, \beta_1 \beta_2 \beta_3 \dots \beta_n \dots$, где
		\begin{align*}
			\beta_1 &\neq \alpha_{1,1} \implies b \neq a_1 \\
			\beta_2 &\neq \alpha_{2,2} \implies b \neq a_2 \\
			\beta_3 &\neq \alpha_{3,3} \implies b \neq a_3 \\
			\cdots & \\
			\beta_n &\neq \alpha_{n,n} \implies b \neq a_n \\
			\cdots & 
		\end{align*}
		При выборе $i$-го знака после запятой в числе $b$ будем избегать цифр 0 и 9, чтобы исключить возможность получения бесконечной дроби с периодом 9, а также совпадения числа $b$ с $0 = 0,(0)$ или $1 = 0,(9)$. Например, можно положить $\beta_n = 3$, если $\alpha_{n,n} \neq 3$, и $\beta_n = 5$ в противном случае
	
		\medskip
		Таким образом, число $b$ будет принадлежать промежутку $(0; 1)$, но не будет совпадать ни с одним из членов последовательности $\{a_1, a_2, a_3, \dots, a_n, \dots\}$.
		Следовательно, множество точек промежутка $(0; 1)$ не является счётным
	\end{proof}
	
	
	\begin{tcolorbox}[colback=red!5, colframe=red!70!black, title={Теорема Кантора (Диагональный метод)}]
	Это доказательство - знаменитый \textbf{диагональный метод Кантора}. Идея гениальна в своей простоте: как бы мы ни старались перечислить все числа интервала, всегда можно построить новое число, которое отличается от каждого числа из списка хотя бы в одном десятичном знаке (на диагонали). Это показывает, что множество действительных чисел принципиально "более бесконечно", чем множество натуральных
	
	\medskip
	Допустим, мы смогли пересчитать все числа в интервале $(0, 1)$. Тогда их можно записать в виде бесконечного списка:
	
	\begin{center}
	\begin{tabular}{r c l}
	$n=1:$ & $z_1 = 0,$ & $\mathbf{d_{11}} d_{12} d_{13} d_{14} \dots$ \\
	$n=2:$ & $z_2 = 0,$ & $d_{21} \mathbf{d_{22}} d_{23} d_{24} \dots$ \\
	$n=3:$ & $z_3 = 0,$ & $d_{31} d_{32} \mathbf{d_{33}} d_{34} \dots$ \\
	$n=4:$ & $z_4 = 0,$ & $d_{41} d_{42} d_{43} \mathbf{d_{44}} \dots$ \\
	$\vdots$ & $\vdots$ & $\vdots$
	\end{tabular}
	\end{center}
	
	Построим новое число $b = 0, b_1 b_2 b_3 b_4 \dots$, выбирая цифры следующим образом:
	\[
	b_n = \begin{cases} 
	1, & \text{если } d_{nn} \neq 1 \\
	2, & \text{если } d_{nn} = 1 
	\end{cases}
	\]
	
	Число $b$ отличается от $z_1$ в первом знаке, от $z_2$ — во втором, и вообще **отличается от любого $z_n$** в $n$-м десятичном знаке 
	
	\medskip
	\textbf{Вывод:} Числа $b$ нет в нашем «полном» списке. Противоречие. Множество действительных чисел \textit{несчётно}
	\end{tcolorbox}
	
	\bigskip
	\myremark
	Обратите внимание на техническую деталь: мы исключили цифры 0 и 9 при построении $\beta_n$, чтобы избежать проблем с неоднозначностью десятичной записи (например, 0,1999... = 0,2000...). Это стандартный приём в таких доказательствах
	\bigskip
	
	\noindent \textbf{Определение 5 (Мощность континуума).}
	Будем говорить, что множества, равномощные множеству точек промежутка $(0; 1)$, имеют \textit{мощность континуума}. Обозначается часто $\mathfrak{c}$ или $2^{\aleph_0}$

	\begin{tcolorbox}[colback=yellow!5, colframe=orange!60!black, title=Примеры множеств мощности континуума]
		\begin{itemize}
			\item Вещественная прямая $\mathbb{R}$;
			\item Любой интервал $(a, b)$, отрезок $[a, b]$ (с концами);
			\item Множество иррациональных чисел;
			\item Множество всех подмножеств натуральных чисел $\mathcal{P}(\mathbb{N})$;
			\item Множество всех бесконечных последовательностей из 0 и 1;
			\item Плоскость $\mathbb{R}^2$ (удивительно, но точек на плоскости столько же, сколько на прямой!).
		\end{itemize}
	\end{tcolorbox}
	
	\begin{tcolorbox}[colback=blue!5, colframe=blue!40!black, title=Отображение интервала на прямую]
		Чтобы проверить, что множество всех вещественных чисел имеет мощность континуума, достаточно установить взаимно однозначное соответствие между точками всей числовой оси и точками $(0; 1)$. Это можно сделать многими способами, например:
		\[
		(0; 1) \ni x \mapsto \operatorname{tg}\left[\left(x - \frac{1}{2}\right)\frac{\pi}{2}\right] \in \mathbb{R}
		\]
		или более простой функцией:
		\[
		f(x) = \frac{2x-1}{x(1-x)}
		\]
		Обе функции биективно отображают интервал $(0; 1)$ на всю вещественную прямую
	\end{tcolorbox}
	
	\bigskip
	\myremark
	Обратите внимание на нетривиальный факт: рациональных чисел на прямой "бесконечно много", но их столько же, сколько и натуральных чисел (они счётны). В то же время иррациональных чисел несчётное количество, и их мощность совпадает с мощностью континуума. Таким образом, \textbf{"большинство"} точек на прямой — это иррациональные числа, хотя рациональные расположены всюду плотно!
	
	\bigskip
	\begin{tcolorbox}[colback=purple!5, colframe=purple!50!black, title=Заметочка: Иерархия бесконечностей]
		Итак, мы имеем как минимум две разные бесконечности:
		\begin{itemize}
			\item \textbf{Счётная мощность} $\aleph_0$ (алеф-ноль) — мощность $\mathbb{N}$, $\mathbb{Z}$, $\mathbb{Q}$
			\item \textbf{Мощность континуума} $\mathfrak{c}$ — мощность $\mathbb{R}$, $(0,1)$, $\mathbb{R}^2$
		\end{itemize}
		Вопрос "есть ли мощность, строго промежуточная между $\aleph_0$ и $\mathfrak{c}$?" — это знаменитая \textit{континуум-гипотеза}. Было доказано, что она не зависит от стандартных аксиом теории множеств (ZFC): можно предполагать как её истинность, так и ложность, не впадая в противоречие
	\end{tcolorbox}
	
	
	
	
	
	
	
	
	
	
	
	

	
	
	
	% ==================== РАЗДЕЛ 3: ОТОБРАЖЕНИЯ ====================
	\newpage
	\section{Отображения. Инъективность. Сюръективность. Биективность. Примеры. Композиция отображений (3 вопрос)}
	
	\subsection{Основные определения}
	
	\noindent \textbf{Определение (Отображение, множество всех отображений, начало, конец).}
	Бинарное отношение $f$ из $A$ в $B$ называется \textit{отображением} $A$ в $B$, если
	\[
	\forall x \in A \ \exists! y \in B \colon x f y,
	\]
	то есть, если каждому элементу множества $A$ соответствует единственный элемент множества $B$
	
	\medskip
	Множество всех отображений $A$ в $B$ обозначается $B^A$
	
	\medskip
	Тот факт, что $f$ есть отображение $A$ в $B$ записывается так: $f \colon A \to B$; мы будем называть $A$ \textit{началом} и $B$ — \textit{концом} отображения $f$
	
	\begin{tcolorbox}[colback=gray!5, colframe=gray!50, title=Примеры отображений]
		\begin{itemize}
			\item $f \colon \mathbb{R} \to \mathbb{R}$, $f(x) = x^2$ — отображение, поскольку каждому действительному числу соответствует единственный квадрат
			\item $g \colon \mathbb{N} \to \{0,1\}$, $g(n) = n \mod 2$ — отображение, ставящее в соответствие числу его чётность
			\item Отношение «$x$ является родителем $y$» не является отображением из множества людей в себя, так как у человека может быть несколько детей или не быть совсем
		\end{itemize}
	\end{tcolorbox}
	
	\myremark
	Обозначение $B^A$ для множества всех отображений связано с тем, что если $A$ и $B$ конечны, то $|B^A| = |B|^{|A|}$. Например, отображений из 3-элементного множества в 2-элементное существует $2^3 = 8$
	
	\subsection{Образ и прообраз}
	
	\noindent \textbf{Определение (Образ, прообраз, полный образ, полный прообраз).}
	Элемент множества $B$, соответствующий элементу $x \in A$ при отображении $f \colon A \to B$, называют \textit{образом} элемента $x$ (при отображении $f$) и обозначают $f(x)$
	
	\medskip
	Если $y$ — образ элемента $x$, то каждый такой $x$ называют \textit{прообразом} элемента $y$
	
	\medskip
	Пусть $f \colon A \to B$, $A' \subset A$, $B' \subset B$. Множество $f(A') = \{f(x) \mid x \in A'\}$, состоящее из образов всех элементов множества $A'$, называется \textit{образом} множества $A'$
	
	\medskip
	Образ $f(A)$ всего начала отображения $f$ называется \textit{образом отображения} $f$ и обозначается $\operatorname{Im} f$
	
	\medskip
	Аналогично, множество $f^{-1}(B') = \{x \in A \mid f(x) \in B'\}$ всех прообразов элементов множества $B'$ называется \textit{полным прообразом} множества $B'$
	
	\begin{tcolorbox}[colback=yellow!5, colframe=orange!60!black, title=Важное различие]
		Важно не путать \textit{прообраз} элемента $y$ с его \textit{полным прообразом}:
		\begin{itemize}
			\item Прообраз элемента $y$ есть \textbf{элемент} множества $A'$ (если он существует), и прообразов одного элемента может быть сколько угодно, в том числе ни одного
			\item Полный прообраз элемента $y$ есть \textbf{подмножество} множества $A$, состоящее из всех прообразов элемента $y$
		\end{itemize}
	\end{tcolorbox}
	
	\begin{tcolorbox}[colback=blue!5, colframe=blue!40!black, title=Пример]
		Рассмотрим $f \colon \mathbb{R} \to \mathbb{R}$, $f(x) = x^2$.
		\begin{itemize}
			\item Образ числа $2$: $f(2) = 4$.
			\item Прообраз числа $4$: это $x = 2$ и $x = -2$ (два прообраза)
			\item Полный прообраз числа $4$: $f^{-1}(\{4\}) = \{2, -2\}$.
			\item Образ множества $[-2, -1]$: $f([-2, -1]) = [1, 4]$.
			\item Полный прообраз множества $[1, 4]$: $f^{-1}([1, 4]) = [-2, -1] \cup [1, 2]$
			\item Образ отображения: $\operatorname{Im} f = [0, +\infty)$.
		\end{itemize}
	\end{tcolorbox}
	
	\subsection{Инъективность, сюръективность, биективность}
	
	\noindent \textbf{Определение.} Отображение $f \colon A \to B$ называется:
	\begin{itemize}
		\item \textit{инъективным} (или \textit{инъекцией}), если $\forall a, b \in A \ a \neq b \Rightarrow f(a) \neq f(b)$;
		\item \textit{сюръективным} (или \textit{сюръекцией}), если $\operatorname{Im} f = B$;
		\item \textit{биективным} (или \textit{биекцией}), если оно инъективно и сюръективно
	\end{itemize}
	
	\medskip
	\noindent \textbf{Примеры.}
	
	\begin{tcolorbox}[colback=gray!5, colframe=gray!50, title=Инъекция]
		$f \colon \mathbb{R}_+ \to \mathbb{R}$, $f(x) = x^2$ — \textbf{инъекция}, так как различные неотрицательные числа имеют различные квадраты
	\end{tcolorbox}
	
	\begin{tcolorbox}[colback=gray!5, colframe=gray!50, title=Не инъекция]
		$f \colon \mathbb{R} \to \mathbb{R}$, $f(x) = x^2$ \textbf{не является инъекцией}, поскольку $f(-2) = f(2) = 4$, но $-2 \neq 2$
	\end{tcolorbox}
	
	\begin{tcolorbox}[colback=gray!5, colframe=gray!50, title=Сюръекция]
		$f(x) = \sin x$, $f \colon \mathbb{R} \to [-1; 1]$ — \textbf{сюръекция}, так как любое число из отрезка $[-1, 1]$ достигается как значение синуса
	\end{tcolorbox}
	
	\begin{tcolorbox}[colback=gray!5, colframe=gray!50, title=Не сюръекция]
		$f(x) = \sin x$, $f \colon \mathbb{R} \to \mathbb{R}$ \textbf{не является сюръекцией}, поскольку, например, $2 \notin \operatorname{Im} f$ (синус принимает значения только от $-1$ до $1$)
	\end{tcolorbox}
	
	\begin{tcolorbox}[colback=gray!5, colframe=gray!50, title=Биекция]
		$f \colon \mathbb{R} \to \mathbb{R}$, $f(x) = 2x + 1$ — \textbf{биекция} (и инъективна, и сюръективна)
	\end{tcolorbox}
	
	\begin{tcolorbox}[colback=gray!5, colframe=gray!50, title=Биекция]
		$f \colon \mathbb{R} \to \mathbb{R}$, $f(x) = x^3$ — \textbf{биекция} (кубическая парабола строго монотонна и принимает все действительные значения)
	\end{tcolorbox}
	
	\myremark
	Инъективность можно проверить так: если $f(a) = f(b)$, то обязательно $a = b$. Сюръективность означает, что уравнение $f(x) = y$ имеет решение для любого $y \in B$. Биекция устанавливает взаимно однозначное соответствие между элементами множеств $A$ и $B$
	
	\begin{tcolorbox}[colback=purple!5, colframe=purple!50!black, title=Связь с мощностями]
		Если существует биекция $f \colon A \to B$, то множества $A$ и $B$ равномощны: $|A| = |B|$. Это ещё один способ определять равенство мощностей
	\end{tcolorbox}
	
	\subsection{Композиция отображений}
	
	\noindent \textbf{Определение (Композиция).}
	Пусть $f \colon B \to C$ и $g \colon A \to B$ — пара отображений. \textit{Произведением} этой пары называется отображение $fg \colon A \to C$ такое, что $\forall x \in A \ fg(x) = f(g(x))$. Произведение отображений называют также \textit{композицией} или \textit{суперпозицией}, обозначают также $f \circ g$
	
	\begin{tcolorbox}[colback=blue!5, colframe=blue!40!black, title=Пример композиции]
		Пусть $g \colon \mathbb{R} \to \mathbb{R}$, $g(x) = 2x$, и $f \colon \mathbb{R} \to \mathbb{R}$, $f(y) = y^2 + 1$.
		Тогда $f \circ g \colon \mathbb{R} \to \mathbb{R}$:
		\[
		(f \circ g)(x) = f(g(x)) = f(2x) = (2x)^2 + 1 = 4x^2 + 1.
		\]
		А $g \circ f$ (если определить соответствующим образом) было бы $g(f(x)) = 2(x^2 + 1) = 2x^2 + 2$. Композиция, вообще говоря, не коммутативна!
	\end{tcolorbox}
	
	\myremark
	Важно следить за порядком: сначала применяется отображение, стоящее справа от знака композиции. Запись $f \circ g$ означает «сначала $g$, потом $f$»
	
	\subsection{Свойства композиции}
	
	\noindent \textbf{Предложение (Композиция ассоциативна).}
	То есть, если есть три отображения — $h \colon U \to V$, $g \colon V \to W$, $f \colon W \to T$, то
	\[
	f(gh) = (fg)h.
	\]
	
	\medskip
	\noindent \textbf{Доказательство.}
	$\forall u \in U$:
	\[
	(f(gh))(u) = f(gh(u)) = f(g(h(u))) = (fg)(h(u)) = ((fg)h)(u).
	\]
	Таким образом, значения отображений совпадают для любого $u$, следовательно, отображения равны.
	
	\begin{tcolorbox}[colback=green!5, colframe=green!50!black, title=Пояснение]
		Ассоциативность означает, что скобки можно расставлять произвольно: результат композиции трёх отображений не зависит от того, скомпонуем ли мы сначала $g$ и $h$, а потом $f$, или сначала $f$ и $g$, а потом применим к $h$. Благодаря этому свойству запись $f \circ g \circ h$ однозначна без скобок
	\end{tcolorbox}
	
	\myremark
	Ассоциативность — фундаментальное свойство, которое позволяет рассматривать композицию как бинарную операцию на множестве отображений из множества в себя, превращая это множество в полугруппу (а если есть тождественное отображение — то в моноид)
	
	\subsection{Ограничение, срезка, сужение}
	
	\noindent \textbf{Определение (Ограничение, срезка, сужение).}
	Пусть $f \colon A \to B$, $A' \subset A$, $B' \subset B$, $\operatorname{Im} f \subset B'$
	
	\begin{itemize}
		\item \textit{Ограничением} отображения $f$ на $A'$ называется отображение $f|_{A'} \colon A' \to B$ такое, что $\forall x \in A' \ f|_{A'}(x) = f(x)$
		
		\item \textit{Срезкой} отображения $f$ на $B'$ будем называть отображение $A$ в $B'$, действующее на элементы из $A$ так же, как $f$. Стандартного обозначения срезка не имеет
		
		\item \textit{Сужением} отображения $f$ будем называть срезку ограничения $f$, т.е. отображение $A'$ в $B'$, где $A' \subset A$, $f(A') \subset B' \subset B$
	\end{itemize}
	
	\begin{tcolorbox}[colback=gray!5, colframe=gray!50, title=Примеры]
		Рассмотрим $f \colon \mathbb{R} \to \mathbb{R}$, $f(x) = x^2$.
		
		\begin{itemize}
			\item \textbf{Ограничение} на $A' = [0, +\infty)$: $f|_{[0, +\infty)} \colon [0, +\infty) \to \mathbb{R}$, $f|_{[0, +\infty)}(x) = x^2$. Это отображение уже инъективно, в отличие от исходного
			
			\item \textbf{Срезка} на $B' = [0, +\infty)$: $g \colon \mathbb{R} \to [0, +\infty)$, $g(x) = x^2$. Это отображение сюръективно (в отличие от исходного), но не инъективно
			
			\item \textbf{Сужение} на $A' = [0, +\infty)$ и $B' = [0, +\infty)$: $h \colon [0, +\infty) \to [0, +\infty)$, $h(x) = x^2$. Это отображение уже биективно!
		\end{itemize}
	\end{tcolorbox}
	
	\myremark
	Ограничение, срезка и сужение — полезные конструкции, позволяющие «исправлять» отображения, делая их инъективными (ограничивая начало) или сюръективными (сужая конец). Сужение позволяет получить биекцию, если исходное отображение было инъективно на некотором подмножестве и его образ совпадает с выбранным подмножеством конца
	
	\begin{tcolorbox}[colback=yellow!5, colframe=orange!60!black, title=Заметочка: Связь с предыдущими темами]
		Понятия инъекции, сюръекции и биекции тесно связаны с мощностями множеств:
		\begin{itemize}
			\item Существование инъекции $A \to B$ означает $|A| \leq |B|$
			\item Существование сюръекции $A \to B$ означает $|A| \geq |B|$ (для конечных множеств; для бесконечных это тоже верно при определённых условиях)
			\item Существование биекции $A \to B$ означает $|A| = |B|$.
		\end{itemize}
		Таким образом, отображения — это инструмент для сравнения мощностей
	\end{tcolorbox}
	
	
	
	
	
	
	
	
	
	% ==================== РАЗДЕЛ 4: ОБРАТНОЕ ОТОБРАЖЕНИЕ ====================
	\newpage
	\section{Обратное отображение. Теорема об обратном отображении (4 вопрос)}
	
	\subsection{Тождественное отображение}
	
	\noindent \textbf{Определение (Тождественное отображение).}
	Пусть $A$ — произвольное множество. \textit{Тождественным отображением} в $A$ называется отображение $\operatorname{id}_A \colon A \to A$ такое, что $\forall x \in A \ \operatorname{id}_A(x) = x$
	
	\begin{tcolorbox}[colback=blue!5, colframe=blue!40!black, title=Пример]
		Для множества $A = \{1, 2, 3\}$ тождественное отображение действует так:
		$\operatorname{id}_A(1) = 1$, $\operatorname{id}_A(2) = 2$, $\operatorname{id}_A(3) = 3$.
		На графике это просто диагональ $y = x$
	\end{tcolorbox}
	
	\myremark
	Тождественное отображение играет роль «нейтрального элемента» относительно композиции: при композиции с любым отображением оно не меняет результат
	
	
	
	
		
	\subsection{Свойство тождественного отображения}
	
	\noindent \textbf{Свойство.} Пусть $f \colon A \to B$. Тогда $f \circ \operatorname{id}_A = \operatorname{id}_B \circ f = f$.
	
	\begin{center}
	\begin{tikzpicture}[>=stealth]
	
	% Узлы в одну линию
	\node (B1) at (0,0) {$B$};
	\node (B2) at (2.5,0) {$B$};
	\node (A1) at (5,0) {$A$};
	\node (A2) at (7.5,0) {$A$};
	
	% Горизонтальные стрелки (влево)
	\draw[->] (B2) -- node[above] {$\operatorname{id}_B$} (B1);
	\draw[->] (A1) -- node[above] {$f$} (B2);
	\draw[->] (A2) -- node[above] {$\operatorname{id}_A$} (A1);
	
	% Верхняя дуга
	\draw[->, bend left=35]
	(A2) to node[above] {$f \circ \operatorname{id}_A$} (B2);
	
	% Нижняя дуга
	\draw[->, bend right=35]
	(A1) to node[below] {$\operatorname{id}_B \circ f$} (B1);
	
	\end{tikzpicture}
	\end{center}
	
		Из диаграммы видно, что отображения $\operatorname{id}_B \circ f$ и $f \circ \operatorname{id}_A$ определены, причём их начала и концы совпадают с началом и концом исходного отображения $f$:
	\begin{itemize}
	\item Отображение $f \circ \operatorname{id}_A$ начинается в $A$ (нижнем) и заканчивается в $B$ (верхнем)
	\item Отображение $\operatorname{id}_B \circ f$ начинается в $A$ (верхнем) и заканчивается в $B$ (нижнем)
	\item Само отображение $f$ определено из верхнего $A$ в верхний $B$, а также (пунктиром) из нижнего $A$ в нижний $B$ после применения тождественных
	\end{itemize}
	
	Проверим равенство отображений поэлементно. Для любого элемента $x \in A$:
	\[
	(f \circ \operatorname{id}_A)(x) = f(\operatorname{id}_A(x)) = f(x),
	\]
	\[
	(\operatorname{id}_B \circ f)(x) = \operatorname{id}_B(f(x)) = f(x).
	\]
	
	Таким образом, значения всех трёх отображений на любом элементе совпадают. Совпадение областей определения, областей значений и графиков (т.е. поточечное совпадение) означает, что эти отображения равны как функции:
	\[
	f \circ \operatorname{id}_A = f = \operatorname{id}_B \circ f.
	\]
	
	\noindent \textbf{Итог.} Тождественное отображение является нейтральным элементом относительно операции композиции функций: композиция любой функции с тождественной (слева или справа) даёт исходную функцию
	
	
	
	\medskip	
	\begin{tcolorbox}[colback=green!5, colframe=green!50!black, title=Пояснение]
		Это свойство аналогично умножению на единицу в арифметике: $a \cdot 1 = 1 \cdot a = a$. Тождественное отображение — это «единица» в мире отображений относительно операции композиции
	\end{tcolorbox}
	
	\subsection{Обратные отображения}
	
	\noindent \textbf{Определение (Обратные отображения).}
	Пусть $f \colon A \to B$. Отображение $g$ называется:
	\begin{itemize}
		\item \textit{Левым обратным} отображению $f$, если $gf = \operatorname{id}_A$;
		\item \textit{Правым обратным} отображению $f$, если $fg = \operatorname{id}_B$;
		\item \textit{Обратным} отображению $f$, если $g$ есть левое и правое обратное
	\end{itemize}
	
	\begin{tcolorbox}[colback=gray!5, colframe=gray!50, title=Пример]
		Рассмотрим $f \colon \mathbb{R} \to [0, +\infty)$, $f(x) = x^2$.
		\begin{itemize}
			\item Отображение $g \colon [0, +\infty) \to \mathbb{R}$, $g(y) = \sqrt{y}$ является \textbf{правым обратным}, потому что $f(g(y)) = (\sqrt{y})^2 = y = \operatorname{id}_{[0, +\infty)}(y)$. Но оно не является левым обратным, так как $g(f(x)) = \sqrt{x^2} = |x| \neq x$ при $x < 0$.
			\item Отображение $h \colon [0, +\infty) \to \mathbb{R}_+$, $h(y) = \sqrt{y}$ будет \textbf{обратным} к $f|_{\mathbb{R}_+} \colon \mathbb{R}_+ \to [0, +\infty)$ (ограничению на неотрицательные числа).
		\end{itemize}
	\end{tcolorbox}
	
	\begin{tcolorbox}[colback=yellow!5, colframe=orange!60!black, title=Замечание]
		Ясно, что $g$ — левое обратное $f$ тогда и только тогда, когда $f$ — правое обратное $g$
	\end{tcolorbox}
	
	\myremark
	Левое и правое обратное — это не одно и то же! Левое обратное «отменяет» действие $f$ после того, как оно произошло, а правое обратное — до того. В общем случае они могут существовать по отдельности и быть разными
	
	\subsection{Единственность обратного отображения}
	
	\noindent \textbf{Предложение (О единственности обратного отображения).}
	Если у отображения $f$ существует левое и правое обратные отображения, то они совпадают, и в этом случае у $f$ имеется единственное левое обратное, оно же единственное правое обратное, оно же единственное обратное отображение. В частности, если у отображения существует обратное, то оно единственно
	
	\medskip
	\noindent \textbf{Доказательство.}
	Пусть $g$ — левое, $h$ — правое обратные отображению $f \colon A \to B$. Тогда
	\[
	gfh = 
	\begin{cases}
		(gf)h = \operatorname{id}_A \circ h = h, \\
		g(fh) = g \circ \operatorname{id}_B = g,
	\end{cases}
	\]
	откуда $g = h$.
	
	Если $g'$ — также левое обратное $f$, то, по доказанному, $g' = h$, т.е. $g' = g$. Аналогично показывается единственность правого обратного. Последнее утверждение предложения очевидно
	
	\medskip
	Обратное $f$ отображение обычно обозначается $f^{-1}$
	
	\begin{tcolorbox}[colback=green!5, colframe=green!50!black, title=Пояснение к доказательству]
		Ключевой момент: мы рассматриваем тройную композицию $gfh$ и группируем скобки двумя способами. Благодаря ассоциативности композиции, оба способа дают одно и то же. Из равенства $g = h$ следует, что левое и правое обратные — это одно и то же отображение
	\end{tcolorbox}
	
	\subsection{Обратимость}
	
	\noindent \textbf{Определение (Обратимость).}
	Отображение называется:
	\begin{itemize}
		\item \textit{Обратимым слева}, если существует левое обратное ему отображение;
		\item \textit{Обратимым справа}, если существует правое обратное ему отображение;
		\item \textit{Обратимым}, если существует обратное ему отображение
	\end{itemize}
	
	\subsection{Критерии обратимости}
	
	\noindent \textbf{Теорема (Критерии обратимости).}
	Отображение
	\begin{enumerate}
		\item обратимо слева тогда и только тогда, когда оно инъективно,
		\item обратимо справа тогда и только тогда, когда оно сюръективно,
		\item обратимо тогда и только тогда, когда оно биективно
	\end{enumerate}
	
	\medskip
	\noindent \textbf{Доказательство.}
	
	\textbf{1.} Пусть $f \colon A \to B$ обратимо слева, т.е. $\exists g \colon B \to A$ такое, что $gf = \operatorname{id}_A$. Пусть $x, x' \in A$ и $f(x) = f(x')$. Тогда
	\[
	x = \operatorname{id}_A(x) = (gf)(x) = g(f(x)) = g(f(x')) = (gf)(x') = \operatorname{id}_A(x') = x',
	\]
	то есть $f$ — инъекция
	
	\medskip
	Обратно, пусть $f$ инъективно. Тогда $\forall y \in \operatorname{Im} f \subset B \ \exists! x(y) \in A \colon f(x) = y$. Положим $g(y) = x(y)$ для $y \in \operatorname{Im} f$, а для $y \in B \setminus \operatorname{Im} f$ значение $g(y)$ выберем в $A$ совершенно произвольно. Тогда, очевидно, получим $gf = \operatorname{id}_A$
	
	\begin{tcolorbox}[colback=blue!5, colframe=blue!40!black, title=Пояснение к пункту 1]
		В обратную сторону мы строим левое обратное, используя инъективность: для каждого образа выбираем его единственный прообраз, а элементы вне образа отображаем куда угодно (например, в фиксированный элемент). Произвол в определении на дополнении показывает, что левых обратных может быть много
	\end{tcolorbox}
	
	\textbf{2.} Пусть $f \colon A \to B$ обратимо справа, т.е. $\exists g \colon B \to A$ такое, что $fg = \operatorname{id}_B$. Тогда $\forall y \in B$:
	\[
	y = \operatorname{id}_B(y) = (fg)(y) = f(g(y)),
	\]
	т.е. $f$ — сюръекция.
	
	\medskip
	Обратно, пусть $f$ сюръективно. У каждого элемента $y \in B$ есть его полный прообраз при отображении $f$. Эти прообразы у разных элементов $y_1, y_2 \in B$ не пересекаются. В качестве $g(y)$ выберем любой $x$ из полного прообраза $y$, т.е. $x \in f^{-1}(y) \subset A$. Тогда, очевидно, $fg = \operatorname{id}_B$
	
	\begin{tcolorbox}[colback=blue!5, colframe=blue!40!black, title=Пояснение к пункту 2]
		В обратную сторону мы строим правое обратное, используя сюръективность: для каждого $y$ выбираем произвольный прообраз из непустого множества $f^{-1}(y)$. Здесь требуется аксиома выбора (для бесконечных множеств), так как мы делаем бесконечно много произвольных выборов одновременно
	\end{tcolorbox}
	
	\textbf{3.} Очевидным образом следует из пп. 1 и 2.
	
	\begin{tcolorbox}[colback=red!5, colframe=red!70!black, title=Следствие теоремы]
		Очевидным образом вытекают из теоремы и замечания следующие утверждения:
		\begin{itemize}
			\item Левое обратное какому-либо отображению (необходимо инъективному) сюръективно
			\item Правое обратное какому-либо отображению (необходимо сюръективному) инъективно
			\item Обратное какому-либо отображению (необходимо биективному) биективно
		\end{itemize}
	\end{tcolorbox}
	
	\myremark
	Эти критерии дают глубокую связь между алгебраическими свойствами (существование обратного относительно композиции) и теоретико-множественными свойствами (инъективность и сюръективность). Биективные отображения — это в точности те, для которых существует двусторонний обратный элемент
	
	\begin{tcolorbox}[colback=purple!5, colframe=purple!50!black, title=Заметочка: Связь с предыдущими темами]
		\begin{itemize}
			\item Инъективность $\Leftrightarrow$ существование левого обратного $\Leftrightarrow$ $|A| \leq |B|$ (для конечных множеств)
			\item Сюръективность $\Leftrightarrow$ существование правого обратного $\Leftrightarrow$ $|A| \geq |B|$ (для конечных множеств)
			\item Биективность $\Leftrightarrow$ существование обратного $\Leftrightarrow$ $|A| = |B|$
		\end{itemize}
		Таким образом, понятие обратного отображения замыкает круг идей, связывая мощности, свойства отображений и возможность «обратить» действие
	\end{tcolorbox}
	
	\begin{tcolorbox}[colback=yellow!5, colframe=orange!60!black, title=Важное замечание о единственности]
		Из предложения 4.4 и теоремы 4.6 следует, что если отображение биективно, то обратное к нему существует и единственно. Для инъективных (но не сюръективных) отображений левых обратных может быть много (из-за произвола на дополнении к образу). Для сюръективных (но не инъективных) отображений правых обратных также может быть много (из-за произвола в выборе прообраза)
	\end{tcolorbox}
	
	
	
	
	
	
	
	
	
	
	
	
	
	
	% ==================== РАЗДЕЛ 5: ОТНОШЕНИЯ. ЭКВИВАЛЕНТНОСТЬ ====================
	\newpage
	\section{Отношения. Эквивалентность. Разбиение множества на классы эквивалентности (5 вопрос)}
	
	\subsection{Бинарные отношения}
	
	\noindent \textbf{Определение (Бинарное отношение, график, соответствие).}
	Пусть $(A, B)$ — пара непустых множеств. \textit{Бинарным отношением} в этой паре (или из $A$ в $B$) называется тройка $(A, B, R)$, где $R \subset A \times B$
	
	\medskip
	Если $\rho = (A, B, R)$ — бинарное отношение, то $R$ называется его \textit{графиком}. Вместо $(a, b) \in R$ пишут $a \rho b$ и говорят, что $a$ соответствует $b$ (по отношению $\rho$), или что $a$ и $b$ находятся в этом отношении (соответствии). Если $A = B$, то говорят, что $\rho$ — отношение в $A$
	
	\begin{tcolorbox}[colback=gray!5, colframe=gray!50, title=Примеры бинарных отношений]
		\begin{itemize}
			\item Отношение «$\leq$» на множестве действительных чисел: $(A, B) = (\mathbb{R}, \mathbb{R})$, график $R = \{(x, y) \mid x \leq y\}$
			\item Отношение «быть братом» на множестве людей: $a \rho b$, если $a$ — брат $b$
			\item Отношение «делит» на множестве натуральных чисел: $a \rho b$, если $a \mid b$.
			\item Отношение «жить в одном городе» на множестве людей.
		\end{itemize}
	\end{tcolorbox}
	
	\myremark
	Бинарное отношение — это просто правило, которое для каждой упорядоченной пары $(a, b)$ говорит, связаны они этим отношением или нет. График $R$ — это подмножество декартова произведения, состоящее из всех пар, для которых отношение выполняется
	
	\subsection{Свойства бинарных отношений}
	
	\noindent \textbf{Определение (Рефлексивность, симметричность, транзитивность).}
	Отношение $\rho$ в множестве $A$ называется
	\begin{itemize}
		\item \textit{Рефлексивным}, если $\forall a \in A \ a \rho a$;
		\item \textit{Симметричным}, если $a \rho b \Rightarrow b \rho a$;
		\item \textit{Транзитивным}, если $a \rho b, b \rho c \Rightarrow a \rho c$
	\end{itemize}
	
	\begin{tcolorbox}[colback=blue!5, colframe=blue!40!black, title=Примеры свойств]
		\begin{itemize}
			\item Отношение «$\leq$» на $\mathbb{R}$: \textbf{рефлексивно} ($x \leq x$), \textbf{не симметрично} ($x \leq y$ не влечёт $y \leq x$), \textbf{транзитивно} ($x \leq y$ и $y \leq z$ влечёт $x \leq z$).
			\item Отношение «$=$» на любом множестве: \textbf{рефлексивно}, \textbf{симметрично}, \textbf{транзитивно}.
			\item Отношение «брат» на множестве мужчин: \textbf{не рефлексивно} (человек не может быть братом самому себе), \textbf{симметрично} (если $a$ брат $b$, то $b$ брат $a$), \textbf{не транзитивно} (брат брата может быть тем же самым человеком или не быть братом).
			\item Отношение «делит» на $\mathbb{N}$: \textbf{рефлексивно} ($a \mid a$), \textbf{не симметрично} ($a \mid b$ и $b \mid a$ возможно только при $a = b$), \textbf{транзитивно} ($a \mid b$ и $b \mid c$ влечёт $a \mid c$)
		\end{itemize}
	\end{tcolorbox}
	
	\myremark
	Рефлексивность означает, что каждый элемент связан с самим собой. Симметричность — что отношение не зависит от порядка пары. Транзитивность — что отношение можно «продолжать» по цепочке
	
	\subsection{Отношение эквивалентности}
	
	\noindent \textbf{Определение (Эквивалентность, класс эквивалентности, представитель).}
	Отношение в некотором множестве называется \textit{отношением эквивалентности}, или \textit{эквивалентностью}, если оно рефлексивно, симметрично и транзитивно
	
	\medskip
	Пусть $\sim$ обозначает некоторую эквивалентность в $A$, и $a \in A$. Множество $\widetilde{a} = \{x \in A \mid x \sim a\}$ называется \textit{классом эквивалентности} $\sim$, порождённым $a$ (ясно, что $\widetilde{a} \subset A$), его элементы называются \textit{представителями} этого класса. В силу рефлексивности отношения $\forall a \in A \ a \in \widetilde{a}$
	
	\begin{tcolorbox}[colback=gray!5, colframe=gray!50, title=Примеры отношений эквивалентности]
		\begin{itemize}
			\item Отношение равенства $=$ на любом множестве: каждый элемент образует свой собственный класс эквивалентности $\widetilde{a} = \{a\}$
			\item Отношение «иметь одинаковый остаток при делении на $m$» на $\mathbb{Z}$ (сравнение по модулю $m$): классы эквивалентности — это классы вычетов $\{0, 1, \dots, m-1\}$
			\item Отношение «быть похожим по цвету» на множестве предметов: классы — это группы предметов одного цвета
			\item Отношение «жить в одной стране» на множестве людей: классы — это группы граждан одной страны
		\end{itemize}
	\end{tcolorbox}
	
	\begin{tcolorbox}[colback=yellow!5, colframe=orange!60!black, title=Замечание о классах]
		Класс эквивалентности $\widetilde{a}$ — это множество всех элементов, эквивалентных $a$. Поскольку отношение симметрично и транзитивно, все элементы внутри одного класса эквивалентны друг другу, а элементы из разных классов — нет
	\end{tcolorbox}
	
	\subsection{Свойства классов эквивалентности}
	
	\noindent \textbf{Немного свойств классов эквивалентности.}
	
	Пусть $a' \in \widetilde{a}$, т.е. $a' \sim a$. Тогда $\forall x \in \widetilde{a'} \ x \sim a' \sim a$, и по транзитивности $x \sim a$. Значит, $\widetilde{a'} \subset \widetilde{a}$. В силу симметричности $a \sim a'$ и таким же рассуждением получим $\widetilde{a} \subset \widetilde{a'}$. Итак, если $a' \in \widetilde{a}$, то $\widetilde{a'} = \widetilde{a}$, т.е. каждый класс эквивалентности порождается любым своим элементом (представителем)
	
	Пусть $\widetilde{a}$, $\widetilde{b}$ — пара классов эквивалентности $\sim$. Если $\widetilde{a} \cap \widetilde{b} \neq \varnothing$, то $\exists x \in \widetilde{a} \cap \widetilde{b}$, т.е. $x \in \widetilde{a}$, $x \in \widetilde{b}$, а потому $\widetilde{x} = \widetilde{a}$, $\widetilde{x} = \widetilde{b}$, то есть $\widetilde{a} = \widetilde{b}$. Таким образом, классы эквивалентности либо не пересекаются, либо совпадают
	
	\begin{tcolorbox}[colback=green!5, colframe=green!50!black, title=Пояснение]
		Эти два свойства — ключевые:
		\begin{enumerate}
			\item Любой элемент класса может служить его представителем (класс не зависит от выбора представителя).
			\item Классы образуют разбиение множества: они попарно не пересекаются и покрывают всё множество
		\end{enumerate}
	\end{tcolorbox}
	
	\myremark
	Второе свойство часто формулируют так: отношение эквивалентности разбивает множество на непересекающиеся классы. Это фундаментальный факт, лежащий в основе многих математических конструкций (факторгруппы, факторкольца, факторпространства и т.д.)
	
	\subsection{Фактор-множество и разбиение}
	
	\noindent \textbf{Определение (Фактор-множество, разбиение / дизъюнктное объединение).}
	Множество всех классов эквивалентности $\sim$ в $A$ называется \textit{фактор-множеством} множества $A$ по этой эквивалентности и обозначается $A/{\sim}$
	
	$A$ есть объединение множества всех классов данной эквивалентности: $A = \bigcup_{\widetilde{a} \in A/{\sim}} \widetilde{a}$. При этом если $\widetilde{a} \neq \widetilde{b}$, то $\widetilde{a} \cap \widetilde{b} = \varnothing$. Представление некоторого множества как объединение его попарно непересекающихся непустых подмножеств называется \textit{разбиением} этого множества (говорят также, что это \textit{дизъюнктное объединение}). Таким образом, всякая эквивалентность в некотором множестве определяет некоторое разбиение этого множества, а именно разбиение его на классы этой эквивалентности
	
	\begin{tcolorbox}[colback=blue!5, colframe=blue!40!black, title=Пример фактор-множества]
		Рассмотрим $A = \mathbb{Z}$ и отношение $\sim$: $a \sim b$, если $a - b$ делится на $3$ (сравнение по модулю $3$).
		\begin{itemize}
			\item Классы эквивалентности: 
			$\widetilde{0} = \{\dots, -3, 0, 3, 6, \dots\}$,
			$\widetilde{1} = \{\dots, -2, 1, 4, 7, \dots\}$,
			$\widetilde{2} = \{\dots, -1, 2, 5, 8, \dots\}$.
			\item Фактор-множество $\mathbb{Z}/{\sim} = \{\widetilde{0}, \widetilde{1}, \widetilde{2}\}$ состоит из трёх элементов.
			\item Это разбиение $\mathbb{Z}$ на три непересекающихся класса
		\end{itemize}
	\end{tcolorbox}
	
	\subsection{Взаимно однозначное соответствие}
	
	\noindent \textbf{Теорема (Соответствие между эквивалентностями и разбиениями).}
	Обратно, пусть $A = \bigcup_i A_i$ — произвольное разбиение множества $A$. Зададим в $A$ отношение $\sim$ следующим образом: $a \sim a' \iff \exists i \colon a \in A_i, a' \in A_i$. Тогда $\sim$ — эквивалентность, а все $A_i$ — все её классы
	
	Таким образом, имеется взаимно однозначное соответствие между множеством всех эквивалентностей в $A$ и множеством всех разбиений этого множества. Другими словами, понятия эквивалентности и разбиения по существу равнозначны, и всё равно, говорить ли о некоторой эквивалентности или о некотором разбиении
	
	
	\begin{tcolorbox}[colback=purple!5, colframe=purple!50!black, title=Заметочка: Значение эквивалентности]
		Понятие эквивалентности — одно из важнейших в математике. Оно позволяет:
		\begin{itemize}
			\item Отождествлять объекты, обладающие общим свойством (например, дроби $\frac{1}{2}$ и $\frac{2}{4}$ считаются одним рациональным числом)
			\item Строить новые математические объекты (фактор-конструкции) — факторгруппы, факторкольца, факторпространства
			\item Классифицировать объекты по некоторому признаку (например, в геометрии классификация фигур по группам симметрии)
		\end{itemize}
		Фактор-множество $A/{\sim}$ «склеивает» эквивалентные элементы в один объект, позволяя работать с классами как с новыми элементами
	\end{tcolorbox}
	
	\begin{tcolorbox}[colback=yellow!5, colframe=orange!60!black, title=Важное замечание]
		Каноническая проекция $\pi \colon A \to A/{\sim}$, заданная формулой $\pi(a) = \widetilde{a}$, является сюръективным отображением. Она каждому элементу ставит в соответствие его класс эквивалентности. Это отображение играет ключевую роль в теории категорий и алгебраических конструкциях
	\end{tcolorbox}
	
	\myremark
	Таким образом, задать эквивалентность на множестве — это то же самое, что разбить его на непересекающиеся подмножества. Это простая, но глубокая идея: когда мы говорим «эти объекты одинаковы с точностью до некоторого признака», мы неявно задаём эквивалентность и работаем с классами
	
	
	
	
	
	
	% ==================== РАЗДЕЛ 6: ПЕРЕСТАНОВКИ ====================
	\newpage
	\section{Перестановки. Разложение перестановки в произведение независимых циклов и транспозиций (6 вопрос)}
	
	\subsection{Определение перестановки}
	
	\noindent \textbf{Определение (Перестановка).}
	Пусть $\Omega$ — конечное множество из $n$ элементов. Удобно считать, что $\Omega = \{1, 2, \dots, n\}$, поскольку природа его элементов для нас несущественна. Элементы множества $S_n = S(\Omega)$ всех взаимно однозначных преобразований $\Omega \to \Omega$ называются \textit{перестановками} и обычно обозначаются строчными буквами греческого алфавита. Лишь за единичным преобразованием $e = e_\Omega$ сохранилась буква латинского алфавита
	
	\begin{tcolorbox}[colback=blue!5, colframe=blue!40!black, title=Пример перестановки]
		Для $n = 4$ перестановка $\pi$ может действовать так:
		$\pi(1) = 2$, $\pi(2) = 4$, $\pi(3) = 1$, $\pi(4) = 3$.
	\end{tcolorbox}
	
	\myremark
	Перестановка — это биекция конечного множества на себя. Такие отображения автоматически являются биективными, поскольку для конечного множества инъективность, сюръективность и биективность равносильны
	
	\subsection{Развёрнутая запись}
	
	\noindent \textbf{Определение (Развёрнутая запись).}
	В развёрнутой и наглядной форме произвольную перестановку $\pi \colon i \to \pi(i)$, $i = 1, 2, \dots, n$ изображают в виде
	\[
	\pi = \begin{pmatrix}
		1 & 2 & \cdots & n \\
		i_1 & i_2 & \cdots & i_n
	\end{pmatrix}
	\]
	
	полностью указывая все образы:
	\[
	\pi \colon \begin{pmatrix}
	1 & 2 & \dots & n \\
	\downarrow & \downarrow & & \downarrow \\
	i_1 & i_2 & \dots & i_n
	\end{pmatrix}
	\]
	
	где $i_k = \pi(k)$ ($k = 1, 2, \dots, n$) — переставленные символы $1, 2, \dots, n$.
	
	\begin{tcolorbox}[colback=gray!5, colframe=gray!50, title=Пример]
		Перестановка из предыдущего примера запишется как
		\[
		\pi = \begin{pmatrix}
			1 & 2 & 3 & 4 \\
			2 & 4 & 1 & 3
		\end{pmatrix}
		\]
	\end{tcolorbox}
	
	\subsection{Композиция перестановок}
	
	\noindent \textbf{Определение (Композиция перестановок).}
	Перестановки перемножаются в соответствии с общим правилом композиции отображений: $(\sigma \tau)(i) = \sigma(\tau(i))$
	
	
	
	\begin{tcolorbox}[colback=gray!5, colframe=gray!50, title={Пример умножения перестановок}]
	Для перестановок
	\[
	\sigma = \begin{pmatrix}
	1 & 2 & 3 & 4 \\
	2 & 3 & 4 & 1
	\end{pmatrix}, \quad
	\tau = \begin{pmatrix}
	1 & 2 & 3 & 4 \\
	4 & 3 & 2 & 1
	\end{pmatrix}
	\]
	имеем
	\[
	\sigma \tau = \begin{pmatrix}
	1 & 2 & 3 & 4 \\
	2 & 3 & 4 & 1
	\end{pmatrix}
	\begin{pmatrix}
	1 & 2 & 3 & 4 \\
	4 & 3 & 2 & 1
	\end{pmatrix} =
	\begin{array}{cccc}
	1 & 2 & 3 & 4 \\
	\downarrow & \downarrow & \downarrow & \downarrow \\
	4 & 3 & 2 & 1 \\
	\downarrow & \downarrow & \downarrow & \downarrow \\
	1 & 4 & 3 & 2
	\end{array} = \begin{pmatrix}
	1 & 2 & 3 & 4 \\
	1 & 4 & 3 & 2
	\end{pmatrix}
	\]
	\end{tcolorbox}
	
	
	\subsection{Свойства умножения перестановок}
	
	\noindent \textbf{Свойства.}
	Ясно, что умножение перестановок подчиняется правилам композиции отображений:
	\begin{enumerate}
		\item Умножение ассоциативно.
		\item $S_n$ обладает единичным элементом: $\pi e = e \pi = \pi \ \ \forall \pi \in S_n$
		\item Для каждой перестановки $\pi \in S_n$ существует обратная перестановка $\pi^{-1}$: $\pi \pi^{-1} = \pi^{-1} \pi = e$.
	\end{enumerate}
	
	Эти 3 свойства позволяют говорить о \textit{группе} $S_n$. Множество $S_n$, рассматриваемое вместе с естественной операцией умножения его элементов (композицией перестановок) называется \textit{симметрической группой} степени $n$
	
	Ясно, что порядок группы $\operatorname{Card} S_n = |S_n| = n!$.
	
	\begin{tcolorbox}[colback=purple!5, colframe=purple!50!black, title=Заметочка: Симметрическая группа]
		Симметрическая группа $S_n$ — один из важнейших объектов в алгебре. Она содержит $n!$ элементов и играет ключевую роль в теории групп, комбинаторике и приложениях. Уже при $n = 3$ группа $S_3$ некоммутативна
	\end{tcolorbox}
	
	\subsection{Циклы и произведение циклов}
	
	\noindent \textbf{Определение (Цикл, произведение циклов).}
	Перестановка $\sigma$, кратко записываемая в виде $\sigma = (1\ 2\ 3\ 4)$, или, что то же самое,
	в виде $\sigma = (2\ 3\ 4\ 1) = (3\ 4\ 1\ 2) = (4\ 1\ 2\ 3)$, называется \textit{циклом длины 4}, а
	перестановка $\tau = (1\ 4)(2\ 3)$ — \textit{произведением двух независимых (непересекающихся) циклов} $(1\ 4)$ и $(2\ 3)$ длины 2
	
	
	
	
	\begin{center}
	\begin{tikzpicture}[>=stealth, scale=1.2, every node/.style={circle, minimum size=0.5cm}]
	
	% -------- сигма (цикл из 4 элементов) --------
	\node at (-1.2,0.9) {$\sigma =$};
	
	\node (1) at (0,0.9) {3};
	\node (2) at (0.9,1.8) {2};
	\node (3) at (1.8,0.9) {1};
	\node (4) at (0.9,0) {4};
	
	\draw[->, bend right=25] (4) to (3);
	\draw[->, bend right=25] (1) to (4);
	\draw[->, bend right=25] (2) to (1);
	\draw[->, bend right=25] (3) to (2);
	
	% -------- тау (два независимых цикла) --------
	\begin{scope}[shift={(4,0)}]
	\node at (-0.8,0.9) {$, \quad \tau =$};
	
	% Цикл (1 4)
	\node (1a) at (0,0.9) {4};
	\node (4a) at (1.2,0.9) {1};
	\draw[->, bend right=40] (1a) to (4a);
	\draw[->, bend right=40] (4a) to (1a);
	
	% Цикл (2 3)
	\begin{scope}[shift={(2.5,0)}]
	\node (2a) at (0,0.9) {3};
	\node (3a) at (1.2,0.9) {2};
	\draw[->, bend right=40] (2a) to (3a);
	\draw[->, bend right=40] (3a) to (2a);
	\end{scope}
	\end{scope}
	\end{tikzpicture}
	\end{center}
	
	
	
	
	
	
	Заметим, что $\sigma^2 = (1\ 3)(2\ 4)$, $\sigma^4 = (\sigma^2)^2 = e$, $\tau^2 = e$
	
	\noindent \textbf{Определение (Степень перестановки).}
	Степень перестановки $\pi^s$ определяется по индукции
	\[
	\pi^s = \begin{cases}
		\pi(\pi^{s-1}) & s > 0 \\
		e & s = 0 \\
		\pi^{-1}((\pi^{-1})^{(-s-1)}) & s < 0
	\end{cases}
	\]
	Ясно, что $\pi^s \pi^t = \pi^{s+t} = \pi^t \pi^s$, $s,t \in \mathbb{Z}$.
	
	\subsection{Порядок перестановки, орбиты}
	
	\noindent \textbf{Определение (Порядок перестановки, эквивалентность, орбита).}
	Так как $|\Omega| < \infty$, то для каждой перестановки $\pi \in S_n$ найдётся однозначно определённое натуральное число $q = q(\pi)$ такое, что все различные степени содержатся во множестве $\langle \pi \rangle = \{e, \pi, \dots, \pi^{q-1}\}$ и $\pi^q = e$. Это число $q$ называется \textit{порядком} перестановки $\pi$
	
	Две точки (два элемента) $i, j \in \Omega$ назовём \textit{$\pi$-эквивалентными}, если $j = \pi^s(i)$ для некоторого $s \in \mathbb{Z}$
	
	Так как
	\begin{itemize}
		\item $i = \pi^0(i)$ (рефлексивность)
		\item $j = \pi^s(i) \Rightarrow i = \pi^{-s}(j)$ (симметричность)
		\item $j = \pi^s(i), k = \pi^t(j) \Rightarrow k = \pi^{s+t}(i)$ (транзитивность)
	\end{itemize}
	то мы имеем дело с рефлексивным, симметричным и транзитивным отношением на $\Omega$
	
	В соответствии с общим свойством эквивалентности получаем разбиение
	\[
	\Omega = \Omega_1 \cup \cdots \cup \Omega_p \tag{1}
	\]
	множества $\Omega$ на попарно непересекающиеся классы $\Omega_1, \Omega_2, \dots, \Omega_p$, которые принято называть \textit{$\pi$-орбитами}. Каждая точка $i \in \Omega$ принадлежит в точности одной орбите, и если $i \in \Omega_k$, то $\Omega_k$ состоит из образов точки $i$ при действии степеней элемента $\pi$:
	\[
	i, \pi(i), \pi^2(i), \dots, \pi^{l_k-1}(i)
	\]
	здесь $l_k = |\Omega_k|$ — длина $\pi$-орбиты $\Omega_k$.
	
	Ясно, что $l_k \leq q = \operatorname{Card} \pi$, $\pi^{l_k}(i) = i$, причём $l_k$ — наименьшее число, обладающее этим свойством
	
	\begin{tcolorbox}[colback=green!5, colframe=green!50!black, title=Пояснение]
		Орбита точки $i$ — это множество всех точек, в которые можно попасть из $i$, многократно применяя перестановку $\pi$ (или обратную к ней). Орбиты образуют разбиение множества $\Omega$, и на каждой орбите перестановка действует как цикл
	\end{tcolorbox}
	
	\subsection{Разложение перестановки в произведение циклов}
	
	Положив
	\[
	\pi_k = (\pi(i) \ \pi^2(i) \ \dots \ \pi^{l_k-1}(i)) = \begin{pmatrix}
		i & \pi(i) & \dots & \pi^{l_k-2}(i) \\
		\pi(i) & \pi^2(i) & \dots & \pi^{l_k-1}(i)
	\end{pmatrix}
	\]
	мы придём как раз к перестановке, называемой \textit{циклом длины $l_k$}
	
	Цикл $\pi_k$ оставляет на месте все точки из множества $\Omega \setminus \Omega_k$, а $\pi(j) = \pi_k(j)$ для любой точки $j \in \Omega_k$. Это даёт нам основание называть $\pi_s$, $\pi_t$, $s \neq t$ \textit{независимыми}, или \textit{непересекающимися}, циклами. Так как $\pi^{l_k}(i) = i$ для $i \in \Omega_k$, то $\pi_k^{l_k} = e$
	
	Итак, с разбиением (1) ассоциируется разложение перестановки $\pi$ в произведение
	\[
	\pi = \pi_1 \pi_2 \dots \pi_p \tag{2}
	\]
	где все циклы перестановочны. Если цикл $\pi_k = (i)$ имеет длину 1, то он действует как единичная перестановка. Естественно такие циклы в произведении (2) опускать
	
	\begin{tcolorbox}[colback=gray!5, colframe=gray!50, title=Пример]
		Перестановку
		\[
		\pi = \begin{pmatrix}
			1 & 2 & 3 & 4 & 5 & 6 & 7 & 8 \\
			2 & 3 & 4 & 5 & 1 & 7 & 6 & 8
		\end{pmatrix} \in S_8
		\]
		мы запишем в виде
		\[
		\pi = (1\ 2\ 3\ 4\ 5)(6\ 7)(8) = (1\ 2\ 3\ 4\ 5)(6\ 7)
		\]
	\end{tcolorbox}
	
	\subsection{Теорема о единственности разложения}
	
	\noindent \textbf{Теорема (О разложении перестановки в произведение независимых циклов).}
	Каждая перестановка $\pi \neq e$ в $S_n$ является произведением независимых циклов длины $\geq 2$. Это разложение в произведение определено однозначно с точностью до порядка следования циклов
	
	\medskip
	\noindent \textbf{Доказательство.}
	Итак,
	\[
	\pi = \pi_1 \pi_2 \dots \pi_m, \quad l_k > 1, \ 1 \leq k \leq m \tag{3}
	\]
	Пусть наряду с разложением (3) мы имеем ещё одно разложение $\pi = \alpha_1 \alpha_2 \dots \alpha_r$ в произведение независимых циклов, и пусть $i$ — символ, не остающийся на месте при действии $\pi$.
	
	Тогда $\pi_s(i) \neq i$, $\alpha_t(i) \neq i$ для одного (и только одного) из циклов $\pi_1, \pi_2, \dots, \pi_m$ и одного из $\alpha_1, \alpha_2, \dots, \alpha_r$
	
	Имеем
	\[
	\pi_s(i) = \pi(i) = \alpha_t(i)
	\]
	
	Если мы уже знаем, что
	\[
	\pi_s^k(i) = \pi^k(i) = \alpha_t^k(i) \tag{4}
	\]
	то (индукционный переход), применяя к этим неравенствам перестановку $\pi$ и используя перестановочность $\pi$ с $\pi_s^k$ и с $\alpha_t^k$, получаем
	\[
	\pi \pi_s^k(i) = \pi^{k+1}(i) = \pi \alpha_t^k(i)
	\]
	откуда $\pi_s^k \pi(i) = \pi^{k+1}(i) = \alpha_t^k \pi(i)$ и, наконец,
	\[
	\pi_s^{k+1}(i) = \pi^{k+1}(i) = \alpha_t^{k+1}(i)
	\]
	
	Значит, равенства (4) справедливы при любом $k = 0, 1, 2, \dots$. Но цикл однозначно определяется действием его степеней на любой символ, который не остаётся на месте. Следовательно, $\pi_s = \alpha_t$. Дальше очевидная индукция по $m$ или $r$
	
	\begin{tcolorbox}[colback=red!5, colframe=red!70!black, title=Итог теоремы]
		Таким образом, разложение перестановки в произведение независимых циклов единственно с точностью до перестановки циклов. Это фундаментальный факт, позволяющий классифицировать перестановки по их цикловой структуре
	\end{tcolorbox}
	
	\subsection{Транспозиции}
	
	\noindent \textbf{Определение (Транспозиция).}
	Цикл длины 2 называется \textit{транспозицией}.
	
	\begin{tcolorbox}[colback=blue!5, colframe=blue!40!black, title=Пример]
		$(1\ 5)$, $(2\ 7)$, $(3\ 8)$ — транспозиции. Транспозиция меняет местами два элемента, оставляя все остальные на месте
	\end{tcolorbox}
	
	\subsection{Разложение перестановки в произведение транспозиций}
	
	\noindent \textbf{Следствие.}
	Каждая перестановка является произведением транспозиций
	
	В самом деле, в силу теоремы достаточно записать в виде произведения транспозиций каждый из циклов. Это можно сделать, например, так:
	\[
	(1\ 2\ \dots\ l-1\ l) = (1\ l)(1\ l-1)\dots(1\ 3)(1\ 2)
	\]
	
	
	
	
	
	
	\begin{tcolorbox}[colback=gray!5, colframe=gray!50, title=Пример]
	Для цикла $(1\ 2\ 3\ 4\ 5)$:
	\[
	(1\ 2\ 3\ 4\ 5) = \begin{pmatrix}
	1 & 2 & 3 & 4 & 5 \\
	2 & 3 & 4 & 5 & 1
	\end{pmatrix} =
	\]
	
	\[
	\begin{array}{ccccc}
	1 & 2 & 3 & 4 & 5 \\
	\downarrow & \downarrow & \downarrow & \downarrow & \downarrow \\
	2 & 1 & 3 & 4 & 5 \\
	\downarrow & \downarrow & \downarrow & \downarrow & \downarrow \\
	2 & 3 & 1 & 4 & 5 \\
	\downarrow & \downarrow & \downarrow & \downarrow & \downarrow \\
	2 & 3 & 4 & 1 & 5 \\
	\downarrow & \downarrow & \downarrow & \downarrow & \downarrow \\
	2 & 3 & 4 & 5 & 1 \\
	\end{array} = (1\ 5)(1\ 4)(1\ 3)(1\ 2)
	\]
	
	
	\textbf{Проверка:} Покажем, как действует это произведение (напоминаем: транспозиции применяются справа налево).
	\begin{align*}
	\text{Исходное множество:} &\quad 1, 2, 3, 4, 5 \\
	\text{После } (1\ 2): &\quad 2, 1, 3, 4, 5 \\
	\text{После } (1\ 3): &\quad 2, 3, 1, 4, 5 \\
	\text{После } (1\ 4): &\quad 2, 3, 4, 1, 5 \\
	\text{После } (1\ 5): &\quad 2, 3, 4, 5, 1
	\end{align*}
	В результате элемент 1 перешёл в конец, а все остальные сдвинулись влево, что соответствует действию цикла $(1\ 2\ 3\ 4\ 5)$
	
	
	\end{tcolorbox}
	
	
	
	
	
	
	\myremark
	Разложение в произведение транспозиций не единственно. Например, $(1\ 2\ 3) = (1\ 3)(1\ 2) = (2\ 3)(1\ 3) = (1\ 2)(2\ 3)$. Однако чётность числа транспозиций в разложении (чётность перестановки) является инвариантом — это фундаментальный факт, лежащий в основе определения знака перестановки
	
	\begin{tcolorbox}[colback=purple!5, colframe=purple!50!black, title=Заметочка: Чётность перестановки]
		Несмотря на неединственность разложения в произведение транспозиций, количество транспозиций в любом разложении данной перестановки имеет одну и ту же чётность. Это позволяет определить \textit{знак перестановки}:
		\[
		\operatorname{sgn}(\pi) = (-1)^{\text{число транспозиций}}
		\]
		Перестановки с $\operatorname{sgn}(\pi) = 1$ называются чётными, с $\operatorname{sgn}(\pi) = -1$ — нечётными. Чётные перестановки образуют подгруппу $A_n$ — знакопеременную группу
	\end{tcolorbox}
	
	\begin{tcolorbox}[colback=yellow!5, colframe=orange!60!black, title=Важное замечание]
		Из теоремы о разложении в произведение независимых циклов следует, что порядок перестановки равен наименьшему общему кратному длин циклов в её разложении. Например, для перестановки $(1\ 2\ 3)(4\ 5)$ порядок равен $\operatorname{НОК}(3, 2) = 6$
	\end{tcolorbox}
	
	
	
	% ==================== РАЗДЕЛ 7: ПОНЯТИЕ ГРУППЫ ====================
	\newpage
	\section{Понятие группы. Пример симметричной и знакопеременной группы. Их порядки (7 вопрос)}
	
	\subsection{Определение группы}
	
	\noindent \textbf{Определение (Группа, абелева группа).}
	\textit{Группой} называется структура, заданная на некотором множестве $G$ одним всюду определённым внутренним действием $*$, удовлетворяющим следующим аксиомам:
	\begin{enumerate}
		\item $*$ ассоциативно,
		\item В $G$ существует нейтральный относительно $*$ элемент,
		\item Каждый элемент из $G$ обратим относительно $*$
	\end{enumerate}
	
	Если действие в группе коммутативно, то группа называется \textit{коммутативной}, или \textit{абелевой}
	
	\begin{tcolorbox}[colback=blue!5, colframe=blue!40!black, title=Примеры групп]
		\begin{itemize}
			\item $(\mathbb{Z}, +)$ — аддитивная группа целых чисел (абелева)
			\item $(\mathbb{R}\setminus\{0\}, \cdot)$ — мультипликативная группа ненулевых действительных чисел (абелева)
			\item $S_n$ — симметрическая группа перестановок (неабелева при $n \geq 3$)
			\item Множество всех обратимых матриц $n \times n$ с операцией умножения (неабелева при $n \geq 2$)
		\end{itemize}
	\end{tcolorbox}
	
	\subsection{Порядок группы}
	
	\noindent \textbf{Определение (Порядок группы, единица / ноль).}
	Наличие в $G$ действия означает, в частности, что $G \neq \varnothing$. Мощность $|G|$ множества $G$ называют \textit{порядком} группы $G$. Нейтральный элемент называется, в зависимости от записи, \textit{единицей} (при мультипликативной записи) или \textit{нулём} (при аддитивной записи) группы $G$
	
	\begin{tcolorbox}[colback=yellow!5, colframe=orange!60!black, title=Замечание о записи]
		При мультипликативной записи нейтральный элемент обозначают $1$ или $e$, обратный к $g$ обозначают $g^{-1}$, а групповую операцию называют умножением.
		При аддитивной записи нейтральный элемент обозначают $0$, обратный к $g$ обозначают $-g$, а операцию называют сложением
	\end{tcolorbox}
	
	\subsection{Симметрическая группа}
	
	\noindent \textbf{Определение (Симметрическая группа степени $n$).}
	Пусть $\Omega = \{1, 2, \dots, n\}$ — конечное множество элементов. Множество всех взаимно однозначных преобразований $\Omega \to \Omega$ образуют группу перестановок $n$ элементов, которую называют \textit{симметрической группой степени $n$} и обозначают $S_n$
	
	Ясно, что порядок этой группы $|S_n| = n!$.
	
	\begin{tcolorbox}[colback=gray!5, colframe=gray!50, title=Примеры]
		\begin{itemize}
			\item $S_1 = \{e\}$, $|S_1| = 1$
			\item $S_2 = \{e, (1\ 2)\}$, $|S_2| = 2$ (коммутативна)
			\item $S_3 = \{e, (1\ 2), (1\ 3), (2\ 3), (1\ 2\ 3), (1\ 3\ 2)\}$, $|S_3| = 6$ (некоммутативна)
		\end{itemize}
	\end{tcolorbox}
	
	\myremark
	Симметрическая группа $S_n$ — важнейший пример конечной группы. Она играет ключевую роль в теории групп, комбинаторике и приложениях
	
	\subsection{Знак перестановки}
	
	\noindent \textbf{Теорема (Определение знака перестановки).}
	Пусть $\pi$ — перестановка из $S_n$,
	\[
	\pi = \tau_1 \tau_2 \dots \tau_k \tag{1}
	\]
	— произвольное разложение $\pi$ в произведение транспозиций
	
	Тогда число $\varepsilon_\pi = (-1)^k$, называемое \textit{знаком} $\pi$ (иначе: сигнатурой или чётностью), полностью определяется перестановкой $\pi$ и не зависит от способа разложения (1), т.е. чётность целого числа $k$ для данной перестановки $\pi$ всегда одна и та же. Кроме того,
	\[
	\varepsilon_{\alpha \beta} = \varepsilon_\alpha \varepsilon_\beta
	\]
	для всех $\alpha, \beta \in S_n$
	
	\medskip
	\noindent \textbf{Доказательство независимости знака от способа разложения.}
	
	Предположим, что наряду с (1) мы имеем также разложение
	\[
	\pi = \tau'_1 \tau'_2 \dots \tau'_{k'},
	\]
	причём чётности $k$ и $k'$ различны. Тогда целое число $k + k'$ нечётно. Так как $(\tau'_s)^2 = e$, то, следовательно, умножая справа обе части равенства
	\[
	\tau_1 \tau_2 \dots \tau_k = \tau'_1 \tau'_2 \dots \tau'_{k'}
	\]
	на $\tau'_{k'}, \dots, \tau'_2, \tau'_1$, получим
	\[
	\tau_1 \tau_2 \dots \tau_k \tau'_{k'} \dots \tau'_2 \tau'_1 = e.
	\]
	
	Мы свели нашу задачу к следующей. Пусть
	\[
	e = \sigma_1 \sigma_2 \dots \sigma_{m-1} \sigma_m, \quad m > 0 \tag{2}
	\]
	— запись единичной перестановки в виде произведения $m > 0$ транспозиций. Нужно показать, что $m$ — обязательно чётное число
	
	Нам нужно обосновать переход от записи (2) к записи $e$ в виде $m-2$ транспозиций
	
	Пусть $s$, $1 \leq s \leq n$ — любое фиксированное натуральное число, входящее в одну из транспозиций. Для определённости считаем, что
	\[
	e = \sigma_1 \dots \sigma_{p-1} \sigma_p \sigma_{p+1} \dots \sigma_m
	\]
	где $\sigma_p = (s\ t)$, а $\sigma_{p+1}, \dots, \sigma_m$ не содержат $s$
	
	Для $\sigma_{p-1}$ имеются 4 возможности:
	
	\begin{enumerate}
		\item[(а)] $\sigma_{p-1} = (s\ t)$; тогда отрезок $(s\ t)(s\ t)$ из записи $e$ удаляется и мы приходим к $m-2$ транспозициям;
		
		
		\item[б)] $\sigma_{p-1} = (s\ r)$, $r \neq s, t$; здесь
		\[
		\sigma_{p-1}\sigma_p = (s\ r)(s\ t) = (s\ t)(r\ t)
		\]
		
		\[
		\begin{array}{ccc}
		r & s & t \\
		\downarrow & \downarrow & \downarrow \\
		r & t & s \\
		\downarrow & \downarrow & \downarrow \\
		s & t & r
		\end{array}
		\quad \quad
		\begin{array}{ccc}
		r & s & t \\
		\downarrow & \downarrow & \downarrow \\
		t & s & r \\
		\downarrow & \downarrow & \downarrow \\
		s & t & r
		\end{array}
		\]
		
		и мы сдвинули вхождение $s$ на одну позицию влево, не изменив $m$;
			

		\item[(в)] $\sigma_{p-1} = (t\ r)$, $r \neq s, t$; здесь
		\[
		\sigma_{p-1}\sigma_p = (t\ r)(s\ t) = (s\ r)(t\ r)
		\]
		\[
		\begin{array}{ccc}
		r & s & t \\
		\downarrow & \downarrow & \downarrow \\
		r & t & s \\
		\downarrow & \downarrow & \downarrow \\
		t & r & s
		\end{array}
		\quad \quad
		\begin{array}{ccc}
		r & s & t \\
		\downarrow & \downarrow & \downarrow \\
		t & s & r \\
		\downarrow & \downarrow & \downarrow \\
		t & r & s
		\end{array}
		\]
		и снова, как и в случае (б), произошёл сдвиг $s$ влево без изменения $m$;
		
		\item[(г)] $\sigma_{p-1}\sigma_p = (q\ r)(s\ t) = (s\ t)(q\ r)$ (не содержат общих символов)
	\end{enumerate}
	
	В случае (а) наша цель достигнута. В случаях (б)-(г) повторяем процесс, сдвигая $s$ на одну позицию влево. В конечном счёте мы придём либо к случаю (а), либо к экстремальному случаю, когда $e = \sigma'_1 \sigma'_2 \dots \sigma'_m$, причём $\sigma'_1 = (s\ t')$ и $s$ не имеет вхождений в $\sigma'_2 \dots \sigma'_m$. Значит, $\forall k > 1 \ \sigma'_k(s) = s$ и
	\[
	s = e(s) = \sigma'_1(s) = t' \neq s,
	\]
	что является противоречием. Полученное противоречие доказывает утверждение об инвариантности $\varepsilon_\pi$, потому что, продолжая этот спуск от $m$ к $m-2$ множителей, мы пришли бы при нечётном $m$ к одной транспозиции $\tau$. Но, очевидно, $e \neq \tau$
	
	\medskip
	\noindent \textbf{Доказательство мультипликативности.}
	
	Если $\alpha = \tau_1 \tau_2 \dots \tau_k$, $\beta = \tau_{k+1} \dots \tau_{k+l}$, то $\alpha \beta = \tau_1 \tau_2 \dots \tau_k \tau_{k+1} \dots \tau_{k+l}$ и
	\[
	\varepsilon_\alpha = (-1)^k,\ \varepsilon_\beta = (-1)^l,\ \varepsilon_{\alpha \beta} = (-1)^k (-1)^l = (-1)^{k+l} = \varepsilon_\alpha \varepsilon_\beta.
	\]
	
	\begin{tcolorbox}[colback=green!5, colframe=green!50!black, title=Пояснение к доказательству]
		Доказательство независимости знака от разложения — это классический комбинаторный аргумент, показывающий, что чётность числа транспозиций в разложении перестановки является инвариантом. Основная идея: если бы существовало два разложения с разной чётностью, то единичную перестановку можно было бы представить как произведение нечётного числа транспозиций, что приводит к противоречию
	\end{tcolorbox}
	
	\subsection{Чётные и нечётные перестановки}
	
	\noindent \textbf{Определение (Чётная, нечётная перестановка).}
	Перестановка $\pi \in S_n$ называется \textit{чётной}, если $\varepsilon_\pi = 1$, и \textit{нечётной}, если $\varepsilon_\pi = -1$.
	
	\begin{tcolorbox}[colback=gray!5, colframe=gray!50, title=Примеры]
		\begin{itemize}
			\item Тождественная перестановка $e$ — чётная (можно представить как $0$ транспозиций, $(-1)^0 = 1$)
			\item Транспозиция $(i\ j)$ — нечётная (представляется одной транспозицией, $(-1)^1 = -1$)
			\item Цикл длины $3$, например $(1\ 2\ 3) = (1\ 3)(1\ 2)$ — произведение двух транспозиций, значит чётный
			\item Цикл длины $4$ — нечётный, так как раскладывается в три транспозиции
		\end{itemize}
	\end{tcolorbox}
	
	\subsection{Знак перестановки через цикловое разложение}
	
	\noindent \textbf{Следствие.}
	Пусть перестановка $\pi \in S_n$ разложена в произведение независимых циклов длин $l_1, l_2, \dots, l_m$. Тогда
	\[
	\varepsilon_\pi = (-1)^{\sum_{k=1}^m (l_k - 1)}
	\]
	
	Действительно, $\varepsilon_\pi = \varepsilon_{\pi_1 \dots \pi_m} = \varepsilon_{\pi_1} \dots \varepsilon_{\pi_m}$.
	
	Кроме того,
	\[
	\varepsilon_{\pi_k} = (-1)^{l_k - 1}
	\]
	поскольку $\pi_k$ записывается в виде произведения $l_k - 1$ транспозиций:
	\[
	(1\ 2\ \dots\ l_k-1\ l_k) = (1\ l_k)(1\ l_k-1)\dots(1\ 3)(1\ 2)
	\]
	
	Окончательно,
	\[
	\varepsilon_\pi = (-1)^{l_1-1} \dots (-1)^{l_m-1} = (-1)^{\sum_{k=1}^m (l_k - 1)}
	\]
	
	\begin{tcolorbox}[colback=blue!5, colframe=blue!40!black, title=Пример]
		Для перестановки $\pi = (1\ 2\ 3)(4\ 5) \in S_5$:
		\begin{itemize}
			\item Длины циклов: $l_1 = 3$, $l_2 = 2$
			\item $\sum (l_k - 1) = (3-1) + (2-1) = 2 + 1 = 3$.
			\item $\varepsilon_\pi = (-1)^3 = -1$ — перестановка нечётная
		\end{itemize}
		Проверим: $(1\ 2\ 3)$ — чётный ($2$ транспозиции), $(4\ 5)$ — нечётный ($1$ транспозиция). Произведение чётного на нечётное даёт нечётное
	\end{tcolorbox}
	
	\subsection{Количество чётных и нечётных перестановок}
	
	\noindent \textbf{Теорема (Количество чётных / нечётных перестановок).}
	Запишем $S_n$ в виде объединения $S_n = A_n \cup \overline{A_n}$, где $A_n = \{\pi \in S_n \mid \varepsilon_\pi = 1\}$ — множество всех чётных перестановок, $\overline{A_n} = S_n \setminus A_n$ — множество нечётных перестановок. Пусть $\tau = (i\ j)$ — любая транспозиция. Отображение $S_n$ в себя, определённое правилом $L_\tau \colon \pi \to \tau \pi$, очевидно, биективно. ($L_\tau^2$ — единичное отображение, $L_\tau^{-1} = L_\tau$).
	
	Заметим, что
	\[
	L_\tau(A_n) = \overline{A_n}, \quad L_\tau(\overline{A_n}) = A_n
	\]
	
	Значит, число чётных перестановок совпадает с числом нечётных перестановок, откуда
	\[
	|A_n| = \frac{1}{2} |S_n| = \frac{n!}{2}
	\]
	
	\begin{tcolorbox}[colback=green!5, colframe=green!50!black, title=Пояснение]
		Умножение на фиксированную транспозицию $\tau$ переводит чётные перестановки в нечётные и наоборот, причём это отображение биективно. Следовательно, чётных и нечётных перестановок поровну. Это простое, но элегантное рассуждение
	\end{tcolorbox}
	
	\subsection{Знакопеременная группа}
	
	\noindent \textbf{Определение (Знакопеременная группа).}
	Поскольку произведение двух чётных перестановок является чётной перестановкой, единичное преобразование является чётной перестановкой, и обратное преобразование для любой чётной перестановки также является чётной перестановкой (так как раскладывается в произведение транспозиций, или проще: $1 = \varepsilon(e) = \varepsilon(\varsigma \varsigma^{-1}) \Rightarrow \varepsilon(\varsigma) = \varepsilon(\varsigma^{-1})$, так как $\varepsilon = \pm 1$), множество чётных перестановок тоже образует группу
	
	Она называется \textit{знакопеременной группой} степени $n$ и обозначается обычно $A_n$. Порядок знакопеременной группы степени $n$ равен $\frac{1}{2} |S_n| = \frac{n!}{2}$
	
	\begin{tcolorbox}[colback=gray!5, colframe=gray!50, title=Примеры]
		\begin{itemize}
			\item $A_2 = \{e\}$, $|A_2| = 1$
			\item $A_3 = \{e, (1\ 2\ 3), (1\ 3\ 2)\}$, $|A_3| = 3$
			\item $A_4$ содержит $12$ элементов (половина от $24$)
		\end{itemize}
	\end{tcolorbox}
	
	\myremark
	Знакопеременная группа $A_n$ является нормальной подгруппой в $S_n$ индекса $2$. При $n \geq 5$ группа $A_n$ простая (не имеет нетривиальных нормальных подгрупп), что играет ключевую роль в теории разрешимости уравнений в радикалах
	
	\begin{tcolorbox}[colback=purple!5, colframe=purple!50!black, title=Заметочка: Связь с предыдущими темами]
		\begin{itemize}
			\item $S_n$ — пример конечной группы порядка $n!$, которую мы уже изучали как группу перестановок
			\item Знак перестановки — это гомоморфизм $\varepsilon \colon S_n \to \{\pm 1\}$.
			\item Ядро этого гомоморфизма — в точности $A_n$ (чётные перестановки)
			\item $|S_n : A_n| = 2$, поэтому $A_n$ — нормальная подгруппа индекса $2$
		\end{itemize}
	\end{tcolorbox}
	
	\begin{tcolorbox}[colback=yellow!5, colframe=orange!60!black, title=Важное замечание]
		Знакопеременная группа $A_n$ при $n \geq 5$ играет фундаментальную роль в теории Галуа: именно её простота является причиной того, что общее уравнение степени $n \geq 5$ не разрешимо в радикалах
	\end{tcolorbox}
	
	
	
	
	
	
	
	% ==================== РАЗДЕЛ 8: ПОНЯТИЕ КОЛЬЦА ====================
	\newpage
	\section{Понятие кольца. Примеры колец. Область целостности. Кольцо классов вычетов и область целостности (8 вопрос)}
	
	\subsection{Определение кольца}
	
	\noindent \textbf{Определение (Кольцо, коммутативное кольцо, кольцо с единицей, тело, поле, кольцо с нулевым умножением).}
	\textit{Кольцом} называется структура $(A, +, \cdot)$, заданная на некотором множестве $A$ парой всюду определённых внутренних действий — сложением и умножением, — удовлетворяющих следующим аксиомам:
	\begin{enumerate}
		\item Сложение ассоциативно;
		\item Сложение коммутативно;
		\item Существует $0$ — нейтральный элемент по сложению: $\forall a \in A \ a + 0 = 0 + a = a$;
		\item Каждый элемент из $A$ обратим по сложению: $\forall a \in A \ \exists (-a) \colon a + (-a) = 0$;
		\item Умножение ассоциативно;
		\item Умножение дистрибутивно относительно сложения: $a(b+c) = ab + ac$, $(a+b)c = ac + bc$
	\end{enumerate}
	
	Кроме аксиом 1-6 действия в кольце могут удовлетворять и другим требованиям:
	\begin{enumerate}
		\item[7.] Умножение коммутативно;
		\item[8.] Существует $1$ — нейтральный элемент по умножению: $\forall a \in A \ a \neq 0 \Rightarrow a \cdot 1 = 1 \cdot a = a$, причём $1 \neq 0$;
		\item[9.] Каждый отличный от нуля элемент кольца обратим
	\end{enumerate}
	
	Кольцо, дополнительно удовлетворяющее:
	\begin{itemize}
		\item аксиоме 7 называется \textit{коммутативным кольцом};
		\item аксиоме 8 — \textit{кольцом с единицей};
		\item аксиомам 7 и 8 — \textit{коммутативным кольцом с единицей};
		\item аксиомам 8 и 9 — \textit{телом};
		\item аксиомам 7, 8 и 9 — \textit{полем}
	\end{itemize}
	
	\begin{tcolorbox}[colback=blue!5, colframe=blue!40!black, title=Базовые примеры колец]
		\begin{itemize}
			\item $\mathbb{Z}$ — коммутативное кольцо с единицей (но не поле, так как обратимые элементы только $\pm 1$)
			\item $\mathbb{Q}$, $\mathbb{R}$, $\mathbb{C}$ — поля (коммутативные кольца с единицей, где каждый ненулевой элемент обратим)
			\item Множество $2\mathbb{Z}$ всех чётных чисел с обычными действиями есть кольцо коммутативное, но без единицы
		\end{itemize}
	\end{tcolorbox}
	
	\subsection{Кольцо с нулевым умножением}
	
	\noindent \textbf{Определение (Кольцо с нулевым умножением).}
	Пусть $A$ — абелева группа, а умножение задано так: $\forall a, b \in A \ ab = 0$. Получается коммутативное кольцо без единицы, так называемое \textit{кольцо с нулевым умножением}
	
	\begin{tcolorbox}[colback=gray!5, colframe=gray!50, title=Пример кольца с нулевым умножением]
		Рассмотрим $\mathbb{Z}_2 \times \mathbb{Z}_2$ с покоординатным сложением и умножением, заданным как $(a,b)(c,d) = (0,0)$. Это кольцо с нулевым умножением
	\end{tcolorbox}
	
	\myremark
	В кольце с нулевым умножением произведение любых двух элементов равно нулю. Это экстремальный пример, показывающий, что умножение может быть "тривиальным"
	
	\subsection{Делители нуля}
	
	\noindent \textbf{Определение (Делитель нуля).}
	Элемент $a \in R$ называется \textit{делителем нуля}, если существует такое $x \in R$, $x \neq 0$, что $ax = 0$
	
	\begin{tcolorbox}[colback=blue!5, colframe=blue!40!black, title=Примеры делителей нуля]
		\begin{itemize}
			\item В кольце $\mathbb{Z}_6$ элемент $\bar{2}$ — делитель нуля, так как $\bar{2} \cdot \bar{3} = \bar{6} = \bar{0}$, но $\bar{2} \neq \bar{0}$, $\bar{3} \neq \bar{0}$
			\item В кольце матриц $M_2(\mathbb{R})$ ненулевая матрица $\begin{pmatrix} 1 & 0 \\ 0 & 0 \end{pmatrix}$ — делитель нуля, так как $\begin{pmatrix} 1 & 0 \\ 0 & 0 \end{pmatrix} \cdot \begin{pmatrix} 0 & 0 \\ 0 & 1 \end{pmatrix} = \begin{pmatrix} 0 & 0 \\ 0 & 0 \end{pmatrix}$
			\item В поле $\mathbb{Q}$ нет делителей нуля
		\end{itemize}
	\end{tcolorbox}
	
	\subsection{Кольцо без делителей нуля}
	
	\noindent \textbf{Определение (Кольцо без делителей нуля).}
	Кольцо $R$ называется \textit{кольцом без делителей нуля}, если $0$ — единственный делитель нуля в $R$
	
	\begin{tcolorbox}[colback=yellow!5, colframe=orange!60!black, title=Замечание]
		В кольце без делителей нуля из $ab = 0$ следует $a = 0$ или $b = 0$. Это свойство обобщает привычное свойство целых чисел
	\end{tcolorbox}
	
	\subsection{Область целостности}
	
	\noindent \textbf{Определение (Область целостности).}
	Коммутативное кольцо с единицей без делителей нуля называется \textit{областью целостности}
	
	\begin{tcolorbox}[colback=gray!5, colframe=gray!50, title=Примеры областей целостности]
		\begin{itemize}
			\item $\mathbb{Z}$ — область целостности (но не поле)
			\item $\mathbb{Z}[x]$ — кольцо многочленов с целыми коэффициентами
			\item Любое поле (например, $\mathbb{Q}$, $\mathbb{R}$, $\mathbb{C}$) является областью целостности
			\item $\mathbb{Z}_p$ при простом $p$ — поле, следовательно, область целостности
		\end{itemize}
	\end{tcolorbox}
	
	\begin{tcolorbox}[colback=yellow!5, colframe=orange!60!black, title=Контрпримеры]
		\begin{itemize}
			\item $\mathbb{Z}_6$ не является областью целостности, так как содержит делители нуля
			\item $2\mathbb{Z}$ не является областью целостности, так как не имеет единицы
			\item Кольцо матриц $M_n(\mathbb{R})$ не является областью целостности (оно некоммутативно и содержит делители нуля)
		\end{itemize}
	\end{tcolorbox}
	
	\subsection{Кольцо классов вычетов}
	
	\noindent \textbf{Предложение (О взаимной простоте, обратимости и делителе нуля в $\mathbb{Z}_m$).}
	Пусть $a, m \in \mathbb{Z}$.
	\begin{enumerate}
		\item Если $a$, $m$ взаимно просты, то $\bar{a}$ обратим в $\mathbb{Z}_m$
		\item Если $a$, $m$ не взаимно просты, то $\bar{a}$ — делитель нуля в $\mathbb{Z}_m$
	\end{enumerate}
	
	\medskip
	\noindent \textbf{Доказательство.}
	
	\textbf{1.} По теореме о линейном представлении НОД $1 = au + m\upsilon$ при некоторых $u, \upsilon \in \mathbb{Z}$. Отсюда в $\mathbb{Z}_m$:
	\[
	\bar{1} = \bar{a}\bar{u} + \bar{m}\bar{\upsilon} = \bar{a}\bar{u},
	\]
	т. е. $\bar{u} = \bar{a}^{-1}$, значит $\bar{a}$ обратим
	
	
	
		
	\textbf{2.} Пусть $d =$ НОД$(a, m) \neq 1$. Запишем $a = a_1d$, $m = m_1d$, где $a_1, m_1 \in \mathbb{Z}$. Так как $d > 1$, то $m_1 = m/d$ не кратно $m$, следовательно, его класс вычетов $\overline{m}_1 \neq \overline{0}$ в кольце $\mathbb{Z}_m$
	
	Однако:
	\[
	\overline{a} \cdot \overline{m}_1 = \overline{am_1} = \overline{a_1 d m_1} = \overline{a_1 m} = \overline{0},
	\]
	поскольку $a_1 m$ делится на $m$. Таким образом, $\overline{a}$ является делителем нуля в кольце $\mathbb{Z}_m$
	
	
	
	
	
	
	\begin{tcolorbox}[colback=green!5, colframe=green!50!black, title=Пояснение]
		Это предложение даёт полную классификацию элементов в кольце классов вычетов $\mathbb{Z}_m$:
		\begin{itemize}
			\item Если НОД$(a, m) = 1$, то $\bar{a}$ обратим (имеет мультипликативный обратный)
			\item Если НОД$(a, m) = d > 1$, то $\bar{a}$ — делитель нуля
		\end{itemize}
	\end{tcolorbox}
	
	\subsection{Следствие о структуре $\mathbb{Z}_m$}
	
	\noindent \textbf{Следствие.}
	\begin{enumerate}
		\item Если $p \in \mathbb{N}$ простое, то $\mathbb{Z}_p$ — поле
		\item Если $m \in \mathbb{N}$ не простое и не равно 1, то $\mathbb{Z}_m$ содержит ненулевые делители нуля
	\end{enumerate}
	
	\medskip
	\noindent \textbf{Доказательство.}
	
	\textbf{1.} Пусть $\bar{a} \in \mathbb{Z}_p$. Поскольку $p$ простое, либо $p \mid a$, либо $a$, $p$ взаимно просты (свойства простых чисел), т. е. либо $\bar{a} = \bar{0}$, либо $\bar{a}$ обратим. Таким образом, каждый ненулевой класс из $\mathbb{Z}_p$ обратим, значит $\mathbb{Z}_p$ — поле
	
	\textbf{2.} Поскольку $m$ не простое и не равно единице, $m$ имеет нетривиальный делитель $a$, причём $(m, a) = a \neq 1$, т. е. $m$, $a$ не взаимно просты. Тогда по предложению $\bar{a}$ — делитель нуля в $\mathbb{Z}_m$. При этом $a \neq m$, т. е. $\bar{a} \neq \bar{0}$
	
	\begin{tcolorbox}[colback=blue!5, colframe=blue!40!black, title=Итоговая классификация $\mathbb{Z}_m$]
		Таким образом, при $m \geq 2$ кольцо $\mathbb{Z}_m$ либо суть поле, либо не есть даже область целостности
		\begin{itemize}
			\item При $m = 0$ кольцо $\mathbb{Z}_m = \mathbb{Z}_0$, которое можно отождествить с $\mathbb{Z}$, есть область целостности, но не поле
			\item При $m = 1$ кольцо $\mathbb{Z}_1 = \{0\}$ не имеет ненулевых делителей нуля, но не есть область целостности, ибо в нём $1 = 0$
		\end{itemize}
		Поля $\mathbb{Z}_p$ (при простых $p$) являют собой примеры конечных полей (их называют \textit{полями Галуа} и обозначают $\mathbb{F}_p$ или $GF(p)$)
	\end{tcolorbox}
	
	\myremark
	Кольца классов вычетов $\mathbb{Z}_m$ дают прекрасные примеры, иллюстрирующие все введённые понятия:
	\begin{itemize}
		\item $\mathbb{Z}_p$ (при простом $p$) — поле (и область целостности)
		\item $\mathbb{Z}_m$ (при составном $m$) — коммутативное кольцо с единицей, но не область целостности (есть делители нуля)
		\item $\mathbb{Z}_0 \cong \mathbb{Z}$ — область целостности, но не поле
		\item $\mathbb{Z}_1$ — тривиальное кольцо (в нём $1 = 0$), которое не является областью целостности по определению
	\end{itemize}
	
	\begin{tcolorbox}[colback=purple!5, colframe=purple!50!black, title=Заметочка: Поля Галуа]
		Конечные поля $\mathbb{F}_p$ (где $p$ — простое) являются простейшими примерами конечных полей. Они играют фундаментальную роль в теории чисел, криптографии и алгебраической геометрии. Обозначение $GF(p)$ происходит от "Galois Field" (поле Галуа) в честь великого математика Эвариста Галуа
	\end{tcolorbox}
	
	\begin{tcolorbox}[colback=yellow!5, colframe=orange!60!black, title=Важное замечание о $\mathbb{Z}_1$]
		Кольцо $\mathbb{Z}_1 = \{0\}$ часто называют \textit{нулевым кольцом}. В нём $1 = 0$, поэтому оно не удовлетворяет требованию $1 \neq 0$ из определения кольца с единицей (хотя формально нейтральный элемент по умножению существует, но он совпадает с нулём). Обычно нулевое кольцо исключают из рассмотрения как тривиальный случай
	\end{tcolorbox}
	
	
	
	
	
	% ==================== РАЗДЕЛ 9: ДЕЛИМОСТЬ В КОЛЬЦЕ ЦЕЛЫХ ЧИСЕЛ ====================
	\newpage
	\section{Свойства делимости в кольце целых чисел. Теорема о делении с остатком (9 вопрос)}
	
	\subsection{Определение делимости}
	
	\noindent \textbf{Определение (Делимость).}
	Пусть $a, b \in \mathbb{Z}$. Говорят, что $a$ \textit{делится} на $b$ ($a \vdots b$), или что $b$ \textit{делит} $a$ ($b \mid a$), если существует такое $x \in \mathbb{Z}$, что $a = bx$.
	
	\begin{tcolorbox}[colback=blue!5, colframe=blue!40!black, title=Примеры]
		\begin{itemize}
			\item $15 \vdots 3$, так как $15 = 3 \cdot 5$
			\item $20 \vdots 4$, $20 \vdots 5$, $20 \vdots 10$
			\item $7 \vdots 1$ и $7 \vdots 7$ (любое число делится на себя и на 1)
			\item $0 \vdots 5$, так как $0 = 5 \cdot 0$
			\item $8$ не делится на $3$, так как нет целого $x$ с $8 = 3x$
		\end{itemize}
	\end{tcolorbox}
	
	\subsection{Свойства делимости}
	
	\noindent \textbf{Свойства делимости.}
	
	\begin{enumerate}
		\item $0$ делится на любое число: $\forall a \in \mathbb{Z} \ 0 \vdots a$
		\item $\forall a \in \mathbb{Z} \ a \vdots 1$ и $a \vdots a$
		\item Делимость транзитивна: $a \vdots b$ и $b \vdots c \Rightarrow a \vdots c$
		\item Если все члены семейства $(a_i)$ делятся на $b$, то и любая линейная комбинация этого семейства делится на $b$
	\end{enumerate}
	
	\begin{tcolorbox}[colback=gray!5, colframe=gray!50, title=Примеры на свойства]
		\begin{itemize}
			\item $12 \vdots 3$ и $3 \vdots 1$ (транзитивность не обязательно выполняется в обратную сторону)
			\item $15 \vdots 5$ и $20 \vdots 5$, тогда $15+20=35 \vdots 5$, и $2\cdot 15 - 3\cdot 20 = 30-60 = -30 \vdots 5$ (линейная комбинация делится на 5)
		\end{itemize}
	\end{tcolorbox}
	
	\myremark
	Свойство 4 особенно важно: если два числа делятся на $b$, то их сумма, разность и любая целочисленная комбинация также делятся на $b$. Это свойство широко используется в теории чисел
	
	\subsection{Теорема о делении целых чисел с остатком}
	
	\noindent \textbf{Теорема (О делении с остатком).}
	$\forall a, b \in \mathbb{Z}$, $b \neq 0$ $\exists! q, r \in \mathbb{Z}$ такие, что
	\[
	a = bq + r \quad \text{и} \quad 0 \leq r < |b|
	\]
	
	\begin{tcolorbox}[colback=red!5, colframe=red!70!black, title=Формулировка]
		Для любых целых $a$ и $b \neq 0$ существуют единственные целые числа $q$ (частное) и $r$ (остаток), такие что $a = bq + r$ и остаток $r$ удовлетворяет неравенству $0 \leq r < |b|$
	\end{tcolorbox}
	
	\subsection{Доказательство существования}
	
	\noindent \textbf{Доказательство существования.}
	Рассмотрим 3 случая
	
	\textbf{Случай 1.} Пусть $a \geq 0$, $b > 0$. Считая $b$ фиксированным (но любым), проведём индукцию по $a$
	
	\textit{База индукции.} $a < b$. Очевидно, что $a = b \cdot 0 + a$, т. е. пара $q = 0$, $r = a$ такая, как требуется в формулировке теоремы
	
	\textit{Индукционный переход.} Пусть $a \geq b$. Тогда $0 \leq a-b < a$, т. е. пара $a-b$, $b$ удовлетворяет условиям теоремы ($b \neq 0$), условиям Случая 1 ($a-b \geq 0$, $b > 0$) и условиям индукционного предположения ($a-b < a$). Поэтому $\exists q', r \colon a-b = bq' + r$, $0 \leq r < |b|$, а тогда $a = b(q' + 1) + r$, и пара $q = q' + 1$, $r$ — такая, как надо в формулировке теоремы
	
	\begin{tcolorbox}[colback=blue!5, colframe=blue!40!black, title=Иллюстрация случая 1]
		Для $a \geq b$ мы уменьшаем задачу, вычитая $b$ из $a$, пока не получим число, меньшее $b$. Например, для $a=17$, $b=5$:
		\begin{align*}
			17 &= 5 \cdot 3 + 2 \quad \text{(непосредственно)}\\
			\text{или по индукции: } & 17-5=12, 12-5=7, 7-5=2 < 5, \text{ тогда } 17 = 5 \cdot (3) + 2
		\end{align*}
	\end{tcolorbox}
	
	\textbf{Случай 2.} Пусть $a < 0$, $b > 0$. Тогда $-a > 0$, и, как доказано в Случае 1,
	\[
	\exists q', r' \colon -a = bq' + r', \quad 0 \leq r' < |b| = b
	\]
	
	Если $r' = 0$, то поскольку тогда $a = (-q')b$, пара $q = -q'$, $r = 0$ такая, как надо. Если $r' \neq 0$, т. е. $r' > 0$, то, поскольку $a = b(-q' - 1) + (b - r')$ и $0 < b - r' < b = |b|$, пара $q = -q' - 1$, $r = b - r'$ — такая, как надо
	
	\begin{tcolorbox}[colback=blue!5, colframe=blue!40!black, title=Пример случая 2]
		$a = -17$, $b = 5$:
		\begin{itemize}
			\item $-a = 17 = 5 \cdot 3 + 2$, т.е. $q' = 3$, $r' = 2$
			\item $r' \neq 0$, поэтому $q = -3 - 1 = -4$, $r = 5 - 2 = 3$
			\item Проверка: $5 \cdot (-4) + 3 = -20 + 3 = -17$, и $0 \leq 3 < 5$
		\end{itemize}
	\end{tcolorbox}
	
	\textbf{Случай 3.} Пусть $b < 0$ (при произвольном $a$). Тогда $-b > 0$, и, как показано в Случае 1 и Случае 2, существует $q', r$ такие, что
	\[
	a = (-b)q' + r = b(-q') + r \quad \text{и} \quad 0 \leq r < -b = |b|
	\]
	т. е. пара $q = -q'$, $r$ такая, как надо
	
	\begin{tcolorbox}[colback=blue!5, colframe=blue!40!black, title=Пример случая 3]
		$a = 17$, $b = -5$:
		\begin{itemize}
			\item $-b = 5$, применяем случай 1: $17 = 5 \cdot 3 + 2$, т.е. $q' = 3$, $r = 2$
			\item Тогда $q = -3$, $r = 2$
			\item Проверка: $(-5) \cdot (-3) + 2 = 15 + 2 = 17$, и $0 \leq 2 < 5 = |-5|$
		\end{itemize}
		
		$a = -17$, $b = -5$:
		\begin{itemize}
			\item $-b = 5$, применяем случай 2: $-17 = 5 \cdot (-4) + 3$? Нет, случай 2 даёт для $-17$ при $b=5$: $q = -4$, $r = 3$
			\item Тогда $q = -(-4) = 4$, $r = 3$
			\item Проверка: $(-5) \cdot 4 + 3 = -20 + 3 = -17$, и $0 \leq 3 < 5 = |-5|$
		\end{itemize}
	\end{tcolorbox}
	
	\subsection{Доказательство единственности}
	
	\noindent \textbf{Доказательство единственности.}
	Пусть $a = bq + r = bq' + r'$, где $0 \leq r \leq r' < |b|$. Тогда $0 \leq r' - r < |b|$ и $b(q - q') = r' - r$, откуда
	\[
	|b| \cdot |q - q'| = |r' - r| = r' - r < |b|,
	\]
	что возможно лишь при $q = q'$, $r = r'$
	
	\begin{tcolorbox}[colback=green!5, colframe=green!50!black, title=Пояснение единственности]
		Если бы $q \neq q'$, то $|q - q'| \geq 1$, и тогда $|b| \cdot |q - q'| \geq |b|$, но мы получили строгое неравенство $|b| \cdot |q - q'| < |b|$. Противоречие. Следовательно, $q = q'$, а тогда и $r = r'$
	\end{tcolorbox}
	
	\myremark
	Теорема о делении с остатком — фундаментальный результат, на котором базируется вся теория делимости в целых числах. Она гарантирует, что мы всегда можем разделить любое целое число на любое ненулевое целое с остатком, причём остаток будет меньше делителя по абсолютной величине
	
	\begin{tcolorbox}[colback=purple!5, colframe=purple!50!black, title=Заметочка: Алгоритм Евклида]
		Теорема о делении с остатком лежит в основе алгоритма Евклида для нахождения наибольшего общего делителя двух чисел. Последовательно деля с остатком, мы получаем цепочку равенств:
		\[
		a = bq_1 + r_1,\ b = r_1q_2 + r_2,\ r_1 = r_2q_3 + r_3,\ \dots
		\]
		Последний ненулевой остаток и будет НОД(a, b)
	\end{tcolorbox}
	
	\begin{tcolorbox}[colback=yellow!5, colframe=orange!60!black, title=Важное замечание]
		В некоторых учебниках остаток определяют в пределах $0 \leq r < |b|$, в других — в пределах $-|b|/2 < r \leq |b|/2$ (так называемое "деление с остатком с симметрией"). Мы используем классическое определение с неотрицательным остатком
	\end{tcolorbox}
	
	
	
	
	
	
	
	% ==================== РАЗДЕЛ 10: НАИБОЛЬШИЙ ОБЩИЙ ДЕЛИТЕЛЬ ====================
	\newpage
	\section{Наибольший общий делитель. Свойства (без алгоритма Евклида) (10 вопрос)}
	
	\subsection{Определение наибольшего общего делителя}
	
	\noindent \textbf{Определение (Наибольший общий делитель пары, семейства).}
	Пусть $(a, b)$ — пара чисел из $\mathbb{Z}$. \textit{Наибольшим общим делителем} этой пары называется число $d \in \mathbb{Z}$ такое, что $d \mid a$, $d \mid b$ и $d \vdots x$ для всякого $x \in \mathbb{Z}$ такого, что $x \mid a$, $x \mid b$
	
	Аналогично, \textit{наибольший общий делитель семейства} — это общий делитель этого семейства, делящийся на любой его общий делитель
	
	\begin{tcolorbox}[colback=blue!5, colframe=blue!40!black, title=Примеры]
		\begin{itemize}
			\item Для пары $(12, 18)$: общие делители: $\pm 1, \pm 2, \pm 3, \pm 6$. Наибольший общий делитель: $6$ (и $-6$, но обычно берут положительный)
			\item Для пары $(7, 13)$: общие делители: $\pm 1$. НОД: $1$
			\item Для пары $(24, 0)$: любое число, делящее $24$, делит и $0$. Наибольший общий делитель: $24$ (и $-24$)
			\item Для семейства $(6, 10, 15)$: общие делители: $\pm 1$. НОД: $1$
		\end{itemize}
	\end{tcolorbox}
	
	\myremark
	Определение НОД через свойство "делится на любой общий делитель" эквивалентно обычному определению "наибольший по величине", но работает и для отрицательных чисел, и для нуля
	
	\subsection{Ассоциированность НОД}
	
	\noindent \textbf{Предложение (Об ассоциированности НОД).}
	Пусть $d$ — наибольший общий делитель некоторой пары (семейства) в $\mathbb{Z}$. Число $d'$ есть наибольший общий делитель этой пары (семейства) тогда и только тогда, когда $d \sim d'$, т. е. $d \vdots d'$ и $d' \vdots d$
	
	\medskip
	\noindent \textbf{Доказательство.}
	
	\textit{Прямое.} Если $d$ и $d'$ — наибольшие общие делители пары (семейства), то $d \vdots d'$ (поскольку $d'$ — общий делитель, а $d$ делится на любой общий делитель) и $d' \vdots d$ (симметрично). Значит, $d \sim d'$
	
	\textit{Обратное.} Если $d \vdots d'$, то $d'$ — общий делитель исходной пары (семейства) (так как $d$ делится на $d'$, а $d$ сам делит все числа семейства). Если $d \sim d'$, то $d' \vdots d$, и значит $d'$ делится на любой общий делитель (поскольку $d$ на него делится). Следовательно, $d'$ — наибольший общий делитель исходной пары (семейства)
	
	\begin{tcolorbox}[colback=green!5, colframe=green!50!black, title=Пояснение]
		Два числа $a$ и $b$ называются \textit{ассоциированными} ($a \sim b$), если они отличаются лишь множителем $\pm 1$, т.е. $a = \pm b$. В кольце целых чисел это означает, что они делят друг друга
	\end{tcolorbox}
	
	\subsection{Итог о существовании и единственности НОД}
	
	\noindent \textbf{Итог.}
	Таким образом,
	\begin{itemize}
		\item если $n \geq 2$ и $a_1 = a_2 = \dots = a_n = 0$, то для чисел $a_1, \dots, a_n$ существует единственный наибольший общий делитель, равный $0$;
		\item если $a_1, a_2, \dots, a_n$ все не равны $0$, то они имеют ровно два наибольших общих делителя, которые отличаются только знаком
	\end{itemize}
	
	Множество всех наибольших общих делителей семейства (или пары) $(a_i)$ мы будем обозначать НОД$(a_i)$ или НОД$((a_i))$, а для пары $(a, b)$ — НОД$(a, b)$
	
	\begin{tcolorbox}[colback=yellow!5, colframe=orange!60!black, title=Замечание об обозначениях]
		Обычно под $\operatorname{НОД}(a, b)$ понимают \textit{положительный} наибольший общий делитель. Это удобно, так как он определён однозначно. В теоретических рассуждениях иногда удобно рассматривать множество $\{\pm d\}$
	\end{tcolorbox}
	
	\subsection{Лемма о добавлении кратного}
	
	\noindent \textbf{Лемма.}
	$\forall a, b, q \in \mathbb{Z}: \medspace$ НОД(a, b) = НОД(a + bq, b)
	
	
	\medskip
	\noindent \textbf{Доказательство.}
	Пусть $M$ — множество всех общих делителей пары $(a, b)$, а $M'$ — множество всех общих делителей пары $(a + bq, b)$. Тогда
	\[
	x \in M \iff (a \vdots x, b \vdots x) \iff (a + bq \vdots x, b \vdots x) \iff x \in M'
	\]
	(поскольку если $a \vdots x$ и $b \vdots x$, то $a + bq \vdots x$, и обратно: если $a + bq \vdots x$ и $b \vdots x$, то $(a + bq) - bq = a \vdots x$)
	
	Таким образом, множества общих делителей совпадают, следовательно, совпадают и наибольшие общие делители
	
	\begin{tcolorbox}[colback=blue!5, colframe=blue!40!black, title=Пример]
		Рассмотрим $a = 18$, $b = 12$, $q = 2$:
		\[
		a + bq = 18 + 12 \cdot 2 = 18 + 24 = 42
		\]
		Проверим: $\operatorname{НОД}(18, 12) = 6$, $\operatorname{НОД}(42, 12) = 6$ (так как $42 = 6 \cdot 7$, $12 = 6 \cdot 2$, общих делителей больше нет). Лемма подтверждается.
	\end{tcolorbox}
	
	\myremark
	Эта лемма — ключевой шаг в алгоритме Евклида. Она позволяет уменьшать числа, не меняя их НОД, заменяя большее число на остаток от деления
	
	\subsection{Свойства НОД}
	
	\noindent \textbf{Свойства НОД.}
	
	\begin{enumerate}
		\item НОД$(ca, cb) = c \cdot \operatorname{НОД}(a, b)$;
		\item Если $a$ и $b$ взаимно просты, то НОД$(a, c) \cdot $НОД$(b, c) =$ НОД$(ab, c)$;
		\item НОД$(a, 0) = a$;
		\item Если $a \vdots b$, то НОД$(a, b) = b$;
		\item НОД$(ab, c) =$ НОД$(b, c)$, если $a$, $c$ взаимно просты
	\end{enumerate}
	
	\begin{tcolorbox}[colback=gray!5, colframe=gray!50, title=Примеры на свойства]
		\begin{itemize}
			\item[1.] НОД$(12, 18) = 6$, тогда НОД$(24, 36) = 12 = 2 \cdot 6$
			\item[2.] $a=3$, $b=5$ взаимно просты, $c=15$. НОД$(3,15)=3$, НОД$(5,15)=5$, произведение $15$, НОД$(3\cdot5,15)=$ НОД$(15,15)=15$
			\item[3.] НОД$(12,0)=12$
			\item[4.] $12 \vdots 3$, поэтому НОД$(12,3)=3$
			\item[5.] НОД$(6\cdot 5, 7) = $ НОД$(30,7)=1$, и НОД$(5,7)=1$, так как $6$ и $7$ взаимно просты
		\end{itemize}
	\end{tcolorbox}
	
	\begin{tcolorbox}[colback=purple!5, colframe=purple!50!black, title=Заметочка: Взаимно простые числа]
		Числа $a$ и $b$ называются \textit{взаимно простыми}, если НОД$(a, b) = 1$. Это понятие постоянно используется в свойствах НОД. Например, свойство 2 показывает, что для взаимно простых сомножителей НОД мультипликативен
	\end{tcolorbox}
	
	\myremark
	Свойство 5 особенно полезно: если один из сомножителей взаимно прост с делителем, то его можно "отбросить" при вычислении НОД. Например, НОД$(14 \cdot 9, 25) =$ НОД$(9, 25)$, так как $14$ и $25$ взаимно просты
	
	\begin{tcolorbox}[colback=yellow!5, colframe=orange!60!black, title=Важное замечание]
		Все эти свойства выводятся непосредственно из определения НОД и леммы о добавлении кратного, без использования алгоритма Евклида. Они составляют базовый аппарат теории делимости
	\end{tcolorbox}
	
	
	
	
	
	% ==================== РАЗДЕЛ 11: АЛГОРИТМ ЕВКЛИДА ====================
	\newpage
	\section{Алгоритм Евклида. Линейное представление наибольшего общего делителя (11 вопрос)}
	
	\subsection{Лемма о добавлении кратного}
	
	\noindent \textbf{Лемма.}
	$\forall a, b, q \in \mathbb{Z}$:
	\[
	\NOD(a, b) = \NOD(a + bq, b)
	\]
	
	\medskip
	\noindent \textbf{Доказательство.}
	Пусть $M$ — множество всех общих делителей пары $(a, b)$, а $M'$ — множество всех общих делителей пары $(a + bq, b)$. Тогда
	\[
	x \in M \iff (a \vdots x, b \vdots x) \iff (a + bq \vdots x, b \vdots x) \iff x \in M'
	\]
	(поскольку если $a \vdots x$ и $b \vdots x$, то $a + bq \vdots x$, и обратно: если $a + bq \vdots x$ и $b \vdots x$, то $(a + bq) - bq = a \vdots x$)
	
	Таким образом, множества общих делителей совпадают, следовательно, совпадают и наибольшие общие делители
	
	\begin{tcolorbox}[colback=blue!5, colframe=blue!40!black, title=Пример]
	Для $a = 18$, $b = 12$, $q = 2$:
	\[
	a + bq = 18 + 12 \cdot 2 = 42,\quad \NOD(18,12)=6,\quad \NOD(42,12)=6
	\]
	\end{tcolorbox}
	
	\myremark
	Эта лемма — основа алгоритма Евклида. Она позволяет заменять пару чисел на другую пару с тем же НОД, уменьшая при этом числа
	
	\subsection{Алгоритм Евклида}
	
	\noindent \textbf{Алгоритм Евклида.}
	Пусть $(a, b)$ — пара чисел из $\mathbb{Z}$. Если $b = 0$, то $\NOD(a, 0) = a$. Пусть $b \neq 0$
	
	Разделим с остатком $a$ на $b$:
	\[
	a = bq_1 + r_1, \quad 0 \leq r_1 < |b|
	\]
	При этом $r_1 = a - bq_1$ и по лемме $\NOD(a, b) = \NOD(b, r_1)$, т. е. мы можем заменить пару $(a, b)$ парой $(b, r_1)$. Это и есть один шаг алгоритма. Делая цепочку таких шагов, мы будем иметь строго ограниченную снизу цепочку целых чисел
	\[
	|b| > r_1 > r_2 > \dots \geq 0
	\]
	Такая последовательность обязательно конечна, поэтому конечно и число шагов алгоритма. Очередной шаг невозможен лишь при условии, что при некотором $k$ мы получим $r_{k+1} = 0$, потому что только в этом случае $r_{k+1}$ невозможно разделить с остатком. Т. е. мы получим $(r_k, 0)$. Значит, $\NOD(a, b) = \NOD(r_k, 0) = r_k$. (Если уже $r_1 = 0$, то полагаем $r_0 = b$)

	Таким образом, алгоритм Евклида можно представить такой схемой:
	\begin{align*}
	a &= bq_1 + r_1 \\
	b &= r_1q_2 + r_2 \\
	r_1 &= r_2q_3 + r_3 \\
	\cdots &\cdots \cdots \cdots \cdots \\
	r_{k-2} &= r_{k-1}q_k + r_k \\
	r_{k-1} &= r_kq_{k+1}
	\end{align*}
	
	\begin{tcolorbox}[colback=gray!5, colframe=gray!50, title=Пример работы алгоритма Евклида]
	Найдём $\NOD(123, 27)$:
	\begin{align*}
	123 &= 27 \cdot 4 + 15 \quad (q_1=4, r_1=15) \\
	27 &= 15 \cdot 1 + 12 \quad (q_2=1, r_2=12) \\
	15 &= 12 \cdot 1 + 3 \quad (q_3=1, r_3=3) \\
	12 &= 3 \cdot 4 + 0 \quad (q_4=4, r_4=0)
	\end{align*}
	Последний ненулевой остаток $r_3 = 3$, значит $\NOD(123, 27) = 3$.
	\end{tcolorbox}
	
	\begin{tcolorbox}[colback=green!5, colframe=green!50!black, title=Пояснение]
	На каждом шаге мы заменяем большее число на остаток от деления на меньшее. Поскольку остаток строго меньше делителя, числа строго убывают, и процесс обязательно закончится получением нулевого остатка. Последний ненулевой остаток и есть НОД
	\end{tcolorbox}
	
	\subsection{Теорема о линейном представлении наибольшего общего делителя}
	
	\noindent \textbf{Теорема (О линейном представлении НОД).}
	У любого семейства $(a_i)_{i \in I}$ в $\mathbb{Z}$ существует наибольший общий делитель, и если $d = \NOD(a_i)$, то в $\mathbb{Z}$ существует семейство $(u_i)_{i \in I}$ такое, что
	\[
	d = \sum a_i u_i
	\]
	
	\medskip
	\noindent \textbf{Доказательство.}
	
	Если $\forall i \ a_i = 0$, то $d = 0$ и можно взять $\forall i \ u_i = 0$, так что утверждения теоремы очевидны
	
	Пусть в семействе $(a_i)$ есть ненулевые члены и пусть
	\[
	M = \left\{ \sum a_i y_i \mid \forall i \ y_i \in \mathbb{Z} \text{ и почти все } y_i = 0 \right\}.
	\]
	Очевидно, что $\forall i \ a_i \in M$ (можно взять $y_i = 1$, $\forall j \neq i \ y_j = 0$) и в частности $M \neq \{0\}$
	
	Множество $\{|x| \mid x \in M \setminus \{0\}\}$ не пусто, состоит из целых чисел и ограничено снизу, поэтому в нём есть наименьшее число, которое есть абсолютная величина некоторого $d' \in M$, $d' \neq 0$, $d' = \sum a_i u'_i$ при некоторых $u'_i \in \mathbb{Z}$. Покажем, что $d'$ — наибольший общий делитель семейства $(a_i)$
	
	Для произвольного $j \in I$ разделим с остатком $a_j$ на $d'$: $a_j = d' q + r$. Тогда
	\[
	r = a_j - d' q = a_j - q \sum a_i u'_i = (1 - u'_j q) a_j + \sum_{i \neq j} (-u'_i q) a_i \in M
	\]
	и либо $r = 0$, либо $|r| < |d'|$
	
	Поскольку $r \in M$, а $|d'|$ — наименьшее значение абсолютной величины на $M$, неравенство $|r| < |d'|$ невозможно, поэтому $r = 0$ и $a_j \vdots d'$. Таким образом, $d'$ — общий делитель семейства $(a_i)$

	Если $z$ — произвольный общий делитель $(a_i)$, то $d' = \sum a_i u'_i \vdots z$, т. е. $d'$ — наибольший общий делитель семейства $(a_i)$. Существование наибольшего общего делителя доказано
	
	Если $d$ — любой наибольший общий делитель семейства $(a_i)$, то $d \vdots d'$, поэтому $d = \nu d'$, $d = \nu \sum a_i u'_i = \sum a_i (\nu u'_i) = \sum a_i u_i$, где $u_i = \nu u'_i$
	
	\begin{tcolorbox}[colback=red!5, colframe=red!70!black, title=Итог теоремы]
	Любой наибольший общий делитель чисел $a_1, \dots, a_n$ может быть представлен в виде линейной комбинации этих чисел с целыми коэффициентами. Это фундаментальный факт теории чисел
	\end{tcolorbox}
	
	\begin{tcolorbox}[colback=blue!5, colframe=blue!40!black, title=Пример линейного представления]
	Для чисел $123$ и $27$ из предыдущего примера:
	\begin{align*}
	3 &= 15 - 12 \cdot 1 \\
	12 &= 27 - 15 \cdot 1 \\
	15 &= 123 - 27 \cdot 4
	\end{align*}
	Подставляем обратно:
	\begin{align*}
	3 &= 15 - (27 - 15) = 2 \cdot 15 - 27 \\
	3 &= 2 \cdot (123 - 27 \cdot 4) - 27 = 2 \cdot 123 - 8 \cdot 27 - 27 = 2 \cdot 123 - 9 \cdot 27
	\end{align*}
	Таким образом, $3 = 123 \cdot 2 + 27 \cdot (-9)$.
	\end{tcolorbox}
	
	\myremark
	Коэффициенты в линейном представлении можно найти, последовательно выражая остатки из алгоритма Евклида и подставляя их обратно. Этот метод называется \textit{обратным ходом алгоритма Евклида}
	
	\begin{tcolorbox}[colback=purple!5, colframe=purple!50!black, title=Заметочка: Следствия теоремы]
	Из теоремы о линейном представлении вытекают важные следствия:
	\begin{itemize}
	\item Если $d = \NOD(a, b)$, то существуют целые $u, v$ такие, что $au + bv = d$.
	\item Числа $a$ и $b$ взаимно просты тогда и только тогда, когда существуют целые $u, v$ с $au + bv = 1$.
	\item Если $a \mid bc$ и $\NOD(a, b) = 1$, то $a \mid c$ (лемма Евклида)
	\end{itemize}
	\end{tcolorbox}
	
	\begin{tcolorbox}[colback=yellow!5, colframe=orange!60!black, title=Важное замечание]
	Теорема о линейном представлении справедлива для любого конечного семейства целых чисел. Для бесконечного семейства утверждение также верно, но сумма становится бесконечной, поэтому требуется, чтобы почти все коэффициенты были нулевыми (что мы и предполагали в определении $M$)
	\end{tcolorbox}








	
	

	
	% ==================== РАЗДЕЛ 12: ВЗАИМНО ПРОСТЫЕ ЧИСЛА ====================
	\newpage
	\section{Взаимно простые числа. Связь НОД и НОК. Наименьшее общее кратное взаимно простых чисел (12 вопрос)}
	
	\subsection{Определение взаимной простоты}
	
	\noindent \textbf{Определение (Взаимная простота).}
	Пусть $a, b \in \mathbb{Z}$. Говорят, что $a$ и $b$ \textit{взаимно просты}, если $\NOD(a, b) = 1$
	
	Говорят, что семейство $(a)$ \textit{взаимно просто}, если $\NOD((a)) = 1$
	
	\begin{tcolorbox}[colback=blue!5, colframe=blue!40!black, title=Примеры]
	\begin{itemize}
	\item $6$ и $35$ взаимно просты, так как $\NOD(6,35)=1$
	\item $8$ и $15$ взаимно просты ($\NOD=1$)
	\item $12$ и $18$ не взаимно просты ($\NOD=6$)
	\item Семейство $(6, 10, 15)$ взаимно просто, так как $\NOD(6,10,15)=1$
	\item Семейство $(6, 10, 14)$ не взаимно просто ($\NOD=2$)
	\end{itemize}
	\end{tcolorbox}
	
	\subsection{Свойства взаимной простоты}
	
	\noindent \textbf{Свойство 1.}
	Семейство $(a_i)_{i \in I}$ взаимно просто тогда и только тогда, когда существует семейство $(u_i)_{i \in I}$ в $\mathbb{Z}$ такое, что $\sum_{i \in I} a_i u_i = 1$
	
	\medskip
	\noindent \textbf{Доказательство.}
	
	\textit{Прямое.} Это частный случай теоремы о линейном представлении НОД. Если $\NOD((a_i)) = 1$, то по теореме о линейном представлении существуют коэффициенты $u_i$ с $\sum a_i u_i = 1$
	
	\textit{Обратное.} Если $\forall i \in I \ a_i \vdots d$ и $\sum a_i u_i = 1$, то $1 \vdots d$, поэтому $\NOD((a_i)) = 1$
	
	\begin{tcolorbox}[colback=green!5, colframe=green!50!black, title=Пояснение]
	Это свойство даёт эквивалентное определение взаимной простоты: числа взаимно просты тогда и только тогда, когда некоторая их целочисленная линейная комбинация равна единице
	\end{tcolorbox}
	
	\begin{tcolorbox}[colback=gray!5, colframe=gray!50, title=Пример]
	Для $a=6$, $b=35$: $6 \cdot 6 + 35 \cdot (-1) = 36 - 35 = 1$.
	Для $a=8$, $b=15$: $8 \cdot 2 + 15 \cdot (-1) = 16 - 15 = 1$.
	\end{tcolorbox}
	
	\medskip
	\noindent \textbf{Свойство 2.}
	Если $d = \NOD$ семейства $(a_i)$ и $\forall i \ b_i = \frac{a_i}{d}$, то семейство $(b_i)$ взаимно просто
	
	\medskip
	\noindent \textbf{Доказательство.}
	По теореме о линейном представлении НОД $d = \sum a_i u_i = \sum d b_i u_i \Rightarrow 1 = \sum b_i u_i$, значит $(b_i)$ взаимно просто по свойству 1
	
	\begin{tcolorbox}[colback=gray!5, colframe=gray!50, title=Пример]
	Для чисел $12$ и $18$: $d = 6$, тогда $b_1 = 12/6 = 2$, $b_2 = 18/6 = 3$. Числа $2$ и $3$ взаимно просты
	\end{tcolorbox}
	
	\medskip
	\noindent \textbf{Свойство 3.}
	Если каждый член конечного семейства $(a_i)$ взаимно прост с $b$, то и $\prod a_i$ взаимно просто с $b$
	
	\medskip
	\noindent \textbf{Доказательство.}
	По свойству 1 имеем $\forall i \ a_i u_i + b \upsilon_i = 1$ для некоторых $u_i, \upsilon_i$. Тогда $a_i u_i = 1 - b \upsilon_i$;
	\[
	(\prod a_i) u = \prod (a_i u_i) = \prod (1 - b \upsilon_i) = 1 - b \upsilon,
	\]
	где $u = \prod u_i$, а $\upsilon$ — некоторое целое число (получающееся при раскрытии скобок). Итак, $(\prod a_i) u + b \upsilon = 1$ и по свойству 1 $\prod a_i$ и $b$ взаимно просты
	
	\begin{tcolorbox}[colback=blue!5, colframe=blue!40!black, title=Пример]
	$6$ и $35$ взаимно просты, $7$ и $35$ взаимно просты ($\NOD(7,35)=7$? Нет, $\NOD(7,35)=7$, поэтому условие не выполнено). 
	Возьмём корректный пример: $4$ взаимно просто с $9$, $5$ взаимно просто с $9$. Тогда $4 \cdot 5 = 20$ взаимно просто с $9$ ($\NOD(20,9)=1$)
	\end{tcolorbox}
	
	\medskip
	\noindent \textbf{Свойство 4.}
	Если $ab \vdots c$ и $a$, $c$ взаимно просты, то $b \vdots c$
	
	\medskip
	\noindent \textbf{Доказательство.}
	При некоторых $u, \upsilon \in \mathbb{Z}$ $au + c\upsilon = 1$ (по свойству 1). Умножая на $b$, получаем
	\[
	b = b \cdot 1 = abu + cb\upsilon.
	\]
	Поскольку $ab \vdots c$ и $cb\upsilon \vdots c$, то и $b \vdots c$
	
	\begin{tcolorbox}[colback=red!5, colframe=red!70!black, title=Лемма Евклида]
	Это свойство известно как \textit{лемма Евклида}: если простое число делит произведение, то оно делит один из сомножителей. В обобщённой форме: если число, взаимно простое с одним из сомножителей, делит произведение, то оно делит другой сомножитель
	\end{tcolorbox}
	
	\begin{tcolorbox}[colback=gray!5, colframe=gray!50, title=Пример]
	$6 \cdot 10 = 60$ делится на $15$. Проверим: $\NOD(6,15)=3 \neq 1$, условие не выполнено.
	Возьмём $6 \cdot 10 = 60$ и $c = 5$. $\NOD(6,5)=1$, и $60 \vdots 5$. Тогда по свойству $10 \vdots 5$ — верно
	\end{tcolorbox}
	
	\medskip
	\noindent \textbf{Свойство 5.}
	Если $a \vdots b$, $a \vdots c$ и $b$, $c$ взаимно просты, то $a \vdots bc$
	
	\medskip
	\noindent \textbf{Доказательство.}
	$a \vdots b \Rightarrow a = bx$. Так как $a \vdots c$, то $bx \vdots c$. Поскольку $b$ и $c$ взаимно просты, по свойству 4 получаем $x \vdots c$, т. е. $x = cy$, и тогда $a = bcy \vdots bc$
	
	\begin{tcolorbox}[colback=gray!5, colframe=gray!50, title=Пример]
	$a = 60$, $b = 6$, $c = 5$. $60 \vdots 6$, $60 \vdots 5$, $6$ и $5$ взаимно просты. Тогда $60 \vdots 6 \cdot 5 = 30$ — верно
	\end{tcolorbox}
	
	\myremark
	Свойства 4 и 5 являются ключевыми в теории делимости. Они показывают, как взаимная простота позволяет "разделять" условия делимости
	
	\subsection{Наименьшее общее кратное}
	
	\noindent \textbf{Определение (Наименьшее общее кратное).}
	\textit{Наименьшим общим кратным} целых чисел $a_1, a_2, \dots, a_n$ называется любое целое число $k$, удовлетворяющее следующим условиям:
	\begin{enumerate}
	\item $k$ есть общее кратное чисел $a_1, a_2, \dots, a_n$, т. е. $k \vdots a_1, \dots, k \vdots a_n$;
	\item $k$ делит любое общее кратное чисел $a_1, a_2, \dots, a_n$, т. е. $\forall k_1 \in \mathbb{Z}$ $(k_1 \vdots a_1, \dots, k_1 \vdots a_n) \Rightarrow k \mid k_1$
	\end{enumerate}
	
	Множество всех наименьших общих кратных чисел $a_1, \dots, a_n$ обозначают $\NOK\{a_1, \dots, a_n\}$
	
	\begin{tcolorbox}[colback=blue!5, colframe=blue!40!black, title=Замечание о существовании]
	Ясно, что если $n \geq 2$ и хотя бы одно из целых чисел равно $0$, то для них существует единственное наименьшее общее кратное, равное $0$. Если же целые числа $a_1, \dots, a_n$ отличны от $0$, то для них существует ровно два наименьших общих кратных, отличающихся только знаком. Неотрицательное наименьшее общее кратное часто обозначают $[a_1, \dots, a_n]$
	\end{tcolorbox}
	
	\begin{tcolorbox}[colback=gray!5, colframe=gray!50, title=Примеры]
	\begin{itemize}
	\item $[6, 8] = 24$ (общие кратные: $24, 48, 72, \dots$)
	\item $[12, 18] = 36$
	\item $[7, 11] = 77$ (взаимно простые)
	\item $[4, 6, 15] = 60$
	\end{itemize}
	\end{tcolorbox}
	
	\subsection{Связь НОД и НОК через каноническое разложение}
	
	Из канонического разложения чисел
	\[
	n = \pm p_1^{\alpha_1} \dots p_k^{\alpha_k}, \quad m = \pm p_1^{\beta_1} \dots p_k^{\beta_k}
	\]
	если условиться допускать нулевые показатели, следует, что
	\[
	(n, m) = p_1^{\gamma_1} \dots p_k^{\gamma_k}, \quad [n, m] = p_1^{\delta_1} \dots p_k^{\delta_k}
	\]
	где $\gamma_i = \min(\alpha_i, \beta_i)$, $\delta_i = \max(\alpha_i, \beta_i)$, $i = 1, 2, \dots, k$.
	
	\begin{tcolorbox}[colback=green!5, colframe=green!50!black, title=Пояснение]
	НОД берёт минимальные показатели для каждого простого, НОК — максимальные. Это прямое следствие определений
	\end{tcolorbox}
	
	\begin{tcolorbox}[colback=gray!5, colframe=gray!50, title=Пример]
	$n = 12 = 2^2 \cdot 3^1$, $m = 18 = 2^1 \cdot 3^2$.
	Тогда $\NOD(12,18) = 2^{\min(2,1)} \cdot 3^{\min(1,2)} = 2^1 \cdot 3^1 = 6$.
	$\NOK(12,18) = 2^{\max(2,1)} \cdot 3^{\max(1,2)} = 2^2 \cdot 3^2 = 36$.
	\end{tcolorbox}
	
	\subsection{Основное соотношение между НОД и НОК}
	
	\noindent \textbf{Теорема (Связь НОД и НОК).}
	При $n > 0$, $m > 0$ выполнено соотношение:
	\[
	\NOD(n, m) \cdot \NOK(n, m) = n m
	\]
	
	\medskip
	\noindent \textbf{Доказательство.}
	Из канонических разложений:
	\[
	\NOD(n, m) \cdot \NOK(n, m) = p_1^{\min(\alpha_1,\beta_1) + \max(\alpha_1,\beta_1)} \dots p_k^{\min(\alpha_k,\beta_k) + \max(\alpha_k,\beta_k)}
	\]
	Но $\min(\alpha_i, \beta_i) + \max(\alpha_i, \beta_i) = \alpha_i + \beta_i$, поэтому произведение равно $p_1^{\alpha_1+\beta_1} \dots p_k^{\alpha_k+\beta_k} = n m$
	
	\begin{tcolorbox}[colback=blue!5, colframe=blue!40!black, title=Пример]
	Для $n=12$, $m=18$: $\NOD=6$, $\NOK=36$, $6 \cdot 36 = 216 = 12 \cdot 18$.
	\end{tcolorbox}
	
	\subsection{НОК взаимно простых чисел}
	
	\noindent \textbf{Следствие (НОК взаимно простых чисел).}
	Если $n$, $m$ взаимно просты, то
	\[
	\NOK(n, m) = n \cdot m
	\]
	
	\medskip
	\noindent \textbf{Доказательство.}
	Из основного соотношения: $\NOD(n, m) \cdot \NOK(n, m) = n m$. Так как $\NOD(n, m) = 1$, получаем $\NOK(n, m) = n m$
	
	\begin{tcolorbox}[colback=gray!5, colframe=gray!50, title=Примеры]
	\begin{itemize}
	\item $[7, 11] = 77 = 7 \cdot 11$
	\item $[8, 15] = 120 = 8 \cdot 15$
	\item $[6, 35] = 210 = 6 \cdot 35$
	\end{itemize}
	\end{tcolorbox}
	
	\myremark
	Это свойство часто используется: для взаимно простых чисел их произведение является наименьшим общим кратным. Для общих кратных это не так: например, $[4,6] = 12$, а $4 \cdot 6 = 24$ (не наименьшее)
	
	\begin{tcolorbox}[colback=purple!5, colframe=purple!50!black, title=Заметочка: Обобщение на несколько чисел]
	Для нескольких чисел $\NOK(a_1, \dots, a_n)$ можно вычислять последовательно:
	\[
	\NOK(a_1, a_2, a_3) = \NOK(\NOK(a_1, a_2), a_3)
	\]
	Аналогично НОД. Для взаимно простых в совокупности чисел (т.е. $\NOD(a_1, \dots, a_n) = 1$) произведение не обязано быть НОК (например, $(6,10,15)$ взаимно просты в совокупности, но $6 \cdot 10 \cdot 15 = 900$, а $\NOK(6,10,15)=30$). Для НОК требуется попарная взаимная простота
	\end{tcolorbox}
	
	
	
	
	
	
	% ==================== РАЗДЕЛ 13: ПРОСТЫЕ ЧИСЛА ====================
	\newpage
	\section{Простые числа. Наименьший делитель числа. Бесконечность простых чисел (13 вопрос)}
	
	\subsection{Определение простого числа}
	
	\noindent \textbf{Определение (Простое число).}
	Натуральное число $p \neq 1$ называют \textit{простым}, если оно не имеет натуральных делителей, отличных от $1$ и $p$, в противном случае оно называется \textit{составным}. Число $1$ не относится ни к простым, ни к составным числам
	
	\begin{tcolorbox}[colback=blue!5, colframe=blue!40!black, title=Примеры]
		\begin{itemize}
			\item Простые числа: $2, 3, 5, 7, 11, 13, 17, 19, 23, \dots$
			\item Составные числа: $4, 6, 8, 9, 10, 12, 14, 15, 16, \dots$
			\item $1$ — не простое и не составное
			\item $2$ — единственное чётное простое число
		\end{itemize}
	\end{tcolorbox}
	
	\myremark
	Определение простого числа требует, чтобы число было больше $1$ и имело ровно два натуральных делителя: $1$ и само себя. Число $1$ имеет только один делитель, поэтому оно не подходит под определение
	
	\subsection{Наименьший делитель числа}
	
	\noindent \textbf{Предложение (О простоте наименьшего натурального делителя).}
	Для любого натурального числа $n > 1$ наименьший отличный от $1$ натуральный делитель всегда является простым числом
	
	\medskip
	\noindent \textbf{Доказательство.}
	Рассмотрим множество $M$, состоящее из натуральных отличных от $1$ делителей числа $n$. Множество $M$ не пустое, так как $n \in M$. Значит, в множестве $M$ существует наименьшее число $q > 1$
	
	Пусть $q$ не простое. Тогда существует такое $a$, что $1 < a < q$ и $q \vdots a$ (по определению составного числа). Поскольку $q \vdots a$ и $q \vdots n$ (так как $q$ — делитель $n$), то $a \vdots n$ (транзитивность делимости). Значит, $a \in M$ и $a < q$, что противоречит минимальности $q$
	
	Следовательно, $q$ — простое число
	
	\begin{tcolorbox}[colback=green!5, colframe=green!50!black, title=Пояснение]
		Идея доказательства: берём наименьший делитель числа $n$, больший $1$. Если бы он был составным, то его собственный делитель был бы ещё меньшим делителем $n$, что противоречит минимальности
	\end{tcolorbox}
	
	\begin{tcolorbox}[colback=gray!5, colframe=gray!50, title=Пример]
		Рассмотрим $n = 45$. Делители числа $45$, большие $1$: $3, 5, 9, 15, 45$. Наименьший из них — $3$, и $3$ — простое число.
		Рассмотрим $n = 91$. Делители: $7, 13, 91$. Наименьший — $7$, простое.
		Рассмотрим $n = 97$ (простое). Делители: $97$. Наименьший — $97$, простое
	\end{tcolorbox}
	
	\myremark
	Это предложение даёт простой способ проверки числа на простоту: если у числа $n$ нет делителей, не превосходящих $\sqrt{n}$, то оно простое. Действительно, если бы у $n$ был делитель $d > \sqrt{n}$, то соответствующий парный делитель $n/d$ был бы меньше $\sqrt{n}$
	
	\subsection{Бесконечность множества простых чисел}
	
	\noindent \textbf{Теорема (О бесконечности множества простых чисел).}
	Множество простых чисел бесконечно
	
	\medskip
	\noindent \textbf{Доказательство.}
	Предположим, что множество простых чисел конечно. Выписав их все в порядке возрастания, получим ряд чисел:
	\[
	2, 3, 5, \dots, p_r \tag{1}
	\]
	
	Рассмотрим число
	\[
	N = (2 \cdot 3 \cdot 5 \cdots p_r) + 1
	\]
	
	Так как каждое число из (1) делит произведение $2 \cdot 3 \cdot 5 \cdots p_r$, но не делит $1$, то число $N$ не делится ни на одно из чисел (1), т. е. ни на одно простое число
	
	А так как число $N$ больше единицы, то это противоречит основной теореме арифметики (о том, что всякое натуральное число, большее $1$, разлагается на простые множители). Следовательно, наше предположение о конечности множества простых чисел неверно
	
	\begin{tcolorbox}[colback=red!5, colframe=red!70!black, title=Классическое доказательство Евклида]
		Это знаменитое доказательство принадлежит Евклиду (III век до н.э.) и содержится в его "Началах". Оно считается одним из самых элегантных в математике
	\end{tcolorbox}
	
	\begin{tcolorbox}[colback=blue!5, colframe=blue!40!black, title=Иллюстрация]
		Пусть простые числа: $2, 3, 5$.
		\[
		N = 2 \cdot 3 \cdot 5 + 1 = 30 + 1 = 31
		\]
		Число $31$ не делится на $2, 3, 5$. Оно само простое, причём новое.
		
		Пусть простые числа: $2, 3, 5, 7, 11, 13$.
		\[
		N = 2 \cdot 3 \cdot 5 \cdot 7 \cdot 11 \cdot 13 + 1 = 30030 + 1 = 30031
		\]
		$30031 = 59 \cdot 509$ — составное, но его простые делители ($59$ и $509$) не входят в исходный список
	\end{tcolorbox}
	
	\begin{tcolorbox}[colback=green!5, colframe=green!50!black, title=Пояснение]
		Число $N$ может быть как простым, так и составным. Но в любом случае оно имеет простой делитель, отличный от всех $p_1, \dots, p_r$:
		\begin{itemize}
			\item Если $N$ простое, то это новый простой делитель
			\item Если $N$ составное, то любой его простой делитель не может равняться ни одному из $p_i$, так как $N$ при делении на $p_i$ даёт остаток $1$
		\end{itemize}
		Таким образом, всегда находится новое простое число
	\end{tcolorbox}
	
	\myremark
	Из этой теоремы следует, что простых чисел бесконечно много. Более того, существуют различные уточнения: например, что простых чисел вида $4k+3$ также бесконечно много, что простых чисел в любой арифметической прогрессии (при выполнении некоторых условий) бесконечно много (теорема Дирихле)
	
	\begin{tcolorbox}[colback=purple!5, colframe=purple!50!black, title=Заметочка: Основная теорема арифметики]
		Основная теорема арифметики утверждает, что любое натуральное число $n > 1$ единственным образом (с точностью до порядка сомножителей) представляется в виде произведения простых чисел. Именно на неё опирается доказательство Евклида
	\end{tcolorbox}
	
	\begin{tcolorbox}[colback=yellow!5, colframe=orange!60!black, title=Важное замечание]
		Существуют и другие доказательства бесконечности простых чисел. Например, аналитическое доказательство Эйлера использует расходимость гармонического ряда и свойства дзета-функции. Однако доказательство Евклида остаётся самым простым и наглядным
	\end{tcolorbox}
	
	
	
	
	

	
	% ==================== РАЗДЕЛ 14: ДВЕ ЛЕММЫ О ПРОСТЫХ ЧИСЛАХ ====================
	\newpage
	\section{«Две» леммы о простых числах (НОД и деление произведения) (14 вопрос)}
	
	\subsection{Лемма о наибольшем общем делителе}
	
	\noindent \textbf{Лемма 1 (О НОД суммы с кратным).}
	Для любых целых чисел $a, b, q \in \mathbb{Z}$:
	\[
	\NOD(a, b) = \NOD(a + bq, b)
	\]
	
	\medskip
	\noindent \textbf{Доказательство.}
	Пусть $M$ и $M'$ — множества всех общих делителей пар $(a, b)$ и $(a + bq, b)$ соответственно.
	Тогда
	\[
	x \in M \iff (a \vdots x, b \vdots x) \iff (a + bq \vdots x, b \vdots x) \iff x \in M'
	\]
	Действительно:
	\begin{itemize}
	\item Если $a \vdots x$ и $b \vdots x$, то $a + bq \vdots x$ (линейная комбинация чисел, кратных $x$, также кратна $x$).
	\item Обратно, если $a + bq \vdots x$ и $b \vdots x$, то $a = (a + bq) - bq \vdots x$ (разность двух чисел, кратных $x$, кратна $x$).
	\end{itemize}
	Таким образом, $M = M'$. Поскольку множества общих делителей совпадают, совпадают и их наибольшие элементы, т. е. $\NOD(a, b) = \NOD(a + bq, b)$
	
	\begin{tcolorbox}[colback=blue!5, colframe=blue!40!black, title=Пример]
	Возьмём $a = 18$, $b = 12$, $q = 2$:
	\[
	a + bq = 18 + 12 \cdot 2 = 18 + 24 = 42
	\]
	$\NOD(18, 12) = 6$, $\NOD(42, 12) = 6$ — совпадают
	\end{tcolorbox}
	
	\myremark
	Эта лемма уже использовалась в алгоритме Евклида. Она показывает, что прибавление к числу кратного другого числа не меняет множества общих делителей, а значит, и НОД
	
	\subsection{Лемма о делении произведения (первая форма)}
	
	\noindent \textbf{Лемма 2 (О делении произведения).}
	Если $ab \vdots c$ ($a, b, c \in \mathbb{Z}$) и $a$, $c$ взаимно просты, то $b \vdots c$
	
	\medskip
	\noindent \textbf{Доказательство.}
	Поскольку $a$ и $c$ взаимно просты, по теореме о линейном представлении НОД существуют такие $u, \upsilon \in \mathbb{Z}$, что
	\[
	au + c\upsilon = 1
	\]
	
	Умножая это равенство на $b$, имеем:
	\[
	b = b \cdot 1 = b(au + c\upsilon) = abu + cb\upsilon
	\]
	
	По условию $ab \vdots c$, значит $abu \vdots c$. Очевидно, $cb\upsilon \vdots c$. Следовательно, сумма $abu + cb\upsilon = b$ также делится на $c$
	
	\begin{tcolorbox}[colback=green!5, colframe=green!50!black, title=Пояснение]
	Ключевая идея: представить $b$ как линейную комбинацию $ab$ и $c$, используя соотношение $au + c\upsilon = 1$. Это классический приём в теории чисел.
	\end{tcolorbox}
	
	\begin{tcolorbox}[colback=gray!5, colframe=gray!50, title=Пример]
	Пусть $a = 6$, $c = 5$ (взаимно просты), и известно, что $6b \vdots 5$. Докажем, что $b \vdots 5$.
	Находим линейное представление: $6 \cdot 1 + 5 \cdot (-1) = 1$.
	Тогда $b = b(6 \cdot 1 + 5 \cdot (-1)) = 6b \cdot 1 + 5 \cdot (-b)$.
	Поскольку $6b \vdots 5$ и $5 \cdot (-b) \vdots 5$, то $b \vdots 5$.
	\end{tcolorbox}
	
	\myremark
	Эта лемма является обобщением известного факта: если простое число делит произведение, то оно делит один из сомножителей. Здесь роль простого числа играет $c$, взаимно простой с $a$
	
	\subsection{Лемма о делении произведения (вторая форма)}
	
	\noindent \textbf{Лемма 3 (О делении произведения на простое число).}
	Если $ab \vdots p$ ($a, b, p \in \mathbb{Z}$) и $p$ простое, то $a \vdots p$ или $b \vdots p$
	
	\medskip
	\noindent \textbf{Доказательство.}
	Если $a$ не делится на $p$, то $\NOD(a, p) = 1$ (поскольку $p$ простое, его единственные делители — $1$ и $p$, а $a$ на $p$ не делится, значит общий делитель только $1$). Тогда $a$ и $p$ взаимно просты, и по предыдущей лемме $b \vdots p$
	
	Таким образом, $a \vdots p$ или $b \vdots p$
	
	\begin{tcolorbox}[colback=red!5, colframe=red!70!black, title=Лемма Евклида]
	Это утверждение известно как \textit{лемма Евклида}. Оно является ключевым при доказательстве основной теоремы арифметики
	\end{tcolorbox}
	
	\begin{tcolorbox}[colback=gray!5, colframe=gray!50, title=Пример]
	Пусть $p = 7$, $a = 14$, $b = 5$: $14 \cdot 5 = 70$ делится на $7$. Здесь $a \vdots 7$, условие выполнено
	
	Пусть $p = 7$, $a = 8$, $b = 21$: $8 \cdot 21 = 168$ делится на $7$. Здесь $a$ не делится на $7$, но $b \vdots 7$
	
	Пусть $p = 7$, $a = 8$, $b = 9$: $8 \cdot 9 = 72$ не делится на $7$, так что условие посылки не выполнено
	\end{tcolorbox}
	
	\begin{tcolorbox}[colback=blue!5, colframe=blue!40!black, title=Контрпример для составного числа]
	Для составного числа лемма не работает. Например, $4 \cdot 3 = 12$ делится на $6$, но ни $4$, ни $3$ не делятся на $6$
	\end{tcolorbox}
	
	\myremark
	Лемма Евклида — фундаментальное свойство простых чисел. Оно отличает простые числа от составных: для составного числа $m$ может существовать произведение $ab$, делящееся на $m$, при том что ни один из сомножителей не делится на $m$
	
	\begin{tcolorbox}[colback=purple!5, colframe=purple!50!black, title=Заметочка: Связь лемм]
	Три леммы этого раздела тесно связаны:
	\begin{itemize}
	\item Лемма 1 — техническая, используется в алгоритме Евклида
	\item Лемма 2 — обобщение леммы Евклида на случай, когда делитель взаимно прост с одним из сомножителей
	\item Лемма 3 — частный случай леммы 2 для простого делителя (поскольку если $p$ простое и не делит $a$, то $a$ и $p$ автоматически взаимно просты)
	\end{itemize}
	Вместе они составляют основу теории делимости
	\end{tcolorbox}
	
	\begin{tcolorbox}[colback=yellow!5, colframe=orange!60!black, title=Важное замечание]
	Лемма 2 часто формулируется так: если $a$ и $c$ взаимно просты и $c \mid ab$, то $c \mid b$. Это эквивалентная формулировка, которая иногда удобнее в применениях
	\end{tcolorbox}
	
	
	
	
	
	
	
	
	
	
	% ==================== РАЗДЕЛ 15: ОСНОВНАЯ ТЕОРЕМА АРИФМЕТИКИ ====================
	\newpage
	\section{Основная теорема арифметики. Каноническое разложение наибольшего общего делителя и наименьшего общего кратного (15 вопрос)}
	
	\subsection{Основная теорема арифметики}
	
	\noindent \textbf{Теорема (Основная теорема арифметики).}
	Каждое положительное целое число $n \neq 1$ может быть записано в виде произведения простых чисел:
	\[
	n = p_1 p_2 \dots p_s \tag{1}
	\]
	Эта запись единственна с точностью до порядка множителей
	
	\begin{tcolorbox}[colback=red!5, colframe=red!70!black, title=Формулировка]
	Любое натуральное число $n > 1$ единственным образом (с точностью до перестановки сомножителей) раскладывается в произведение простых чисел
	\end{tcolorbox}
	
	\subsection{Доказательство существования разложения}
	
	\noindent \textbf{Доказательство существования.}
	Так как $2$ — простое число, то для $n = 2$ утверждение теоремы верно. Допустим, что оно верно для всех целых положительных чисел $n \colon 2 \leq n \leq m$ при любом фиксированном натуральном числе $m \geq 2$. Докажем существование разложения (1) для $m + 1$
	
	Пусть $m + 1$ не простое число. Тогда оно делится на некоторое число $a$ такое, что $1 < a < m + 1$. Следовательно, $m + 1 = a \cdot b$, где $1 < b < m + 1$
	
	По предположению индукции каждое из чисел $a$, $b$ разлагается в произведение простых чисел:
	\[
	a = p_1 \dots p_k, \quad b = q_1 \dots q_l
	\]
	где $p_1, \dots, p_k, q_1, \dots, q_l$ — простые числа, $k, l \in \mathbb{N}$.
	
	Отсюда
	\[
	m + 1 = p_1 \dots p_k q_1 \dots q_l.
	\]
	
	\begin{tcolorbox}[colback=green!5, colframe=green!50!black, title=Пояснение]
	Доказательство использует индукцию по $n$. База: $n=2$ (простое). Переход: если $n$ простое, то разложение тривиально (оно само простое); если $n$ составное, то представляем его в виде произведения двух меньших чисел, для которых разложение уже есть по индукционному предположению
	\end{tcolorbox}
	
	\subsection{Доказательство единственности разложения}
	
	\noindent \textbf{Доказательство единственности.}
	Пусть $n = p_1 \dots p_r = p'_1 \dots p'_{r'}$ и все $p_i$, $p'_j$ просты. Пусть для определённости $r \leq r'$
	
	\textit{База индукции:} $r = 1$. Тогда $p_1 = p'_1 \dots p'_{r'} \vdots p'_1$, значит $p'_1 = p_1$ (поскольку $p_1$ простое и $p'_1$ — простой делитель). После сокращения на $p_1$ получаем $1 = p'_2 \dots p'_{r'}$, что возможно только при $r' = 1$. Таким образом, $r = r' = 1$ и $p_1 = p'_1$
	
	\textit{Индукционный переход.} Пусть $r \geq 2$. Поскольку $p'_1 \dots p'_{r'} \vdots p_r$, то по лемме Евклида (если простое число делит произведение, то оно делит один из сомножителей) существует $k$ такое, что $p'_k = p_r$. Перенумеровав, если надо, сомножители, можем считать, что $k = r'$. Сокращая исходное равенство на $p_r$, получаем
	\[
	p_1 \dots p_{r-1} = p'_1 \dots p'_{r'-1}.
	\]
	
	По предположению индукции (применимому, поскольку $r-1 < r$) имеем $r-1 = r'-1$ (т. е. $r = r'$) и $p_i = p'_i$ для $i = 1, 2, \dots, r-1$
	
	\begin{tcolorbox}[colback=blue!5, colframe=blue!40!black, title=Иллюстрация единственности]
	Рассмотрим $n = 60$. Два разложения:
	\[
	60 = 2 \cdot 2 \cdot 3 \cdot 5 = 2 \cdot 3 \cdot 2 \cdot 5
	\]
	Они отличаются лишь порядком множителей. Единственность гарантирует, что набор простых множителей (с учётом кратностей) определён однозначно
	\end{tcolorbox}
	
	\myremark
	Единственность разложения — ключевое свойство натуральных чисел. Оно не выполняется, например, в некоторых кольцах целых алгебраических чисел, что привело к развитию теории идеалов
	
	\subsection{Разложение на простые множители}
	
	\noindent \textbf{Определение (Разложение на простые множители).}
	Указанное в теореме разложение называют \textit{разложением на простые множители}
	
	\subsection{Каноническое разложение}
	
	\noindent \textbf{Определение (Каноническое разложение).}
	Перенумеровав сомножители так, чтобы они были попарно различны, мы получим
	\[
	n = p_1^{k_1} p_2^{k_2} \dots p_m^{k_m}
	\]
	Такое представление называется \textit{каноническим разложением} числа $n$
	
	\begin{tcolorbox}[colback=gray!5, colframe=gray!50, title=Примеры канонического разложения]
	\begin{align*}
	60 &= 2^2 \cdot 3^1 \cdot 5^1 \\
	84 &= 2^2 \cdot 3^1 \cdot 7^1 \\
	100 &= 2^2 \cdot 5^2 \\
	17 &= 17^1 \quad (\text{простое число})
	\end{align*}
	\end{tcolorbox}
	
	\subsection{Каноническое разложение наибольшего общего делителя и наименьшего общего кратного}
	
	\noindent \textbf{Теорема (НОД и НОК через каноническое разложение).}
	Любые два целых числа можно записать в виде произведения степеней одних и тех же простых чисел,
	\[
	n = \pm p_1^{\alpha_1} \dots p_k^{\alpha_k}, \quad m = \pm p_1^{\beta_1} \dots p_k^{\beta_k}
	\]
	если условиться допускать нулевые показатели, считая $p_i^0 = 1$. Тогда, очевидно,
	\[
	\NOD(n, m) = p_1^{\gamma_1} \dots p_k^{\gamma_k}, \quad \NOK(n, m) = p_1^{\delta_1} \dots p_k^{\delta_k} \tag{2}
	\]
	где $\gamma_i = \min(\alpha_i, \beta_i)$, $\delta_i = \max(\alpha_i, \beta_i)$, $i = 1, 2, \dots, k$
	
	\begin{tcolorbox}[colback=blue!5, colframe=blue!40!black, title=Пример]
	Для $n = 60 = 2^2 \cdot 3^1 \cdot 5^1$ и $m = 84 = 2^2 \cdot 3^1 \cdot 7^1$ запишем с общим набором простых $2,3,5,7$:
	\begin{align*}
	n &= 2^2 \cdot 3^1 \cdot 5^1 \cdot 7^0 \\
	m &= 2^2 \cdot 3^1 \cdot 5^0 \cdot 7^1
	\end{align*}
	Тогда:
	\begin{align*}
	\NOD &= 2^{\min(2,2)} \cdot 3^{\min(1,1)} \cdot 5^{\min(1,0)} \cdot 7^{\min(0,1)} = 2^2 \cdot 3^1 \cdot 5^0 \cdot 7^0 = 12 \\
	\NOK &= 2^{\max(2,2)} \cdot 3^{\max(1,1)} \cdot 5^{\max(1,0)} \cdot 7^{\max(0,1)} = 2^2 \cdot 3^1 \cdot 5^1 \cdot 7^1 = 420
	\end{align*}
	Проверка: $12 \cdot 420 = 5040 = 60 \cdot 84$.
	\end{tcolorbox}
	
	\subsection{Свойства, вытекающие из канонического разложения}
	
	Из (2) вытекают следующие утверждения:
	
	\begin{enumerate}
	\item $\NOD(n, m) \mid n$, $\NOD(n, m) \mid m$, и если $k \mid n$, $k \mid m$, то $k \mid \NOD(n, m)$.
	\item $n \mid \NOK(n, m)$, $m \mid \NOK(n, m)$, и если $n \mid u$, $m \mid u$, то $\NOK(n, m) \mid u$
	\end{enumerate}
	
	\begin{tcolorbox}[colback=green!5, colframe=green!50!black, title=Пояснение]
	Эти свойства выражают характеристическое свойство НОД как наибольшего по делимости общего делителя, и НОК как наименьшего по делимости общего кратного. Каноническое разложение делает их очевидными: показатель в НОД — минимум, поэтому любой общий делитель имеет показатели не больше минимумов; показатель в НОК — максимум, поэтому любое общее кратное имеет показатели не меньше максимумов
	\end{tcolorbox}
	
	\subsection{Основное соотношение между НОД и НОК}
	
	\noindent \textbf{Следствие.}
	При $n > 0$, $m > 0$ выполнено соотношение:
	\[
	\NOD(n, m) \cdot \NOK(n, m) = n m
	\]
	
	\medskip
	\noindent \textbf{Доказательство.}
	Из канонических разложений:
	\[
	\NOD(n, m) \cdot \NOK(n, m) = p_1^{\min(\alpha_1,\beta_1) + \max(\alpha_1,\beta_1)} \dots p_k^{\min(\alpha_k,\beta_k) + \max(\alpha_k,\beta_k)}
	\]
	Но $\min(\alpha_i, \beta_i) + \max(\alpha_i, \beta_i) = \alpha_i + \beta_i$, поэтому произведение равно $p_1^{\alpha_1+\beta_1} \dots p_k^{\alpha_k+\beta_k} = n m$
	
	\begin{tcolorbox}[colback=gray!5, colframe=gray!50, title=Пример]
	Для $n=12=2^2\cdot 3^1$, $m=18=2^1\cdot 3^2$:
	$\NOD=2^1\cdot 3^1=6$, $\NOK=2^2\cdot 3^2=36$, $6\cdot 36=216=12\cdot 18$
	\end{tcolorbox}
	
	\subsection{НОК взаимно простых чисел}
	
	\noindent \textbf{Следствие.}
	Если $n$, $m$ взаимно просты, то
	\[
	\NOK(n, m) = n \cdot m
	\]
	
	\medskip
	\noindent \textbf{Доказательство.}
	Из основного соотношения: $\NOD(n, m) \cdot \NOK(n, m) = n m$. Так как $\NOD(n, m) = 1$, получаем $\NOK(n, m) = n m$
	
	\begin{tcolorbox}[colback=gray!5, colframe=gray!50, title=Примеры]
	\begin{itemize}
	\item $[7, 11] = 77 = 7 \cdot 11$
	\item $[8, 15] = 120 = 8 \cdot 15$
	\item $[6, 35] = 210 = 6 \cdot 35$
	\end{itemize}
	\end{tcolorbox}
	
	\myremark
	Основная теорема арифметики является фундаментом всей элементарной теории чисел. Из неё следуют все свойства делимости, включая свойства НОД и НОК, лемму Евклида и другие важные факты
	
	\begin{tcolorbox}[colback=purple!5, colframe=purple!50!black, title=Заметочка: История]
	Основная теорема арифметики была известна ещё Евклиду, хотя в явном виде сформулирована позже. Её доказательство опирается на лемму Евклида (если простое число делит произведение, то оно делит один из сомножителей)
	\end{tcolorbox}
	
	
	
	
	
	
	
	% ==================== РАЗДЕЛ 16: НЕПРЕРЫВНЫЕ ДРОБИ ====================
	\newpage
	\section{Непрерывные дроби. Вычисление подходящих дробей (16 вопрос)}
	
	\subsection{Определение непрерывной дроби}
	
	\noindent \textbf{Определение (Числовая непрерывная дробь).}
	\textit{Числовой непрерывной дробью} называется выражение
	\[
	\frac{a_1}{b_1 + \frac{a_2}{b_2 + \frac{a_3}{b_3 + \dots}}} \tag{1}
	\]
	где $a_1, b_1, a_2, b_2, \dots$ — некоторые отличные от нуля числа.
	
	\begin{tcolorbox}[colback=blue!5, colframe=blue!40!black, title=Замечание об обозначениях]
		Обычно рассматривают непрерывные дроби вида
		\[
		b_0 + \frac{1}{b_1 + \frac{1}{b_2 + \frac{1}{b_3 + \dots}}}
		\]
		где $b_0 \in \mathbb{Z}$, $b_i \in \mathbb{N}$ при $i \geq 1$. Такие дроби называются \textit{правильными} или \textit{обыкновенными непрерывными дробями}. В общем случае числители $a_i$ могут быть произвольными
	\end{tcolorbox}
	
	\subsection{Конечные и бесконечные непрерывные дроби}
	
	\noindent \textbf{Определение (Конечная/бесконечная непрерывная дробь).}
	Непрерывная дробь называется \textit{конечной}, если она имеет лишь конечное число звеньев. В противном случае она называется \textit{бесконечной}
	
	\begin{tcolorbox}[colback=gray!5, colframe=gray!50, title=Примеры]
		\begin{itemize}
			\item $1 + \frac{1}{2 + \frac{1}{3}}$ — конечная непрерывная дробь (три звена)
			\item $1 + \frac{1}{2 + \frac{1}{2 + \frac{1}{2 + \dots}}}$ — бесконечная периодическая непрерывная дробь
		\end{itemize}
	\end{tcolorbox}
	
	\subsection{Подходящие дроби}
	
	\noindent \textbf{Определение (Подходящая дробь).}
	$n$-й \textit{подходящей дробью} непрерывной дроби (1) называется конечная дробь
	\[
	\frac{a_1}{b_1 + \frac{a_2}{b_2 + \dots + \frac{a_n}{b_n}}}
	\]
	получаемая из непрерывной дроби (1) её обрыванием на $n$-ом шаге.
	
	В этом определении предполагается, что число звеньев в непрерывной дроби (1) не меньше $n$
	
	\subsection{Пример построения непрерывной дроби для рационального числа}
	
	\begin{tcolorbox}[colback=gray!5, colframe=gray!50, title=Пример: $\frac{10}{7}$]
		Рассмотрим дробь $\frac{10}{7}$. Наибольшее целое число, не превосходящее эту дробь — это $1$.
		\[
		\frac{10}{7} = 1 + \frac{3}{7}
		\]
		Перевернём дробь $\frac{3}{7}$:
		\[
		\frac{10}{7} = 1 + \frac{1}{\frac{7}{3}}
		\]
		Наибольшее целое число, не превосходящее $\frac{7}{3}$ — это $2$.
		\[
		\frac{10}{7} = 1 + \frac{1}{2 + \frac{1}{3}}
		\]
		Таким образом, $\frac{10}{7} = [1; 2, 3]$ — непрерывная дробь для $\frac{10}{7}$
	\end{tcolorbox}
	
	\subsection{Связь с дробно-линейными преобразованиями}
	
	\noindent \textbf{Композиция дробно-линейных преобразований.}
	Пусть $(S_1 \circ \dots \circ S_n)(w)$ — композиция дробно-линейных преобразований, где
	\[
	S_n(w) = \frac{a_n}{b_n + w}, \quad n = 1, 2, \dots
	\]
	Ясно, что $(S_1 \circ \dots \circ S_n)(0)$ равно значению $n$-й подходящей дроби непрерывной дроби (1)
	
	\begin{tcolorbox}[colback=green!5, colframe=green!50!black, title=Пояснение]
		Это важное наблюдение позволяет рассматривать непрерывные дроби как результат применения последовательности дробно-линейных преобразований к нулю. Каждое преобразование добавляет очередной этаж дроби
	\end{tcolorbox}
	
	\subsection{Примеры разложения иррациональных чисел в непрерывные дроби}
	
	\begin{tcolorbox}[colback=gray!5, colframe=gray!50, title=Пример: $\sqrt{2}$]
		\begin{align*}
			\sqrt{2} &= 1 + (\sqrt{2} - 1) = 1 + \frac{1}{\sqrt{2} + 1} \\
			&= 1 + \frac{1}{2 + (\sqrt{2} - 1)} = 1 + \frac{1}{2 + \frac{1}{\sqrt{2} + 1}} \\
			&= 1 + \frac{1}{2 + \frac{1}{2 + \frac{1}{\sqrt{2} + 1}}} = \dots
		\end{align*}
		Таким образом, $\sqrt{2} = [1; 2, 2, 2, \dots]$ — бесконечная периодическая непрерывная дробь
	\end{tcolorbox}
	
	\begin{tcolorbox}[colback=gray!5, colframe=gray!50, title=Пример: $\frac{\sqrt{5} - 1}{2}$ (золотое сечение)]
		\begin{align*}
			\frac{\sqrt{5} - 1}{2} &= 0 + \frac{1}{\frac{2}{\sqrt{5} - 1}} = \frac{1}{\frac{2(\sqrt{5} + 1)}{4}} = \frac{1}{1 + \frac{\sqrt{5} + 1}{2} - 1} = \frac{1}{1 + \frac{\sqrt{5} - 1}{2}} \\
			&= \frac{1}{1 + \frac{1}{\frac{2}{\sqrt{5} - 1}}} = \frac{1}{1 + \frac{1}{\frac{2(\sqrt{5} + 1)}{4}}} = \frac{1}{1 + \frac{1}{1 + \frac{\sqrt{5} + 1}{2} - 1}} \\
			&= \frac{1}{1 + \frac{1}{1 + \frac{\sqrt{5} - 1}{2}}} = \dots
		\end{align*}
		Таким образом, $\frac{\sqrt{5} - 1}{2} = [0; 1, 1, 1, \dots]$ — тоже периодическая дробь (чисто периодическая)
	\end{tcolorbox}
	
	\myremark
	Число $\frac{\sqrt{5} - 1}{2}$ связано с золотым сечением $\varphi = \frac{1+\sqrt{5}}{2}$: это $\varphi - 1 = \frac{1}{\varphi}$. Его непрерывная дробь состоит из одних единиц
	
	\subsection{Основная теорема о представлении чисел непрерывными дробями}
	
	\noindent \textbf{Теорема.}
	Рациональному числу $\frac{p}{q}$ ставится в соответствие \textit{конечная} непрерывная дробь $[b_0; b_1, b_2, \dots, b_n]$. Иррациональному числу $\alpha$ ставится в соответствие \textit{бесконечная} непрерывная дробь $[b_0; b_1, b_2, \dots, b_n, \dots]$. При этом знак соответствия в случае рационального числа можно заменить знаком равенства ($\alpha \sim [b_0; b_1, b_2, \dots]$; $b_0 \in \mathbb{Z}$)
	
	\begin{tcolorbox}[colback=blue!5, colframe=blue!40!black, title=Алгоритм разложения]
		Всякому положительному вещественному числу $\alpha$ ставится в соответствие непрерывная дробь с натуральными коэффициентами по следующему алгоритму:
		\[
		\alpha = [\alpha] + \{\alpha\} = [\alpha] + \frac{1}{\frac{1}{\{\alpha\}}} = \dots
		\]
		То есть:
		\begin{enumerate}
			\item $b_0 = [\alpha]$ (целая часть);
			\item $\alpha_1 = \frac{1}{\{\alpha\}}$;
			\item $b_1 = [\alpha_1]$;
			\item $\alpha_2 = \frac{1}{\{\alpha_1\}}$;
			\item и так далее.
		\end{enumerate}
		Процесс обрывается тогда и только тогда, когда $\alpha$ рационально.
	\end{tcolorbox}
	
	\begin{tcolorbox}[colback=red!5, colframe=red!70!black, title=Итог теоремы]
		\begin{itemize}
			\item Рациональные числа $\longleftrightarrow$ конечные непрерывные дроби
			\item Иррациональные числа $\longleftrightarrow$ бесконечные непрерывные дроби
			\item Подходящие дроби $[b_0; b_1, \dots, b_n]$ дают последовательные рациональные приближения к числу
		\end{itemize}
	\end{tcolorbox}
	
	\myremark
	Непрерывные дроби дают наилучшие приближения действительных чисел рациональными. Для квадратичных иррациональностей (корней квадратных уравнений с целыми коэффициентами) непрерывные дроби являются периодическими (теорема Лагранжа)
	
	\begin{tcolorbox}[colback=purple!5, colframe=purple!50!black, title=Заметочка: Применения непрерывных дробей]
		Непрерывные дроби используются в:
		\begin{itemize}
			\item теории приближений (наилучшие рациональные приближения);
			\item решении диофантовых уравнений (уравнение Пелля);
			\item теории чисел (квадратичные иррациональности);
			\item криптографии (атаки на RSA с помощью цепных дробей)
		\end{itemize}
	\end{tcolorbox}
	
	
	
	% ==================== РАЗДЕЛ 17: СРАВНЕНИЯ ====================
	\newpage
	\section{Сравнения. Свойства сравнений (17 вопрос)}
	
	\subsection{Кольцо классов вычетов}
	
	\noindent Пусть $m$ — фиксированное натуральное число, $m > 1$. Множество $m\mathbb{Z}$, очевидно, замкнуто не только относительно операции сложения, но и относительно операции умножения и удовлетворяет всем аксиомам кольца
	
	\begin{tcolorbox}[colback=blue!5, colframe=blue!40!black, title=Замечание]
	Кольцо $m\mathbb{Z}$ состоит из всех чисел, кратных $m$. Оно является идеалом в кольце $\mathbb{Z}$. Факторкольцо $\mathbb{Z}/m\mathbb{Z}$ — это кольцо классов вычетов по модулю $m$
	\end{tcolorbox}
	
	\subsection{Определение сравнения}
	
	\noindent \textbf{Определение (Сравнимые числа, модуль сравнения).}
	Два целых числа $n$, $n'$ называются \textit{сравнимыми по модулю $m$}, если при делении на $m$ они дают одинаковые остатки. При этом пишут $n \equiv n' \ (m)$ или $n \equiv n' \pmod{m}$, а число $m$ называют \textit{модулем сравнения}
	
	\begin{tcolorbox}[colback=gray!5, colframe=gray!50, title=Примеры]
	\begin{itemize}
	\item $17 \equiv 5 \pmod{12}$, так как $17 = 12 \cdot 1 + 5$, $5 = 12 \cdot 0 + 5$ (остаток $5$)
	\item $100 \equiv 2 \pmod{7}$, так как $100 = 7 \cdot 14 + 2$, $2 = 7 \cdot 0 + 2$
	\item $-3 \equiv 4 \pmod{7}$, так как $-3 = 7 \cdot (-1) + 4$, $4 = 7 \cdot 0 + 4$ (остаток должен быть неотрицательным)
	\item $25 \equiv 0 \pmod{5}$, так как $25$ делится на $5$ (остаток $0$)
	\end{itemize}
	\end{tcolorbox}
	
	\myremark
	Сравнение $a \equiv b \pmod{m}$ равносильно тому, что $a - b$ делится на $m$. Это эквивалентное определение часто удобнее в использовании
	
	\subsection{Свойства сравнений}
	
	\noindent \textbf{Свойство 1.}
	Для любых целых чисел $a, b, c, m$ ($m > 1$):
	\[
	ac \equiv bc \pmod{m} \iff a \equiv b \pmod{\frac{m}{\NOD(c,m)}}
	\]
	В частности, если $c$, $m$ взаимно просты, то
	\[
	ac \equiv bc \pmod{m} \iff a \equiv b \pmod{m}
	\]
	
	\medskip
	\noindent \textbf{Доказательство.}
	Пусть $d = \NOD(c, m)$, $c = d c_1$, $m = d m_1$, тогда $c_1$, $m_1$ взаимно просты.
	\[
	ac \equiv bc \pmod{m} \iff ac - bc \vdots m \iff (a - b)c = m x
	\]
	\[
	\iff (a - b) d c_1 = d m_1 x \iff (a - b) c_1 = m_1 x
	\]
	Поскольку $c_1$ и $m_1$ взаимно просты, из $(a - b) c_1 \vdots m_1$ следует $a - b \vdots m_1$ (по лемме о делении произведения: если $c_1$ и $m_1$ взаимно просты и $m_1 \mid (a-b)c_1$, то $m_1 \mid a-b$). Таким образом,
	\[
	a - b \vdots m_1 \iff a \equiv b \pmod{m_1}, \text{ где } m_1 = \frac{m}{\NOD(c, m)}
	\]
	
	\begin{tcolorbox}[colback=blue!5, colframe=blue!40!black, title=Пример]
	Рассмотрим сравнение $6x \equiv 12 \pmod{15}$.
	Здесь $c = 6$, $m = 15$, $\NOD(6,15) = 3$, $m_1 = 15/3 = 5$.
	Сокращаем на $6$ с учётом свойства:
	\[
	6x \equiv 12 \pmod{15} \iff x \equiv 2 \pmod{5}
	\]
	Действительно, решения: $x = 2, 7, 12, \dots$ (проверка: $6\cdot2=12$, $6\cdot7=42\equiv12$, $6\cdot12=72\equiv12$ по модулю $15$).
	\end{tcolorbox}
	
	\begin{tcolorbox}[colback=green!5, colframe=green!50!black, title=Пояснение]
	Это свойство — аналог сокращения в обычных уравнениях, но с осторожностью: сокращать на число, не взаимно простое с модулем, можно только с изменением модуля. Если же число взаимно просто с модулем, то сокращение происходит без изменения модуля
	\end{tcolorbox}
	
	\medskip
	\noindent \textbf{Свойство 2.}
	Если $m$, $n$ взаимно просты, то
	\[
	a \equiv a' \pmod{mn} \iff 
	\begin{cases}
	a \equiv a' \pmod{m} \\
	a \equiv a' \pmod{n}
	\end{cases}
	\]
	
	\medskip
	\noindent \textbf{Доказательство.}
	
	\textit{Прямое.} Если $a \equiv a' \pmod{mn}$, то $a - a'$ делится на $mn$, следовательно, делится и на $m$, и на $n$. Значит, $a \equiv a' \pmod{m}$ и $a \equiv a' \pmod{n}$
	
	\textit{Обратное.} Если $a \equiv a' \pmod{m}$ и $a \equiv a' \pmod{n}$, то $a - a'$ делится на $m$ и на $n$. Поскольку $m$ и $n$ взаимно просты, то $a - a'$ делится на их произведение $mn$ (свойство 5 из раздела 12: если число делится на два взаимно простых числа, то оно делится на их произведение). Следовательно, $a \equiv a' \pmod{mn}$
	
	\begin{tcolorbox}[colback=blue!5, colframe=blue!40!black, title=Пример]
	Рассмотрим $a = 23$, $a' = 5$, $m = 6$, $n = 7$ (взаимно просты).
	$23 \equiv 5 \pmod{6}$? $23-5=18$, $18 \vdots 6$ — да.
	$23 \equiv 5 \pmod{7}$? $18 \vdots 7$? нет, $18$ не делится на $7$. Значит, по модулю $42$ сравнения нет
	
	Возьмём $a = 23$, $a' = 5$, $m = 6$, $n = 9$ (не взаимно просты).
	$23 \equiv 5 \pmod{6}$ — да, $23 \equiv 5 \pmod{9}$? $18 \vdots 9$ — да, но по модулю $54$ сравнения нет, так как условие на взаимную простоту не выполнено (и действительно, $18$ не делится на $54$)
	\end{tcolorbox}
	
	\myremark
	Свойство 2 показывает, что сравнение по составному модулю эквивалентно системе сравнений по взаимно простым множителям этого модуля. Это лежит в основе китайской теоремы об остатках
	
	\begin{tcolorbox}[colback=purple!5, colframe=purple!50!black, title=Заметочка: Обобщение]
	Для попарно взаимно простых модулей $m_1, m_2, \dots, m_k$:
	\[
	a \equiv b \pmod{m_1 m_2 \dots m_k} \iff a \equiv b \pmod{m_i} \text{ для всех } i = 1, \dots, k
	\]
	Это свойство широко используется в теории чисел
	\end{tcolorbox}
	
	\begin{tcolorbox}[colback=yellow!5, colframe=orange!60!black, title=Важное замечание]
	Сравнения обладают также очевидными свойствами, аналогичными свойствам равенств:
	\begin{itemize}
	\item Рефлексивность: $a \equiv a \pmod{m}$
	\item Симметричность: если $a \equiv b \pmod{m}$, то $b \equiv a \pmod{m}$
	\item Транзитивность: если $a \equiv b \pmod{m}$ и $b \equiv c \pmod{m}$, то $a \equiv c \pmod{m}$
	\item Если $a \equiv b \pmod{m}$ и $c \equiv d \pmod{m}$, то $a \pm c \equiv b \pm d \pmod{m}$ и $ac \equiv bd \pmod{m}$
	\end{itemize}
	Эти свойства следуют непосредственно из определения $a \equiv b \pmod{m} \iff a-b \vdots m$
	\end{tcolorbox}
	
	
	
	
	
	
	% ==================== РАЗДЕЛ 18: ПОЛНАЯ СИСТЕМА ВЫЧЕТОВ ====================
	\newpage
	\section{Полная система вычетов. Кольцо классов вычетов (корректность операций) (18 вопрос)}
	
	\subsection{Определение сравнимых чисел}
	
	\noindent \textbf{Определение (Сравнимые числа).}
	Пусть $a, a', m \in \mathbb{Z}$; $m > 1$. Говорят, что $a$ \textit{сравнимо} с $a'$ по модулю $m$ (и пишут $a \equiv a' \pmod{m}$), если $a - a' \vdots m$. Запись $a \equiv a' \pmod{m}$ называют \textit{сравнением} (по модулю $m$)
	
	\begin{tcolorbox}[colback=blue!5, colframe=blue!40!black, title=Пример]
		$17 \equiv 5 \pmod{12}$, так как $17-5 = 12 \vdots 12$
		$100 \equiv 2 \pmod{7}$, так как $100-2 = 98 = 7 \cdot 14$
		$-3 \equiv 4 \pmod{7}$, так как $-3-4 = -7 \vdots 7$
	\end{tcolorbox}
	
	\subsection{Сравнимость как отношение эквивалентности}
	
	\noindent \textbf{Определение (Сравнимость, класс вычетов).}
	Данное определение для каждого фиксированного $m$ задаёт в множестве $\mathbb{Z}$ бинарное отношение — \textit{сравнимость} (по модулю $m$). Ясно, что сравнимость рефлексивна, симметрична и транзитивна и, значит, разбивает $\mathbb{Z}$ на классы эквивалентности. Эти классы называются \textit{классами вычетов} (по модулю $m$). В каждом классе сравнимости по модулю $m$ можно выбрать представителя
	
	\begin{tcolorbox}[colback=green!5, colframe=green!50!black, title=Проверка свойств отношения сравнимости]
		\begin{itemize}
			\item \textbf{Рефлексивность:} $a - a = 0 \vdots m$, значит $a \equiv a \pmod{m}$
			\item \textbf{Симметричность:} если $a \equiv b \pmod{m}$, то $a-b \vdots m$, тогда $b-a = -(a-b) \vdots m$, значит $b \equiv a \pmod{m}$
			\item \textbf{Транзитивность:} если $a \equiv b \pmod{m}$ и $b \equiv c \pmod{m}$, то $a-b \vdots m$ и $b-c \vdots m$, тогда $a-c = (a-b)+(b-c) \vdots m$, значит $a \equiv c \pmod{m}$
		\end{itemize}
	\end{tcolorbox}
	
	\myremark
	Класс вычетов, содержащий число $a$, обычно обозначают $\bar{a}$ или $[a]_m$. Он состоит из всех чисел, сравнимых с $a$ по модулю $m$:
	\[
	\bar{a} = \{x \in \mathbb{Z} \mid x \equiv a \pmod{m}\} = a + m\mathbb{Z}.
	\]
	
	\subsection{Полная система вычетов}
	
	\noindent \textbf{Определение (Полная система вычетов).}
	Множество представителей, выбранных по одному из каждого класса вычетов, называют \textit{полной системой вычетов} (по данному модулю)
	
	\begin{tcolorbox}[colback=gray!5, colframe=gray!50, title=Примеры полных систем вычетов по модулю 5]
		\begin{itemize}
			\item $\{0, 1, 2, 3, 4\}$ — наименьшие неотрицательные вычеты
			\item $\{-2, -1, 0, 1, 2\}$ — симметричная система (абсолютно наименьшие вычеты)
			\item $\{5, 6, 7, 8, 9\}$ — тоже полная система (сдвиг на 5)
			\item $\{0, 1, 2, 3, 4\}$ — наиболее часто используемая
		\end{itemize}
	\end{tcolorbox}
	
	\myremark
	Всего существует ровно $m$ классов вычетов по модулю $m$, поэтому любая полная система вычетов состоит из $m$ чисел, попарно не сравнимых по модулю $m$
	
	\subsection{Корректность операций на классах вычетов}
	
	\noindent \textbf{Свойство (Согласованность сравнений с арифметическими операциями).}
	Если $a \equiv a' \pmod{m}$ и $b \equiv b' \pmod{m}$, то
	\[
	a \pm b \equiv a' \pm b' \pmod{m}, \quad ab \equiv a'b' \pmod{m}
	\]
	(сравнения по одному модулю можно почленно складывать, вычитать и перемножать)
	
	\medskip
	\noindent \textbf{Доказательство.}
	Из $a \equiv a'$ и $b \equiv b'$ следует $a - a' \vdots m$ и $b - b' \vdots m$.
	Тогда $(a \pm b) - (a' \pm b') = (a - a') \pm (b - b') \vdots m$, значит $a \pm b \equiv a' \pm b' \pmod{m}$.
	Для умножения: $ab - a'b' = ab - ab' + ab' - a'b' = a(b - b') + (a - a')b' \vdots m$, так как каждое слагаемое делится на $m$
	
	\begin{tcolorbox}[colback=blue!5, colframe=blue!40!black, title=Пример]
		По модулю $7$: $15 \equiv 1$, $22 \equiv 1$. Тогда $15+22 = 37 \equiv 1+1 = 2 \pmod{7}$ (проверка: $37-35=2$). $15 \cdot 22 = 330$, $1 \cdot 1 = 1$, $330 \equiv 1 \pmod{7}$? $7 \cdot 47 = 329$, остаток $1$ — верно
	\end{tcolorbox}
	
	\myremark
	Другими словами, сравнимость согласована со сложением, вычитанием и умножением, т.е. является \textit{конгруэнцией} в кольце $\mathbb{Z}$
	
	\subsection{Построение кольца классов вычетов}
	
	Согласованность сравнений с операциями позволяет перенести действия кольца на множество классов конгруэнции, т.е. классов вычетов
	
	\noindent \textbf{Определение операций на классах вычетов.}
	Для классов вычетов $\bar{a}$ и $\bar{b}$ (по модулю $m$) определим:
	\[
	\bar{a} + \bar{b} = \overline{a + b}, \quad \bar{a} \cdot \bar{b} = \overline{ab}.
	\]
	
	\noindent \textbf{Проверка корректности (независимости от выбора представителей).}
	Пусть $\bar{a} = \bar{a'}$ и $\bar{b} = \bar{b'}$. Тогда $a \equiv a' \pmod{m}$ и $b \equiv b' \pmod{m}$. По доказанному выше,
	\[
	a + b \equiv a' + b' \pmod{m} \Rightarrow \overline{a + b} = \overline{a' + b'},
	\]
	\[
	ab \equiv a'b' \pmod{m} \Rightarrow \overline{ab} = \overline{a'b'}.
	\]
	Таким образом, класс-результат не зависит от выбора представителей классов-операндов
	
	Ясно, что эти действия в множестве классов вычетов всюду определены
	
	\subsection{Структура кольца классов вычетов}
	
	\noindent \textbf{Теорема.}
	Введённые выше операции задают на множестве классов вычетов по модулю $m$ структуру коммутативного кольца с единицей
	
	\medskip
	\noindent \textbf{Проверка аксиом кольца.}
	Ясно, что все аксиомы кольца выполняются для множества классов вычетов:
	\begin{itemize}
		\item $\bar{0}$ — нуль сложения классов: $\bar{a} + \bar{0} = \overline{a+0} = \bar{a}$
		\item $-\bar{a} = \overline{-a}$ — противоположный элемент
		\item $\bar{1}$ — единица умножения классов: $\bar{a} \cdot \bar{1} = \overline{a \cdot 1} = \bar{a}$
		\item Коммутативность сложения и умножения, ассоциативность и дистрибутивность наследуются из соответствующих свойств целых чисел
	\end{itemize}
	
	\begin{tcolorbox}[colback=red!5, colframe=red!70!black, title=Итог]
		Таким образом, введённые нами действия задают на множестве классов вычетов по модулю $m$ структуру коммутативного кольца с единицей. Это кольцо называется \textit{кольцом вычетов} (или \textit{кольцом классов вычетов}) по модулю $m$ и обозначается $\mathbb{Z}_m$ (или $\mathbb{Z}/m\mathbb{Z}$, $\mathbb{Z}/(m)$).
	\end{tcolorbox}
	
	\begin{tcolorbox}[colback=gray!5, colframe=gray!50, title=Пример: $\mathbb{Z}_6$]
		Классы вычетов по модулю $6$: $\bar{0}, \bar{1}, \bar{2}, \bar{3}, \bar{4}, \bar{5}$.
		Таблицы сложения и умножения:
		\begin{align*}
			\bar{2} + \bar{3} &= \bar{5} \\
			\bar{2} + \bar{4} &= \bar{0} \quad (\text{так как } 2+4=6 \equiv 0) \\
			\bar{2} \cdot \bar{3} &= \bar{0} \quad (\text{делители нуля!}) \\
			\bar{2} \cdot \bar{5} &= \bar{4} \quad (2\cdot5=10 \equiv 4)
		\end{align*}
	\end{tcolorbox}
	
	\myremark
	Кольцо $\mathbb{Z}_m$ является полем тогда и только тогда, когда $m$ — простое число. При составном $m$ в нём есть делители нуля (как, например, $\bar{2}$ и $\bar{3}$ в $\mathbb{Z}_6$)
	
	\begin{tcolorbox}[colback=purple!5, colframe=purple!50!black, title=Заметочка: Обозначения]
		В разных учебниках используются различные обозначения: $\mathbb{Z}_m$, $\mathbb{Z}/m\mathbb{Z}$, $\mathbb{Z}/(m)$. Первое обозначение иногда путают с кольцом $m$-адических чисел, поэтому в серьёзной литературе чаще используют $\mathbb{Z}/m\mathbb{Z}$
	\end{tcolorbox}
	
	
	
	
	% ==================== РАЗДЕЛ 19: ПРИВЕДЁННАЯ СИСТЕМА ВЫЧЕТОВ ====================
	\newpage
	\section{Приведённая система вычетов. Функция Эйлера. Вычисление функции Эйлера (без доказательства мультипликативности) (19 вопрос)}
	
	\subsection{Приведённая система вычетов}
	
	\noindent \textbf{Определение (Приведённая система вычетов).}
	Группа $(\mathbb{Z}/m\mathbb{Z})^*$ состоит из всех тех классов вычетов, чьи представители взаимно просты с модулем $m$; мы будем говорить, что такой класс \textit{взаимно прост с модулем}
	
	Множество представителей, выбранных по одному из каждого такого класса, называют \textit{приведённой системой вычетов} (по данному модулю). Таким образом, $\{a_1, a_2, \dots\}$ тогда и только тогда есть приведённая система вычетов по модулю $m$, когда $\{\bar{a}_1, \bar{a}_2, \dots\} = (\mathbb{Z}/m\mathbb{Z})^*$
	
	\begin{tcolorbox}[colback=blue!5, colframe=blue!40!black, title=Примеры]
		\begin{itemize}
			\item По модулю $6$: числа, взаимно простые с $6$ в пределах от $1$ до $6$: $1$ и $5$. Приведённая система вычетов: $\{1, 5\}$
			\item По модулю $7$ (простое): все ненулевые классы взаимно просты с $7$. Приведённая система: $\{1,2,3,4,5,6\}$ (совпадает с полной системой без нуля)
			\item По модулю $8$: числа, взаимно простые с $8$: $1,3,5,7$. Приведённая система: $\{1,3,5,7\}$
			\item По модулю $9$: числа, не кратные $3$: $1,2,4,5,7,8$. Приведённая система: $\{1,2,4,5,7,8\}$
		\end{itemize}
	\end{tcolorbox}
	
	\myremark
	Приведённая система вычетов состоит из $\varphi(m)$ элементов, где $\varphi$ — функция Эйлера. Она играет ключевую роль в теории чисел, так как классы вычетов, взаимно простые с модулем, образуют мультипликативную группу
	
	\subsection{Функция Эйлера}
	
	\noindent \textbf{Определение (Функция Эйлера).}
	\textit{Функцией Эйлера} называется отображение $\varphi \colon \mathbb{N} \to \mathbb{N}$ такое, что $\forall n \in \mathbb{N}$ $\varphi(n)$ есть количество натуральных чисел, не превосходящих $n$ и взаимно простых с $n$
	
	\begin{tcolorbox}[colback=gray!5, colframe=gray!50, title=Примеры]
		\begin{itemize}
			\item $\varphi(1) = 1$ (число $1$ взаимно просто с $1$)
			\item $\varphi(2) = 1$ (число $1$)
			\item $\varphi(3) = 2$ (числа $1,2$)
			\item $\varphi(4) = 2$ (числа $1,3$)
			\item $\varphi(5) = 4$ (числа $1,2,3,4$)
			\item $\varphi(6) = 2$ (числа $1,5$)
			\item $\varphi(7) = 6$ (все числа от $1$ до $6$)
			\item $\varphi(8) = 4$ (числа $1,3,5,7$)
			\item $\varphi(9) = 6$ (числа $1,2,4,5,7,8$)
			\item $\varphi(10) = 4$ (числа $1,3,7,9$)
		\end{itemize}
	\end{tcolorbox}
	
	\subsection{Свойство мультипликативности (без доказательства)}
	
	\noindent \textbf{Свойство.}
	Функция Эйлера \textit{мультипликативна}, т.е. для любой взаимно простой пары $a, b \in \mathbb{N}$
	\[
	\varphi(ab) = \varphi(a) \varphi(b)
	\]
	
	\begin{tcolorbox}[colback=green!5, colframe=green!50!black, title=Пояснение]
		Мультипликативность означает, что значение функции на произведении взаимно простых чисел равно произведению значений на каждом сомножителе. Это свойство позволяет вычислять функцию Эйлера для любого числа, зная её значения для степеней простых чисел
	\end{tcolorbox}
	
	\begin{tcolorbox}[colback=blue!5, colframe=blue!40!black, title=Пример]
		$\varphi(6) = \varphi(2 \cdot 3) = \varphi(2) \cdot \varphi(3) = 1 \cdot 2 = 2$ — совпадает с вычисленным ранее.
		$\varphi(12) = \varphi(3 \cdot 4) = \varphi(3) \cdot \varphi(4) = 2 \cdot 2 = 4$ (действительно, числа $1,5,7,11$).
	\end{tcolorbox}
	
	\subsection{Вычисление функции Эйлера}
	
	\noindent \textbf{Теорема (Формула для функции Эйлера).}
	Если $n = p_1^{k_1} p_2^{k_2} \dots p_m^{k_m}$ — каноническое разложение числа $n \in \mathbb{N}$, то
	\[
	\varphi(n) = p_1^{k_1-1}(p_1 - 1) \cdot p_2^{k_2-1}(p_2 - 1) \cdots p_m^{k_m-1}(p_m - 1)
	\]
	
	\medskip
	\noindent \textbf{Вывод для степени простого числа.}
	Если $p \in \mathbb{N}$ простое, то не взаимно просты с $p^k$ те и только те числа, которые кратны $p$. Ясно, что в интервале $[1, p^k]$ имеется $p^{k-1}$ таких чисел, и поэтому
	\[
	\varphi(p^k) = p^k - p^{k-1} = p^{k-1}(p - 1)
	\]
	
	\begin{tcolorbox}[colback=red!5, colframe=red!70!black, title=Доказательство для степени простого]
		Среди чисел от $1$ до $p^k$ на $p$ делятся числа: $p, 2p, 3p, \dots, p^{k-1} \cdot p = p^k$. Всего их $p^{k-1}$. Все остальные числа взаимно просты с $p^k$. Следовательно,
		\[
		\varphi(p^k) = p^k - p^{k-1} = p^{k-1}(p-1).
		\]
	\end{tcolorbox}
	
	\begin{tcolorbox}[colback=blue!5, colframe=blue!40!black, title=Примеры]
		\begin{itemize}
			\item $\varphi(2^3) = \varphi(8) = 2^{3-1}(2-1) = 2^2 \cdot 1 = 4$ (числа $1,3,5,7$)
			\item $\varphi(3^2) = \varphi(9) = 3^{2-1}(3-1) = 3^1 \cdot 2 = 6$ (числа $1,2,4,5,7,8$)
			\item $\varphi(5^3) = \varphi(125) = 5^{2} \cdot 4 = 25 \cdot 4 = 100$
		\end{itemize}
	\end{tcolorbox}
	
	\subsection{Примеры вычисления для составных чисел}
	
	Используя мультипликативность и формулу для степени простого, получаем общую формулу
	
	\begin{tcolorbox}[colback=gray!5, colframe=gray!50, title=Примеры]
		\begin{itemize}
			\item $n = 12 = 2^2 \cdot 3^1$:
			\[
			\varphi(12) = 2^{2-1}(2-1) \cdot 3^{1-1}(3-1) = 2^1 \cdot 1 \cdot 3^0 \cdot 2 = 2 \cdot 1 \cdot 1 \cdot 2 = 4
			\]
			\item $n = 30 = 2 \cdot 3 \cdot 5$:
			\[
			\varphi(30) = 2^{0}(2-1) \cdot 3^{0}(3-1) \cdot 5^{0}(5-1) = 1 \cdot 1 \cdot 2 \cdot 1 \cdot 4 = 8
			\]
			(числа $1,7,11,13,17,19,23,29$ — действительно 8 чисел)
			\item $n = 100 = 2^2 \cdot 5^2$:
			\[
			\varphi(100) = 2^{1}(2-1) \cdot 5^{1}(5-1) = 2 \cdot 1 \cdot 5 \cdot 4 = 40
			\]
			\item $n = 36 = 2^2 \cdot 3^2$:
			\[
			\varphi(36) = 2^{1} \cdot 1 \cdot 3^{1} \cdot 2 = 2 \cdot 3 \cdot 2 = 12
			\]
		\end{itemize}
	\end{tcolorbox}
	
	\myremark
	Функция Эйлера имеет огромное значение в теории чисел. Она используется в теореме Эйлера ($a^{\varphi(m)} \equiv 1 \pmod{m}$ для $a$, взаимно простых с $m$), в RSA-криптосистеме, в различных комбинаторных задачах
	
	\begin{tcolorbox}[colback=purple!5, colframe=purple!50!black, title=Заметочка: Связь с приведённой системой]
		Количество элементов в приведённой системе вычетов по модулю $m$ равно $\varphi(m)$. Действительно, приведённая система выбирает по одному представителю из каждого класса, взаимно простого с модулем, а таких классов ровно $\varphi(m)$
	\end{tcolorbox}
	
	\begin{tcolorbox}[colback=yellow!5, colframe=orange!60!black, title=Важное замечание]
		Формула $\varphi(p^k) = p^k - p^{k-1}$ имеет простой комбинаторный смысл: из всех $p^k$ чисел от $1$ до $p^k$ нужно выбросить те, которые делятся на $p$ (их $p^{k-1}$). Для составного числа используется мультипликативность, которая здесь принимается без доказательства
	\end{tcolorbox}
	
	
	
	
	% ==================== РАЗДЕЛ 20: СИСТЕМА СРАВНЕНИЙ ====================
	\newpage
	\section{Система сравнений. Китайская теорема об остатках (20 вопрос)}
	
	\subsection{Китайская теорема об остатках}
	
	\noindent \textbf{Теорема (Китайская теорема об остатках).}
	Пусть $(m_i)_{i \in I}$ — конечное семейство попарно простых натуральных чисел. Тогда для любого семейства $(a_i)_{i \in I}$ существует $x \in \mathbb{Z}$ такое, что
	\[
	\begin{cases}
		x \equiv a_1 \pmod{m_1} \\
		x \equiv a_2 \pmod{m_2} \\
		\cdots\cdots\cdots\cdots\cdots\cdots \\
		x \equiv a_k \pmod{m_k}
	\end{cases}
	\qquad (|I| = k)
	\]
	и это решение единственно по модулю $m_1 m_2 \dots m_k$.
	
	\begin{tcolorbox}[colback=red!5, colframe=red!70!black, title=Формулировка]
		Система линейных сравнений с попарно взаимно простыми модулями всегда имеет решение, и все решения сравнимы между собой по модулю произведения модулей
	\end{tcolorbox}
	
	\subsection{Доказательство}
	
	\noindent \textbf{Доказательство.}
	Проведём индукцию по $k$.
	
	\textit{База индукции.} При $k = 1$ утверждение очевидно: $x \equiv a_1 \pmod{m_1}$ имеет решение (например, $x = a_1$) и единственно по модулю $m_1$
	
	\textit{Индукционный переход.} Пусть $k > 1$. По предположению индукции система, составленная из первых $k-1$ сравнений, имеет единственное по модулю $m' = m_1 m_2 \dots m_{k-1}$ решение. Пусть это решение есть
	\[
	x \equiv a \pmod{m'}
	\]
	
	Тогда класс вычетов $a \pmod{m'}$ — это множество всех чисел вида
	\[
	x = a + m' y, \quad y \in \mathbb{Z} \tag{1}
	\]
	
	Для нахождения всех решений исходной системы сравнений осталось найти все те значения $y$, при которых числа вида (1) удовлетворяют последнему сравнению исходной системы. Подставим вместо $x$ его выражение (1) и решим полученное сравнение относительно $y$:
	\[
	a + m' y \equiv a_k \pmod{m_k}
	\]
	\[
	m' y \equiv a_k - a \pmod{m_k}
	\]
	
	Так как $m'$ и $m_k$ взаимно просты (по условию попарной простоты всех модулей), то $m'$ обратим по модулю $m_k$. Следовательно, это сравнение имеет единственное решение по модулю $m_k$. Пусть это будет класс
	\[
	y \equiv b \pmod{m_k}, \quad \text{т.е. } y = b + m_k t, \ t \in \mathbb{Z}
	\]
	
	Подставляя это выражение для $y$ в (1), получаем:
	\[
	x = a + m'(b + m_k t) = a + m'b + m'm_k t
	\]
	
	Итак, множество решений исходной системы сравнений совпадает с классом
	\[
	x \equiv a + b m' \pmod{m_1 m_2 \dots m_k}
	\]
	
	\begin{tcolorbox}[colback=green!5, colframe=green!50!black, title=Пояснение]
		Индукция позволяет свести задачу к последовательному решению сравнений с двумя модулями. На каждом шаге мы используем обратимость произведения предыдущих модулей относительно нового модуля
	\end{tcolorbox}
	
	\subsection{Алгоритм решения системы сравнений}
	
	\noindent \textbf{Замечание (Алгоритм).}
	Из доказательства китайской теоремы об остатках виден и алгоритм решения системы сравнений:
	
	\begin{enumerate}
		\item Из первого сравнения находим $x = a_1 + m_1 y$
		\item Подставив $x$ во второе и решив полученное сравнение относительно $y$, получим $y = b_1 + m_2 z$, следовательно $x = a_1 + m_1 b_1 + m_1 m_2 z$
		\item Подставив найденные выражения для $x$ в третье сравнение системы, находим $z$ и т.д.
	\end{enumerate}
	
	Продолжая этот процесс, после $k$ шагов получим решение в виде $x \equiv X \pmod{m_1 m_2 \dots m_k}$.
	
	\begin{tcolorbox}[colback=gray!5, colframe=gray!50, title=Пример]
		Решим систему:
		\[
		\begin{cases}
			x \equiv 2 \pmod{3} \\
			x \equiv 3 \pmod{5} \\
			x \equiv 2 \pmod{7}
		\end{cases}
		\]
		
		\textbf{Шаг 1:} Из первого сравнения $x = 2 + 3y$
		
		\textbf{Шаг 2:} Подставляем во второе: $2 + 3y \equiv 3 \pmod{5} \Rightarrow 3y \equiv 1 \pmod{5}$.
		Умножаем на обратный к $3$ по модулю $5$ (число $2$, так как $3\cdot2=6\equiv1$): $y \equiv 2 \pmod{5}$.
		Значит $y = 2 + 5z$, и $x = 2 + 3(2+5z) = 2 + 6 + 15z = 8 + 15z$
		
		\textbf{Шаг 3:} Подставляем в третье: $8 + 15z \equiv 2 \pmod{7}$.
		$15 \equiv 1 \pmod{7}$, поэтому $8 + z \equiv 2 \pmod{7} \Rightarrow z \equiv 2 - 8 = -6 \equiv 1 \pmod{7}$.
		Значит $z = 1 + 7t$, и $x = 8 + 15(1+7t) = 8 + 15 + 105t = 23 + 105t$
		
		\textbf{Ответ:} $x \equiv 23 \pmod{105}$
		
		Проверка: $23 \div 3 = 7\cdot3 + 2$, $23 \div 5 = 4\cdot5 + 3$, $23 \div 7 = 3\cdot7 + 2$ — верно
	\end{tcolorbox}
	
	\subsection{Пример: линейное сравнение с одним неизвестным}
	
	\noindent \textbf{Пример (Единственность решения сравнения $ax \equiv b \pmod{m}$ при взаимно простых $a$ и $m$).}
	Рассмотрим сравнение
	\[
	ax \equiv b \pmod{m}
	\]
	где $a$ и $m$ взаимно просты.
	
	\textit{Единственность решения.}
	Пусть $x_1$ и $x_2$ — два решения этого сравнения. Тогда
	\[
	ax_1 \equiv b \pmod{m}, \quad ax_2 \equiv b \pmod{m}
	\]
	Вычитая, получаем:
	\[
	a(x_1 - x_2) \equiv 0 \pmod{m}
	\]
	
	Так как $a$ и $m$ взаимно просты, $a$ обратим по модулю $m$. Умножим обе части сравнения на $a^{-1}$ (существующий в силу взаимной простоты):
	\[
	x_1 - x_2 \equiv 0 \pmod{m}
	\]
	Следовательно, $x_1 \equiv x_2 \pmod{m}$.
	
	Более строго:
	\[
	ab_1 \equiv b \pmod{m}, \quad ab_2 \equiv b \pmod{m} \Rightarrow a(b_1 - b_2) \equiv 0 \pmod{m}
	\]
	Поскольку $a \not\equiv 0 \pmod{m}$ (точнее, $a$ не делится на $m$ и взаимно прост с $m$), умножим на $a^{-1}$ и получим $b_1 \equiv b_2 \pmod{m}$
	
	\begin{tcolorbox}[colback=blue!5, colframe=blue!40!black, title=Замечание]
		Этот факт является частным случаем китайской теоремы об остатках для одного сравнения. Он показывает, что при взаимно простых коэффициенте и модуле решение линейного сравнения существует и единственно по модулю $m$
	\end{tcolorbox}
	
	\myremark
	Китайская теорема об остатках имеет множество применений:
	\begin{itemize}
		\item В криптографии (например, в RSA для ускорения вычислений)
		\item В теории чисел для решения систем сравнений
		\item В компьютерной арифметике для представления больших чисел через остатки от деления на взаимно простые модули (система остаточных классов)
	\end{itemize}
	
	\begin{tcolorbox}[colback=purple!5, colframe=purple!50!black, title=Заметочка: Название]
		Теорема называется "китайской", поскольку была известна в древнем Китае. Первая известная формулировка встречается в труде Сунь Цзы "Математические руководства" (III век н.э.), где рассматривалась задача: "Найти число, дающее остатки 2, 3, 2 при делении на 3, 5, 7" (как в нашем примере)
	\end{tcolorbox}
	
	
	
	
	
	
	
	% ==================== РАЗДЕЛ 21: МУЛЬТИПЛИКАТИВНОСТЬ ФУНКЦИИ ЭЙЛЕРА ====================
	\newpage
	\section{Мультипликативность функции Эйлера (21 вопрос)}
	
	\noindent \textbf{Теорема (Мультипликативность функции Эйлера).}
	Функция Эйлера мультипликативна, т.е. для любой взаимно простой пары $a, b \in \mathbb{N}$
	\[
	\varphi(ab) = \varphi(a) \varphi(b)
	\]
	
	\begin{tcolorbox}[colback=red!5, colframe=red!70!black, title=Формулировка]
		Если числа $a$ и $b$ взаимно просты, то количество чисел, взаимно простых с их произведением, равно произведению количеств чисел, взаимно простых с каждым сомножителем.
	\end{tcolorbox}
	
	\subsection{Доказательство}
	
	\noindent \textbf{Доказательство.}
	Пусть $a_1, a_2, \dots, a_{\varphi(a)}$ и $b_1, b_2, \dots, b_{\varphi(b)}$ — приведённые системы вычетов по модулям $a$ и $b$ соответственно. Это означает, что:
	\begin{itemize}
		\item $\{a_1, \dots, a_{\varphi(a)}\}$ — полный набор представителей классов вычетов по модулю $a$, взаимно простых с $a$;
		\item $\{b_1, \dots, b_{\varphi(b)}\}$ — полный набор представителей классов вычетов по модулю $b$, взаимно простых с $b$.
	\end{itemize}
	
	\begin{tcolorbox}[colback=blue!5, colframe=blue!40!black, title=Пример для пояснения]
		Пусть $a = 3$, $b = 5$ (взаимно просты). 
		Приведённая система по модулю $3$: $\{1, 2\}$ ($\varphi(3)=2$).
		Приведённая система по модулю $5$: $\{1, 2, 3, 4\}$ ($\varphi(5)=4$).
		По теореме должно быть $\varphi(15)=2\cdot4=8$.
	\end{tcolorbox}
	
	\medskip
	По Китайской теореме об остатках для каждой пары $(i, j)$ существует единственное по модулю $ab$ число $c_{ij}$ такое, что
	\[
	\begin{cases}
		c_{ij} \equiv a_i \pmod{a} \\
		c_{ij} \equiv b_j \pmod{b}
	\end{cases}
	\]
	
	\begin{tcolorbox}[colback=green!5, colframe=green!50!black, title=Пояснение]
		Китайская теорема гарантирует, что для любого набора остатков по взаимно простым модулям существует число, дающее эти остатки, и все такие числа образуют один класс по модулю $ab$.
	\end{tcolorbox}
	
	\medskip
	\textit{Свойство 1.} Число $c_{ij}$ взаимно просто с $a$ и с $b$.
	Действительно:
	\begin{itemize}
		\item $c_{ij} \equiv a_i \pmod{a}$, а $a_i$ взаимно просто с $a$. Если бы $c_{ij}$ имело общий делитель $d > 1$ с $a$, то $a_i$ также делилось бы на $d$ (поскольку $c_{ij} = a_i + a \cdot t$), что противоречит взаимной простоте $a_i$ и $a$.
		\item Аналогично, $c_{ij} \equiv b_j \pmod{b}$ и $b_j$ взаимно просто с $b$, значит $c_{ij}$ взаимно просто с $b$.
	\end{itemize}
	
	По свойству взаимной простоты (если число взаимно просто с каждым из двух взаимно простых между собой сомножителей, то оно взаимно просто с их произведением), из взаимной простоты $a$ и $b$ следует, что $c_{ij}$ взаимно просто с $ab$.
	
	\textit{Свойство 2.} Все числа $c_{ij}$ лежат в разных классах по модулю $ab$.
	Действительно, если $(i,j) \neq (k,l)$, то либо $i \neq k$, либо $j \neq l$ (или оба случая). Если $i \neq k$, то $c_{ij} \equiv a_i \not\equiv a_k \equiv c_{kl} \pmod{a}$, следовательно $c_{ij} \not\equiv c_{kl} \pmod{ab}$ (так как несравнимость по модулю $a$ влечёт несравнимость по модулю $ab$). Аналогично для случая $j \neq l$.
	
	\textit{Свойство 3.} Любое число $x$, взаимно простое с $ab$, сравнимо с некоторым $c_{ij}$ по модулю $ab$.
	Если $x$ взаимно просто с $ab$, то оно взаимно просто и с $a$, и с $b$ (так как любой общий делитель $x$ и $a$ был бы также делителем $ab$). Следовательно, существуют такие $i$ и $j$, что
	\[
	x \equiv a_i \pmod{a}, \quad x \equiv b_j \pmod{b}
	\]
	(поскольку $\{a_i\}$ и $\{b_j\}$ — полные системы вычетов для взаимно простых классов). Тогда по построению $c_{ij}$ является решением этой же системы сравнений, а значит $x \equiv c_{ij} \pmod{ab}$ (поскольку решение системы по модулю $ab$ единственно).
	
		
	\begin{tcolorbox}[colback=blue!5, colframe=blue!40!black, title=Итог]
		Таким образом, числа $c_{ij}$ ($i = 1,\dots,\varphi(a)$, $j = 1,\dots,\varphi(b)$) представляют все классы вычетов по модулю $ab$, взаимно простые с $ab$, причём каждый класс представлен ровно один раз.
		
		Следовательно, количество таких классов равно количеству чисел $c_{ij}$, т.е. произведению количеств $i$ и $j$:
		\[
		\varphi(ab) = \varphi(a) \cdot \varphi(b)
		\]
	\end{tcolorbox}
	
	\begin{tcolorbox}[colback=gray!5, colframe=gray!50, title={Пример для $a=3, b=5$}]
		Приведённые системы: $a_i = \{1,2\}$, $b_j = \{1,2,3,4\}$.
		По китайской теореме находим $c_{ij}$ (по модулю $15$):
		\begin{align*}
			c_{11}: & \begin{cases} x \equiv 1 \pmod{3} \\ x \equiv 1 \pmod{5} \end{cases} \Rightarrow x = 1 \\
			c_{12}: & \begin{cases} x \equiv 1 \pmod{3} \\ x \equiv 2 \pmod{5} \end{cases} \Rightarrow x = 7 \\
			c_{13}: & \begin{cases} x \equiv 1 \pmod{3} \\ x \equiv 3 \pmod{5} \end{cases} \Rightarrow x = 13 \\
			c_{14}: & \begin{cases} x \equiv 1 \pmod{3} \\ x \equiv 4 \pmod{5} \end{cases} \Rightarrow x = 4 \\
			c_{21}: & \begin{cases} x \equiv 2 \pmod{3} \\ x \equiv 1 \pmod{5} \end{cases} \Rightarrow x = 11 \\
			c_{22}: & \begin{cases} x \equiv 2 \pmod{3} \\ x \equiv 2 \pmod{5} \end{cases} \Rightarrow x = 2 \\
			c_{23}: & \begin{cases} x \equiv 2 \pmod{3} \\ x \equiv 3 \pmod{5} \end{cases} \Rightarrow x = 8 \\
			c_{24}: & \begin{cases} x \equiv 2 \pmod{3} \\ x \equiv 4 \pmod{5} \end{cases} \Rightarrow x = 14
		\end{align*}
		Получили ровно $\varphi(15)=8$ чисел: $\{1,2,4,7,8,11,13,14\}$ — все числа от $1$ до $14$, не кратные $3$ и $5$.
	\end{tcolorbox}
	
	\myremark
	Это доказательство наглядно показывает, как китайская теорема об остатках позволяет установить биекцию между парами остатков по взаимно простым модулям и остатками по произведению. Мультипликативность функции Эйлера является ключевым свойством, используемым при вычислении $\varphi(n)$ через каноническое разложение.
	
	\begin{tcolorbox}[colback=purple!5, colframe=purple!50!black, title=Заметочка]
		Из мультипликативности и формулы для степени простого числа $\varphi(p^k) = p^k - p^{k-1}$ немедленно получается общая формула:
		\[
		\varphi(n) = n \prod_{p \mid n} \left(1 - \frac{1}{p}\right)
		\]
		где произведение берётся по всем простым делителям $n$.
	\end{tcolorbox}
	
	\begin{tcolorbox}[colback=yellow!5, colframe=orange!60!black, title=Важное замечание]
		Мультипликативность функции Эйлера не следует путать с полной мультипликативностью. Для произвольных (не обязательно взаимно простых) $a$ и $b$ равенство $\varphi(ab) = \varphi(a)\varphi(b)$, вообще говоря, неверно. Например, $\varphi(4) = 2$, $\varphi(2) = 1$, но $\varphi(8) = 4 \neq 2 \cdot 1$.
	\end{tcolorbox}
	
	
	
	
	
	
	% ==================== РАЗДЕЛ 22: ТЕОРЕМА ЭЙЛЕРА ====================
	\newpage
	\section{Теорема Эйлера. Малая теорема Ферма (22 вопрос)}
	
	\subsection{Теорема Эйлера}
	
	\noindent \textbf{Теорема (Эйлера).}
	Если $a \in \mathbb{Z}$ и $m \in \mathbb{N}$ взаимно просты, то
	\[
	a^{\varphi(m)} \equiv 1 \pmod{m} \quad \Longleftrightarrow \quad \bar{a}^{\varphi(m)} = \bar{1} \ \text{в} \ \mathbb{Z}/m\mathbb{Z}
	\]
	
	\begin{tcolorbox}[colback=red!5, colframe=red!70!black, title=Формулировка]
		Для любого числа $a$, взаимно простого с $m$, возведение в степень $\varphi(m)$ даёт число, сравнимое с $1$ по модулю $m$.
	\end{tcolorbox}
	
	\subsection{Доказательство теоремы Эйлера}
	
	\noindent \textbf{Доказательство.}
	Рассмотрим мультипликативную группу обратимых элементов кольца $\mathbb{Z}/m\mathbb{Z}$:
	\[
	(\mathbb{Z}/m\mathbb{Z})^* = \{\bar{a}_1, \bar{a}_2, \dots, \bar{a}_{\varphi(m)}\}
	\]
	где $\bar{a}_i$ — классы вычетов, взаимно простые с $m$. Количество таких классов равно $\varphi(m)$.
	
	\begin{tcolorbox}[colback=blue!5, colframe=blue!40!black, title=Ключевое наблюдение]
		Умножение на фиксированный обратимый элемент $\bar{a}$ является биекцией группы $(\mathbb{Z}/m\mathbb{Z})^*$ на себя. Действительно, отображение $x \mapsto \bar{a}x$ инъективно (так как $\bar{a}$ обратим) и сюръективно (образ совпадает с группой, поскольку группа конечна).
	\end{tcolorbox}
	
	Из этого наблюдения следует, что
	\[
	\bar{a} \cdot (\mathbb{Z}/m\mathbb{Z})^* = (\mathbb{Z}/m\mathbb{Z})^*
	\]
	то есть множество $\{\bar{a}\bar{a}_1, \bar{a}\bar{a}_2, \dots, \bar{a}\bar{a}_{\varphi(m)}\}$ совпадает с множеством $\{\bar{a}_1, \bar{a}_2, \dots, \bar{a}_{\varphi(m)}\}$ (возможно, в другом порядке).
	
	Перемножим все элементы множества $(\mathbb{Z}/m\mathbb{Z})^*$ двумя способами. С одной стороны, это просто произведение всех $\bar{a}_i$:
	\[
	P = \prod_{i=1}^{\varphi(m)} \bar{a}_i
	\]
	
	С другой стороны, это произведение всех $\bar{a}\bar{a}_i$:
	\[
	\prod_{i=1}^{\varphi(m)} (\bar{a}\bar{a}_i) = \bar{a}^{\varphi(m)} \prod_{i=1}^{\varphi(m)} \bar{a}_i = \bar{a}^{\varphi(m)} P
	\]
	
	Поскольку эти два произведения равны (оба представляют произведение всех элементов группы, взятых в разном порядке), получаем:
	\[
	P = \bar{a}^{\varphi(m)} P
	\]
	
	Так как все $\bar{a}_i$ обратимы, их произведение $P$ также обратимо (произведение обратимых элементов обратимо). Следовательно, на $P$ можно сократить (умножить обе части равенства на $P^{-1}$):
	\[
	\bar{1} = \bar{a}^{\varphi(m)}
	\]
	
	Возвращаясь к языку сравнений, это означает:
	\[
	a^{\varphi(m)} \equiv 1 \pmod{m}
	\]
	
	\begin{tcolorbox}[colback=green!5, colframe=green!50!black, title=Пояснение]
		Доказательство использует тот факт, что умножение на фиксированный элемент переставляет элементы конечной группы. Перемножая все элементы группы двумя способами, получаем требуемое равенство.
	\end{tcolorbox}
	
	\begin{tcolorbox}[colback=gray!5, colframe=gray!50, title=Пример]
		Возьмём $m = 9$, $\varphi(9) = 6$, $a = 2$ (взаимно прост с $9$).
		Группа $(\mathbb{Z}/9\mathbb{Z})^* = \{\bar{1}, \bar{2}, \bar{4}, \bar{5}, \bar{7}, \bar{8}\}$.
		Умножим все элементы на $\bar{2}$: получаем $\{\bar{2}, \bar{4}, \bar{8}, \bar{1}, \bar{5}, \bar{7}\}$ — та же группа.
		Произведение всех элементов: $1 \cdot 2 \cdot 4 \cdot 5 \cdot 7 \cdot 8 = 2240 \equiv ?$ по модулю $9$: $2240/9 = 248\cdot9 + 8$, остаток $8$. Но нам это не нужно — важно, что после умножения на $2^6$ получится то же произведение.
		Проверим теорему: $2^6 = 64$, $64 \div 9 = 7\cdot9 + 1$, значит $64 \equiv 1 \pmod{9}$ — верно.
	\end{tcolorbox}
	
	\myremark
	Теорема Эйлера — одно из фундаментальных утверждений теории чисел. Она обобщает малую теорему Ферма на случай составного модуля.
	
	\subsection{Малая теорема Ферма}
	
	\noindent \textbf{Следствие (Малая теорема Ферма).}
	Пусть $p \in \mathbb{N}$ простое, $a \in \mathbb{Z}$. Если $p$ простое и $a$ не делится на $p$ ($a \not\vdots p$), то
	\[
	a^{p-1} \equiv 1 \pmod{p},
	\]
	или, что то же, $\bar{a}^{p-1} = \bar{1}$ в поле $\mathbb{F}_p$.
	
	\medskip
	\noindent \textbf{Доказательство.}
	Это частный случай теоремы Эйлера. Если $p$ простое и $a \not\vdots p$, то $a$ и $p$ взаимно просты. Так как $p$ простое, $\varphi(p) = p - 1$. Подставляя $m = p$ в теорему Эйлера, получаем:
	\[
	a^{\varphi(p)} = a^{p-1} \equiv 1 \pmod{p}
	\]
	
	\begin{tcolorbox}[colback=blue!5, colframe=blue!40!black, title=Пример]
		$p = 7$, $a = 3$ (не делится на $7$). $3^{6} = 729$. $729 \div 7 = 104\cdot7 + 1$, значит $3^6 \equiv 1 \pmod{7}$.
		$p = 13$, $a = 5$. $5^{12} = 244140625$. Проверим: по теореме должно быть $\equiv 1 \pmod{13}$.
	\end{tcolorbox}
	
	\begin{tcolorbox}[colback=red!5, colframe=red!70!black, title=Другая форма малой теоремы Ферма]
		Часто используется эквивалентная формулировка: для любого целого $a$ и простого $p$
		\[
		a^p \equiv a \pmod{p}
		\]
		Эта форма верна для всех $a$ (включая кратные $p$), так как при $a$, делящемся на $p$, обе части сравнимы с $0$ по модулю $p$.
	\end{tcolorbox}
	
	\myremark
	Малая теорема Ферма имеет огромное значение:
	\begin{itemize}
		\item Она лежит в основе многих тестов на простоту.
		\item Используется в криптографии (например, в RSA для проверки ключей).
		\item Применяется при решении сравнений и в теории делимости.
	\end{itemize}
	
	\begin{tcolorbox}[colback=purple!5, colframe=purple!50!black, title=Заметочка: История]
		Пьер Ферма сформулировал эту теорему в 1640 году без доказательства. Первое доказательство было дано Эйлером в 1736 году. Эйлер же в 1763 году обобщил её на случай произвольного модуля (теорема Эйлера).
	\end{tcolorbox}
	
	\begin{tcolorbox}[colback=yellow!5, colframe=orange!60!black, title=Важное замечание]
		Условие взаимной простоты $a$ и $m$ в теореме Эйлера существенно. Например, для $a=2$, $m=6$: $\varphi(6)=2$, но $2^2=4 \not\equiv 1 \pmod{6}$ (и действительно, $2$ и $6$ не взаимно просты).
	\end{tcolorbox}
	
	
	
	
	
	
	
	% ==================== РАЗДЕЛ 23: ПОЛЕ КОМПЛЕКСНЫХ ЧИСЕЛ ====================
	\newpage
	\section{Построение поля комплексных чисел. Алгебраическая форма записи (23 вопрос)}
	
	\subsection{Построение поля комплексных чисел}
	
	\noindent Пусть $\mathbb{R}$ — поле вещественных чисел. Рассмотрим множество всех упорядоченных пар вещественных чисел:
	\[
	\mathbb{R}^2 = \{(a, b) \mid a, b \in \mathbb{R}\}
	\]
	
	Определим на этом множестве операции сложения и умножения следующим образом:
	
	\begin{tcolorbox}[colback=blue!5, colframe=blue!40!black, title=Определение операций]
		\[
		(a, b) + (c, d) = (a + c, b + d)
		\]
		\[
		(a, b) \cdot (c, d) = (ac - bd, ad + bc)
		\]
	\end{tcolorbox}
	
	Непосредственно из определения следует, что:
	\begin{itemize}
		\item Пара $(0, 0)$ является нейтральным элементом по сложению (нулём).
		\item Пара $(1, 0)$ является нейтральным элементом по умножению (единицей).
		\item Сложение ассоциативно и коммутативно (очевидно из определения).
		\item Умножение ассоциативно и коммутативно (проверяется вычислениями).
		\item Выполняется дистрибутивность умножения относительно сложения.
	\end{itemize}
	
	\subsection{Существование обратного элемента}
	
	\noindent \textbf{Обратный элемент по умножению.}
	Если $(a, b) \neq (0, 0)$, то существует обратный элемент. Действительно, рассмотрим пару
	\[
	\left( \frac{a}{a^2 + b^2}, \frac{-b}{a^2 + b^2} \right)
	\]
	
	Умножим её на $(a, b)$:
	\[
	(a, b) \cdot \left( \frac{a}{a^2 + b^2}, \frac{-b}{a^2 + b^2} \right) = 
	\left( a \cdot \frac{a}{a^2 + b^2} - b \cdot \frac{-b}{a^2 + b^2},\; 
	a \cdot \frac{-b}{a^2 + b^2} + b \cdot \frac{a}{a^2 + b^2} \right) =
	\]
	\[
	= \left( \frac{a^2 + b^2}{a^2 + b^2}, \frac{-ab + ab}{a^2 + b^2} \right) = (1, 0) = 1
	\]
	
	Таким образом, $\left( \frac{a}{a^2+b^2}, \frac{-b}{a^2+b^2} \right)$ является обратным к $(a, b)$.
	
	\begin{tcolorbox}[colback=green!5, colframe=green!50!black, title=Вывод]
		Множество $\mathbb{R}^2$ с введёнными операциями образует поле. Это поле называется \textit{полем комплексных чисел} и обозначается $\mathbb{C}$.
	\end{tcolorbox}
	
	\subsection{Вложение вещественных чисел}
	
	\noindent Рассмотрим отображение $\varphi \colon \mathbb{R} \to \mathbb{R}^2$, заданное формулой:
	\[
	\varphi(a) = (a, 0)
	\]
	
	Это отображение обладает следующими свойствами:
	\begin{itemize}
		\item Инъективность: $\varphi(a) = \varphi(a') \Rightarrow (a,0) = (a',0) \Rightarrow a = a'$.
		\item Сохранение сложения: $\varphi(a + b) = (a+b, 0) = (a,0) + (b,0) = \varphi(a) + \varphi(b)$.
		\item Сохранение умножения: $\varphi(ab) = (ab, 0) = (a,0)(b,0) = \varphi(a)\varphi(b)$.
	\end{itemize}
	
	Таким образом, $\varphi$ — инъективный гомоморфизм полей, т.е. \textit{вложение} $\mathbb{R}$ в $\mathbb{C}$. Образ этого вложения
	\[
	\varphi(\mathbb{R}) = \{(a, 0) \mid a \in \mathbb{R}\}
	\]
	отождествляется с полем вещественных чисел. Обычно мы пишем просто $a$ вместо $(a,0)$.
	
	\subsection{Мнимая единица}
	
	\noindent \textbf{Определение.} Обозначим
	\[
	i = (0, 1)
	\]
	
	Вычислим квадрат этой пары:
	\[
	i^2 = (0, 1)(0, 1) = (0 \cdot 0 - 1 \cdot 1,\; 0 \cdot 1 + 1 \cdot 0) = (-1, 0) = -1
	\]
	
	Таким образом, $i^2 = -1$. Элемент $i$ называется \textit{мнимой единицей}.
	
	\subsection{Алгебраическая форма записи}
	
	\noindent \textbf{Теорема (Алгебраическая форма).}
	Любое комплексное число $z \in \mathbb{C}$ единственным образом представляется в виде
	\[
	z = a + bi
	\]
	где $a, b \in \mathbb{R}$, а $i$ — мнимая единица.
	
	\medskip
	\noindent \textbf{Доказательство.}
	Для произвольной пары $(a, b) \in \mathbb{R}^2$ имеем:
	\[
	(a, b) = (a, 0) + (0, b) = (a, 0) + (b, 0)(0, 1) = a + bi
	\]
	
	Единственность представления следует из того, что если $a + bi = a' + b'i$, то $(a, b) = (a', b')$, откуда $a = a'$ и $b = b'$.
	
	\begin{tcolorbox}[colback=red!5, colframe=red!70!black, title=Алгебраическая форма]
		\[
		\boxed{z = a + bi,\quad a,b \in \mathbb{R},\ i^2 = -1}
		\]
		Число $a$ называется \textit{действительной частью} комплексного числа $z$ и обозначается $\operatorname{Re} z$. Число $b$ называется \textit{мнимой частью} и обозначается $\operatorname{Im} z$.
	\end{tcolorbox}
	
	\subsection{Свойства алгебраической формы}
	
	В алгебраической форме операции сложения и умножения принимают привычный вид:
	\[
	(a + bi) + (c + di) = (a + c) + (b + d)i
	\]
	\[
	(a + bi)(c + di) = ac + adi + bci + bdi^2 = (ac - bd) + (ad + bc)i
	\]
	что полностью соответствует определению операций на парах.
	
	\begin{tcolorbox}[colback=gray!5, colframe=gray!50, title=Примеры]
		\begin{itemize}
			\item $(3 + 2i) + (1 - 5i) = 4 - 3i$
			\item $(2 + 3i)(1 - 2i) = 2 - 4i + 3i - 6i^2 = 2 - i + 6 = 8 - i$
			\item Обратный элемент: $\frac{1}{a+bi} = \frac{a-bi}{a^2+b^2}$
		\end{itemize}
	\end{tcolorbox}
	
	\subsection{Минимальность поля комплексных чисел}
	
	\noindent \textbf{Замечание.} Поле $\mathbb{C}$ является минимальным полем, содержащим $\mathbb{R}$ и элемент $i$ с условием $i^2 = -1$. Если $P$ — любое поле, содержащее $\mathbb{R}$ и элемент $i$ такой, что $i^2 = -1$, то $\mathbb{C} \subset P$ (точнее, существует вложение $\mathbb{C}$ в $P$). В этом смысле $\mathbb{C}$ — наименьшее поле, содержащее $\mathbb{R}$ и корень уравнения $x^2 + 1 = 0$.
	
	\begin{tcolorbox}[colback=purple!5, colframe=purple!50!black, title=Заметочка: История]
		Мнимая единица $i$ была введена Леонардом Эйлером в XVIII веке. Термин "мнимая единица" отражает первоначальное скептическое отношение математиков к числам, квадрат которых отрицателен. Сейчас комплексные числа являются фундаментальным инструментом во многих областях математики и физики.
	\end{tcolorbox}
	
	\myremark
	Комплексные числа можно представлять геометрически как точки плоскости (плоскость Аргана). Действительная часть отвечает за координату по оси $x$, мнимая — по оси $y$. Такое представление позволяет интерпретировать умножение комплексных чисел как повороты и растяжения плоскости.
	
	
	
	
	
	
	% ==================== РАЗДЕЛ 24: ГЕОМЕТРИЧЕСКОЕ ИЗОБРАЖЕНИЕ ====================
	\newpage
	\section{Геометрическое изображение комплексных чисел. Тригонометрическая форма записи (24 вопрос)}
	
	\subsection{Комплексная плоскость}
	
	\noindent Алгебраическая форма комплексного числа $z = a + bi$ задаёт взаимно однозначное соответствие
	\[
	z \mapsto (\operatorname{Re} z, \operatorname{Im} z) = (a, b)
	\]
	между множеством $\mathbb{C}$ всех комплексных чисел и множеством $\mathbb{R}^2$ всех упорядоченных пар вещественных чисел.
	
	Поскольку точки плоскости также находятся во взаимно однозначном соответствии с парами координат $(x, y)$, мы получаем соответствие между комплексными числами и точками плоскости:
	\[
	z = a + bi \quad \longleftrightarrow \quad \text{точка } (a, b)
	\]
	
	\begin{tcolorbox}[colback=blue!5, colframe=blue!40!black, title=Комплексная плоскость]
		Плоскость, точки которой изображают комплексные числа, называется \textit{комплексной плоскостью}. При этом:
		\begin{itemize}
			\item Ось абсцисс называется \textit{действительной осью} (на ней лежат числа вида $a + 0i$, т.е. вещественные числа).
			\item Ось ординат называется \textit{мнимой осью} (на ней лежат числа вида $0 + bi$, т.е. чисто мнимые числа).
		\end{itemize}
	\end{tcolorbox}
	
	\subsection{Геометрическая интерпретация операций}
	
	\noindent Каждому комплексному числу $z = a + bi$ можно сопоставить также \textit{радиус-вектор} — вектор, идущий из начала координат в точку $(a, b)$.
	
	\begin{itemize}
		\item \textbf{Сложение комплексных чисел} соответствует векторному сложению радиус-векторов: если $z_1 = a_1 + b_1i$, $z_2 = a_2 + b_2i$, то их сумме $z_1 + z_2$ соответствует вектор, являющийся суммой векторов $(a_1, b_1)$ и $(a_2, b_2)$ (правило параллелограмма).
		
		\item \textbf{Умножение на вещественное число} соответствует растяжению (сжатию) вектора: число $\lambda \in \mathbb{R}$ умножает длину вектора на $|\lambda|$, а при $\lambda < 0$ ещё и меняет направление на противоположное.
	\end{itemize}
	
	\myremark
	Геометрическая интерпретация умножения двух произвольных комплексных чисел более сложна и связана с поворотом и растяжением — она естественно выражается в тригонометрической форме.
	
	\subsection{Модуль комплексного числа}
	
	\noindent \textbf{Определение (Модуль).}
	\textit{Модулем} комплексного числа $z = a + bi$ называется неотрицательное вещественное число
	\[
	|z| = \sqrt{a^2 + b^2}
	\]
	
	Геометрически модуль $|z|$ равен расстоянию от точки $z$ до начала координат (длине радиус-вектора).
	
	\begin{tcolorbox}[colback=gray!5, colframe=gray!50, title=Свойства модуля]
		\begin{itemize}
			\item $|z| \geq 0$, причём $|z| = 0 \iff z = 0$.
			\item $|\overline{z}| = |z|$, где $\overline{z} = a - bi$ — сопряжённое число.
			\item $|z_1 z_2| = |z_1| \cdot |z_2|$.
			\item $|z_1 + z_2| \leq |z_1| + |z_2|$ (неравенство треугольника).
		\end{itemize}
	\end{tcolorbox}
	
	\subsection{Аргумент комплексного числа}
	
	\noindent \textbf{Определение (Аргумент).}
	Пусть $z \neq 0$. \textit{Аргументом} комплексного числа $z$ называется угол $\varphi$ между положительным направлением действительной оси и радиус-вектором точки $z$. Аргумент обозначается $\arg z$.
	
	\begin{tcolorbox}[colback=blue!5, colframe=blue!40!black, title=Замечание о многозначности]
		Аргумент определён с точностью до слагаемого $2\pi k$, $k \in \mathbb{Z}$:
		\[
		\arg z = \varphi + 2\pi k,\quad k \in \mathbb{Z}
		\]
		Для однозначности часто выбирают \textit{главное значение аргумента} в промежутке $(-\pi, \pi]$ или $[0, 2\pi)$.
	\end{tcolorbox}
	
	Из геометрических соображений (рис. 1) получаем связь декартовых координат с полярными:
	\[
	a = r \cos \varphi,\quad b = r \sin \varphi,\quad r = |z| = \sqrt{a^2 + b^2}
	\]
	
	При этом $\operatorname{tg} \varphi = \frac{b}{a}$ (при $a \neq 0$), но выбор конкретного значения $\varphi$ требует учёта знаков $a$ и $b$ (четверти).
	
	\begin{tcolorbox}[colback=green!5, colframe=green!50!black, title=Нахождение аргумента]
		\[
		\arg z = 
		\begin{cases}
			\arctg\frac{b}{a}, & a > 0 \\
			\arctg\frac{b}{a} + \pi, & a < 0,\ b \geq 0 \\
			\arctg\frac{b}{a} - \pi, & a < 0,\ b < 0 \\
			\frac{\pi}{2}, & a = 0,\ b > 0 \\
			-\frac{\pi}{2}, & a = 0,\ b < 0
		\end{cases}
		\]
	\end{tcolorbox}
	
	\subsection{Тригонометрическая форма комплексного числа}
	
	\noindent \textbf{Теорема (Тригонометрическая форма).}
	Любое ненулевое комплексное число $z$ может быть представлено в виде
	\[
	z = r(\cos \varphi + i \sin \varphi)
	\]
	где $r = |z|$, а $\varphi = \arg z$ — некоторое значение аргумента.
	
	\begin{tcolorbox}[colback=red!5, colframe=red!70!black, title=Тригонометрическая форма]
		\[
		\boxed{z = r(\cos \varphi + i \sin \varphi),\quad r \geq 0,\ \varphi \in \mathbb{R}}
		\]
		При $z = 0$ модуль равен $0$, а аргумент не определён.
	\end{tcolorbox}
	
	\subsection{Умножение в тригонометрической форме}
	
	\noindent \textbf{Теорема (Умножение).}
	Если $z_1 = r_1(\cos \varphi_1 + i \sin \varphi_1)$, $z_2 = r_2(\cos \varphi_2 + i \sin \varphi_2)$, то
	\[
	z_1 z_2 = r_1 r_2 (\cos(\varphi_1 + \varphi_2) + i \sin(\varphi_1 + \varphi_2))
	\]
	
	\medskip
	\noindent \textbf{Доказательство.}
	\[
	\begin{aligned}
		z_1 z_2 &= r_1 r_2 (\cos \varphi_1 + i \sin \varphi_1)(\cos \varphi_2 + i \sin \varphi_2) \\
		&= r_1 r_2 [(\cos \varphi_1 \cos \varphi_2 - \sin \varphi_1 \sin \varphi_2) + i(\cos \varphi_1 \sin \varphi_2 + \sin \varphi_1 \cos \varphi_2)] \\
		&= r_1 r_2 (\cos(\varphi_1 + \varphi_2) + i \sin(\varphi_1 + \varphi_2))
	\end{aligned}
	\]
	используя формулы косинуса и синуса суммы.
	
	\begin{tcolorbox}[colback=blue!5, colframe=blue!40!black, title=Геометрический смысл]
		При умножении комплексных чисел их модули перемножаются, а аргументы складываются. Геометрически: растяжение вектора в $r_2$ раз и поворот на угол $\varphi_2$.
	\end{tcolorbox}
	
	\subsection{Формула Муавра}
	
	\noindent \textbf{Следствие (Формула Муавра).}
	Для любого целого $n$
	\[
	(\cos \varphi + i \sin \varphi)^n = \cos n\varphi + i \sin n\varphi
	\]
	
	\subsection{Деление и возведение в степень}
	
	Из формулы умножения получаем:
	\[
	\frac{z_1}{z_2} = \frac{r_1}{r_2} (\cos(\varphi_1 - \varphi_2) + i \sin(\varphi_1 - \varphi_2)),\quad z_2 \neq 0
	\]
	\[
	z^n = r^n (\cos n\varphi + i \sin n\varphi)
	\]
	
	\subsection{Извлечение корня}
	
	\noindent \textbf{Корни из комплексного числа.}
	Уравнение $w^n = z$ для $z \neq 0$ имеет ровно $n$ различных решений:
	\[
	w_k = \sqrt[n]{r} \left( \cos\frac{\varphi + 2\pi k}{n} + i \sin\frac{\varphi + 2\pi k}{n} \right),\quad k = 0, 1, \dots, n-1
	\]
	где $\sqrt[n]{r}$ — арифметический корень из положительного числа $r = |z|$.
	
	\begin{tcolorbox}[colback=gray!5, colframe=gray!50, title=Пример]
		Найдём все значения $\sqrt[3]{8i}$.
		\[
		8i = 8\left(\cos\frac{\pi}{2} + i \sin\frac{\pi}{2}\right),\quad \sqrt[3]{8} = 2
		\]
		\[
		w_k = 2\left(\cos\frac{\pi/2 + 2\pi k}{3} + i \sin\frac{\pi/2 + 2\pi k}{3}\right),\quad k = 0,1,2
		\]
		\[
		w_0 = 2\left(\cos\frac{\pi}{6} + i \sin\frac{\pi}{6}\right) = 2\left(\frac{\sqrt{3}}{2} + \frac{i}{2}\right) = \sqrt{3} + i
		\]
		\[
		w_1 = 2\left(\cos\frac{5\pi}{6} + i \sin\frac{5\pi}{6}\right) = 2\left(-\frac{\sqrt{3}}{2} + \frac{i}{2}\right) = -\sqrt{3} + i
		\]
		\[
		w_2 = 2\left(\cos\frac{3\pi}{2} + i \sin\frac{3\pi}{2}\right) = 2(0 - i) = -2i
		\]
	\end{tcolorbox}
	
	\myremark
	Тригонометрическая форма особенно удобна для умножения, деления, возведения в степень и извлечения корней. Алгебраическая форма удобнее для сложения и вычитания.
	
	\begin{tcolorbox}[colback=purple!5, colframe=purple!50!black, title=Заметочка: Связь с экспонентой]
		Используя формулу Эйлера $e^{i\varphi} = \cos\varphi + i\sin\varphi$, тригонометрическую форму можно записать как
		\[
		z = r e^{i\varphi}
		\]
		что ещё более упрощает операции: $z_1 z_2 = r_1 r_2 e^{i(\varphi_1+\varphi_2)}$, $z^n = r^n e^{in\varphi}$.
	\end{tcolorbox}
	
	
	
	
	
	% ==================== РАЗДЕЛ 25: МОДУЛЬ КОМПЛЕКСНОГО ЧИСЛА ====================
	\newpage
	\section{Модуль комплексного числа. Неравенство треугольника (25 вопрос)}
	
	\subsection{Определение модуля}
	
	\noindent \textbf{Определение (Модуль комплексного числа).}
	\textit{Модулем} комплексного числа $z = a + bi$ называется неотрицательное вещественное число
	\[
	|z| = \sqrt{a^2 + b^2}
	\]
	
	Геометрически модуль $|z|$ равен расстоянию от точки $z$ на комплексной плоскости до начала координат (длине радиус-вектора).
	
	\begin{tcolorbox}[colback=blue!5, colframe=blue!40!black, title=Связь с сопряжением]
		\[
		|z| = \sqrt{z \cdot \overline{z}},\quad \text{где } \overline{z} = a - bi
		\]
		Действительно, $z\overline{z} = (a+bi)(a-bi) = a^2 + b^2 = |z|^2$.
	\end{tcolorbox}
	
	\subsection{Свойства модуля}
	
	\noindent \textbf{Основные свойства:}
	\begin{enumerate}
		\item $|z| \geq 0$, причём $|z| = 0 \iff z = 0$.
		\item $|\overline{z}| = |z|$ (модуль сопряжённого числа равен модулю исходного).
		\item $|z_1 z_2| = |z_1| \cdot |z_2|$ (модуль произведения равен произведению модулей).
		\item $\left| \frac{z_1}{z_2} \right| = \frac{|z_1|}{|z_2|}$ для $z_2 \neq 0$.
		\item $|z^n| = |z|^n$ для любого целого $n$.
	\end{enumerate}
	
	\begin{tcolorbox}[colback=gray!5, colframe=gray!50, title=Примеры]
		\begin{itemize}
			\item $|3 + 4i| = \sqrt{3^2 + 4^2} = \sqrt{9 + 16} = \sqrt{25} = 5$.
			\item $|1 - i| = \sqrt{1^2 + (-1)^2} = \sqrt{2}$.
			\item $|5| = 5$, $|-5| = 5$ (для вещественных чисел модуль совпадает с абсолютной величиной).
			\item $|i| = \sqrt{0^2 + 1^2} = 1$.
		\end{itemize}
	\end{tcolorbox}
	
	\subsection{Расстояние между точками}
	
	\noindent \textbf{Определение.}
	Для комплексных чисел $x$ и $y$ величина
	\[
	|x - y|
	\]
	равна расстоянию между точками $x$ и $y$ на комплексной плоскости.
	
	Действительно, если $x = a + bi$, $y = c + di$, то
	\[
	|x - y| = \sqrt{(a-c)^2 + (b-d)^2}
	\]
	что совпадает с евклидовым расстоянием между точками $(a,b)$ и $(c,d)$.
	
	\subsection{Неравенство треугольника}
	
	\noindent \textbf{Теорема (Неравенство треугольника).}
	Для любых комплексных чисел $x, y \in \mathbb{C}$ выполняется неравенство
	\[
	|x| - |y| \leq |x + y| \leq |x| + |y|
	\]
	
	\begin{tcolorbox}[colback=red!5, colframe=red!70!black, title=Неравенство треугольника в двух формах]
		Верхняя оценка: $|x + y| \leq |x| + |y|$ (модуль суммы не превосходит суммы модулей).\\
		Нижняя оценка: $|x| - |y| \leq |x + y|$ (разность модулей не превосходит модуля суммы).
	\end{tcolorbox}
	
	\subsection{Доказательство неравенства треугольника}
	
	\noindent \textbf{Доказательство.}
	Запишем числа $x$ и $y$ в тригонометрической форме:
	\[
	x = r_1(\cos \varphi_1 + i \sin \varphi_1),\quad y = r_2(\cos \varphi_2 + i \sin \varphi_2)
	\]
	где $r_1 = |x|$, $r_2 = |y|$.
	
	Вычислим квадрат модуля суммы:
	\[
	\begin{aligned}
		|x + y|^2 &= (r_1 \cos \varphi_1 + r_2 \cos \varphi_2)^2 + (r_1 \sin \varphi_1 + r_2 \sin \varphi_2)^2 = \\
		&= r_1^2 \cos^2 \varphi_1 + 2r_1 r_2 \cos \varphi_1 \cos \varphi_2 + r_2^2 \cos^2 \varphi_2 \;+ \\
		&\quad\; + r_1^2 \sin^2 \varphi_1 + 2r_1 r_2 \sin \varphi_1 \sin \varphi_2 + r_2^2 \sin^2 \varphi_2 = \\
		&= r_1^2(\cos^2 \varphi_1 + \sin^2 \varphi_1) + r_2^2(\cos^2 \varphi_2 + \sin^2 \varphi_2) \;+ \\
		&\quad\; + 2r_1 r_2 (\cos \varphi_1 \cos \varphi_2 + \sin \varphi_1 \sin \varphi_2) = \\
		&= r_1^2 + r_2^2 + 2r_1 r_2 \cos(\varphi_1 - \varphi_2)
	\end{aligned}
	\]
	
	Поскольку $-1 \leq \cos(\varphi_1 - \varphi_2) \leq 1$, получаем:
	\[
	r_1^2 + r_2^2 - 2r_1 r_2 \leq |x + y|^2 \leq r_1^2 + r_2^2 + 2r_1 r_2
	\]
	
	Левая часть есть $(r_1 - r_2)^2$, правая часть — $(r_1 + r_2)^2$. Следовательно:
	\[
	(r_1 - r_2)^2 \leq |x + y|^2 \leq (r_1 + r_2)^2
	\]
	
	Извлекая арифметический квадратный корень (все величины неотрицательны), получаем:
	\[
	|r_1 - r_2| \leq |x + y| \leq r_1 + r_2
	\]
	
	Учитывая, что $r_1 = |x|$, $r_2 = |y|$, окончательно:
	\[
	\bigl| |x| - |y| \bigr| \leq |x + y| \leq |x| + |y|
	\]
	
	\begin{tcolorbox}[colback=green!5, colframe=green!50!black, title=Пояснение]
		Доказательство использует тригонометрическую форму и свойство косинуса разности. Полученное неравенство является комплексным аналогом неравенства треугольника для векторов на плоскости.
	\end{tcolorbox}
	
	\subsection{Геометрическая интерпретация}
	
	Неравенство треугольника имеет простую геометрическую интерпретацию: длина одной стороны треугольника не превосходит суммы длин двух других сторон и не меньше абсолютной величины их разности.
	
	\begin{tcolorbox}[colback=blue!5, colframe=blue!40!black, title=Иллюстрация]
		Рассмотрим треугольник с вершинами в точках $0$, $x$ и $-y$ (или $x$, $y$ и $x+y$). Стороны этого треугольника имеют длины $|x|$, $|y|$ и $|x+y|$. Неравенство треугольника утверждает, что:
		\[
		|x+y| \leq |x| + |y| \quad \text{(длина стороны не больше суммы двух других)}
		\]
		и
		\[
		|x| \leq |x+y| + |y| \quad \Longleftrightarrow \quad |x| - |y| \leq |x+y|
		\]
		(аналогично для $|y|$).
	\end{tcolorbox}
	
	\subsection{Следствия неравенства треугольника}
	
	\begin{tcolorbox}[colback=gray!5, colframe=gray!50, title=Полезные следствия]
		\begin{itemize}
			\item $|x - y| \geq \bigl| |x| - |y| \bigr|$ (замена $x$ на $x-y$ и $y$ на $y$ в исходном неравенстве).
			\item $|x_1 + x_2 + \dots + x_n| \leq |x_1| + |x_2| + \dots + |x_n|$ (обобщение по индукции).
			\item $|x| - |y| \leq |x - y|$ (частный случай при замене $y$ на $-y$).
		\end{itemize}
	\end{tcolorbox}
	
	\myremark
	Неравенство треугольника является фундаментальным свойством нормы (в данном случае модуля комплексного числа). Оно необходимо для определения метрических пространств и играет ключевую роль в комплексном анализе.
	
	\begin{tcolorbox}[colback=purple!5, colframe=purple!50!black, title=Заметочка]
		Модуль комплексного числа — это частный случай понятия нормы в векторных пространствах. Все свойства модуля аналогичны свойствам длины вектора в евклидовом пространстве, что естественно, поскольку комплексная плоскость изоморфна $\mathbb{R}^2$.
	\end{tcolorbox}
	
	
	
	% ==================== РАЗДЕЛ 26: УМНОЖЕНИЕ И ДЕЛЕНИЕ ====================
	\newpage
	\section{Умножение и деление комплексных чисел в тригонометрической форме. Формула Муавра (26 вопрос)}
	
	\subsection{Умножение в тригонометрической форме}
	
	\noindent \textbf{Теорема (Умножение).}
	Если комплексные числа заданы в тригонометрической форме:
	\[
	x = r_1(\cos \varphi_1 + i \sin \varphi_1),\quad y = r_2(\cos \varphi_2 + i \sin \varphi_2),
	\]
	то их произведение равно
	\[
	x \cdot y = r_1 r_2 \bigl(\cos(\varphi_1 + \varphi_2) + i \sin(\varphi_1 + \varphi_2)\bigr)
	\]
	
	\begin{tcolorbox}[colback=red!5, colframe=red!70!black, title=Правило умножения]
		При умножении комплексных чисел в тригонометрической форме их модули перемножаются, а аргументы складываются:
		\[
		|xy| = |x| \cdot |y|,\quad \arg(xy) = \arg x + \arg y
		\]
	\end{tcolorbox}
	
	\medskip
	\noindent \textbf{Доказательство.}
	Вычислим произведение непосредственно:
	\[
	\begin{aligned}
		xy &= r_1 r_2 (\cos \varphi_1 + i \sin \varphi_1)(\cos \varphi_2 + i \sin \varphi_2) \\
		&= r_1 r_2 \bigl[(\cos \varphi_1 \cos \varphi_2 - \sin \varphi_1 \sin \varphi_2) + i(\cos \varphi_1 \sin \varphi_2 + \sin \varphi_1 \cos \varphi_2)\bigr] \\
		&= r_1 r_2 \bigl[\cos(\varphi_1 + \varphi_2) + i \sin(\varphi_1 + \varphi_2)\bigr]
	\end{aligned}
	\]
	используя формулы косинуса и синуса суммы.
	
	\begin{tcolorbox}[colback=blue!5, colframe=blue!40!black, title=Геометрический смысл]
		Умножение на комплексное число $y = r_2(\cos \varphi_2 + i \sin \varphi_2)$ геометрически означает:
		\begin{itemize}
			\item растяжение (сжатие) радиус-вектора числа $x$ в $r_2$ раз;
			\item поворот этого вектора на угол $\varphi_2$ против часовой стрелки.
		\end{itemize}
	\end{tcolorbox}
	
	\subsection{Деление в тригонометрической форме}
	
	\noindent \textbf{Теорема (Деление).}
	Для ненулевых комплексных чисел в тригонометрической форме:
	\[
	\frac{x}{y} = \frac{r_1}{r_2} \bigl(\cos(\varphi_1 - \varphi_2) + i \sin(\varphi_1 - \varphi_2)\bigr),\quad y \neq 0
	\]
	
	\begin{tcolorbox}[colback=red!5, colframe=red!70!black, title=Правило деления]
		При делении комплексных чисел их модули делятся, а аргументы вычитаются:
		\[
		\left|\frac{x}{y}\right| = \frac{|x|}{|y|},\quad \arg\left(\frac{x}{y}\right) = \arg x - \arg y
		\]
	\end{tcolorbox}
	
	\medskip
	\noindent \textbf{Доказательство.}
	Достаточно проверить, что $\frac{x}{y} \cdot y = x$:
	\[
	\frac{r_1}{r_2} (\cos(\varphi_1 - \varphi_2) + i \sin(\varphi_1 - \varphi_2)) \cdot r_2 (\cos \varphi_2 + i \sin \varphi_2) = r_1 (\cos \varphi_1 + i \sin \varphi_1)
	\]
	по правилу умножения.
	
	\subsection{Частные случаи}
	
	\begin{tcolorbox}[colback=gray!5, colframe=gray!50, title=Обратное число]
		Для $y \neq 0$ обратное число имеет вид:
		\[
		y^{-1} = \frac{1}{r_2} (\cos(-\varphi_2) + i \sin(-\varphi_2)) = \frac{1}{r_2} (\cos \varphi_2 - i \sin \varphi_2)
		\]
		При этом:
		\[
		|y^{-1}| = |y|^{-1},\quad \arg(y^{-1}) = -\arg y
		\]
	\end{tcolorbox}
	
	\subsection{Возведение в целую степень}
	
	\noindent \textbf{Теорема (Возведение в степень).}
	Для любого целого числа $n \in \mathbb{Z}$:
	\[
	x^n = r_1^n (\cos n\varphi_1 + i \sin n\varphi_1)
	\]
	где $x = r_1(\cos \varphi_1 + i \sin \varphi_1)$.
	
	\begin{tcolorbox}[colback=red!5, colframe=red!70!black, title=Правило возведения в степень]
		При возведении комплексного числа в целую степень его модуль возводится в эту степень, а аргумент умножается на показатель степени:
		\[
		|x^n| = |x|^n,\quad \arg(x^n) = n \arg x
		\]
	\end{tcolorbox}
	
	\medskip
	\noindent \textbf{Доказательство.}
	Для $n > 0$ утверждение доказывается индукцией по $n$, используя правило умножения. Для $n = 0$: $x^0 = 1 = r_1^0(\cos 0 + i \sin 0)$. Для $n < 0$: $x^n = (x^{-1})^{-n}$, и утверждение следует из правила для обратного числа и положительной степени.
	
	\subsection{Формула Муавра}
	
	\noindent \textbf{Следствие (Формула Муавра).}
	Для любого целого $n$:
	\[
	(\cos \varphi + i \sin \varphi)^n = \cos n\varphi + i \sin n\varphi
	\]
	
	\begin{tcolorbox}[colback=blue!5, colframe=blue!40!black, title=Примеры]
		\begin{itemize}
			\item $(\cos \varphi + i \sin \varphi)^2 = \cos 2\varphi + i \sin 2\varphi$.
			\item $(\cos 30^\circ + i \sin 30^\circ)^3 = \cos 90^\circ + i \sin 90^\circ = i$.
			\item $(1 + i)^6$: сначала $1+i = \sqrt{2}(\cos 45^\circ + i \sin 45^\circ)$, тогда $(1+i)^6 = (\sqrt{2})^6(\cos 270^\circ + i \sin 270^\circ) = 8(0 - i) = -8i$.
		\end{itemize}
	\end{tcolorbox}
	
	\subsection{Применение формулы Муавра}
	
	Формула Муавра позволяет:
	\begin{itemize}
		\item Выражать $\cos n\varphi$ и $\sin n\varphi$ через степени $\cos \varphi$ и $\sin \varphi$.
		\item Находить корни из комплексных чисел.
		\item Решать уравнения вида $z^n = a$.
	\end{itemize}
	
	\begin{tcolorbox}[colback=green!5, colframe=green!50!black, title=Пример: выражение $\cos 3\varphi$]
		\[
		\cos 3\varphi + i \sin 3\varphi = (\cos \varphi + i \sin \varphi)^3 = \cos^3 \varphi + 3i \cos^2 \varphi \sin \varphi - 3 \cos \varphi \sin^2 \varphi - i \sin^3 \varphi
		\]
		Приравнивая действительные части:
		\[
		\cos 3\varphi = \cos^3 \varphi - 3 \cos \varphi \sin^2 \varphi = 4\cos^3 \varphi - 3\cos \varphi
		\]
	\end{tcolorbox}
	
	\myremark
	Формула Муавра — один из важнейших результатов в теории комплексных чисел. Она показывает, как тригонометрические функции кратных углов связаны со степенями комплексных чисел. В сочетании с формулой Эйлера $e^{i\varphi} = \cos\varphi + i\sin\varphi$ она даёт мощный аппарат для вычислений.
	
	\begin{tcolorbox}[colback=purple!5, colframe=purple!50!black, title=Заметочка: История]
		Формула названа в честь французского математика Абрахама де Муавра (1667-1754), который открыл её в 1707 году. Интересно, что сам Муавр использовал эту формулу для аппроксимации биномиального распределения, что привело к появлению нормального распределения.
	\end{tcolorbox}
	
	\begin{tcolorbox}[colback=yellow!5, colframe=orange!60!black, title=Важное замечание]
		При использовании формулы Муавра для целых отрицательных $n$ нужно помнить, что
		\[
		(\cos \varphi + i \sin \varphi)^{-1} = \cos(-\varphi) + i \sin(-\varphi) = \cos \varphi - i \sin \varphi
		\]
		что согласуется с формулой.
	\end{tcolorbox}
	
	
	
	
	
	% ==================== РАЗДЕЛ 27: КОРНИ ИЗ КОМПЛЕКСНЫХ ЧИСЕЛ ====================
	\newpage
	\section{Извлечение корня из комплексных чисел. Изображение корней из единицы (27 вопрос)}
	
	\subsection{Определение корня n-й степени}
	
	Первое, что необходимо сделать — это строго определить, что понимается под корнем n-й степени из комплексного числа.
	
	\medskip
	
	\noindent \textbf{Определение.} Пусть \( n \in \mathbb{N} \) и \( x \in \mathbb{C} \). \textit{Корнем n-й степени} из числа \( x \) называется такое комплексное число \( y \in \mathbb{C} \), что
	\[
	y^n = x.
	\]
	
	\subsection{Случай \( x = 0 \)}
	
	Рассмотрим тривиальный случай.
	\begin{itemize}
		\item Если \( x = 0 \), то уравнение \( y^n = 0 \) имеет единственное решение \( y = 0 \) для любого натурального \( n \).
	\end{itemize}
	
	\subsection{Случай \( x \neq 0 \)}
	
	Пусть теперь \( x \neq 0 \). Запишем число \( x \) в тригонометрической форме:
	\[
	x = r(\cos \theta + i \sin \theta),
	\]
	где \( r = |x| > 0 \) — модуль числа, а \( \theta = \arg x \) — его аргумент (определенный с точностью до \( 2\pi k, k \in \mathbb{Z} \)).
	
	Будем искать корень \( y \) также в тригонометрической форме:
	\[
	y = \rho(\cos \varphi + i \sin \varphi),
	\]
	где \( \rho = |y| \geq 0 \), а \( \varphi = \arg y \).
	
	\subsection{Вывод формулы для корня}
	
	Воспользуемся формулой Муавра для возведения в степень:
	\[
	y^n = \rho^n (\cos (n\varphi) + i \sin (n\varphi)).
	\]
	
	Для выполнения равенства \( y^n = x \) необходимо и достаточно, чтобы модули были равны, а аргументы совпадали с точностью до периода \( 2\pi \):
	\[
	\begin{cases}
		\rho^n = r, \\
		n\varphi = \theta + 2\pi k, \quad k \in \mathbb{Z}.
	\end{cases}
	\]
	
	Отсюда получаем:
	\[
	\rho = \sqrt[n]{r} \quad (\text{арифметический корень из положительного числа}),
	\]
	\[
	\varphi = \frac{\theta + 2\pi k}{n}, \quad k \in \mathbb{Z}.
	\]
	
	\subsection{Количество различных корней}
	
	Несмотря на то, что параметр \( k \) пробегает все целые числа, различные значения аргумента \( \varphi \) будут давать только \( n \) различных корней. Это происходит из-за периодичности тригонометрических функций (\( \cos \) и \( \sin \) имеют период \( 2\pi \)).
	
	Два аргумента \( \varphi_{k_1} \) и \( \varphi_{k_2} \) задают одно и то же число тогда и только тогда, когда их разность кратна \( 2\pi \):
	\[
	\frac{\theta + 2\pi k_1}{n} - \frac{\theta + 2\pi k_2}{n} = 2\pi m \quad \Rightarrow \quad \frac{2\pi (k_1 - k_2)}{n} = 2\pi m \quad \Rightarrow \quad k_1 - k_2 = m n.
	\]
	Таким образом, \( k_1 \) и \( k_2 \) должны быть сравнимы по модулю \( n \).
	
	\medskip
	Следовательно, все различные корни получаются при взятии \( n \) последовательных значений \( k \), например:
	\[
	k = 0, 1, 2, \dots, n-1.
	\]
	
	\myremark
	При \( x \neq 0 \) существует ровно \( n \) различных комплексных корней n-й степени.
	
	\subsection{Множество всех корней}
	
	Обозначим корни через \( y_k \). Тогда множество всех корней n-й степени из числа \( x \) имеет вид:
	\[
	y_k = \sqrt[n]{r} \left( \cos \frac{\theta + 2\pi k}{n} + i \sin \frac{\theta + 2\pi k}{n} \right), \quad k = 0, 1, \dots, n-1.
	\]
	
	\subsection{Геометрическая интерпретация}
	
	Все корни \( y_k \) из одного числа \( x \) имеют одинаковый модуль \( \sqrt[n]{r} \). На комплексной плоскости они располагаются на окружности радиуса \( \sqrt[n]{r} \) с центром в начале координат.
	
	Аргументы корней отличаются друг от друга на угол \( \frac{2\pi}{n} \). Таким образом, точки, соответствующие корням, делят окружность на \( n \) равных дуг, т.е. являются вершинами правильного \( n \)-угольника, вписанного в эту окружность.
	
	\begin{center}
		\begin{tikzpicture}[scale=2.5]
			% Оси
			\draw[thin,->] (-1.2,0) -- (1.2,0) node[right] {Re};
			\draw[thin,->] (0,-1.2) -- (0,1.2) node[above] {Im};
			
			% Окружность
			\draw[thick,blue!50] (0,0) circle (1);
			
			% Точки - корни 5-й степени из 1 (для примера)
			\foreach \k in {0,1,2,3,4} {
				\filldraw[red] ({cos(72*\k)},{sin(72*\k)}) circle (0.5pt);
			}
			
			% Подписи
			\node[blue] at (0.7,0.7) {$|z| = \sqrt[n]{r}$};
			\node[red] at (1.1,0.2) {$\frac{2\pi}{n}$};
			
			% Дуга для угла
			\draw[thick,red,->] (0.8,0) arc (0:72:0.8);
			
		\end{tikzpicture}
	\end{center}
	
	\subsection{Корни из единицы}
	
	Важным частным случаем являются корни n-й степени из числа 1. Положим \( x = 1 \). В тригонометрической форме:
	\[
	1 = \cos 0 + i \sin 0 \quad \Rightarrow \quad r = 1,\ \theta = 0.
	\]
	
	Тогда корни n-й степени из 1 (обозначаемые \( \varepsilon_k \)) находятся по формуле:
	\[
	\varepsilon_k = \cos \frac{2\pi k}{n} + i \sin \frac{2\pi k}{n}, \quad k = 0, 1, \dots, n-1.
	\]
	
	\begin{tcolorbox}[colback=blue!5, colframe=blue!40!black, title=Свойства корней из единицы]
		\begin{enumerate}
			\item \textbf{Модуль:} Все корни из единицы лежат на единичной окружности, так как \( |\varepsilon_k| = 1 \).
			\item \textbf{Вершины многоугольника:} Они являются вершинами правильного \( n \)-угольника, вписанного в единичную окружность. Одна из вершин всегда находится в точке \( 1 \) (при \( k = 0 \)).
			\item \textbf{Замкнутость относительно умножения:} Произведение любых двух корней n-й степени из 1 также является корнем n-й степени из 1. Это свойство позволяет говорить о группе корней из единицы.
			\item \textbf{Основное тождество:} Если \( \varepsilon \neq 1 \) (т.е. \( \varepsilon = \varepsilon_1 = \cos \frac{2\pi}{n} + i \sin \frac{2\pi}{n} \)), то:
			\[
			1 + \varepsilon + \varepsilon^2 + \dots + \varepsilon^{n-1} = 0.
			\]
			Это легко показать, используя формулу суммы геометрической прогрессии.
		\end{enumerate}
	\end{tcolorbox}
	
	\subsection{Пример для \( n = 3 \) (кубические корни из 1)}
	
	Для \( n = 3 \):
	\[
	\varepsilon_0 = 1,
	\]
	\[
	\varepsilon_1 = \cos \frac{2\pi}{3} + i \sin \frac{2\pi}{3} = -\frac{1}{2} + i\frac{\sqrt{3}}{2},
	\]
	\[
	\varepsilon_2 = \cos \frac{4\pi}{3} + i \sin \frac{4\pi}{3} = -\frac{1}{2} - i\frac{\sqrt{3}}{2}.
	\]
	
	Геометрически они образуют равносторонний треугольник, вписанный в окружность.
	
	\begin{center}
		\begin{tikzpicture}[scale=2.5]
			% Оси
			\draw[thin,->] (-1.2,0) -- (1.2,0) node[right] {Re};
			\draw[thin,->] (0,-1.2) -- (0,1.2) node[above] {Im};
			
			% Окружность
			\draw[thick,blue!50] (0,0) circle (1);
			
			% Точки - корни 3-й степени из 1
			\filldraw[red] (1,0) circle (1pt) node[below right] {$\varepsilon_0 = 1$};
			\filldraw[red] ({cos(120)},{sin(120)}) circle (1pt) node[above left] {$\varepsilon_1$};
			\filldraw[red] ({cos(240)},{sin(240)}) circle (1pt) node[below left] {$\varepsilon_2$};
			
			% Соединим их треугольником
			\draw[red, dashed] (1,0) -- ({cos(120)},{sin(120)}) -- ({cos(240)},{sin(240)}) -- cycle;
			
		\end{tikzpicture}
	\end{center}
	
	\subsection{Связь между корнями произвольного числа и корнями из единицы}
	
	Если известен один из корней n-й степени из числа \( x \), например, \( y_0 = \sqrt[n]{r}(\cos \frac{\theta}{n} + i \sin \frac{\theta}{n}) \), то все остальные корни можно получить умножением этого корня на все корни n-й степени из единицы:
	\[
	y_k = y_0 \cdot \varepsilon_k, \quad k = 0, 1, \dots, n-1.
	\]
	
	\myremark
	Это свойство вытекает из того факта, что \( (y_0 \varepsilon_k)^n = y_0^n \varepsilon_k^n = x \cdot 1 = x \). Именно это и доказывается в конце предоставленного вами фрагмента.
	
	
	
	
	
	
	
	% ==================== РАЗДЕЛ 28: КОРНИ ИЗ ЕДИНИЦЫ ====================
	\newpage
	\section{Корни из единицы. Свойства. Группа корней из единицы (28 вопрос)}
	
	\subsection{Определение корней из единицы}
	
	Важнейшим частным случаем корней n-й степени из комплексного числа являются корни из числа 1.
	
	\medskip
	
	\noindent \textbf{Определение.} \textit{Корнем n-й степени из единицы} называется такое комплексное число \( \varepsilon \in \mathbb{C} \), что
	\[
	\varepsilon^n = 1.
	\]
	
	Так как \( |1| = 1 \), а \( \arg 1 = 0 \), то, подставляя \( r = 1 \) и \( \theta = 0 \) в общую формулу корня, получаем множество всех различных корней n-й степени из единицы:
	\[
	\varepsilon_k = \cos \frac{2\pi k}{n} + i \sin \frac{2\pi k}{n}, \quad k = 0, 1, \dots, n-1.
	\]
	
	\subsection{Свойства корней из единицы}
	
	Корни из единицы обладают рядом замечательных алгебраических свойств.
	
	\begin{tcolorbox}[colback=green!5, colframe=green!50!black, title=Основные свойства]
		
		\textbf{Свойство 1 (Замкнутость относительно умножения).} 
		Если \( \varepsilon \) и \( \delta \) — корни n-й степени из единицы (т.е. \( \varepsilon^n = 1 \) и \( \delta^n = 1 \)), то их произведение \( \varepsilon \delta \) также является корнем n-й степени из единицы.
		
		\textit{Доказательство.} 
		\[
		(\varepsilon \delta)^n = \varepsilon^n \delta^n = 1 \cdot 1 = 1.
		\]
		
		\medskip
		\textbf{Свойство 2 (Существование обратного элемента).}
		Если \( \varepsilon \) — корень n-й степени из единицы, то обратный к нему элемент \( \varepsilon^{-1} \) также является корнем n-й степени из единицы. Более того, для корней из единицы обратный элемент совпадает с сопряженным: \( \varepsilon^{-1} = \overline{\varepsilon} \).
		
		\textit{Доказательство.} 
		Так как \( |\varepsilon| = 1 \), то \( \varepsilon \overline{\varepsilon} = |\varepsilon|^2 = 1 \), откуда \( \overline{\varepsilon} = \varepsilon^{-1} \). Далее,
		\[
		(\varepsilon^{-1})^n = (\varepsilon^n)^{-1} = 1^{-1} = 1.
		\]
		
		\medskip
		\textbf{Свойство 3 (Единица).}
		Число 1 всегда является корнем n-й степени из единицы (при \( k = 0 \)), так как \( 1^n = 1 \).
		
		\medskip
		\textbf{Свойство 4 (Сумма всех корней).}
		Сумма всех различных корней n-й степени из единицы равна нулю:
		\[
		\sum_{k=0}^{n-1} \varepsilon_k = 0.
		\]
		
		\textit{Доказательство.} 
		Все корни \( \varepsilon_k \) являются корнями уравнения \( z^n - 1 = 0 \). По теореме Виета для многочлена \( z^n - 1 \) коэффициент при \( z^{n-1} \) равен нулю, что и означает равенство нулю суммы всех корней.
		
		Альтернативное доказательство: пусть \( \varepsilon = \varepsilon_1 = \cos \frac{2\pi}{n} + i \sin \frac{2\pi}{n} \neq 1 \). Тогда сумма представляет собой геометрическую прогрессию:
		\[
		1 + \varepsilon + \varepsilon^2 + \dots + \varepsilon^{n-1} = \frac{\varepsilon^n - 1}{\varepsilon - 1} = \frac{1 - 1}{\varepsilon - 1} = 0.
		\]
		
	\end{tcolorbox}
	
	\subsection{Группа корней из единицы}
	
	Множество всех корней n-й степени из единицы образует алгебраическую структуру, называемую \textit{циклической группой}.
	
	\medskip
	\noindent \textbf{Определение.} Множество
	\[
	U_n = \{ \varepsilon_k \mid k = 0, 1, \dots, n-1 \}
	\]
	с операцией умножения комплексных чисел называется \textit{группой корней n-й степени из единицы}.
	
	\medskip
	Проверим выполнение аксиом группы:
	\begin{itemize}
		\item \textbf{Замкнутость:} следует из Свойства 1.
		\item \textbf{Ассоциативность:} умножение комплексных чисел ассоциативно.
		\item \textbf{Нейтральный элемент:} \( \varepsilon_0 = 1 \).
		\item \textbf{Обратный элемент:} для \( \varepsilon_k \) обратным является \( \varepsilon_{n-k} = \overline{\varepsilon_k} \) (следует из Свойства 2).
	\end{itemize}
	
	\subsection{Циклическая природа группы}
	
	Группа \( U_n \) является \textit{циклической}. Это означает, что существует элемент, называемый \textit{первообразным корнем} или \textit{образующей}, степени которого порождают всю группу.
	
	\medskip
	\noindent \textbf{Определение.} Корень \( \varepsilon \) из единицы называется \textit{первообразным} (или примитивным) корнем n-й степени, если его порядок в группе равен \( n \), т.е. наименьшее положительное \( m \) такое, что \( \varepsilon^m = 1 \), есть \( n \).
	
	\medskip
	Элемент \( \varepsilon_1 = \cos \frac{2\pi}{n} + i \sin \frac{2\pi}{n} \) является первообразным корнем, так как:
	\[
	\varepsilon_1^k = \cos \frac{2\pi k}{n} + i \sin \frac{2\pi k}{n} = \varepsilon_k,
	\]
	и \( \varepsilon_1^k = 1 \) только тогда, когда \( k \) кратно \( n \).
	
	\myremark
	Все первообразные корни — это \( \varepsilon_k \), где \( k \) взаимно просто с \( n \) (т.е. \( \gcd(k, n) = 1 \)). Их количество равно \( \varphi(n) \) (функция Эйлера).
	
	\subsection{Геометрическая интерпретация}
	
	Геометрически элементы группы \( U_n \) располагаются на единичной окружности комплексной плоскости и являются вершинами правильного \( n \)-угольника.
	
	\begin{center}
		\begin{tikzpicture}[scale=2.5]
			% Оси
			\draw[thin,->] (-1.2,0) -- (1.2,0) node[right] {Re};
			\draw[thin,->] (0,-1.2) -- (0,1.2) node[above] {Im};
			
			% Единичная окружность
			\draw[thick,blue!50] (0,0) circle (1);
			
			% Точки - корни 6-й степени из 1
			\foreach \k in {0,1,2,3,4,5} {
				\filldraw[red] ({cos(60*\k)},{sin(60*\k)}) circle (1.5pt);
			}
			
			% Подписи
			\node[red] at (1,0) [below right] {$\varepsilon_0$};
			\node[red] at (0.5,0.866) [above right] {$\varepsilon_1$};
			\node[red] at (-0.5,0.866) [above left] {$\varepsilon_2$};
			\node[red] at (-1,0) [above left] {$\varepsilon_3$};
			\node[red] at (-0.5,-0.866) [below left] {$\varepsilon_4$};
			\node[red] at (0.5,-0.866) [below right] {$\varepsilon_5$};
			
			% Правильный шестиугольник
			\draw[red, dashed] (1,0) -- (0.5,0.866) -- (-0.5,0.866) -- (-1,0) -- (-0.5,-0.866) -- (0.5,-0.866) -- cycle;
			
		\end{tikzpicture}
	\end{center}
	
	\subsection{Связь с общей формулой корней}
	
	Как было показано в вопросе 27, если \( y_0 \) — один из корней n-й степени из произвольного числа \( x \), то все остальные корни получаются умножением \( y_0 \) на все элементы группы \( U_n \):
	\[
	\{ y_0, y_0\varepsilon_1, y_0\varepsilon_2, \dots, y_0\varepsilon_{n-1} \}.
	\]
	
	Таким образом, группа корней из единицы играет роль структурного «множителя», описывающего симметрию множества корней произвольного числа.
	
	
	
	
	
	% ==================== РАЗДЕЛ 29: КОЛЬЦО МНОГОЧЛЕНОВ ====================
	\newpage
	\section{Построение кольца многочленов над коммутативным кольцом (29 вопрос)}
	
	\subsection{Мотивация и предварительные понятия}
	
	Пусть \( R \) — коммутативное кольцо с единицей \( 1_R \). Перед тем как строить кольцо многочленов, необходимо ввести понятие элемента, трансцендентного (первичного) над кольцом.
	
	\medskip
	
	\noindent \textbf{Определение (Первичный элемент).} 
	Пусть \( R \) — коммутативное кольцо с 1, и пусть \( A \) — некоторое расширение кольца \( R \) (т.е. \( R \subseteq A \) и единица \( 1_R \) является единицей в \( A \)). Элемент \( \xi \in A \) называется \textit{первичным} (или трансцендентным) над \( R \), если:
	\begin{enumerate}
		\item \( \forall a \in R: a\xi = \xi a \) (элемент коммутирует со всеми элементами кольца \( R \));
		\item Равенство 
		\[
		\sum_{i=0}^{n} b_i \xi^i = 0, \quad b_i \in R,
		\]
		возможно только при условии, что все коэффициенты \( b_i = 0 \).
	\end{enumerate}
	
	\myremark
	Из определения следует, что никакой элемент самого кольца \( R \) не может быть первичным над \( R \), так как для \( \xi = a_0 \in R \) существует нетривиальная линейная комбинация \( 1 \cdot \xi - a_0 \cdot 1 = 0 \).
	
	\subsection{Определение кольца многочленов}
	
	Существует два эквивалентных подхода к построению кольца многочленов: аксиоматический и конструктивный (через последовательности).
	
	\subsubsection*{Аксиоматическое определение}
	
	\noindent \textbf{Определение.} Пусть \( R \) — коммутативное кольцо с 1. Кольцо \( S \) называется \textit{кольцом многочленов от одной переменной над \( R \)}, если:
	\begin{itemize}
		\item \( S \) является расширением кольца \( R \) (т.е. \( R \subseteq S \) и единица совпадает);
		\item В \( S \) существует первичный над \( R \) элемент \( t \) такой, что \( S = R[t] \), где
		\[
		R[t] = \left\{ \sum_{i \in \mathbb{N}_0} a_i t^i \mid a_i \in R \text{ и лишь конечное число } a_i \neq 0 \right\}.
		\]
	\end{itemize}
	
	\subsection{Конструктивное построение через последовательности}
	
	Дадим явную конструкцию кольца многочленов \( R[x] \), не опирающуюся на существование расширения \( A \).
	
	\medskip
	\noindent \textbf{Шаг 1: Множество последовательностей.}
	Рассмотрим множество
	\[
	P = \{ (a_0, a_1, a_2, \dots) \mid a_i \in R \text{ и лишь конечное число } a_i \neq 0 \}.
	\]
	Элементы этого множества — бесконечные последовательности элементов из \( R \), в которых все члены, начиная с некоторого номера, равны нулю.
	
	\medskip
	\noindent \textbf{Шаг 2: Операции сложения и умножения.}
	Определим на \( P \) две бинарные операции:
	
	\begin{itemize}
		\item \textbf{Сложение (покомпонентное):}
		\[
		(a_0, a_1, a_2, \dots) + (b_0, b_1, b_2, \dots) = (a_0 + b_0, a_1 + b_1, a_2 + b_2, \dots).
		\]
		
		\item \textbf{Умножение (свертка):}
		\[
		(a_0, a_1, a_2, \dots) \cdot (b_0, b_1, b_2, \dots) = (c_0, c_1, c_2, \dots),
		\]
		где 
		\[
		c_k = \sum_{i+j = k} a_i b_j = a_0 b_k + a_1 b_{k-1} + \dots + a_k b_0.
		\]
	\end{itemize}
	
	\medskip
	\noindent \textbf{Шаг 3: Проверка аксиом кольца.}
	Можно проверить, что относительно этих операций \( P \) образует коммутативное кольцо с единицей:
	\begin{itemize}
		\item Коммутативность и ассоциативность сложения очевидны.
		\item Нейтральный элемент по сложению: \( (0, 0, 0, \dots) \).
		\item Противоположный элемент: \( (-a_0, -a_1, -a_2, \dots) \).
		\item Ассоциативность умножения следует из свойств свертки.
		\item Дистрибутивность проверяется непосредственно.
		\item Единица кольца: \( (1_R, 0, 0, \dots) \).
	\end{itemize}
	
	\subsection{Отождествление констант}
	
	Рассмотрим подмножество \( R' \subset P \), состоящее из последовательностей вида \( (a, 0, 0, \dots) \), где \( a \in R \).
	Заметим, что:
	\[
	(a, 0, 0, \dots) + (b, 0, 0, \dots) = (a+b, 0, 0, \dots),
	\]
	\[
	(a, 0, 0, \dots) \cdot (b, 0, 0, \dots) = (ab, 0, 0, \dots).
	\]
	Таким образом, отображение \( a \mapsto (a, 0, 0, \dots) \) является инъективным гомоморфизмом колец. Мы будем отождествлять элемент \( a \in R \) с последовательностью \( (a, 0, 0, \dots) \) и считать \( R \) подкольцом в \( P \).
	
	\subsection{Введение переменной}
	
	Рассмотрим специальный элемент
	\[
	x = (0, 1_R, 0, 0, \dots) \in P.
	\]
	
	Вычислим его степени:
	\[
	x^2 = (0, 0, 1_R, 0, 0, \dots),
	\]
	\[
	x^3 = (0, 0, 0, 1_R, 0, \dots),
	\]
	и так далее. По индукции, \( x^k \) — это последовательность, у которой на \( k \)-й позиции стоит \( 1_R \), а на остальных — нули.
	
	\subsection{Представление произвольного элемента}
	
	Любой элемент \( f \in P \) можно записать в виде конечной линейной комбинации степеней \( x \) с коэффициентами из \( R \). Действительно, пусть
	\[
	f = (a_0, a_1, a_2, \dots, a_n, 0, 0, \dots).
	\]
	Тогда:
	\[
	\begin{aligned}
		f &= (a_0, 0, 0, \dots) + (0, a_1, 0, \dots) + \dots + (0, 0, \dots, a_n, 0, \dots) \\
		&= a_0 + a_1 x + a_2 x^2 + \dots + a_n x^n.
	\end{aligned}
	\]
	
	Таким образом, каждый элемент \( f \in P \) однозначно представляется в виде
	\[
	f = \sum_{i=0}^{n} a_i x^i, \quad a_i \in R,
	\]
	где \( n \) — наибольший индекс с ненулевым коэффициентом (если такой существует).
	
	\subsection{Обозначения и терминология}
	
	Построенное кольцо обозначается \( R[x] \) и называется \textit{кольцом многочленов} от переменной \( x \) над кольцом \( R \).
	
	\begin{tcolorbox}[colback=yellow!5, colframe=orange!50!black, title=Терминология]
		\begin{itemize}
			\item Элементы кольца \( R[x] \) называются \textit{многочленами} (или полиномами).
			\item Кольцо \( R \) называется \textit{кольцом коэффициентов}.
			\item Элементы \( R \), рассматриваемые как многочлены вида \( a = a x^0 \), называются \textit{константами}.
			\item Переменная \( x \) является первичным (трансцендентным) элементом над \( R \).
			\item Если \( f = \sum_{i=0}^{n} a_i x^i \) и \( a_n \neq 0 \), то число \( n \) называется \textit{степенью} многочлена и обозначается \( \deg f \), а \( a_n \) называется \textit{старшим коэффициентом}.
		\end{itemize}
	\end{tcolorbox}
	
	\subsection{Универсальное свойство}
	
	Построенное кольцо многочленов обладает важным универсальным свойством.
	
	\begin{tcolorbox}[colback=blue!5, colframe=blue!40!black, title=Теорема (Универсальное свойство)]
		Для любого коммутативного кольца \( S \) с единицей, любого гомоморфизма колец \( \varphi: R \to S \) и любого элемента \( s \in S \), коммутирующего со всеми \( \varphi(a) \) (т.е. \( s \cdot \varphi(a) = \varphi(a) \cdot s \) для всех \( a \in R \)), существует единственный гомоморфизм колец \( \Phi: R[x] \to S \) такой, что:
		\begin{enumerate}
			\item \( \Phi(a) = \varphi(a) \) для всех \( a \in R \) (т.е. \( \Phi \) продолжает \( \varphi \));
			\item \( \Phi(x) = s \).
		\end{enumerate}
		При этом
		\[
		\Phi\left( \sum_{i=0}^{n} a_i x^i \right) = \sum_{i=0}^{n} \varphi(a_i) s^i.
		\]
	\end{tcolorbox}
	
	Это свойство означает, что подстановка элемента \( s \) вместо переменной \( x \) является корректно определенным гомоморфизмом.
	
	\subsection{Изоморфизм колец многочленов}
	
	Из универсального свойства следует, что кольцо многочленов определено однозначно с точностью до изоморфизма: если \( t \) и \( u \) — два первичных элемента над \( R \) в некоторых расширениях, то существует изоморфизм
	\[
	R[t] \to R[u], \quad \sum a_i t^i \mapsto \sum a_i u^i.
	\]
	
	\myremark
	Таким образом, конструкция через последовательности дает конкретную реализацию абстрактно определенного кольца многочленов.
	
	
	
	
	
	
	
	% ==================== РАЗДЕЛ 30: ДЕЛИМОСТЬ В КОЛЬЦЕ МНОГОЧЛЕНОВ ====================
	\newpage
	\section{Делимость в кольце многочленов над областью целостности (30 вопрос)}
	
	\subsection{Область целостности кольца многочленов}
	
	Напомним, что кольцо \( R \) называется \textit{областью целостности} (или целостным кольцом), если оно коммутативно, имеет единицу \( 1 \neq 0 \) и не содержит делителей нуля.
	
	\medskip
	\noindent \textbf{Теорема 1.} Если \( R \) — область целостности, то кольцо многочленов \( R[x] \) также является областью целостности.
	
	\medskip
	\textit{Доказательство.}
	Пусть \( f(x), g(x) \in R[x] \) — ненулевые многочлены. Запишем их в виде:
	\[
	f(x) = a_0 + a_1x + \dots + a_m x^m, \quad a_m \neq 0,
	\]
	\[
	g(x) = b_0 + b_1x + \dots + b_n x^n, \quad b_n \neq 0.
	\]
	
	Старший член произведения \( f(x)g(x) \) равен \( a_m b_n x^{m+n} \). Так как \( R \) — область целостности, произведение старших коэффициентов \( a_m b_n \neq 0 \). Следовательно, \( f(x)g(x) \neq 0 \). Таким образом, в \( R[x] \) нет делителей нуля. \hfill \(\square\)
	
	\medskip
	\noindent \textbf{Следствие 1.} В условиях теоремы для ненулевых многочленов \( f, g \in R[x] \) выполняется:
	\[
	\deg(f \cdot g) = \deg f + \deg g.
	\]
	
	\myremark
	Это свойство является характеристическим для областей целостности: если в кольце многочленов степень произведения равна сумме степеней, то кольцо коэффициентов не имеет делителей нуля.
	
	\subsection{Делимость и ассоциированность}
	
	\noindent \textbf{Определение.} Пусть \( R \) — область целостности, \( f, g \in R[x] \). Говорят, что \( g \) \textit{делит} \( f \) (обозначение \( g \mid f \)), если существует такой \( h \in R[x] \), что \( f = gh \).
	
	\noindent \textbf{Определение.} Многочлены \( f \) и \( g \) называются \textit{ассоциированными}, если \( f \mid g \) и \( g \mid f \). В области целостности это эквивалентно тому, что \( f = ug \) для некоторого обратимого элемента \( u \in R^* \).
	
	\subsection{Деление с остатком в \( R[x] \)}
	
	Для существования алгоритма деления с остатком требуется обратимость старшего коэффициента делителя.
	
	\medskip
	\noindent \textbf{Лемма (о делении с остатком).} 
	Пусть \( R \) — коммутативное кольцо с 1, и пусть \( f, g \in R[x] \), причем \( g \neq 0 \) и старший коэффициент многочлена \( g \) обратим в \( R \) (т.е. является единицей кольца). Тогда существуют единственные многочлены \( q, r \in R[x] \) такие, что:
	\[
	f = gq + r,
	\]
	где либо \( r = 0 \), либо \( \deg r < \deg g \).
	
	\medskip
	\textit{Доказательство.}
	
	\textit{1. Существование.} Проведем индукцию по степени \( \deg f = n \).
	
	\textit{База индукции:} Если \( \deg f < \deg g \) или \( f = 0 \), то можно взять \( q = 0 \), \( r = f \), так как \( f = g \cdot 0 + f \).
	
	\textit{Индукционный переход:} Пусть \( \deg f = n \geq m = \deg g \). Запишем многочлены:
	\[
	f = a_n x^n + a_{n-1} x^{n-1} + \dots + a_0, \quad a_n \neq 0,
	\]
	\[
	g = b_m x^m + b_{m-1} x^{m-1} + \dots + b_0, \quad b_m \in R^*.
	\]
	
	Рассмотрим многочлен
	\[
	h = f - a_n b_m^{-1} x^{n-m} g.
	\]
	По построению, старшие члены сокращаются, поэтому \( \deg h < n \) или \( h = 0 \).
	
	Применим к \( h \) предположение индукции: существуют \( q' \) и \( r \) такие, что \( h = g q' + r \), где \( \deg r < \deg g \) или \( r = 0 \). Тогда
	\[
	f = h + a_n b_m^{-1} x^{n-m} g = g q' + r + a_n b_m^{-1} x^{n-m} g = g (q' + a_n b_m^{-1} x^{n-m}) + r.
	\]
	Таким образом, искомая пара: \( q = q' + a_n b_m^{-1} x^{n-m} \), \( r \) — остаток.
	
	\textit{2. Единственность.} Предположим, что существуют две пары \( (q_1, r_1) \) и \( (q_2, r_2) \), удовлетворяющие условиям:
	\[
	f = g q_1 + r_1 = g q_2 + r_2.
	\]
	Тогда
	\[
	g (q_1 - q_2) = r_2 - r_1.
	\]
	
	Если \( q_1 \neq q_2 \), то левая часть имеет степень не ниже \( \deg g \) (так как старший коэффициент \( g \) обратим, а значит, не является делителем нуля). Правая же часть имеет степень меньше \( \deg g \) (как разность остатков). Противоречие. Следовательно, \( q_1 = q_2 \), и тогда \( r_1 = r_2 \). \hfill \(\square\)
	
	\subsection{Случай поля коэффициентов}
	
	Если кольцо коэффициентов является полем \( K \), то любой ненулевой элемент поля обратим. Поэтому условие обратимости старшего коэффициента выполняется для любого ненулевого делителя \( g \in K[x] \).
	
	\medskip
	\noindent \textbf{Следствие 2.} Для любой пары многочленов \( f, g \in K[x] \), \( g \neq 0 \), существуют единственные \( q, r \in K[x] \) такие, что
	\[
	f = gq + r,
	\]
	где \( \deg r < \deg g \) или \( r = 0 \).
	
	\subsection{Евклидовость кольца многочленов над полем}
	
	Наличие деления с остатком позволяет ввести понятие евклидовой нормы.
	
	\medskip
	\noindent \textbf{Определение.} Область целостности \( R \) называется \textit{евклидовым кольцом}, если существует функция \( \nu: R \setminus \{0\} \to \mathbb{N}_0 \) (евклидова норма) такая, что для любых \( a, b \in R \), \( b \neq 0 \), найдутся \( q, r \in R \) со свойствами:
	\[
	a = bq + r, \quad \text{где либо } r = 0, \text{ либо } \nu(r) < \nu(b).
	\]
	
	\medskip
	\noindent \textbf{Теорема 2.} Кольцо многочленов \( K[x] \) над полем \( K \) является евклидовым кольцом с евклидовой нормой \( \nu(f) = \deg f \).
	
	\textit{Доказательство.} Проверим аксиомы евклидова кольца:
	\begin{itemize}
		\item \( K[x] \) — область целостности (по теореме 1).
		\item Для любых \( f, g \in K[x] \), \( g \neq 0 \), деление с остатком существует и единственно по следствию 2.
		\item Условие на норму выполнено: \( \deg r < \deg g \) или \( r = 0 \).
	\end{itemize}
	\hfill \(\square\)
	
	\myremark
	Евклидовость кольца многочленов над полем влечет за собой выполнение основной теоремы арифметики (однозначное разложение на неприводимые множители) и существование алгоритма Евклида для нахождения наибольшего общего делителя.
	
	\subsection{Наибольший общий делитель многочленов}
	
	В евклидовом кольце \( K[x] \) можно определить понятие наибольшего общего делителя.
	
	\medskip
	\noindent \textbf{Определение.} Пусть \( f, g \in K[x] \). Многочлен \( d \in K[x] \) называется \textit{наибольшим общим делителем} (НОД) многочленов \( f \) и \( g \), если:
	\begin{enumerate}
		\item \( d \mid f \) и \( d \mid g \);
		\item Если \( h \mid f \) и \( h \mid g \), то \( h \mid d \).
	\end{enumerate}
	
	\medskip
	\noindent \textbf{Алгоритм Евклида.} Благодаря делению с остатком, НОД можно находить последовательностью делений:
	\[
	\begin{aligned}
		f &= g q_1 + r_1, \quad \deg r_1 < \deg g, \\
		g &= r_1 q_2 + r_2, \quad \deg r_2 < \deg r_1, \\
		r_1 &= r_2 q_3 + r_3, \quad \deg r_3 < \deg r_2, \\
		&\vdots \\
		r_{k-2} &= r_{k-1} q_k + r_k, \\
		r_{k-1} &= r_k q_{k+1} + 0.
	\end{aligned}
	\]
	Последний ненулевой остаток \( r_k \) и будет НОД(\( f, g \)) с точностью до умножения на константу.
	
	\subsection{Линейное представление НОД}
	
	Из алгоритма Евклида следует важнейшее свойство: НОД двух многочленов всегда можно выразить в виде их линейной комбинации.
	
	\medskip
	\noindent \textbf{Теорема 3 (Тождество Безу).} Для любых \( f, g \in K[x] \) существуют такие \( u, v \in K[x] \), что
	\[
	u f + v g = \operatorname{НОД}(f, g).
	\]
	
	\myremark
	Это представление не является единственным: если \( (u, v) \) — одна пара, то \( (u + t g, v - t f) \) при любом \( t \in K[x] \) также дает представление того же НОД.
	
	
	
	% ==================== РАЗДЕЛ 31: ТЕОРЕМА БЕЗУ И ЧИСЛО КОРНЕЙ ====================
	\newpage
	\section{Теорема Безу. Теорема о числе корней многочлена (31 вопрос)}
	
	\subsection{Значение многочлена в точке}
	
	Прежде чем формулировать теорему Безу, необходимо определить, что значит вычислить многочлен в точке, если коэффициенты и точка принадлежат некоторому кольцу.
	
	\medskip
	
	\noindent \textbf{Определение.} Пусть \( R \) — коммутативное кольцо с единицей, и пусть \( A \) — расширение кольца \( R \) (т.е. \( R \subseteq A \) и единица \( 1_R \) является единицей в \( A \)). Пусть \( \xi \in A \) таков, что \( \forall a \in R: a\xi = \xi a \) (элемент \( \xi \) коммутирует со всеми элементами \( R \)).
	
	Для многочлена \( f \in R[t] \), записанного в виде
	\[
	f = \sum_{i=0}^{n} a_i t^i = a_0 + a_1 t + a_2 t^2 + \dots + a_n t^n,
	\]
	определим его \textit{значение} в точке \( \xi \) как элемент
	\[
	f(\xi) = \sum_{i=0}^{n} a_i \xi^i = a_0 + a_1 \xi + a_2 \xi^2 + \dots + a_n \xi^n \in A.
	\]
	
	\medskip
	\noindent \textbf{Определение.} Элемент \( \xi \in A \) называется \textit{корнем} многочлена \( f \in R[t] \), если \( f(\xi) = 0 \).
	
	\subsection{Теорема Безу}
	
	Теорема Безу устанавливает связь между делимостью многочлена на линейный двучлен и значением многочлена в соответствующей точке.
	
	\begin{tcolorbox}[colback=blue!5, colframe=blue!40!black, title=Теорема 1 (Безу)]
		Остаток от деления многочлена \( f \in R[t] \) на двучлен \( t - d \) (где \( d \in R \)) равен \( f(d) \).
	\end{tcolorbox}
	
	\medskip
	\textit{Доказательство.}
	Старший коэффициент двучлена \( t - d \) равен 1, который обратим в любом кольце с единицей. Следовательно, по лемме о делении с остатком (см. вопрос 30), существуют единственные многочлены \( q \in R[t] \) и \( r \in R \) такие, что:
	\[
	f = (t - d) q + r,
	\]
	где \( r = \text{const} \) (так как степень остатка должна быть меньше степени делителя \( \deg(t-d) = 1 \)).
	
	Подставим \( t = d \) в это равенство. Учитывая, что \( (d - d) = 0 \), получаем:
	\[
	f(d) = (d - d) q(d) + r = 0 \cdot q(d) + r = r.
	\]
	Таким образом, \( r = f(d) \). \hfill \(\square\)
	
	\medskip
	\noindent \textbf{Следствие 1 (Признак корня).}
	Многочлен \( f \in R[t] \) делится на двучлен \( t - d \) (где \( d \in R \)) без остатка тогда и только тогда, когда \( d \) является корнем \( f \), т.е. \( f(d) = 0 \).
	
	\myremark
	Это следствие часто также называют теоремой Безу. Оно дает простой способ проверки, является ли данный элемент кольца корнем многочлена.
	
	\subsection{Теорема о числе корней многочлена}
	
	Перейдем к вопросу о максимально возможном количестве корней многочлена. Здесь принципиальным является требование, чтобы кольцо коэффициентов было областью целостности.
	
	\begin{tcolorbox}[colback=green!5, colframe=green!50!black, title=Теорема 2 (о числе корней)]
		Пусть \( R \) — область целостности, и пусть \( f \in R[x] \) — ненулевой многочлен степени \( n \). Тогда \( f \) имеет в \( R \) не более \( n \) различных корней.
	\end{tcolorbox}
	
	\medskip
	\textit{Доказательство.}
	Применим метод математической индукции по степени многочлена \( n \).
	
	\textit{База индукции:} \( n = 0 \). Многочлен нулевой степени — это ненулевая константа \( f(x) = c \neq 0 \). Такой многочлен не имеет корней, и \( 0 \leq 0 \) — верно.
	
	\textit{Индукционный переход:} Предположим, что утверждение верно для всех многочленов степени \( n-1 \). Докажем его для многочлена \( f(x) \) степени \( n \).
	
	Если \( f(x) \) не имеет корней в \( R \), то утверждение (\( 0 \leq n \)) очевидно выполнено.
	
	Пусть \( f(x) \) имеет хотя бы один корень \( c_1 \in R \). Тогда по теореме Безу (следствие 1) многочлен \( f(x) \) делится на \( (x - c_1) \):
	\[
	f(x) = (x - c_1) \, h(x),
	\]
	где \( h(x) \in R[x] \) и, так как \( R \) — область целостности, \( \deg h = n-1 \).
	
	Рассмотрим любой другой корень \( c_2 \in R \), \( c_2 \neq c_1 \). Тогда:
	\[
	0 = f(c_2) = (c_2 - c_1) h(c_2).
	\]
	Поскольку \( R \) — область целостности и \( c_2 - c_1 \neq 0 \), произведение может равняться нулю только если \( h(c_2) = 0 \). Таким образом, любой корень \( f \), отличный от \( c_1 \), является корнем \( h(x) \).
	
	По предположению индукции, многочлен \( h(x) \) степени \( n-1 \) имеет не более \( n-1 \) корней в \( R \). Следовательно, \( f(x) \) имеет не более \( 1 + (n-1) = n \) корней. \hfill \(\square\)
	
	\subsection{Существенность условия области целостности}
	
	Условие, что \( R \) является областью целостности, является существенным. В кольцах с делителями нуля многочлен может иметь корней больше, чем его степень.
	
	\medskip
	\noindent \textbf{Пример.} Рассмотрим кольцо \( \mathbb{Z}_8 \) (вычеты по модулю 8). Это кольцо не является областью целостности, так как, например, \( 2 \cdot 4 = 0 \). Многочлен \( f(x) = x^2 - 1 \) степени 2 имеет в \( \mathbb{Z}_8 \) четыре корня:
	\[
	f(1) = 1 - 1 = 0, \quad f(3) = 9 - 1 = 8 \equiv 0 \pmod{8},
	\]
	\[
	f(5) = 25 - 1 = 24 \equiv 0 \pmod{8}, \quad f(7) = 49 - 1 = 48 \equiv 0 \pmod{8}.
	\]
	Таким образом, \( x^2 - 1 \) имеет 4 корня, что больше его степени.
	
	\subsection{Следствия из теоремы о числе корней}
	
	\begin{tcolorbox}[colback=yellow!5, colframe=orange!50!black, title=Следствия]
		\begin{enumerate}
			\item Если два многочлена \( f, g \in R[x] \) степени не выше \( n \) совпадают в \( n+1 \) различных точках кольца \( R \), то \( f = g \).
			
			\item Многочлен над бесконечной областью целостности (например, над полем \( \mathbb{R} \) или \( \mathbb{C} \)) однозначно определяется своими значениями.
			
			\item Если многочлен \( f \in R[x] \) имеет более \( \deg f \) корней, то \( f = 0 \) (нулевой многочлен).
		\end{enumerate}
	\end{tcolorbox}
	
	\myremark
	Последнее следствие часто используется для доказательства равенства многочленов: если разность двух многочленов имеет бесконечно много корней (или больше, чем степень), то она тождественно равна нулю.
	
	
	
	
	
	
	
	
	% ==================== РАЗДЕЛ 32: НЕПРИВОДИМЫЕ МНОГОЧЛЕНЫ ====================
	\newpage
	\section{Неприводимые многочлены. Свойства (32 вопрос)}
	
	\subsection{Определения}
	
	Пусть \( K \) — поле (или, в более общем случае, область целостности). Многочлены, которые нельзя разложить в произведение многочленов меньшей степени, играют роль, аналогичную простым числам в арифметике целых чисел.
	
	\medskip
	
	\noindent \textbf{Определение 1 (Нормированный многочлен).}
	Многочлен, старший коэффициент которого равен 1, называется \textit{нормированным} (или \textit{унитарным}).
	
	\medskip
	
	\noindent \textbf{Определение 2 (Неприводимый многочлен).}
	Пусть \( K \) — поле. Многочлен \( p \in K[x] \) называется \textit{неприводимым} над \( K \), если:
	\begin{enumerate}
		\item \( \deg p \geq 1 \) (т.е. \( p \) не является константой);
		\item Из равенства \( p = f \cdot g \) с \( f, g \in K[x] \) следует, что либо \( f \), либо \( g \) является константой (обратимым элементом \( K \)).
	\end{enumerate}
	
	Если многочлен степени \( \geq 1 \) не является неприводимым, он называется \textit{приводимым}.
	
	\myremark
	Обратимыми элементами в \( K[x] \) являются ненулевые константы (элементы поля \( K^* \)). Два многочлена называются \textit{ассоциированными}, если они отличаются умножением на ненулевую константу.
	
	\subsection{Основные свойства неприводимых многочленов над полем}
	
	Далее \( K \) обозначает поле, а все многочлены рассматриваются в кольце \( K[x] \).
	
	\begin{tcolorbox}[colback=blue!5, colframe=blue!40!black, title=Свойство 1 (Взаимная простота с неприводимым)]
		Если \( f \in K[x] \) и \( p \in K[x] \) неприводим над \( K \), то либо \( p \mid f \), либо \( \operatorname{НОД}(f, p) = 1 \) (т.е. \( f \) и \( p \) взаимно просты).
	\end{tcolorbox}
	
	\medskip
	\textit{Доказательство.}
	Рассмотрим нормированный наибольший общий делитель \( d = \operatorname{НОД}(f, p) \). Так как \( d \mid p \) и \( p \) неприводим, то \( d \) либо ассоциирован с \( p \) (т.е. \( d = c \cdot p \) для некоторого \( c \in K^* \)), либо является константой (обратимым элементом).
	
	\begin{itemize}
		\item В первом случае \( p \mid d \), а так как \( d \mid f \), то \( p \mid f \).
		\item Во втором случае \( d \) — константа, следовательно, \( f \) и \( p \) взаимно просты. \hfill \(\square\)
	\end{itemize}
	
	\medskip
	\begin{tcolorbox}[colback=blue!5, colframe=blue!40!black, title=Свойство 2 (Свойство простого делителя)]
		Если произведение \( f_1 f_2 \) делится на неприводимый многочлен \( p \), то хотя бы один из сомножителей делится на \( p \).
	\end{tcolorbox}
	
	\medskip
	\textit{Доказательство.}
	Предположим, что \( p \nmid f_1 \). Тогда по свойству 1, \( \operatorname{НОД}(f_1, p) = 1 \). Следовательно, существуют многочлены \( u, v \in K[x] \) такие, что (тождество Безу):
	\[
	u f_1 + v p = 1.
	\]
	Умножим это равенство на \( f_2 \):
	\[
	u f_1 f_2 + v p f_2 = f_2.
	\]
	По условию, \( p \mid f_1 f_2 \), значит, \( p \) делит левую часть (первое слагаемое делится на \( p \) по условию, второе — очевидно). Следовательно, \( p \mid f_2 \). \hfill \(\square\)
	
	\medskip
	\begin{tcolorbox}[colback=blue!5, colframe=blue!40!black, title=Свойство 3 (Обобщение на произведение)]
		Если произведение \( f_1 f_2 \dots f_k \) делится на неприводимый многочлен \( p \), то хотя бы один из множителей \( f_i \) делится на \( p \).
	\end{tcolorbox}
	
	\medskip
	\textit{Доказательство.}
	Применяем индукцию по \( k \). База \( k = 2 \) доказана в свойстве 2. Пусть утверждение верно для \( k-1 \). Если \( p \mid (f_1 f_2 \dots f_k) \), то по свойству 2 либо \( p \mid f_k \), либо \( p \mid (f_1 \dots f_{k-1}) \). Во втором случае по предположению индукции \( p \) делит один из \( f_1, \dots, f_{k-1} \). \hfill \(\square\)
	
	\medskip
	\begin{tcolorbox}[colback=blue!5, colframe=blue!40!black, title=Свойство 4 (Взаимная простота двух неприводимых)]
		Если \( p_1 \) и \( p_2 \) — два неприводимых многочлена над \( K \), не являющихся ассоциированными, то они взаимно просты.
	\end{tcolorbox}
	
	\medskip
	\textit{Доказательство.}
	Пусть \( d = \operatorname{НОД}(p_1, p_2) \). Так как \( d \mid p_1 \) и \( p_1 \) неприводим, то \( d \) либо ассоциирован с \( p_1 \), либо является константой. Если \( d \) ассоциирован с \( p_1 \), то \( p_1 \mid p_2 \). Но тогда в силу неприводимости \( p_2 \) и того, что \( \deg p_1 \geq 1 \), получаем, что \( p_1 \) и \( p_2 \) ассоциированы, что противоречит условию. Следовательно, \( d \) — константа. \hfill \(\square\)
	
	\medskip
	\begin{tcolorbox}[colback=blue!5, colframe=blue!40!black, title=Свойство 5 (Корень и делимость)]
		Пусть \( p \in K[x] \) неприводим над \( K \). Если в некотором расширении поля \( K \) (например, в алгебраическом замыкании) многочлен \( p \) имеет корень \( \alpha \), и этот корень является также корнем многочлена \( f \in K[x] \), то \( p \mid f \).
	\end{tcolorbox}
	
	\medskip
	\textit{Доказательство.}
	Рассмотрим \( d = \operatorname{НОД}(f, p) \) в кольце \( K[x] \). Так как \( \alpha \) является корнем и \( f \), и \( p \), то минимальный многочлен элемента \( \alpha \) над \( K \) делит оба многочлена. Но в силу неприводимости \( p \), этот минимальный многочлен ассоциирован с \( p \). Следовательно, \( p \mid d \), а так как \( d \mid f \), то \( p \mid f \). \hfill \(\square\)
	
	\subsection{Примеры и замечания}
	
	\begin{itemize}
		\item Над полем комплексных чисел \( \mathbb{C} \) (в силу основной теоремы алгебры) неприводимы только многочлены первой степени.
		\item Над полем действительных чисел \( \mathbb{R} \) неприводимы многочлены первой степени и второй степени с отрицательным дискриминантом (например, \( x^2 + 1 \)).
		\item Над полем рациональных чисел \( \mathbb{Q} \) существуют неприводимые многочлены любой степени (например, \( x^n - 2 \) неприводим по критерию Эйзенштейна).
	\end{itemize}
	
	\myremark
	Свойства 1–5 показывают, что неприводимые многочлены над полем ведут себя аналогично простым числам в кольце целых чисел. Это позволяет строить теорию делимости и однозначность разложения на множители в \( K[x] \).
	
	
	
	
	
	% ==================== РАЗДЕЛ 33: ФАКТОРИАЛЬНОСТЬ КОЛЬЦА МНОГОЧЛЕНОВ ====================
	\newpage
	\section{Факториальность кольца многочленов (33 вопрос)}
	
	\subsection{Введение}
	
	Напомним, что кольцо называется \textit{факториальным} (или кольцом с однозначным разложением на множители), если в нем каждый ненулевой необратимый элемент может быть представлен в виде произведения неприводимых элементов, и такое представление единственно с точностью до порядка сомножителей и умножения на обратимые элементы.
	
	В данном разделе мы докажем, что кольцо многочленов \( K[x] \) над полем \( K \) является факториальным.
	
	\subsection{Существование разложения на неприводимые множители}
	
	\begin{tcolorbox}[colback=blue!5, colframe=blue!40!black, title=Предложение 1 (Существование разложения)]
		Каждый многочлен \( f \in K[x] \) степени \( \deg f \geq 1 \) делится по крайней мере на один неприводимый над \( K \) многочлен. Более того, \( f \) может быть представлен в виде произведения неприводимых многочленов.
	\end{tcolorbox}
	
	\medskip
	\textit{Доказательство.}
	Применим индукцию по степени многочлена \( n = \deg f \).
	
	\textit{База индукции:} \( n = 1 \). Многочлен первой степени всегда неприводим над полем (так как его нельзя разложить в произведение многочленов меньшей положительной степени). Следовательно, он сам является искомым произведением (из одного множителя).
	
	\textit{Индукционный переход:} Предположим, что утверждение верно для всех многочленов степени меньше \( n \). Рассмотрим многочлен \( f \) степени \( n \).
	
	Если \( f \) неприводим, то он сам является своим разложением. Если \( f \) приводим, то по определению существуют многочлены \( f_1, f_2 \in K[x] \) такие, что \( f = f_1 f_2 \), причем \( 1 \leq \deg f_1 < n \) и \( 1 \leq \deg f_2 < n \).
	
	По предположению индукции, \( f_1 \) и \( f_2 \) раскладываются в произведения неприводимых многочленов:
	\[
	f_1 = p_1 p_2 \dots p_k, \quad f_2 = q_1 q_2 \dots q_m,
	\]
	где все \( p_i, q_j \) неприводимы. Тогда
	\[
	f = p_1 p_2 \dots p_k q_1 q_2 \dots q_m
	\]
	есть искомое разложение. \hfill \(\square\)
	
	\subsection{Единственность разложения}
	
	Для доказательства единственности нам потребуются свойства неприводимых многочленов, установленные в вопросе 32.
		
	\begin{tcolorbox}[colback=green!5, colframe=green!50!black, title={Теорема (Факториальность $K[x]$)}]
		Любой многочлен $f \in K[x]$ степени $\deg f \geq 1$ может быть представлен в виде произведения неприводимых над $K$ многочленов. Такое представление единственно с точностью до:
		\begin{itemize}
			\item порядка следования сомножителей;
			\item умножения сомножителей на обратимые элементы поля $K$ (т.е. на ненулевые константы).
		\end{itemize}
	\end{tcolorbox}
	
	\medskip
	\textit{Доказательство.}
	Существование разложения уже доказано в предложении 1. Докажем единственность индукцией по степени многочлена.
	
	\textit{База индукции:} \( \deg f = 1 \). Многочлен первой степени неприводим, поэтому любое его разложение может состоять только из одного множителя (с точностью до умножения на константу).
	
	\textit{Индукционный переход:} Предположим, что утверждение верно для всех многочленов степени меньше \( n \). Пусть многочлен \( f \) степени \( n \) имеет два разложения на неприводимые множители:
	\[
	f = p_1 p_2 \dots p_k = q_1 q_2 \dots q_m,
	\]
	где все \( p_i, q_j \) неприводимы над \( K \).
	
	Рассмотрим неприводимый множитель \( p_1 \). Так как \( p_1 \mid f \), то \( p_1 \mid (q_1 q_2 \dots q_m) \). По свойству простого делителя (свойство 3 из вопроса 32), \( p_1 \) делит хотя бы один из множителей \( q_j \). Перенумеруем множители так, чтобы \( p_1 \mid q_1 \).
	
	Поскольку \( q_1 \) неприводим, а \( p_1 \) не является константой (иначе разложение было бы тривиальным), они ассоциированы. Следовательно, существует такая ненулевая константа \( c_1 \in K^* \), что
	\[
	q_1 = c_1 p_1.
	\]
	
	Подставим это выражение в исходное равенство:
	\[
	p_1 p_2 \dots p_k = c_1 p_1 q_2 \dots q_m.
	\]
	Сократим на \( p_1 \) (кольцо \( K[x] \) — область целостности, сокращение допустимо):
	\[
	p_2 \dots p_k = (c_1 q_2) \dots q_m.
	\]
	
	Многочлен в левой части имеет степень \( n - \deg p_1 < n \). По предположению индукции, его разложение на неприводимые множители единственно с точностью до порядка и ассоциированности. Это означает, что после соответствующей перенумерации и, возможно, умножения на константы, множества множителей \( \{p_2, \dots, p_k\} \) и \( \{c_1 q_2, q_3, \dots, q_m\} \) совпадают.
	
	Таким образом, исходные разложения \( f \) совпадают с точностью до порядка и умножения на константы. \hfill \(\square\)
	
	\subsection{Каноническое разложение}
	
	Если привести все неприводимые множители к нормированному виду (т.е. умножить на подходящие константы так, чтобы старший коэффициент стал равен 1) и сгруппировать одинаковые множители, то получим \textit{каноническое разложение}.
	
	\medskip
	\noindent \textbf{Следствие (Каноническое разложение).}
	Любой ненулевой многочлен \( f \in K[x] \) может быть единственным образом записан в виде:
	\[
	f = a_0 \cdot p_1^{n_1} p_2^{n_2} \dots p_m^{n_m},
	\]
	где:
	\begin{itemize}
		\item \( a_0 \in K^* \) — старший коэффициент многочлена \( f \);
		\item \( p_1, p_2, \dots, p_m \) — различные нормированные неприводимые многочлены над \( K \);
		\item \( n_1, n_2, \dots, n_m \in \mathbb{N} \) — положительные целые числа (показатели степеней).
	\end{itemize}
	Единственность понимается с точностью до порядка сомножителей.
	
	\myremark
	Каноническое разложение является аналогом основной теоремы арифметики для кольца многочленов над полем. Оно широко используется при исследовании делимости, нахождении наибольшего общего делителя и наименьшего общего кратного многочленов.
	
	\subsection{Пример}
	
	Рассмотрим многочлен \( f(x) = 2x^4 - 2 \) над полем рациональных чисел \( \mathbb{Q} \).
	
	Сначала вынесем старший коэффициент:
	\[
	f(x) = 2(x^4 - 1).
	\]
	
	Разложим \( x^4 - 1 \):
	\[
	x^4 - 1 = (x^2 - 1)(x^2 + 1) = (x - 1)(x + 1)(x^2 + 1).
	\]
	
	Многочлен \( x^2 + 1 \) неприводим над \( \mathbb{Q} \). Таким образом, каноническое разложение:
	\[
	f(x) = 2 \cdot (x - 1) \cdot (x + 1) \cdot (x^2 + 1).
	\]
	
	Здесь \( a_0 = 2 \), а нормированные неприводимые множители: \( (x - 1) \), \( (x + 1) \), \( (x^2 + 1) \), каждый в первой степени.
	
	\subsection{Обобщение}
	
	Факториальность кольца многочленов над полем является частным случаем более общей теоремы: если \( R \) — факториальное кольцо, то и \( R[x] \) факториально. Однако доказательство этого факта требует более тонких рассмотрений (лемма Гаусса и т.д.). В нашем случае \( R = K \) — поле, что значительно упрощает ситуацию.
	
	
	
	
	
	% ==================== РАЗДЕЛ 34: ПОЛЯ ====================
	\newpage
	\section{Поля. Алгебраически замкнутое поле. Примеры. Изоморфизм полей. Автоморфизмы поля рациональных чисел, комплексных чисел (34 вопрос)}
	
	\subsection{Определение поля}
	
	\begin{tcolorbox}[colback=blue!5, colframe=blue!40!black, title=Определение поля]
		Множество \( K \) с двумя бинарными операциями \( + \) (сложение) и \( \cdot \) (умножение) называется \textit{полем}, если выполнены следующие аксиомы:
		\begin{enumerate}
			\item Сложение ассоциативно: \( \forall a,b,c \in K: (a+b)+c = a+(b+c) \).
			\item Сложение коммутативно: \( \forall a,b \in K: a+b = b+a \).
			\item Существует нейтральный элемент по сложению (нуль): \( \exists 0 \in K: \forall a \in K, a+0 = 0+a = a \).
			\item Существует противоположный элемент: \( \forall a \in K, \exists (-a) \in K: a+(-a) = (-a)+a = 0 \).
			\item Умножение ассоциативно: \( \forall a,b,c \in K: (ab)c = a(bc) \).
			\item Умножение коммутативно: \( \forall a,b \in K: ab = ba \).
			\item Умножение дистрибутивно относительно сложения: \( \forall a,b,c \in K: a(b+c) = ab+ac \).
			\item Существует нейтральный элемент по умножению (единица), причём \( 1 \neq 0 \): \( \exists 1 \in K: \forall a \in K, a \cdot 1 = 1 \cdot a = a \).
			\item Каждый ненулевой элемент имеет обратный по умножению: \( \forall a \in K, a \neq 0, \exists a^{-1} \in K: a a^{-1} = a^{-1} a = 1 \).
		\end{enumerate}
	\end{tcolorbox}
	
	\subsection{Примеры полей}
	
	Приведём основные примеры полей:
	
	\begin{itemize}
		\item Поле рациональных чисел \( \mathbb{Q} \);
		\item Поле действительных чисел \( \mathbb{R} \);
		\item Поле комплексных чисел \( \mathbb{C} \);
		\item Поля вычетов по простому модулю \( \mathbb{F}_p = \mathbb{Z}/p\mathbb{Z} \), где \( p \) — простое число;
		\item Поля рациональных функций \( K(x) \) над полем \( K \).
	\end{itemize}
	
	\subsection{Алгебраически замкнутое поле}
	
	\begin{tcolorbox}[colback=green!5, colframe=green!50!black, title=Определение]
		Поле \( K \) называется \textit{алгебраически замкнутым}, если любой многочлен \( f(x) \in K[x] \) степени \( \geq 1 \) имеет хотя бы один корень в \( K \).
	\end{tcolorbox}
	
	Из этого определения следует, что в алгебраически замкнутом поле любой многочлен раскладывается на линейные множители.
	
	\medskip
	\noindent \textbf{Теорема (Основная теорема алгебры).}
	Поле комплексных чисел \( \mathbb{C} \) алгебраически замкнуто: любой многочлен с комплексными коэффициентами, отличный от константы, имеет по крайней мере один комплексный корень.
	
	\medskip
	\noindent \textbf{Примеры алгебраически замкнутых полей:}
	\begin{itemize}
		\item \( \mathbb{C} \) — поле комплексных чисел;
		\item Поле алгебраических чисел \( \overline{\mathbb{Q}} \) (алгебраическое замыкание \( \mathbb{Q} \) в \( \mathbb{C} \));
		\item Алгебраическое замыкание \( \overline{\mathbb{F}}_p \) конечного поля \( \mathbb{F}_p \).
	\end{itemize}
	
	\medskip
	\noindent \textbf{Примеры полей, не являющихся алгебраически замкнутыми:}
	\begin{itemize}
		\item \( \mathbb{Q} \) (многочлен \( x^2 - 2 \) не имеет корней в \( \mathbb{Q} \));
		\item \( \mathbb{R} \) (многочлен \( x^2 + 1 \) не имеет корней в \( \mathbb{R} \));
		\item Конечные поля \( \mathbb{F}_p \) (многочлен \( x^p - x + 1 \) не имеет корней в \( \mathbb{F}_p \), что следует из малой теоремы Ферма).
	\end{itemize}
	
	\subsection{Изоморфизм полей}
	
	\begin{tcolorbox}[colback=yellow!5, colframe=orange!50!black, title=Определение]
		Пусть \( (K, +, \cdot) \) и \( (K', \oplus, \odot) \) — поля. Отображение \( \varphi: K \to K' \) называется \textit{гомоморфизмом полей}, если оно сохраняет операции:
		\[
		\varphi(a + b) = \varphi(a) \oplus \varphi(b), \quad \varphi(a \cdot b) = \varphi(a) \odot \varphi(b) \quad \forall a,b \in K.
		\]
	\end{tcolorbox}
	
	Из определения автоматически следует, что \( \varphi(0) = 0' \), \( \varphi(1) = 1' \), и для любых целых \( n \in \mathbb{Z} \) и \( a \in K \):
	\[
	\varphi(na) = n \varphi(a), \quad \varphi(a^n) = [\varphi(a)]^n.
	\]
	
	\medskip
	\noindent \textbf{Определение.}
	\begin{itemize}
		\item \textit{Ядро} гомоморфизма: \( \ker \varphi = \{ a \in K \mid \varphi(a) = 0' \} \).
		\item Гомоморфизм полей всегда инъективен, так как его ядро — идеал в поле, а поле не имеет нетривиальных идеалов. Следовательно, \( \ker \varphi = \{0\} \) или \( \ker \varphi = K \). Так как \( \varphi(1) = 1' \neq 0' \), то \( \ker \varphi = \{0\} \). Таким образом, любой гомоморфизм полей является \textit{мономорфизмом} (инъективным).
		\item \textit{Образ} гомоморфизма: \( \operatorname{Im} \varphi = \{ a' \in K' \mid a' = \varphi(a) \text{ для некоторого } a \in K \} \) является подполем в \( K' \).
		\item Гомоморфизм полей называется \textit{изоморфизмом}, если он биективен (сюръективен). В этом случае поля \( K \) и \( K' \) называются \textit{изоморфными} (обозначение: \( K \cong K' \)).
	\end{itemize}
	
	\subsection{Автоморфизмы поля}
	
	\begin{tcolorbox}[colback=red!5, colframe=red!50!black, title=Определение]
		\textit{Автоморфизмом} поля \( K \) называется изоморфизм \( \varphi: K \to K \) поля на себя.
	\end{tcolorbox}
	
	Множество всех автоморфизмов поля \( K \) образует группу относительно композиции. Она обозначается \( \operatorname{Aut}(K) \).
	
	\subsubsection*{Автоморфизмы поля рациональных чисел \( \mathbb{Q} \)}
	
	\begin{tcolorbox}[colback=blue!5, colframe=blue!40!black, title=Теорема]
		Поле рациональных чисел \( \mathbb{Q} \) имеет единственный автоморфизм — тождественный:
		\[
		\operatorname{Aut}(\mathbb{Q}) = \{ \operatorname{id} \}.
		\]
	\end{tcolorbox}
	
	\medskip
	\textit{Доказательство.}
	Пусть \( \varphi: \mathbb{Q} \to \mathbb{Q} \) — автоморфизм. Тогда \( \varphi(1) = 1 \). Для любого натурального \( n \in \mathbb{N} \):
	\[
	\varphi(n) = \varphi(1 + 1 + \dots + 1) = \varphi(1) + \varphi(1) + \dots + \varphi(1) = n \cdot 1 = n.
	\]
	Для целых чисел: \( \varphi(-n) = -\varphi(n) = -n \). Для рациональных чисел \( \frac{p}{q} \) (\( q \neq 0 \)):
	\[
	\varphi\left(\frac{p}{q}\right) = \varphi(p) \cdot \varphi(q^{-1}) = \varphi(p) \cdot (\varphi(q))^{-1} = p \cdot q^{-1} = \frac{p}{q}.
	\]
	Таким образом, \( \varphi(x) = x \) для всех \( x \in \mathbb{Q} \). \hfill \(\square\)
	
	\subsubsection*{Автоморфизмы поля комплексных чисел \( \mathbb{C} \)}
	
	Поле комплексных чисел имеет значительно больше автоморфизмов. Непрерывных (относительно обычной топологии) автоморфизмов всего два.
	
	\begin{tcolorbox}[colback=blue!5, colframe=blue!40!black, title=Теорема]
		Непрерывные автоморфизмы поля комплексных чисел \( \mathbb{C} \) (как расширения \( \mathbb{R} \)) суть тождественное отображение и комплексное сопряжение:
		\[
		\varphi_1(z) = z, \quad \varphi_2(z) = \overline{z}.
		\]
	\end{tcolorbox}
	
	\medskip
	\textit{Доказательство.}
	Любой непрерывный автоморфизм \( \varphi: \mathbb{C} \to \mathbb{C} \) должен переводить вещественные числа в вещественные (так как вещественные числа характеризуются тем, что они инвариантны относительно комплексного сопряжения, но это требует отдельного рассуждения). Известно, что любой автоморфизм поля \( \mathbb{C} \) либо тождествен на \( \mathbb{R} \), либо переводит \( i \) в \( -i \).
	
	Действительно, так как \( \varphi \) сохраняет операции, то:
	\[
	\varphi(i)^2 = \varphi(i^2) = \varphi(-1) = -1.
	\]
	Следовательно, \( \varphi(i) = i \) или \( \varphi(i) = -i \). В первом случае \( \varphi(a + bi) = a + bi \) (тождественный автоморфизм). Во втором случае:
	\[
	\varphi(a + bi) = \varphi(a) + \varphi(b)\varphi(i) = a + b(-i) = a - bi = \overline{a + bi}.
	\]
	\hfill \(\square\)
	
	\myremark
	Существуют и другие (непрерывные) автоморфизмы \( \mathbb{C} \), но их существование зависит от аксиомы выбора, и они не сохраняют порядок на \( \mathbb{R} \) и не являются измеримыми. Группа \( \operatorname{Aut}(\mathbb{C}) \) бесконечна и очень сложно устроена.
	
	\subsection{Свойства автоморфизмов}
	
	\begin{itemize}
		\item Любой автоморфизм поля переводит простые элементы (если они определены) в простые.
		\item Автоморфизмы поля оставляют на месте его простое подполе (наименьшее подполе, содержащее 1). Для \( \mathbb{Q} \) это всё поле, для \( \mathbb{C} \) простое подполе — \( \mathbb{Q} \), и любой автоморфизм \( \mathbb{C} \) действует тождественно на \( \mathbb{Q} \).
		\item Автоморфизмы полей играют ключевую роль в теории Галуа.
	\end{itemize}
	
	
	
	% ==================== РАЗДЕЛ 35: АЛГЕБРАИЧЕСКАЯ ЗАМКНУТОСТЬ \mathbb{C} ====================
	\newpage
	\section{Алгебраическая замкнутость поля комплексных чисел (без доказательства). Каноническое разложение многочленов над полем комплексных чисел (35 вопрос)}
	
	\subsection{Алгебраическая замкнутость поля комплексных чисел}
	
	Напомним определение алгебраически замкнутого поля.
	
	\medskip
	\noindent \textbf{Определение.} Поле \( K \) называется \textit{алгебраически замкнутым}, если любой многочлен \( f(x) \in K[x] \) степени \( \geq 1 \) имеет хотя бы один корень в \( K \).
	
	\medskip
	Поле комплексных чисел \( \mathbb{C} \) обладает этим фундаментальным свойством.
	
	\begin{tcolorbox}[colback=red!5, colframe=red!50!black, title=Теорема (Основная теорема алгебры)]
		Поле комплексных чисел \( \mathbb{C} \) алгебраически замкнуто. Иными словами, любой многочлен с комплексными коэффициентами, отличный от постоянной, имеет по крайней мере один комплексный корень.
	\end{tcolorbox}
	
	\myremark
	Доказательство этой теоремы выходит за рамки данного курса. Оно может быть проведено различными методами: с использованием теории функций комплексного переменного (теорема Лиувилля), топологическими методами (теория степени отображения) или алгебраическими (теорема Гаусса). Мы принимаем этот факт без доказательства.
	
	\subsection{Следствия алгебраической замкнутости}
	
	Из алгебраической замкнутости \( \mathbb{C} \) вытекает важнейшее следствие: любой многочлен над \( \mathbb{C} \) раскладывается в произведение линейных множителей.
	
	\begin{tcolorbox}[colback=blue!5, colframe=blue!40!black, title=Теорема (Разложение на линейные множители)]
		Любой многочлен \( f(x) \in \mathbb{C}[x] \) степени \( n \geq 1 \) со старшим коэффициентом \( a_0 \neq 0 \) может быть представлен в виде:
		\[
		f(x) = a_0 (x - c_1)(x - c_2) \dots (x - c_n),
		\]
		где \( c_1, c_2, \dots, c_n \in \mathbb{C} \) — корни многочлена \( f(x) \) (не обязательно различные). Такое представление единственно с точностью до перестановки сомножителей.
	\end{tcolorbox}
	
	\medskip
	\textit{Доказательство.}
	Докажем оба утверждения (существование и единственность) методом математической индукции по степени многочлена \( n \).
	
	\textit{База индукции:} \( n = 1 \). Многочлен первой степени \( f(x) = a_0 x + a_1 \) (\( a_0 \neq 0 \)) имеет единственный корень \( c = -\frac{a_1}{a_0} \). Очевидно,
	\[
	f(x) = a_0 (x - c).
	\]
	
	Единственность: если \( a_0 (x - c) = a_0 (x - c') \), то \( a_0 (c' - c) = 0 \), а так как \( a_0 \neq 0 \) и \( \mathbb{C} \) — поле (нет делителей нуля), то \( c' = c \).
	
	\textit{Индукционный переход:} Предположим, что утверждение верно для всех многочленов степени \( n-1 \). Пусть \( f(x) \) — многочлен степени \( n \geq 2 \). В силу алгебраической замкнутости \( \mathbb{C} \), \( f(x) \) имеет хотя бы один корень \( c_1 \in \mathbb{C} \). По теореме Безу (следствие из неё), \( f(x) \) делится на \( (x - c_1) \):
	\[
	f(x) = (x - c_1) f_1(x),
	\]
	где \( f_1(x) \in \mathbb{C}[x] \) — многочлен степени \( n-1 \). Запишем его со старшим коэффициентом \( a_0 \):
	\[
	f_1(x) = a_0 x^{n-1} + b_1 x^{n-2} + \dots + b_{n-1}.
	\]
	
	По предположению индукции, \( f_1(x) \) раскладывается в произведение линейных множителей:
	\[
	f_1(x) = a_0 (x - c_2)(x - c_3) \dots (x - c_n).
	\]
	Подставляя это в выражение для \( f(x) \), получаем:
	\[
	f(x) = a_0 (x - c_1)(x - c_2) \dots (x - c_n).
	\]
	
	\textit{Единственность разложения.}
	Пусть \( n = 1 \) — база уже разобрана. Для \( n > 1 \) предположим, что существуют два разложения:
	\[
	f(x) = a_0 (x - c_1)(x - c_2) \dots (x - c_n) = a_0 (x - c_1')(x - c_2') \dots (x - c_n').
	\]
	
	Так как \( f(c_1) = 0 \), то
	\[
	0 = a_0 (c_1 - c_1')(c_1 - c_2') \dots (c_1 - c_n').
	\]
	Поскольку \( a_0 \neq 0 \) и в поле нет делителей нуля, один из множителей должен быть равен нулю. Следовательно, \( c_1 \) совпадает с некоторым \( c_i' \). Перенумеруем множители во втором разложении так, чтобы \( c_1 = c_1' \).
	
	Тогда имеем:
	\[
	a_0 (x - c_1)(x - c_2) \dots (x - c_n) = a_0 (x - c_1)(x - c_2') \dots (x - c_n').
	\]
	Сокращая на \( (x - c_1) \) (что допустимо в кольце многочленов над полем), получаем:
	\[
	a_0 (x - c_2) \dots (x - c_n) = a_0 (x - c_2') \dots (x - c_n').
	\]
	Левая и правая части — многочлены степени \( n-1 \). По предположению индукции, их разложения на линейные множители совпадают с точностью до перестановки. Таким образом, множества \( \{c_2, \dots, c_n\} \) и \( \{c_2', \dots, c_n'\} \) совпадают. \hfill \(\square\)
	
	\subsection{Кратность корня}
	
	В разложении многочлена на линейные множители некоторые корни могут повторяться.
	
	\medskip
	\noindent \textbf{Определение.} Пусть \( c \in \mathbb{C} \) — корень многочлена \( f(x) \). Натуральное число \( m \) называется \textit{кратностью} корня \( c \), если
	\[
	f(x) = (x - c)^m g(x),
	\]
	где \( g(x) \in \mathbb{C}[x] \) и \( g(c) \neq 0 \). Корень кратности 1 называется \textit{простым}, корень кратности \( m \geq 2 \) — \textit{кратным}.
	
	\subsection{Каноническое разложение над \( \mathbb{C} \)}
	
	Сгруппировав одинаковые линейные множители, получим \textit{каноническое разложение} многочлена над полем комплексных чисел.
	
	\begin{tcolorbox}[colback=green!5, colframe=green!50!black, title=Каноническое разложение]
		Любой ненулевой многочлен \( f(x) \in \mathbb{C}[x] \) степени \( n \geq 1 \) единственным образом (с точностью до порядка сомножителей) представляется в виде:
		\[
		f(x) = a_0 (x - c_1)^{m_1} (x - c_2)^{m_2} \dots (x - c_k)^{m_k},
		\]
		где:
		\begin{itemize}
			\item \( a_0 \in \mathbb{C}^* \) — старший коэффициент многочлена \( f \);
			\item \( c_1, c_2, \dots, c_k \in \mathbb{C} \) — \textit{различные} комплексные числа (корни многочлена);
			\item \( m_1, m_2, \dots, m_k \in \mathbb{N} \) — кратности соответствующих корней;
			\item \( m_1 + m_2 + \dots + m_k = n = \deg f \) (сумма кратностей равна степени многочлена).
		\end{itemize}
	\end{tcolorbox}
	
	\myremark
	Каноническое разложение полностью описывает многочлен с точностью до умножения на ненулевую константу: набор корней с их кратностями и старший коэффициент определяют многочлен однозначно.
	
	\subsection{Пример}
	
	Рассмотрим многочлен
	\[
	f(x) = 2x^5 - 6x^4 + 6x^3 - 2x^2 = 2x^2 (x^3 - 3x^2 + 3x - 1).
	\]
	Заметим, что \( x^3 - 3x^2 + 3x - 1 = (x - 1)^3 \). Таким образом,
	\[
	f(x) = 2x^2 (x - 1)^3.
	\]
	Здесь:
	\begin{itemize}
		\item старший коэффициент \( a_0 = 2 \);
		\item корень \( c_1 = 0 \) кратности \( m_1 = 2 \);
		\item корень \( c_2 = 1 \) кратности \( m_2 = 3 \);
		\item \( \deg f = 2 + 3 = 5 \).
	\end{itemize}
	
	\subsection{Связь с теоремой Безу и производными}
	
	Кратность корня можно определить с помощью производных: число \( c \) является корнем кратности \( m \) тогда и только тогда, когда
	\[
	f(c) = f'(c) = \dots = f^{(m-1)}(c) = 0, \quad f^{(m)}(c) \neq 0.
	\]
	
	\subsection{Заключение}
	
	Алгебраическая замкнутость поля комплексных чисел является причиной того, что именно в комплексной области многочлены ведут себя наиболее просто: они всегда имеют ровно \( n \) корней с учётом кратности и полностью определяются своим старшим коэффициентом и множеством корней. Это свойство широко используется в алгебре, математическом анализе и их приложениях.
	
	
	
	
	
	% ==================== РАЗДЕЛ 36: КАНОНИЧЕСКОЕ РАЗЛОЖЕНИЕ НАД \mathbb{R} ====================
	\newpage
	\section{Каноническое разложение многочленов над полем вещественных чисел (36 вопрос)}
	
	\subsection{Неприводимые многочлены над \(\mathbb{R}\)}
	
	Над полем действительных чисел \(\mathbb{R}\) ситуация с неприводимостью многочленов сложнее, чем над \(\mathbb{C}\), но всё ещё достаточно проста.
	
	\begin{tcolorbox}[colback=blue!5, colframe=blue!40!black, title=Предложение (Неприводимые многочлены над \(\mathbb{R}\))]
		Над полем вещественных чисел \(\mathbb{R}\) неприводимыми являются в точности:
		\begin{enumerate}
			\item многочлены первой степени: \(ax + b\), \(a \neq 0\);
			\item многочлены второй степени \(ax^2 + bx + c\), не имеющие вещественных корней, т.е. с отрицательным дискриминантом \(D = b^2 - 4ac < 0\).
		\end{enumerate}
		Многочлены степени выше второй всегда приводимы над \(\mathbb{R}\).
	\end{tcolorbox}
	
	\medskip
	\textit{Доказательство.}
	Пусть \(f(x) = a_0x^n + a_1x^{n-1} + \dots + a_n \in \mathbb{R}[x]\) и \(n \geq 1\).
	
	Если \(n = 1\), то многочлен первой степени, очевидно, неприводим (любое его разложение содержало бы множитель нулевой степени, т.е. константу).
	
	Рассмотрим случай \(n \geq 2\). Если \(f\) имеет вещественный корень \(c \in \mathbb{R}\), то по теореме Безу \(f(x)\) делится на \((x - c)\), следовательно, приводим.
	
	Пусть \(f\) не имеет вещественных корней. Так как \(\mathbb{C}\) алгебраически замкнуто, \(f\) имеет комплексный корень. Поскольку коэффициенты \(f\) вещественны, комплексные корни входят сопряжёнными парами. Пусть \(z_0 = a + bi\) (\(b \neq 0\)) — комплексный корень \(f\). Тогда \(\overline{z_0} = a - bi\) также является корнем \(f\).
	
	Рассмотрим квадратичный многочлен с вещественными коэффициентами:
	\[
	g(x) = (x - z_0)(x - \overline{z_0}) = (x - a - bi)(x - a + bi) = (x - a)^2 + b^2 = x^2 - 2ax + (a^2 + b^2).
	\]
	Его дискриминант \(D = (-2a)^2 - 4(a^2 + b^2) = 4a^2 - 4a^2 - 4b^2 = -4b^2 < 0\), поэтому \(g(x)\) неприводим над \(\mathbb{R}\).
	
	Так как \(z_0\) и \(\overline{z_0}\) являются корнями \(f\), то в кольце \(\mathbb{C}[x]\) многочлен \(f\) делится на \((x - z_0)(x - \overline{z_0}) = g(x)\). Но поскольку все коэффициенты \(f\) и \(g\) вещественны, алгоритм деления с остатком даёт частное и остаток также с вещественными коэффициентами. Следовательно, \(f(x) = g(x) h(x)\) для некоторого \(h(x) \in \mathbb{R}[x]\).
	
	Если степень \(f\) равна 2, то \(h\) — константа, и \(f\) ассоциирован с \(g\), т.е. неприводим. Если степень \(f\) больше 2, то мы получили нетривиальное разложение, значит, \(f\) приводим. \hfill \(\square\)
	
	\subsection{Каноническое разложение над \(\mathbb{R}\)}
	
	Из описания неприводимых многочленов следует структура канонического разложения для произвольного многочлена с вещественными коэффициентами.
	
	\begin{tcolorbox}[colback=green!5, colframe=green!50!black, title=Теорема (Каноническое разложение над \(\mathbb{R}\))]
		Любой ненулевой многочлен \(f(x) \in \mathbb{R}[x]\) степени \(n \geq 1\) единственным образом (с точностью до порядка сомножителей) представляется в виде:
		\[
		f(x) = a_0 (x - x_1)^{m_1} (x - x_2)^{m_2} \dots (x - x_k)^{m_k} \cdot (x^2 + p_1 x + q_1)^{n_1} (x^2 + p_2 x + q_2)^{n_2} \dots (x^2 + p_s x + q_s)^{n_s},
		\]
		где:
		\begin{itemize}
			\item \(a_0 \in \mathbb{R}^*\) — старший коэффициент многочлена \(f\);
			\item \(x_1, x_2, \dots, x_k \in \mathbb{R}\) — все различные \textit{вещественные} корни многочлена \(f\);
			\item \(m_1, m_2, \dots, m_k \in \mathbb{N}\) — их кратности;
			\item \(x^2 + p_j x + q_j\) (\(j = 1, \dots, s\)) — различные неприводимые квадратичные множители (т.е. с отрицательными дискриминантами \(p_j^2 - 4q_j < 0\));
			\item \(n_1, n_2, \dots, n_s \in \mathbb{N}\) — кратности этих квадратичных множителей.
		\end{itemize}
		При этом
		\[
		\sum_{i=1}^k m_i + 2\sum_{j=1}^s n_j = n = \deg f.
		\]
	\end{tcolorbox}
	
	\medskip
	\textit{Доказательство (идея).}
	Существование разложения следует из теоремы о факториальности кольца многочленов над полем (вопрос 33) и классификации неприводимых многочленов над \(\mathbb{R}\), данной в предыдущем предложении.
	
	Единственность (с точностью до порядка следования множителей и умножения на константы) вытекает из общей теоремы о факториальности \(K[x]\) для поля \(K\). \hfill \(\square\)
	
	\subsection{Следствие о сопряжённых корнях}
	
	Из канонического разложения вытекает важное свойство многочленов с вещественными коэффициентами.
	
	\begin{tcolorbox}[colback=yellow!5, colframe=orange!50!black, title=Следствие (Сопряжённые корни)]
		Если многочлен \(f(x) \in \mathbb{R}[x]\) имеет комплексный корень \(z_0 = a + bi\) (\(b \neq 0\)), то он имеет и сопряжённый корень \(\overline{z_0} = a - bi\), причём кратности этих корней совпадают.
	\end{tcolorbox}
	
	\medskip
	\textit{Доказательство.}
	Рассмотрим каноническое разложение \(f(x)\) над \(\mathbb{R}\). В этом разложении присутствуют только линейные множители, соответствующие вещественным корням, и неприводимые квадратичные множители. Каждый квадратичный множитель \(x^2 + px + q\) имеет два комплексно-сопряжённых корня. Если \(f\) имеет корень \(a + bi\), то этот корень является корнем некоторого квадратичного множителя \((x^2 + px + q)^{n_j}\) в разложении. Но тогда этот же квадратичный множитель имеет корнем и \(a - bi\). Следовательно, \(\overline{z_0}\) также является корнем \(f\), и его кратность равна \(n_j\), т.е. совпадает с кратностью \(z_0\). \hfill \(\square\)
	
	\subsection{Примеры}
	
	\textbf{Пример 1.} Разложить \(f(x) = x^5 - x^4 + 2x^3 - 2x^2 + x - 1\) над \(\mathbb{R}\).
	
	Сначала найдём вещественные корни. Проверим \(x = 1\):
	\[
	f(1) = 1 - 1 + 2 - 2 + 1 - 1 = 0,
	\]
	значит, \(x = 1\) — корень. Разделим на \((x - 1)\):
	\[
	f(x) = (x - 1)(x^4 + 2x^2 + 1).
	\]
	Заметим, что \(x^4 + 2x^2 + 1 = (x^2 + 1)^2\). Таким образом,
	\[
	f(x) = (x - 1)(x^2 + 1)^2.
	\]
	Здесь:
	\begin{itemize}
		\item вещественный корень \(x_1 = 1\) кратности \(m_1 = 1\);
		\item неприводимый квадратичный множитель \(x^2 + 1\) (дискриминант \(-4 < 0\)) кратности \(n_1 = 2\);
		\item степень многочлена: \(1 + 2 \cdot 2 = 5\).
	\end{itemize}
	
	\textbf{Пример 2.} Разложить \(f(x) = x^4 + 4\) над \(\mathbb{R}\).
	
	Это классический пример. Добавим и вычтем \(4x^2\):
	\[
	x^4 + 4 = (x^4 + 4x^2 + 4) - 4x^2 = (x^2 + 2)^2 - (2x)^2 = (x^2 + 2x + 2)(x^2 - 2x + 2).
	\]
	Оба квадратичных множителя имеют отрицательные дискриминанты (\(4 - 8 = -4\)), поэтому они неприводимы над \(\mathbb{R}\). Каноническое разложение:
	\[
	f(x) = (x^2 + 2x + 2)(x^2 - 2x + 2).
	\]
	
	\subsection{Сравнение с разложением над \(\mathbb{C}\)}
	
	Над \(\mathbb{C}\) разложение выглядит проще (только линейные множители), но требует выхода в комплексную область. Например, для \(f(x) = x^2 + 1\):
	\[
	\text{над } \mathbb{R}: f(x) = x^2 + 1 \quad (\text{неприводимый квадрат}),
	\]
	\[
	\text{над } \mathbb{C}: f(x) = (x - i)(x + i).
	\]
	
	Каноническое разложение над \(\mathbb{R}\) группирует комплексно-сопряжённые корни в вещественные квадратичные множители, что позволяет оставаться в рамках поля действительных чисел.
	
	
	
	
	
	
	% ==================== РАЗДЕЛ 37: ПРОИЗВОДНАЯ МНОГОЧЛЕНА ====================
	\newpage
	\section{Производная многочлена. Свойства (37 вопрос)}
	
	\subsection{Определение производной многочлена}
	
	Пусть \( R \) — коммутативное кольцо (в частности, может быть полем \(\mathbb{R}\) или \(\mathbb{C}\)). Для многочлена с коэффициентами из \( R \) можно формально определить понятие производной, не прибегая к пределам.
	
	\begin{tcolorbox}[colback=blue!5, colframe=blue!40!black, title=Определение]
		Пусть \( f(x) \in R[x] \) — многочлен:
		\[
		f(x) = a_n x^n + a_{n-1} x^{n-1} + \dots + a_1 x + a_0, \quad a_i \in R.
		\]
		\textit{Производной} многочлена \( f \) называется многочлен \( f'(x) \in R[x] \), определяемый формулой:
		\[
		f'(x) = n a_n x^{n-1} + (n-1) a_{n-1} x^{n-2} + \dots + 2a_2 x + a_1.
		\]
		Для константы (многочлена нулевой степени) производная равна нулю.
	\end{tcolorbox}
	
	\myremark
	В отличие от математического анализа, здесь производная вводится чисто алгебраически, без понятия предела. Это формальная производная.
	
	\subsection{Элементарные свойства}
	
	Непосредственно из определения следуют простейшие свойства.
	
	\begin{tcolorbox}[colback=green!5, colframe=green!50!black, title=Свойства 1–4]
		\begin{enumerate}
			\item Производная константы равна нулю: \((c)' = 0\) для любого \(c \in R\).
			\item Производная линейного двучлена: \((x - c)' = 1\).
			\item Производная суммы равна сумме производных: \((f_1 + f_2)' = f_1' + f_2'\).
			\item Производная произведения константы на многочлен: \((c f_1)' = c f_1'\).
		\end{enumerate}
	\end{tcolorbox}
	
	\subsection{Производная произведения}
	
	Важнейшим свойством является правило Лейбница (производная произведения).
	
	\begin{tcolorbox}[colback=blue!5, colframe=blue!40!black, title=Свойство 5 (Правило Лейбница)]
		Для любых многочленов \( f_1, f_2 \in R[x] \) выполняется:
		\[
		(f_1 f_2)' = f_1' f_2 + f_1 f_2'.
		\]
	\end{tcolorbox}
	
	\medskip
	\textit{Доказательство.}
	Доказательство проведём в несколько этапов, последовательно усложняя вид множителей.
	
	\textbf{Этап a)} Пусть \( f_1 = a x^m \), \( f_2 = b x^k \) — одночлены. Тогда
	\[
	f_1 f_2 = a b x^{m+k}.
	\]
	По определению производной:
	\[
	(f_1 f_2)' = (m+k) a b x^{m+k-1}.
	\]
	С другой стороны,
	\[
	f_1' f_2 = (m a x^{m-1}) \cdot (b x^k) = m a b x^{m+k-1},
	\]
	\[
	f_1 f_2' = (a x^m) \cdot (k b x^{k-1}) = k a b x^{m+k-1}.
	\]
	Складывая, получаем \( (m+k) a b x^{m+k-1} \), что совпадает с \((f_1 f_2)'\). Для одночленов правило доказано.
	
	\textbf{Этап б)} Пусть теперь \( f_1 \) — произвольный многочлен, а \( f_2 = b x^k \) — одночлен. Запишем \( f_1 \) как сумму одночленов:
	\[
	f_1 = a_0 x^m + a_1 x^{m-1} + \dots + a_m.
	\]
	Тогда
	\[
	f_1 f_2 = (a_0 x^m)(b x^k) + (a_1 x^{m-1})(b x^k) + \dots + a_m (b x^k).
	\]
	Используя свойства 3, 4 и результат этапа a) для каждого слагаемого, получаем:
	\[
	(f_1 f_2)' = \big[(a_0 x^m)' (b x^k) + (a_0 x^m) (b x^k)'\big] + \big[(a_1 x^{m-1})' (b x^k) + (a_1 x^{m-1}) (b x^k)'\big] + \dots + \big[(a_m)' (b x^k) + a_m (b x^k)'\big].
	\]
	Группируя слагаемые с \((b x^k)'\) и без, получаем:
	\[
	(f_1 f_2)' = \big[(a_0 x^m)' + (a_1 x^{m-1})' + \dots + a_m'\big] (b x^k) + \big[a_0 x^m + a_1 x^{m-1} + \dots + a_m\big] (b x^k)' = f_1' f_2 + f_1 f_2'.
	\]
	
	\textbf{Этап в)} Общий случай: \( f_1 \) и \( f_2 \) — произвольные многочлены. Представим \( f_2 \) как сумму одночленов:
	\[
	f_2 = b_0 x^k + b_1 x^{k-1} + \dots + b_k.
	\]
	Тогда
	\[
	f_1 f_2 = f_1 (b_0 x^k) + f_1 (b_1 x^{k-1}) + \dots + f_1 b_k.
	\]
	Применяя свойства 3, 4 и результат этапа б) к каждому слагаемому, получаем:
	\[
	(f_1 f_2)' = \big[f_1' (b_0 x^k) + f_1 (b_0 x^k)'\big] + \big[f_1' (b_1 x^{k-1}) + f_1 (b_1 x^{k-1})'\big] + \dots + \big[f_1' b_k + f_1 b_k'\big].
	\]
	Группируя слагаемые с \( f_1' \) и с \( f_1 \), находим:
	\[
	(f_1 f_2)' = f_1' (b_0 x^k + b_1 x^{k-1} + \dots + b_k) + f_1 \big[(b_0 x^k)' + (b_1 x^{k-1})' + \dots + b_k'\big] = f_1' f_2 + f_1 f_2'.
	\]
	\hfill \(\square\)
	
	\subsection{Производная произведения нескольких множителей}
	
	Из правила Лейбница по индукции получается формула для производной произведения любого конечного числа многочленов.
	
	\begin{tcolorbox}[colback=blue!5, colframe=blue!40!black, title=Свойство 6]
		Для любых многочленов \( f_1, f_2, \dots, f_k \in R[x] \) выполняется:
		\[
		(f_1 f_2 \dots f_k)' = f_1' f_2 \dots f_k + f_1 f_2' f_3 \dots f_k + \dots + f_1 f_2 \dots f_{k-1} f_k'.
		\]
	\end{tcolorbox}
	
	\medskip
	\textit{Доказательство.}
	Применим индукцию по \( k \). База \( k = 2 \) — это свойство 5. Пусть утверждение верно для \( k-1 \) множителей. Тогда, рассматривая \( f_1 f_2 \dots f_k = (f_1 f_2 \dots f_{k-1}) f_k \) и применяя свойство 5, а затем предположение индукции к первому сомножителю, получаем требуемую формулу. \hfill \(\square\)
	
	\subsection{Производная степени многочлена}
	
	Как частный случай свойства 6, когда все множители равны, получаем важную формулу.
	
	\begin{tcolorbox}[colback=blue!5, colframe=blue!40!black, title=Свойство 7]
		Для любого многочлена \( f \in R[x] \) и натурального \( k \):
		\[
		(f^k)' = k f^{k-1} f'.
		\]
	\end{tcolorbox}
	
	\medskip
	\textit{Доказательство.}
	В свойстве 6 положим \( f_1 = f_2 = \dots = f_k = f \). Тогда каждое слагаемое имеет вид \( f' \cdot f^{k-1} \), и таких слагаемых ровно \( k \). \hfill \(\square\)
	
	\subsection{Производная линейного двучлена в степени}
	
	Частный случай свойства 7 при \( f(x) = x - c \).
	
	\begin{tcolorbox}[colback=blue!5, colframe=blue!40!black, title=Свойство 8]
		Для любого \( c \in R \):
		\[
		\big((x - c)^k\big)' = k (x - c)^{k-1}.
		\]
	\end{tcolorbox}
	
	\medskip
	\textit{Доказательство.}
	Так как \( (x - c)' = 1 \), то по свойству 7 получаем \( ((x - c)^k)' = k (x - c)^{k-1} \cdot 1 = k (x - c)^{k-1} \). \hfill \(\square\)
	
	\subsection{Формула Тейлора для многочленов}
	
	Пусть \( R \) — область целостности характеристики 0 (например, \( \mathbb{Q}, \mathbb{R}, \mathbb{C} \)). Тогда для любого многочлена \( f \in R[x] \) и любого элемента \( d \in R \) справедливо разложение по степеням \( (x - d) \).
	
	\begin{tcolorbox}[colback=yellow!5, colframe=orange!50!black, title=Теорема (Формула Тейлора)]
		Пусть \( f(x) \in R[x] \) — многочлен степени \( n \). Тогда для любого \( d \in R \) существует единственное представление:
		\[
		f(x) = \sum_{k=0}^n \beta_k (x - d)^k, \quad \beta_k \in R.
		\]
		Коэффициенты \( \beta_k \) выражаются через значения производных многочлена в точке \( d \):
		\[
		\beta_k = \frac{f^{(k)}(d)}{k!},
		\]
		где \( f^{(k)} \) обозначает \( k \)-ю производную (\( f^{(0)} = f \)), а \( k! \) — факториал, понимаемый как элемент кольца \( R \) (кратное единицы).
	\end{tcolorbox}
	
	\medskip
	\textit{Доказательство.}
	Существование разложения следует из того, что систему многочленов \( \{1, (x-d), (x-d)^2, \dots, (x-d)^n\} \) можно получить из стандартного базиса \( \{1, x, x^2, \dots, x^n\} \) обратимым линейным преобразованием (треугольная матрица). Следовательно, любой многочлен единственным образом выражается через этот базис.
	
	Найдём коэффициенты \( \beta_k \). Продифференцируем разложение \( k \) раз и подставим \( x = d \). Заметим, что при \( i < k \) производная порядка \( k \) от \( (x-d)^i \) равна нулю. При \( i > k \) производная порядка \( k \) от \( (x-d)^i \) содержит множитель \( (x-d)^{i-k} \), который обращается в ноль при \( x = d \). Единственное слагаемое, дающее вклад, — это \( i = k \):
	\[
	\frac{d^k}{dx^k}\big[(x-d)^k\big] = k!.
	\]
	Таким образом,
	\[
	f^{(k)}(d) = \beta_k \cdot k!.
	\]
	Отсюда \( \beta_k = \frac{f^{(k)}(d)}{k!} \). \hfill \(\square\)
	
	\myremark
	Формула Тейлора для многочленов является точной (не асимптотической) и часто используется для исследования кратных корней: \( d \) является корнем кратности \( m \) тогда и только тогда, когда
	\[
	f(d) = f'(d) = \dots = f^{(m-1)}(d) = 0, \quad f^{(m)}(d) \neq 0.
	\]
	
	
	
	
	
	
	
	
	
	% ==================== РАЗДЕЛ 38: КРИТЕРИЙ КРАТНОСТИ КОРНЯ ====================
	\newpage
	\section{Критерий кратности корня (38 вопрос)}
	
	\subsection{Определение кратного корня}
	
	Пусть \( R \) — область целостности (в частности, может быть полем \( \mathbb{R} \) или \( \mathbb{C} \)). Рассмотрим многочлен \( f(t) \in R[t] \).
	
	\begin{tcolorbox}[colback=blue!5, colframe=blue!40!black, title=Определение]
		Пусть \( d \in R \) является корнем многочлена \( f \in R[t] \). Натуральное число \( k \in \mathbb{N} \) называется \textit{кратностью} корня \( d \), если
		\[
		f(t) = (t - d)^k g(t),
		\]
		где \( g(t) \in R[t] \) и \( g(d) \neq 0 \). 
		
		Если \( k = 1 \), корень называется \textit{простым}. Если \( k \geq 2 \), корень называется \textit{кратным} (или \textit{многократным}).
	\end{tcolorbox}
	
	\myremark
	Условие \( g(d) \neq 0 \) гарантирует, что \( k \) — максимальная степень \( (t - d) \), на которую делится \( f(t) \). Другими словами, \( f(t) \) делится на \( (t - d)^k \), но не делится на \( (t - d)^{k+1} \).
	
	\subsection{Критерий кратности через производную}
	
	Связь между кратностью корня и его поведением относительно производной даёт эффективный способ определения кратности.
	
	\begin{tcolorbox}[colback=green!5, colframe=green!50!black, title=Теорема (Критерий кратного корня)]
		Пусть \( R \) — область целостности характеристики 0 (например, \( \mathbb{Q}, \mathbb{R}, \mathbb{C} \)), и пусть \( d \in R \) — корень многочлена \( f \in R[t] \) кратности \( k \). Тогда:
		
		\begin{enumerate}
			\item Если \( k = 1 \) (простой корень), то \( d \) \textbf{не является} корнем производной \( f' \), т.е. \( f'(d) \neq 0 \).
			\item Если \( k \geq 2 \) (кратный корень), то \( d \) является корнем производной \( f' \). Более того, кратность \( d \) как корня \( f' \) равна в точности \( k-1 \) при условии, что характеристика кольца не делит \( k \) (в частности, это всегда выполнено для полей характеристики 0).
		\end{enumerate}
	\end{tcolorbox}
	
	\medskip
	\textit{Доказательство.}
	По определению кратности, существует представление
	\[
	f(t) = (t - d)^k g(t),
	\]
	где \( g(t) \in R[t] \) и \( g(d) \neq 0 \).
	
	Вычислим производную \( f'(t) \), используя правило Лейбница:
	\[
	f'(t) = k(t - d)^{k-1} g(t) + (t - d)^k g'(t).
	\]
	Вынесем общий множитель \( (t - d)^{k-1} \):
	\[
	f'(t) = (t - d)^{k-1} \big[ k g(t) + (t - d) g'(t) \big].
	\]
	Обозначим \( h(t) = k g(t) + (t - d) g'(t) \). Тогда
	\[
	f'(t) = (t - d)^{k-1} h(t).
	\]
	
	Рассмотрим значение \( h(t) \) в точке \( t = d \):
	\[
	h(d) = k g(d) + (d - d) g'(d) = k g(d).
	\]
	
	\textbf{Случай 1: \( k = 1 \).} Тогда
	\[
	f'(t) = (t - d)^{0} h(t) = h(t),
	\]
	и \( f'(d) = h(d) = 1 \cdot g(d) = g(d) \neq 0 \). Следовательно, \( d \) не является корнем \( f' \).
	
	\textbf{Случай 2: \( k \geq 2 \).} Тогда в выражении \( f'(t) = (t - d)^{k-1} h(t) \) множитель \( (t - d)^{k-1} \) гарантирует, что \( d \) — корень \( f' \) кратности не меньше \( k-1 \).
	
	Для определения точной кратности нужно выяснить, обращается ли \( h(d) \) в ноль. Имеем \( h(d) = k g(d) \). По условию \( g(d) \neq 0 \). Если характеристика кольца \( R \) не делит \( k \) (т.е. \( k \neq 0 \) в \( R \)), то \( k g(d) \neq 0 \), так как \( R \) — область целостности. Следовательно, \( h(d) \neq 0 \), и множитель \( (t - d) \) не входит в \( h(t) \). Таким образом, кратность \( d \) как корня \( f' \) равна в точности \( k-1 \). \hfill \(\square\)
	
	\subsection{Обобщение на старшие производные}
	
	Из доказательства легко получить общий критерий для определения кратности корня с помощью последовательных производных.
	
	\begin{tcolorbox}[colback=yellow!5, colframe=orange!50!black, title=Следствие (Критерий в терминах производных)]
		Пусть \( R \) — область целостности характеристики 0, \( f \in R[t] \), \( d \in R \). Тогда \( d \) является корнем кратности \( k \) тогда и только тогда, когда
		\[
		f(d) = f'(d) = f''(d) = \dots = f^{(k-1)}(d) = 0, \quad \text{но} \quad f^{(k)}(d) \neq 0.
		\]
	\end{tcolorbox}
	
	\medskip
	\textit{Доказательство.}
	Применим формулу Тейлора для многочлена \( f \) в окрестности точки \( d \):
	\[
	f(t) = \sum_{i=0}^n \frac{f^{(i)}(d)}{i!} (t - d)^i.
	\]
	Кратность корня \( d \) равна наименьшему индексу \( i \), для которого коэффициент при \( (t - d)^i \) ненулевой. Этот коэффициент равен \( \frac{f^{(i)}(d)}{i!} \). Условие \( f^{(i)}(d) = 0 \) для всех \( i < k \) и \( f^{(k)}(d) \neq 0 \) в точности означает, что \( k \) — искомая кратность. \hfill \(\square\)
	
	\subsection{Примеры}
	
	\textbf{Пример 1.} Рассмотрим \( f(t) = (t-1)^3 (t+2) \) над \( \mathbb{R} \).
	\begin{itemize}
		\item Корень \( d = 1 \) имеет кратность \( k = 3 \).
		\item Вычислим производную: \( f'(t) = 3(t-1)^2(t+2) + (t-1)^3 \cdot 1 = (t-1)^2[3(t+2) + (t-1)] = (t-1)^2(4t+5) \).
		\item При \( t = 1 \) получаем \( f'(1) = 0 \), \( f''(1) = 0 \), \( f'''(1) \neq 0 \), что соответствует \( k = 3 \).
	\end{itemize}
	
	\textbf{Пример 2.} Рассмотрим \( f(t) = t^3 - 3t + 2 \) над \( \mathbb{R} \).
	\begin{itemize}
		\item Проверим \( t = 1 \): \( f(1) = 1 - 3 + 2 = 0 \), значит, 1 — корень.
		\item \( f'(t) = 3t^2 - 3 \), \( f'(1) = 3 - 3 = 0 \).
		\item \( f''(t) = 6t \), \( f''(1) = 6 \neq 0 \).
		\item Следовательно, \( t = 1 \) — корень кратности 2 (двукратный корень).
		\item Действительно, \( t^3 - 3t + 2 = (t-1)^2(t+2) \).
	\end{itemize}
	
	\subsection{Замечание о характеристике}
	
	Условие на характеристику кольца существенно. Рассмотрим поле \( \mathbb{F}_p \) характеристики \( p \). Пусть \( f(t) = t^p - a \). Если \( a \) имеет корень в \( \mathbb{F}_p \), то может возникнуть ситуация, когда производная \( f'(t) = p t^{p-1} \equiv 0 \) тождественно, и критерий через производные не работает. В частности, в характеристике \( p \) корень может быть кратным, даже если формальная производная не обращается в ноль в точке (или обращается тождественно в ноль). Для полей характеристики 0 (как \( \mathbb{Q}, \mathbb{R}, \mathbb{C} \)) эти проблемы отсутствуют.
	
	
	
	
	
	
	% ==================== РАЗДЕЛ 39: ХАРАКТЕРИСТИКА ПОЛЯ. ОТДЕЛЕНИЕ КРАТНЫХ КОРНЕЙ ====================
	\newpage
	\section{Характеристика поля. Отделение кратных корней (39 вопрос)}
	
	\subsection{Характеристика кольца и поля}
	
	\begin{tcolorbox}[colback=blue!5, colframe=blue!40!black, title=Определение характеристики]
		Пусть \( A \) — коммутативное кольцо с единицей \( 1_A \). Рассмотрим множество
		\[
		M = \{ s \in \mathbb{Z} \mid s \cdot 1_A = 0 \},
		\]
		где \( s \cdot 1_A \) означает \( s \)-кратное суммы единицы: \( 1_A + 1_A + \dots + 1_A \) (\( s \) раз) при \( s > 0 \), и соответствующее определение для отрицательных \( s \) (с использованием противоположного элемента).
		
		\textit{Характеристикой} кольца \( A \) (обозначение \( \operatorname{char} A \)) называется неотрицательное целое число, равное:
		\begin{itemize}
			\item \( 0 \), если \( M = \{0\} \) (т.е. \( s \cdot 1_A = 0 \) только при \( s = 0 \));
			\item наименьшему положительному числу \( m \in M \), если такое существует.
		\end{itemize}
	\end{tcolorbox}
	
	\medskip
	\noindent \textbf{Примеры.}
	\begin{enumerate}
		\item Для кольца целых чисел \( \mathbb{Z} \): \( s \cdot 1 = s \). Равенство \( s = 0 \) возможно только при \( s = 0 \). Следовательно, \( \operatorname{char} \mathbb{Z} = 0 \).
		\item Для кольца вычетов \( \mathbb{Z}_m \) (\( m \geq 1 \)): \( s \cdot 1 = \overline{s} \). Это равно нулю тогда и только тогда, когда \( m \mid s \). Наименьшее положительное \( s \) с этим свойством — само \( m \). Следовательно, \( \operatorname{char} \mathbb{Z}_m = m \).
		\item Для полей \( \mathbb{Q}, \mathbb{R}, \mathbb{C} \) характеристика равна 0.
		\item Для конечного поля \( \mathbb{F}_p \) (\( p \) — простое) характеристика равна \( p \).
	\end{enumerate}
	
	\subsection{Характеристика области целостности}
	
	\begin{tcolorbox}[colback=green!5, colframe=green!50!black, title=Предложение]
		Характеристика области целостности равна либо нулю, либо простому числу.
	\end{tcolorbox}
	
	\medskip
	\textit{Доказательство.}
	Пусть \( R \) — область целостности, и пусть \( \operatorname{char} R = m \neq 0 \). Предположим, что \( m \) — составное число, т.е. \( m = ab \), где \( 1 < a, b < m \). Тогда в \( R \):
	\[
	(a \cdot 1_R)(b \cdot 1_R) = (ab) \cdot 1_R = m \cdot 1_R = 0.
	\]
	Так как \( R \) — область целостности, произведение равно нулю только если один из множителей равен нулю. Следовательно, либо \( a \cdot 1_R = 0 \), либо \( b \cdot 1_R = 0 \). Но это противоречит минимальности \( m \) (так как \( a, b < m \)). Значит, \( m \) не может быть составным, т.е. \( m \) — простое число. \hfill \(\square\)
	
	\myremark
	В частности, характеристика поля равна либо 0, либо простому числу. Поля характеристики 0 содержат подполе, изоморфное \( \mathbb{Q} \); поля простой характеристики \( p \) содержат подполе, изоморфное \( \mathbb{F}_p \).
	
	\subsection{Отделение кратных корней}
	
	Пусть \( K \) — поле характеристики 0 (или просто поле, характеристика которого не влияет на рассматриваемые свойства). Рассмотрим многочлен \( f \in K[x] \). Задача отделения кратных корней состоит в нахождении многочлена, имеющего те же корни, что и \( f \), но все простые (т.е. без кратных корней).
	
	\begin{tcolorbox}[colback=yellow!5, colframe=orange!50!black, title=Теорема (Отделение кратных корней)]
		Пусть \( f \in K[x] \) — ненулевой многочлен над полем характеристики 0. Тогда многочлен
		\[
		h(x) = \frac{f(x)}{\operatorname{НОД}(f(x), f'(x))}
		\]
		имеет те же корни, что и \( f \), но все кратности этих корней равны 1 (т.е. все корни простые).
	\end{tcolorbox}
	
	\medskip
	\textit{Доказательство.}
	Пусть \( d \in \overline{K} \) (алгебраическое замыкание) — корень \( f \) кратности \( k \geq 1 \). Тогда в разложении над алгебраическим замыканием:
	\[
	f(x) = (x - d)^k g(x), \quad g(d) \neq 0.
	\]
	Вычислим производную:
	\[
	f'(x) = k(x - d)^{k-1} g(x) + (x - d)^k g'(x) = (x - d)^{k-1} [k g(x) + (x - d) g'(x)].
	\]
	Обозначим \( h_1(x) = k g(x) + (x - d) g'(x) \). Тогда \( h_1(d) = k g(d) \neq 0 \) (так как \( k \neq 0 \) в поле характеристики 0 и \( g(d) \neq 0 \)). Следовательно, \( (x - d) \) не делит \( h_1(x) \).
	
	Таким образом, в разложении \( f'(x) \) множитель \( (x - d) \) входит ровно в степени \( k-1 \).
	
	Рассмотрим \( \operatorname{НОД}(f, f') \). Из приведённых разложений видно, что \( (x - d) \) входит в НОД в степени \( k-1 \) (так как в \( f \) степень \( k \), в \( f' \) степень \( k-1 \), и общих множителей, кроме степеней \( (x - d) \), нет, поскольку \( g(d) \neq 0 \) и \( h_1(d) \neq 0 \).
	
	Следовательно, в факторе
	\[
	\frac{f}{\operatorname{НОД}(f, f')}
	\]
	множитель \( (x - d) \) будет входить ровно в степени \( k - (k-1) = 1 \). Это верно для каждого корня \( d \). Таким образом, все корни полученного многочлена простые. \hfill \(\square\)
	
	\subsection{Явный вид НОД(\(f, f'\))}
	
	Из доказательства видно, что
	\[
	\operatorname{НОД}(f, f') = \prod_{i=1}^m (x - d_i)^{k_i - 1},
	\]
	где \( d_i \) — различные корни \( f \), а \( k_i \) — их кратности. Таким образом, НОД \( f \) и \( f' \) содержит каждый неприводимый множитель \( f \) в степени на единицу меньше его кратности.
	
	\subsection{Пример отделения кратных корней}
	
	Рассмотрим пример из предоставленного фрагмента.
	
	\begin{tcolorbox}[colback=red!5, colframe=red!50!black, title=Пример]
		Отделить кратные корни многочлена
		\[
		f(x) = x^5 - 10x^3 - 20x^2 - 15x - 4.
		\]
	\end{tcolorbox}
	
	\textit{Решение.}
	Вычислим производную:
	\[
	f'(x) = 5x^4 - 30x^2 - 40x - 15 = 5(x^4 - 6x^2 - 8x - 3).
	\]
	Константный множитель 5 не влияет на НОД (он обратим в \( \mathbb{Q} \)), поэтому можно работать с \( g(x) = x^4 - 6x^2 - 8x - 3 \).
	
	Найдём \( \operatorname{НОД}(f, g) \) алгоритмом Евклида.
	
	\textit{Шаг 1:} Делим \( f \) на \( g \):
	\[
	\frac{x^5 - 10x^3 - 20x^2 - 15x - 4}{x^4 - 6x^2 - 8x - 3} = x + \frac{-4x^3 - 12x^2 - 12x - 4}{x^4 - 6x^2 - 8x - 3}.
	\]
	Остаток: \( r_1(x) = -4x^3 - 12x^2 - 12x - 4 = -4(x^3 + 3x^2 + 3x + 1) \).
	
	\textit{Шаг 2:} Делим \( g \) на \( x^3 + 3x^2 + 3x + 1 \):
	\[
	\frac{x^4 - 6x^2 - 8x - 3}{x^3 + 3x^2 + 3x + 1} = x - 3 + \frac{0 \cdot x^2 + 0 \cdot x + 0}{x^3 + 3x^2 + 3x + 1}?
	\]
	Выполним деление аккуратно:
	\[
	(x^4 - 6x^2 - 8x - 3) \div (x^3 + 3x^2 + 3x + 1):
	\]
	\begin{itemize}
		\item Частное: \( x - 3 \).
		\item Умножаем: \( (x - 3)(x^3 + 3x^2 + 3x + 1) = x^4 + 3x^3 + 3x^2 + x - 3x^3 - 9x^2 - 9x - 3 = x^4 - 6x^2 - 8x - 3 \).
		\item Остаток: \( 0 \).
	\end{itemize}
	Таким образом, \( x^3 + 3x^2 + 3x + 1 \) делит \( g \) нацело.
	
	Следовательно,
	\[
	\operatorname{НОД}(f, f') = x^3 + 3x^2 + 3x + 1.
	\]
	
	\textit{Шаг 3:} Отделяем кратные корни:
	\[
	\frac{f}{\operatorname{НОД}(f, f')} = \frac{x^5 - 10x^3 - 20x^2 - 15x - 4}{x^3 + 3x^2 + 3x + 1}.
	\]
	Выполним деление:
	\[
	(x^5 - 10x^3 - 20x^2 - 15x - 4) \div (x^3 + 3x^2 + 3x + 1):
	\]
	\begin{itemize}
		\item Частное: \( x^2 - 3x - 4 \).
		\item Проверка: \( (x^2 - 3x - 4)(x^3 + 3x^2 + 3x + 1) = x^5 + 3x^4 + 3x^3 + x^2 - 3x^4 - 9x^3 - 9x^2 - 3x - 4x^3 - 12x^2 - 12x - 4 = x^5 - 10x^3 - 20x^2 - 15x - 4 \).
	\end{itemize}
	
	Итак,
	\[
	h(x) = x^2 - 3x - 4
	\]
	— многочлен, имеющий те же корни, что и \( f \), но все простые. Действительно, \( h(x) = (x - 4)(x + 1) \), его корни: \( 4 \) и \( -1 \) (оба простые).
	
	Проверим исходный многочлен: должно быть \( f(x) = (x-4)^2 (x+1)^3 \) или что-то подобное? Разложим:
	\[
	x^3 + 3x^2 + 3x + 1 = (x+1)^3,
	\]
	а \( h(x) = (x-4)(x+1) \). Тогда
	\[
	f(x) = (x+1)^3 \cdot (x-4)(x+1) = (x-4)(x+1)^4.
	\]
	Действительно, корень \( x = -1 \) имеет кратность 4, корень \( x = 4 \) — кратность 1. После отделения кратных корней получили \( (x-4)(x+1) \) — все корни простые.
	
	\subsection{Замечание о характеристике}
	
	В полях характеристики \( p > 0 \) ситуация сложнее: если \( f(x) = g(x^p) \), то производная \( f'(x) = p x^{p-1} g'(x^p) \) может быть тождественно равна нулю, и описанный метод неприменим. В таких случаях используют понятие \textit{сепарабельных} и \textit{чисто несепарабельных} расширений.
	
	
	
	
	
	% ==================== РАЗДЕЛ 40: ИНТЕРПОЛЯЦИОННЫЙ МНОГОЧЛЕН ЛАГРАНЖА ====================
	\newpage
	\section{Интерполяционный многочлен Лагранжа (40 вопрос)}
	
	\subsection{Постановка задачи интерполяции}
	
	Пусть \( K \) — поле (например, \( \mathbb{R}, \mathbb{C}, \mathbb{Q} \)). Рассмотрим кольцо многочленов \( K[t] \).
	
	\begin{tcolorbox}[colback=blue!5, colframe=blue!40!black, title=Задача интерполяции]
		Заданы \( n \) различных точек \( x_1, x_2, \dots, x_n \in K \) и соответствующие им значения \( y_1, y_2, \dots, y_n \in K \). Требуется найти многочлен \( f(t) \in K[t] \) такой, что
		\[
		f(x_i) = y_i \quad \text{для всех } i = 1, 2, \dots, n.
		\]
		Многочлен \( f \), удовлетворяющий этим условиям, называется \textit{интерполяционным многочленом} для данной таблицы значений.
	\end{tcolorbox}
	
	\myremark
	Обычно ищут многочлен наименьшей возможной степени, который, как будет показано, имеет степень не выше \( n-1 \).
	
	\subsection{Теорема существования и единственности}
	
	\begin{tcolorbox}[colback=green!5, colframe=green!50!black, title=Теорема Лагранжа]
		Для любого набора из \( n \) различных узлов \( x_1, \dots, x_n \in K \) и произвольных значений \( y_1, \dots, y_n \in K \) существует единственный многочлен \( f(t) \in K[t] \) степени не выше \( n-1 \), такой что
		\[
		f(x_i) = y_i \quad (i = 1, \dots, n).
		\]
	\end{tcolorbox}
	
	\medskip
	\textit{Доказательство.}
	
	\textbf{1. Существование.} Построим многочлен явно. Для каждого \( k = 1, \dots, n \) определим вспомогательный многочлен
	\[
	g_k(t) = \prod_{\substack{i=1 \\ i \neq k}}^n (t - x_i).
	\]
	Это многочлен степени \( n-1 \). Заметим, что
	\[
	g_k(x_k) = \prod_{i \neq k} (x_k - x_i) \neq 0,
	\]
	так как все узлы различны, а для \( j \neq k \) имеем \( g_k(x_j) = 0 \), поскольку в произведении присутствует множитель \( (x_j - x_j) = 0 \).
	
	Рассмотрим многочлен
	\[
	f(t) = \sum_{k=1}^n \frac{y_k}{g_k(x_k)} \, g_k(t).
	\]
	Проверим, что он удовлетворяет интерполяционным условиям. Для любого \( j \):
	\[
	f(x_j) = \sum_{k=1}^n \frac{y_k}{g_k(x_k)} \, g_k(x_j).
	\]
	В этой сумме все слагаемые с \( k \neq j \) равны нулю, так как \( g_k(x_j) = 0 \). При \( k = j \) получаем \( g_j(x_j) \neq 0 \), и множители сокращаются:
	\[
	f(x_j) = \frac{y_j}{g_j(x_j)} \cdot g_j(x_j) = y_j.
	\]
	Таким образом, многочлен \( f \) степени не выше \( n-1 \) (как сумма многочленов степени \( n-1 \)) является искомым.
	
	\textbf{2. Единственность.} Предположим, что существуют два многочлена \( f_1 \) и \( f_2 \) степени не выше \( n-1 \), удовлетворяющие интерполяционным условиям. Рассмотрим их разность:
	\[
	h(t) = f_1(t) - f_2(t).
	\]
	Тогда \( \deg h \leq n-1 \) (если \( h \neq 0 \)). Для каждого узла \( x_i \) имеем:
	\[
	h(x_i) = f_1(x_i) - f_2(x_i) = y_i - y_i = 0.
	\]
	Таким образом, \( h(t) \) имеет не менее \( n \) различных корней \( x_1, \dots, x_n \) в поле \( K \). Но ненулевой многочлен степени не выше \( n-1 \) не может иметь более \( n-1 \) корней (теорема о числе корней). Следовательно, \( h(t) \equiv 0 \), т.е. \( f_1 = f_2 \). \hfill \(\square\)
	
	\subsection{Интерполяционная формула Лагранжа}
	
	Выражение для \( g_k(x_k) \) можно записать явно. Так как
	\[
	g_k(x_k) = \prod_{\substack{i=1 \\ i \neq k}}^n (x_k - x_i),
	\]
	то интерполяционный многочлен принимает вид:
	
	\begin{tcolorbox}[colback=yellow!5, colframe=orange!50!black, title=Формула Лагранжа]
		\[
		f(t) = \sum_{k=1}^n y_k \cdot \frac{\displaystyle\prod_{\substack{i=1 \\ i \neq k}}^n (t - x_i)}{\displaystyle\prod_{\substack{i=1 \\ i \neq k}}^n (x_k - x_i)}.
		\]
	\end{tcolorbox}
	
	Это классическая запись интерполяционного многочлена Лагранжа. Каждое слагаемое представляет собой многочлен степени \( n-1 \), равный \( y_k \) в точке \( x_k \) и нулю во всех остальных узлах.
	
	\subsection{Компактная форма с использованием производной}
	
	Введём многочлен
	\[
	g(t) = \prod_{i=1}^n (t - x_i).
	\]
	Это многочлен степени \( n \), имеющий корни во всех узлах. Заметим, что
	\[
	g'(x_k) = \prod_{\substack{i=1 \\ i \neq k}}^n (x_k - x_i).
	\]
	Действительно, производная \( g'(t) \) может быть вычислена по формуле Лейбница, но проще заметить, что \( g(t) = (t - x_k) g_k(t) \), где \( g_k(t) = \prod_{i \neq k} (t - x_i) \). Тогда
	\[
	g'(t) = g_k(t) + (t - x_k) g_k'(t),
	\]
	и подстановка \( t = x_k \) даёт \( g'(x_k) = g_k(x_k) \), так как второе слагаемое обращается в ноль.
	
	Таким образом, формула Лагранжа может быть переписана в более компактном виде:
	
	\begin{tcolorbox}[colback=red!5, colframe=red!50!black, title=Компактная форма]
		\[
		f(t) = \sum_{k=1}^n \frac{y_k}{g'(x_k)} \cdot \frac{g(t)}{t - x_k},
		\]
		где \( g(t) = \prod_{i=1}^n (t - x_i) \).
	\end{tcolorbox}
	
	\myremark
	Эта форма удобна тем, что в ней явно выделен общий множитель \( g(t) \), и каждый член \( \frac{g(t)}{t - x_k} \) является многочленом (так как деление на \( t - x_k \) выполняется нацело в \( K[t] \)).
	
	\subsection{Пример}
	
	Построим интерполяционный многочлен Лагранжа для таблицы:
	
	\[
	\begin{array}{c|ccc}
		x & 1 & 2 & 4 \\
		\hline
		y & 1 & 3 & 3
	\end{array}
	\]
	
	\textit{Решение.} Здесь \( n = 3 \). Вычислим базисные многочлены Лагранжа:
	
	\[
	L_1(t) = \frac{(t-2)(t-4)}{(1-2)(1-4)} = \frac{(t-2)(t-4)}{(-1)(-3)} = \frac{(t-2)(t-4)}{3}.
	\]
	
	\[
	L_2(t) = \frac{(t-1)(t-4)}{(2-1)(2-4)} = \frac{(t-1)(t-4)}{1 \cdot (-2)} = -\frac{(t-1)(t-4)}{2}.
	\]
	
	\[
	L_3(t) = \frac{(t-1)(t-2)}{(4-1)(4-2)} = \frac{(t-1)(t-2)}{3 \cdot 2} = \frac{(t-1)(t-2)}{6}.
	\]
	
	Тогда интерполяционный многочлен:
	\[
	f(t) = 1 \cdot L_1(t) + 3 \cdot L_2(t) + 3 \cdot L_3(t).
	\]
	
	Вычислим:
	\[
	f(t) = \frac{(t-2)(t-4)}{3} - \frac{3}{2}(t-1)(t-4) + \frac{1}{2}(t-1)(t-2).
	\]
	
	Приведём к общему знаменателю 6 и упростим:
	\[
	f(t) = \frac{2(t-2)(t-4) - 9(t-1)(t-4) + 3(t-1)(t-2)}{6}.
	\]
	
	Раскроем скобки:
	\[
	\begin{aligned}
		& 2(t^2 - 6t + 8) = 2t^2 - 12t + 16, \\
		& -9(t^2 - 5t + 4) = -9t^2 + 45t - 36, \\
		& 3(t^2 - 3t + 2) = 3t^2 - 9t + 6.
	\end{aligned}
	\]
	
	Суммируем: \( (2 - 9 + 3)t^2 = -4t^2 \); \( (-12 + 45 - 9)t = 24t \); \( (16 - 36 + 6) = -14 \).
	
	Таким образом,
	\[
	f(t) = \frac{-4t^2 + 24t - 14}{6} = -\frac{2}{3}t^2 + 4t - \frac{7}{3}.
	\]
	
	Проверка: \( f(1) = -\frac{2}{3} + 4 - \frac{7}{3} = 4 - \frac{9}{3} = 4 - 3 = 1 \), \( f(2) = -\frac{8}{3} + 8 - \frac{7}{3} = 8 - \frac{15}{3} = 8 - 5 = 3 \), \( f(4) = -\frac{32}{3} + 16 - \frac{7}{3} = 16 - \frac{39}{3} = 16 - 13 = 3 \).
	
	\subsection{Замечания}
	
	\begin{itemize}
		\item Интерполяционный многочлен Лагранжа существует и единствен для любого набора различных узлов. Это следует из того, что система линейных уравнений для коэффициентов многочлена имеет матрицу Вандермонда, определитель которой отличен от нуля.
		\item Если степень искомого многочлена не ограничена сверху, то существует бесконечно много многочленов, удовлетворяющих интерполяционным условиям (можно добавить множитель \( \prod (t - x_i) \), умноженный на любой многочлен). Условие \( \deg f \leq n-1 \) обеспечивает единственность.
		\item Формула Лагранжа удобна для теоретических построений, но для вычислений при больших \( n \) часто используют интерполяционные схемы Ньютона или другие методы.
	\end{itemize}
	
	
	
	
	
	
	Вот текст для 41-го вопроса «Интерполяционный многочлен Ньютона», оформленный в том же стиле. Я объединил и структурировал разрозненные фрагменты в связное изложение: общая идея интерполяции, постановка задачи, метод неопределённых коэффициентов (форма Ньютона), рекуррентное вычисление коэффициентов (разделённые разности) и пример.
	
	```latex
	% ==================== РАЗДЕЛ 41: ИНТЕРПОЛЯЦИОННЫЙ МНОГОЧЛЕН НЬЮТОНА ====================
	\newpage
	\section{Интерполяционный многочлен Ньютона (41 вопрос)}
	
	\subsection{Общая задача интерполяции}
	
	Напомним постановку задачи интерполяции. Пусть заданы \( n \) различных узлов \( x_1, x_2, \dots, x_n \in K \) (где \( K \) — поле, например \( \mathbb{R} \) или \( \mathbb{C} \)) и соответствующие значения \( y_1, y_2, \dots, y_n \in K \). Требуется найти многочлен \( f(t) \in K[t] \), удовлетворяющий условиям
	\[
	f(x_i) = y_i \quad \text{для всех } i = 1, 2, \dots, n.
	\]
	
	В предыдущем вопросе (40) мы рассмотрели форму Лагранжа, которая даёт явное выражение интерполяционного многочлена в виде суммы базисных многочленов. Однако для практических вычислений, особенно при добавлении новых узлов, удобнее использовать форму Ньютона.
	
	\subsection{Интерполяционный многочлен Ньютона}
	
	Идея метода Ньютона состоит в том, чтобы искать интерполяционный многочлен в виде разложения по системе многочленов
	\[
	1, \quad (t - x_1), \quad (t - x_1)(t - x_2), \quad \dots, \quad (t - x_1)(t - x_2)\dots(t - x_{n-1}).
	\]
	
	\begin{tcolorbox}[colback=blue!5, colframe=blue!40!black, title=Определение]
		\textit{Интерполяционным многочленом Ньютона} для узлов \( x_1, x_2, \dots, x_n \) и значений \( y_1, y_2, \dots, y_n \) называется многочлен вида
		\[
		N(t) = a_0 + a_1(t - x_1) + a_2(t - x_1)(t - x_2) + \dots + a_{n-1}(t - x_1)(t - x_2)\dots(t - x_{n-1}),
		\]
		где коэффициенты \( a_0, a_1, \dots, a_{n-1} \in K \) определяются последовательно из интерполяционных условий.
	\end{tcolorbox}
	
	\myremark
	Степень многочлена Ньютона не превосходит \( n-1 \). Базисные многочлены обладают тем свойством, что при подстановке \( t = x_k \) все члены, начиная с \( k \)-го, обращаются в ноль (так как содержат множитель \( (t - x_k) \)). Это позволяет последовательно находить коэффициенты.
	
	\subsection{Метод последовательного определения коэффициентов}
	
	Подставим \( t = x_1 \). Все слагаемые, кроме первого, содержат множитель \( (x_1 - x_1) = 0 \), поэтому:
	\[
	N(x_1) = a_0 = y_1 \quad \Rightarrow \quad a_0 = y_1.
	\]
	
	Подставим \( t = x_2 \). Обратятся в ноль все слагаемые, начиная с третьего (содержат \( (t - x_1)(t - x_2) \) и выше). Получим:
	\[
	N(x_2) = a_0 + a_1(x_2 - x_1) = y_2.
	\]
	Отсюда, зная \( a_0 \), находим:
	\[
	a_1 = \frac{y_2 - a_0}{x_2 - x_1} = \frac{y_2 - y_1}{x_2 - x_1}.
	\]
	
	Подставим \( t = x_3 \). Обратятся в ноль слагаемые, начиная с четвёртого (содержат \( (t - x_1)(t - x_2)(t - x_3) \)). Имеем:
	\[
	N(x_3) = a_0 + a_1(x_3 - x_1) + a_2(x_3 - x_1)(x_3 - x_2) = y_3.
	\]
	Отсюда выражаем \( a_2 \):
	\[
	a_2 = \frac{y_3 - a_0 - a_1(x_3 - x_1)}{(x_3 - x_1)(x_3 - x_2)}.
	\]
	
	Продолжая этот процесс, мы последовательно находим все коэффициенты \( a_0, a_1, \dots, a_{n-1} \). Полученный многочлен \( N(t) \) совпадает с интерполяционным многочленом Лагранжа (в силу единственности).
	
	\subsection{Разделённые разности}
	
	Введённые коэффициенты \( a_k \) естественно выражаются через так называемые \textit{разделённые разности}. Обозначим:
	\[
	f[x_i] = y_i
	\]
	— разделённая разность нулевого порядка.
	
	Разделённая разность первого порядка:
	\[
	f[x_i, x_j] = \frac{f[x_j] - f[x_i]}{x_j - x_i}.
	\]
	
	Разделённая разность второго порядка:
	\[
	f[x_i, x_j, x_k] = \frac{f[x_j, x_k] - f[x_i, x_j]}{x_k - x_i},
	\]
	и так далее.
	
	Тогда коэффициенты \( a_k \) в интерполяционном многочлене Ньютона равны:
	\[
	a_0 = f[x_1], \quad a_1 = f[x_1, x_2], \quad a_2 = f[x_1, x_2, x_3], \quad \dots, \quad a_{n-1} = f[x_1, x_2, \dots, x_n].
	\]
	
	\begin{tcolorbox}[colback=green!5, colframe=green!50!black, title=Формула Ньютона через разделённые разности]
		\[
		N(t) = f[x_1] + f[x_1, x_2](t - x_1) + f[x_1, x_2, x_3](t - x_1)(t - x_2) + \dots + f[x_1, \dots, x_n](t - x_1)\dots(t - x_{n-1}).
		\]
	\end{tcolorbox}
	
	\subsection{Пример}
	
	Построим интерполяционный многочлен Ньютона для таблицы:
	
	\[
	\begin{array}{c|cccc}
		x & -2 & -1 & 2 & 3 \\
		\hline
		y & 2 & 4 & -2 & -8
	\end{array}
	\]
	
	\textit{Решение.} Будем искать многочлен в форме Ньютона с узлами в порядке \( x_1 = -2, x_2 = -1, x_3 = 2, x_4 = 3 \):
	\[
	N(t) = a_0 + a_1(t + 2) + a_2(t + 2)(t + 1) + a_3(t + 2)(t + 1)(t - 2).
	\]
	
	Последовательно находим коэффициенты.
	
	\textit{Шаг 1.} Подставим \( t = -2 \):
	\[
	N(-2) = a_0 = 2 \quad \Rightarrow \quad a_0 = 2.
	\]
	
	\textit{Шаг 2.} Подставим \( t = -1 \):
	\[
	N(-1) = a_0 + a_1(-1 + 2) = 2 + a_1 = 4 \quad \Rightarrow \quad a_1 = 2.
	\]
	
	\textit{Шаг 3.} Подставим \( t = 2 \):
	\[
	N(2) = a_0 + a_1(2 + 2) + a_2(2 + 2)(2 + 1) = 2 + 2 \cdot 4 + a_2 \cdot 4 \cdot 3 = 2 + 8 + 12a_2 = -2.
	\]
	Отсюда \( 12a_2 = -12 \), значит \( a_2 = -1 \).
	
	\textit{Шаг 4.} Подставим \( t = 3 \):
	\[
	N(3) = a_0 + a_1(3 + 2) + a_2(3 + 2)(3 + 1) + a_3(3 + 2)(3 + 1)(3 - 2).
	\]
	Вычислим известные члены:
	\[
	a_0 + a_1 \cdot 5 = 2 + 10 = 12,
	\]
	\[
	a_2 \cdot 5 \cdot 4 = (-1) \cdot 20 = -20,
	\]
	сумма первых трёх слагаемых: \( 12 - 20 = -8 \).
	Таким образом,
	\[
	N(3) = -8 + a_3 \cdot 5 \cdot 4 \cdot 1 = -8 + 20a_3.
	\]
	По условию \( N(3) = -8 \), следовательно, \( 20a_3 = 0 \), откуда \( a_3 = 0 \).
	
	Итак,
	\[
	N(t) = 2 + 2(t + 2) - (t + 2)(t + 1).
	\]
	
	Раскроем скобки и упростим:
	\[
	\begin{aligned}
		N(t) &= 2 + 2t + 4 - (t^2 + 3t + 2) \\
		&= 2t + 6 - t^2 - 3t - 2 \\
		&= -t^2 - t + 4.
	\end{aligned}
	\]
	
	Таким образом, интерполяционный многочлен: \( f(t) = -t^2 - t + 4 \).
	
	Проверим значения:
	\[
	f(-2) = -4 + 2 + 4 = 2, \quad f(-1) = -1 + 1 + 4 = 4, \quad f(2) = -4 - 2 + 4 = -2, \quad f(3) = -9 - 3 + 4 = -8.
	\]
	Все условия выполнены.
	
	\subsection{Преимущества формы Ньютона}
	
	\begin{itemize}
		\item При добавлении нового узла \( (x_{n+1}, y_{n+1}) \) не нужно пересчитывать все коэффициенты заново. Достаточно добавить к уже имеющемуся многочлену слагаемое
		\[
		a_n (t - x_1)(t - x_2)\dots(t - x_n),
		\]
		где \( a_n = f[x_1, \dots, x_{n+1}] \) вычисляется по уже известным разделённым разностям.
		\item Форма Ньютона естественно обобщается на случай кратных узлов (интерполяция с производными), что приводит к интерполяционному многочлену Эрмита.
		\item Вычисления с разделёнными разностями удобно организовывать в табличной форме.
	\end{itemize}
	
	\subsection{Связь с многочленом Лагранжа}
	
	В силу единственности интерполяционного многочлена степени не выше \( n-1 \), многочлены Ньютона и Лагранжа совпадают. Различие лишь в форме представления и способе вычисления коэффициентов.
	
	
	
	% ==================== РАЗДЕЛ 42: ПОСТРОЕНИЕ ПОЛЯ ЧАСТНЫХ ====================
	\newpage
	\section{Построение поля частных (42 вопрос)}
	
	\subsection{Мотивация и определение}
	
	Не каждое кольцо является полем. Например, кольцо целых чисел \( \mathbb{Z} \) не является полем, так как не у каждого ненулевого элемента есть обратный в \( \mathbb{Z} \). Однако целые числа можно вложить в поле рациональных чисел \( \mathbb{Q} \), где каждый ненулевой элемент уже обратим.
	
	Аналогичная конструкция существует для любой области целостности: по ней можно построить поле, в которое она вкладывается, называемое \textit{полем частных} (или полем отношений).
	
	\begin{tcolorbox}[colback=blue!5, colframe=blue!40!black, title=Определение]
		Пусть \( R \) — область целостности (коммутативное кольцо с единицей \( 1 \neq 0 \) без делителей нуля). \textit{Полем частных} кольца \( R \) называется поле \( K \), обладающее следующими свойствами:
		\begin{enumerate}
			\item \( R \subseteq K \) (точнее, существует инъективный гомоморфизм \( R \hookrightarrow K \));
			\item Каждый элемент \( k \in K \) может быть представлен в виде частного \( \frac{a}{b} \) с \( a, b \in R \), \( b \neq 0 \);
			\item \( K \) является наименьшим полем с этими свойствами: если \( K' \) — поле, содержащее \( R \), то \( K \subseteq K' \).
		\end{enumerate}
	\end{tcolorbox}
	
	\myremark
	Поле частных области целостности единственно с точностью до изоморфизма. Для \( \mathbb{Z} \) это \( \mathbb{Q} \), для кольца многочленов \( K[x] \) — поле рациональных функций \( K(x) \).
	
	\subsection{Конструкция поля частных}
	
	Пусть \( R \) — область целостности. Рассмотрим множество формальных дробей:
	\[
	F = \{ (a, b) \mid a, b \in R, \ b \neq 0 \}.
	\]
	
	Введём на \( F \) отношение эквивалентности:
	\[
	(a, b) \sim (c, d) \iff ad = bc.
	\]
	Корректность: это отношение рефлексивно, симметрично и транзитивно (последнее следует из отсутствия делителей нуля в \( R \)).
	
	Обозначим класс эквивалентности пары \( (a, b) \) через \( \frac{a}{b} \). Множество классов эквивалентности обозначим \( K \).
	
	\subsection{Определение операций в \( K \)}
	
	Определим сложение и умножение на \( K \) по обычным правилам действий с дробями:
	
	\[
	\frac{a}{b} + \frac{c}{d} = \frac{ad + bc}{bd},
	\]
	\[
	\frac{a}{b} \cdot \frac{c}{d} = \frac{ac}{bd}.
	\]
	
	\begin{tcolorbox}[colback=green!5, colframe=green!50!black, title=Лемма (Корректность определений)]
		Операции сложения и умножения определены корректно, т.е. не зависят от выбора представителей классов эквивалентности.
	\end{tcolorbox}
	
	\medskip
	\textit{Доказательство для умножения.}
	Пусть \( \frac{a}{b} = \frac{a'}{b'} \) и \( \frac{c}{d} = \frac{c'}{d'} \). Тогда \( ab' = a'b \) и \( cd' = c'd \). Нужно показать, что \( \frac{ac}{bd} = \frac{a'c'}{b'd'} \), т.е. \( (ac)(b'd') = (a'c')(bd) \).
	
	Имеем:
	\[
	(ac)(b'd') = a c b' d' = (a b')(c d') = (a' b)(c' d) = a' c' b d = (a'c')(bd),
	\]
	что и требовалось. \hfill \(\square\)
	
	\medskip
	\textit{Доказательство для сложения} проводится аналогично, используя дистрибутивность и свойства равенств:
	\[
	(ad + bc)(b'd') = ad b'd' + bc b'd' = (a b')(d d') + (c d')(b b') = (a' b)(d d') + (c' d)(b b') = a' d' b d + b' c' b d = (a'd' + b'c')(bd),
	\]
	что соответствует равенству \( \frac{ad+bc}{bd} = \frac{a'd' + b'c'}{b'd'} \). \hfill \(\square\)
	
	\subsection{Проверка аксиом поля}
	
	\begin{itemize}
		\item \textbf{Сложение:} коммутативно, ассоциативно, нейтральный элемент \( \frac{0}{1} \) (где \( 0 \in R \)), противоположный элемент для \( \frac{a}{b} \) есть \( \frac{-a}{b} \).
		\item \textbf{Умножение:} коммутативно, ассоциативно, дистрибутивно относительно сложения, нейтральный элемент \( \frac{1}{1} \).
		\item \textbf{Обратный элемент:} Пусть \( \frac{a}{b} \neq \frac{0}{1} \) в \( K \). Это означает, что \( a \neq 0 \) в \( R \) (иначе \( a \cdot 1 = 0 \cdot b \) дало бы \( 0 = 0 \), но класс \( \frac{0}{b} \) равен \( \frac{0}{1} \), так как \( 0 \cdot 1 = 0 \cdot b \)). Тогда \( \frac{b}{a} \in K \), и
		\[
		\frac{a}{b} \cdot \frac{b}{a} = \frac{ab}{ba} = \frac{ab}{ab} = \frac{1}{1},
		\]
		так как \( (ab) \cdot 1 = 1 \cdot (ab) \). Таким образом, \( \frac{b}{a} \) — обратный элемент.
	\end{itemize}
	
	Следовательно, \( (K, +, \cdot) \) является полем.
	
	\subsection{Вложение \( R \) в \( K \)}
	
	Зафиксируем некоторый ненулевой элемент \( d \in R \). Рассмотрим отображение
	\[
	\varphi_d: R \to K, \quad \varphi_d(x) = \frac{xd}{d}.
	\]
	
	Покажем, что это отображение не зависит от выбора \( d \). Действительно, если \( d' \) — другой ненулевой элемент, то
	\[
	\frac{xd}{d} = \frac{xd'}{d'} \iff (xd)d' = (xd')d \iff x d d' = x d' d,
	\]
	что верно в коммутативном кольце. Таким образом, все \( \varphi_d \) совпадают, и мы можем обозначить образ элемента \( x \) просто как \( \frac{x}{1} \) (выбирая \( d = 1 \)).
	
	Отображение \( x \mapsto \frac{x}{1} \) инъективно: если \( \frac{x}{1} = \frac{y}{1} \), то \( x \cdot 1 = y \cdot 1 \), т.е. \( x = y \). Также оно сохраняет операции:
	\[
	\frac{x}{1} + \frac{y}{1} = \frac{x+y}{1}, \quad \frac{x}{1} \cdot \frac{y}{1} = \frac{xy}{1}.
	\]
	
	Таким образом, мы отождествляем \( R \) с подкольцом \( \overline{R} = \{ \frac{x}{1} \mid x \in R \} \subseteq K \).
	
	\subsection{Представление элементов поля частных}
	
	Пусть \( \frac{a}{b} \in K \). Тогда в поле \( K \) имеем:
	\[
	\frac{a}{b} = \frac{a}{1} \cdot \left( \frac{b}{1} \right)^{-1} = a \cdot b^{-1},
	\]
	где \( a, b \) рассматриваются как элементы \( K \) через вложение. Таким образом, каждый элемент поля частных действительно является частным двух элементов из \( R \).
	
	\subsection{Универсальное свойство}
	
	\begin{tcolorbox}[colback=yellow!5, colframe=orange!50!black, title=Теорема (Универсальное свойство)]
		Пусть \( R \) — область целостности, \( K \) — её поле частных. Тогда для любого поля \( F \) и любого инъективного гомоморфизма колец \( \psi: R \to F \) существует единственное продолжение до гомоморфизма полей \( \tilde{\psi}: K \to F \), задаваемое формулой
		\[
		\tilde{\psi}\left( \frac{a}{b} \right) = \frac{\psi(a)}{\psi(b)}.
		\]
	\end{tcolorbox}
	
	Это свойство означает, что поле частных является наименьшим полем, содержащим \( R \), и любое другое поле, содержащее \( R \), содержит и его (с точностью до изоморфизма).
	
	\subsection{Примеры}
	
	\begin{itemize}
		\item Для \( R = \mathbb{Z} \) поле частных — \( \mathbb{Q} \).
		\item Для \( R = K[x] \) (кольцо многочленов над полем \( K \)) поле частных — \( K(x) \), поле рациональных функций.
		\item Для \( R = \mathbb{Z}[i] \) (кольцо целых гауссовых чисел) поле частных — \( \mathbb{Q}(i) \).
	\end{itemize}
	
	\subsection{Замечание}
	
	Если \( R \) уже является полем, то его поле частных изоморфно самому \( R \). Действительно, для поля каждый ненулевой элемент обратим, поэтому конструкция поля частных не даёт ничего нового
	
	
	
	
	
	
	% ==================== РАЗДЕЛ 43: ПОЛЕ ДРОБНО-РАЦИОНАЛЬНЫХ ФУНКЦИЙ ====================
	\newpage
	\section{Поле дробно-рациональных функций. Правильные и неправильные дроби. Выделение целой части (43 вопрос)}
	
	\subsection{Определение поля рациональных функций}
	
	Пусть \( K \) — поле. Рассмотрим кольцо многочленов \( K[t] \). Поскольку \( K[t] \) является областью целостности, для него можно построить поле частных.
	
	\begin{tcolorbox}[colback=blue!5, colframe=blue!40!black, title=Определение]
		\textit{Полем рациональных функций} (или полем дробно-рациональных функций) от переменной \( t \) над полем \( K \) называется поле частных кольца многочленов \( K[t] \). Оно обозначается \( K(t) \).
	\end{tcolorbox}
	
	Элементами поля \( K(t) \) являются формальные дроби вида \( \frac{f(t)}{g(t)} \), где \( f, g \in K[t] \), \( g \neq 0 \), с обычными правилами сложения и умножения (как в поле частных). Две дроби \( \frac{f_1}{g_1} \) и \( \frac{f_2}{g_2} \) считаются равными, если \( f_1 g_2 = f_2 g_1 \) в \( K[t] \).
	
	\subsection{Каноническая запись}
	
	Каждую рациональную функцию можно представить в виде дроби, у которой числитель и знаменатель взаимно просты (сокращённая дробь). Обычно знаменатель берут нормированным (старший коэффициент равен 1), что обеспечивает единственность такой записи с точностью до знака.
	
	\medskip
	\noindent \textbf{Определение (Правильная дробь).}
	Рациональная функция \( u = \frac{f}{g} \in K(t) \) называется \textit{правильной дробью}, если \( \deg f < \deg g \). Если \( u = 0 \), её также считают правильной. В противном случае дробь называется \textit{неправильной}.
	
	\myremark
	Для нулевой дроби \( 0 = \frac{0}{1} \) условие \( \deg 0 < \deg 1 \) выполнено (степень нулевого многочлена обычно полагают \( -\infty \) или считают, что условие выполняется автоматически).
	
	\subsection{Свойства правильных дробей}
	
	\begin{tcolorbox}[colback=green!5, colframe=green!50!black, title=Свойство 1]
		Сумма, разность и произведение двух правильных дробей являются правильными дробями. Таким образом, множество правильных дробей образует подкольцо в \( K(t) \).
	\end{tcolorbox}
	
	\medskip
	\textit{Доказательство.}
	Если одна из дробей нулевая, утверждение очевидно. Пусть \( u_1 = \frac{f_1}{g_1} \) и \( u_2 = \frac{f_2}{g_2} \) — ненулевые правильные дроби, т.е. \( \deg f_1 < \deg g_1 \) и \( \deg f_2 < \deg g_2 \).
	
	Рассмотрим сумму:
	\[
	u_1 + u_2 = \frac{f_1 g_2 + f_2 g_1}{g_1 g_2}.
	\]
	Оценим степени числителя и знаменателя. Знаменатель: \( \deg(g_1 g_2) = \deg g_1 + \deg g_2 \). Числитель: каждый из двух членов \( f_1 g_2 \) и \( f_2 g_1 \) имеет степень
	\[
	\deg(f_1 g_2) = \deg f_1 + \deg g_2 < \deg g_1 + \deg g_2,
	\]
	\[
	\deg(f_2 g_1) = \deg f_2 + \deg g_1 < \deg g_2 + \deg g_1.
	\]
	Следовательно, степень числителя не превосходит максимума этих двух чисел, который строго меньше \( \deg g_1 + \deg g_2 \). Таким образом, сумма — правильная дробь. Для разности рассуждение аналогично.
	
	Для произведения:
	\[
	u_1 u_2 = \frac{f_1 f_2}{g_1 g_2}.
	\]
	Имеем \( \deg(f_1 f_2) = \deg f_1 + \deg f_2 < \deg g_1 + \deg g_2 = \deg(g_1 g_2) \), следовательно, произведение — правильная дробь. \hfill \(\square\)
	
	\subsection{Выделение целой части}
	
	Любую рациональную функцию можно представить в виде суммы многочлена (целой части) и правильной дроби. Это аналог выделения целой части у неправильных дробей в арифметике.
	
	\begin{tcolorbox}[colback=yellow!5, colframe=orange!50!black, title=Свойство 2 (Выделение целой части)]
		Для любой рациональной функции \( u \in K(t) \) существуют единственные многочлен \( q \in K[t] \) и правильная дробь \( r \in K(t) \) такие, что
		\[
		u = q + r.
		\]
	\end{tcolorbox}
	
	\medskip
	\textit{Доказательство.}
	Запишем \( u \) в виде дроби \( u = \frac{f}{g} \) с \( f, g \in K[t] \), \( g \neq 0 \). Выполним деление числителя на знаменатель с остатком в кольце \( K[t] \) (это возможно, так как \( K \) — поле):
	\[
	f = q g + r, \quad q, r \in K[t], \quad \deg r < \deg g \text{ или } r = 0.
	\]
	Тогда
	\[
	u = \frac{f}{g} = \frac{q g + r}{g} = q + \frac{r}{g}.
	\]
	Здесь \( q \in K[t] \), а \( \frac{r}{g} \) — правильная дробь (так как \( \deg r < \deg g \)).
	
	Единственность: если \( u = q_1 + r_1 = q_2 + r_2 \) с \( q_i \in K[t] \) и правильными \( r_i \), то \( q_1 - q_2 = r_2 - r_1 \). Левая часть — многочлен, правая — правильная дробь (разность правильных дробей правильна по свойству 1). Единственный многочлен, являющийся правильной дробью, — это нулевой (так как ненулевой многочлен имеет степень \( \geq 0 \), а правильная дробь должна иметь степень числителя меньше степени знаменателя, что для многочлена (знаменатель 1) означало бы степень \( < 0 \), т.е. только нуль). Следовательно, \( q_1 = q_2 \) и \( r_1 = r_2 \). \hfill \(\square\)
	
	\subsection{Разложение на простейшие дроби (основная идея)}
	
	Если знаменатель правильной дроби разлагается на взаимно простые множители, то такую дробь можно представить в виде суммы более простых дробей.
	
	\begin{tcolorbox}[colback=blue!5, colframe=blue!40!black, title=Свойство 3 (Разложение по взаимно простым знаменателям)]
		Пусть \( u = \frac{f}{g_1 g_2} \) — правильная дробь, где \( g_1, g_2 \in K[t] \) взаимно просты (т.е. \( \operatorname{НОД}(g_1, g_2) = 1 \)). Тогда существуют такие многочлены \( h_1, h_2 \in K[t] \), что
		\[
		u = \frac{h_1}{g_1} + \frac{h_2}{g_2},
		\]
		причём дроби \( \frac{h_1}{g_1} \) и \( \frac{h_2}{g_2} \) можно выбрать правильными.
	\end{tcolorbox}
	
	\medskip
	\textit{Доказательство (идея).}
	Из взаимной простоты \( g_1 \) и \( g_2 \) следует существование многочленов \( u_1, u_2 \in K[t] \) таких, что \( u_1 g_1 + u_2 g_2 = 1 \) (тождество Безу). Тогда
	\[
	\frac{f}{g_1 g_2} = \frac{f(u_1 g_1 + u_2 g_2)}{g_1 g_2} = \frac{f u_1}{g_2} + \frac{f u_2}{g_1}.
	\]
	Теперь к каждой из полученных дробей можно применить выделение целой части (свойство 2) и "перебросить" целые части в другую дробь, добиваясь правильности. В результате получим представление в виде суммы правильных дробей с знаменателями \( g_1 \) и \( g_2 \). \hfill \(\square\)
	
	\subsection{Обобщение на несколько сомножителей}
	
	Методом математической индукции свойство 3 обобщается на случай произведения любого конечного числа попарно взаимно простых многочленов.
	
	\begin{tcolorbox}[colback=blue!5, colframe=blue!40!black, title=Свойство 4]
		Пусть \( u = \frac{f}{g_1 g_2 \dots g_n} \) — правильная дробь, где многочлены \( g_1, g_2, \dots, g_n \) попарно взаимно просты. Тогда \( u \) может быть представлена в виде суммы правильных дробей со знаменателями \( g_1, g_2, \dots, g_n \):
		\[
		u = \frac{h_1}{g_1} + \frac{h_2}{g_2} + \dots + \frac{h_n}{g_n},
		\]
		где каждая дробь \( \frac{h_i}{g_i} \) правильная.
	\end{tcolorbox}
	
	\medskip
	\textit{Доказательство.}
	Применяем индукцию по \( n \). База \( n = 2 \) — свойство 3. Для \( n > 2 \) полагаем \( G = g_2 g_3 \dots g_n \). Тогда \( g_1 \) и \( G \) взаимно просты, и по свойству 3
	\[
	\frac{f}{g_1 G} = \frac{h_1}{g_1} + \frac{h}{G},
	\]
	где \( \frac{h_1}{g_1} \) и \( \frac{h}{G} \) правильные. Затем применяем предположение индукции к \( \frac{h}{G} \) (знаменатель \( G \) уже разложен в произведение \( n-1 \) попарно взаимно простых множителей). \hfill \(\square\)
	
	\subsection{Примеры}
	
	\textbf{Пример 1 (Выделение целой части).}
	\[
	u = \frac{t^3 + 2t^2 + 3t + 4}{t^2 + 1}.
	\]
	Делим числитель на знаменатель: \( t^3 + 2t^2 + 3t + 4 = (t + 2)(t^2 + 1) + (2t + 2) \). Таким образом,
	\[
	u = t + 2 + \frac{2t + 2}{t^2 + 1},
	\]
	где \( \frac{2t+2}{t^2+1} \) — правильная дробь (степень числителя 1 < степени знаменателя 2).
	
	\textbf{Пример 2 (Разложение на простейшие).}
	Разложим правильную дробь \( \frac{3t+5}{(t-1)(t+2)} \) на простейшие.
	
	Ищем представление \( \frac{3t+5}{(t-1)(t+2)} = \frac{A}{t-1} + \frac{B}{t+2} \). Приводим к общему знаменателю:
	\[
	\frac{A(t+2) + B(t-1)}{(t-1)(t+2)} = \frac{(A+B)t + (2A - B)}{(t-1)(t+2)}.
	\]
	Приравниваем коэффициенты числителя: \( A+B = 3 \), \( 2A - B = 5 \). Решая систему, находим \( A = \frac{8}{3} \), \( B = \frac{1}{3} \). Таким образом,
	\[
	\frac{3t+5}{(t-1)(t+2)} = \frac{8}{3}\cdot\frac{1}{t-1} + \frac{1}{3}\cdot\frac{1}{t+2}.
	\]
	
	\subsection{Замечание}
	
	Разложение на простейшие дроби играет ключевую роль в интегрировании рациональных функций в математическом анализе, а также в теории линейных дифференциальных уравнений и других областях
	
	
	
	
	% ==================== РАЗДЕЛ 44: РАЗЛОЖЕНИЕ ПРАВИЛЬНОЙ ДРОБИ НА ПРОСТЕЙШИЕ ====================
	\newpage
	\section{Разложение правильной дроби на простейшие (44 вопрос)}
	
	\subsection{Определение простейшей дроби}
	
	Пусть \( K \) — поле (в приложениях чаще всего \( K = \mathbb{R} \) или \( \mathbb{C} \)). Рассмотрим поле рациональных функций \( K(t) \).
	
	\begin{tcolorbox}[colback=blue!5, colframe=blue!40!black, title=Определение]
		\textit{Простейшей дробью} над полем \( K \) называется ненулевая рациональная функция, каноническая запись которой имеет вид
		\[
		\frac{f(t)}{p(t)^k},
		\]
		где \( p(t) \in K[t] \) — неприводимый над \( K \) многочлен, \( k \in \mathbb{N} \), а \( f(t) \in K[t] \) таков, что \( \deg f < \deg p \) и дробь несократима (т.е. \( \operatorname{НОД}(f, p) = 1 \)).
	\end{tcolorbox}
	
	\myremark
	В случае \( K = \mathbb{C} \) неприводимые многочлены — это линейные \( (t - d) \), поэтому простейшие дроби над \( \mathbb{C} \) имеют вид \( \frac{A}{(t - d)^k} \) с \( A \in \mathbb{C} \). В случае \( K = \mathbb{R} \) добавляются дроби вида \( \frac{At + B}{(t^2 + pt + q)^k} \) с отрицательным дискриминантом.
	
	\subsection{Теорема о разложении на простейшие}
	
	\begin{tcolorbox}[colback=green!5, colframe=green!50!black, title=Теорема]
		Любая правильная дробь \( u \in K(t) \) может быть представлена в виде суммы конечного числа простейших дробей над \( K \). Такое представление единственно (с точностью до порядка слагаемых).
	\end{tcolorbox}
	
	\medskip
	\textit{Доказательство.}
	
	Пусть \( u = \frac{f}{g} \) — правильная дробь в канонической записи (т.е. \( \operatorname{НОД}(f, g) = 1 \), \( \deg f < \deg g \)). Разложим знаменатель \( g \) на неприводимые множители над \( K \):
	\[
	g = p_1^{k_1} p_2^{k_2} \dots p_m^{k_m},
	\]
	где \( p_i \) — различные неприводимые многочлены, \( k_i \in \mathbb{N} \).
	
	\textbf{Шаг 1. Сведение к степеням одного неприводимого.}
	Воспользуемся свойством разложения правильной дроби по взаимно простым знаменателям (см. вопрос 43). Поскольку многочлены \( p_i^{k_i} \) попарно взаимно просты, существуют такие правильные дроби \( u_i = \frac{h_i}{p_i^{k_i}} \), что
	\[
	u = u_1 + u_2 + \dots + u_m.
	\]
	
	Таким образом, достаточно доказать теорему для дроби вида \( \frac{h}{p^k} \), где \( p \) неприводим, \( \deg h < \deg(p^k) \) и \( \operatorname{НОД}(h, p) = 1 \) (иначе дробь сократима и её можно упростить).
	
	\textbf{Шаг 2. Индукция по \( k \).}
	Докажем, что любую правильную дробь вида \( \frac{h}{p^k} \) можно разложить в сумму простейших дробей со знаменателями \( p, p^2, \dots, p^k \).
	
	\textit{База индукции:} \( k = 1 \). Дробь \( \frac{h}{p} \) сама является простейшей, так как \( \deg h < \deg p \) и \( \operatorname{НОД}(h, p) = 1 \). Утверждение верно.
	
	\textit{Индукционный переход:} Пусть утверждение верно для всех показателей, меньших \( k \). Рассмотрим дробь \( \frac{h}{p^k} \). Разделим числитель \( h \) на \( p \) с остатком в кольце \( K[t] \):
	\[
	h = q p + r, \quad q, r \in K[t], \quad \deg r < \deg p.
	\]
	Тогда
	\[
	\frac{h}{p^k} = \frac{q p + r}{p^k} = \frac{q}{p^{k-1}} + \frac{r}{p^k}.
	\]
	
	Заметим, что \( \frac{q}{p^{k-1}} \) — правильная дробь? Не обязательно. Однако мы можем применить к ней выделение целой части (свойство 2 из вопроса 43) и снова разложить, но это может нарушить правильность. Вместо этого поступим иначе: заметим, что \( \deg q < \deg p^{k-1} \) не гарантировано. Но мы можем записать \( \frac{q}{p^{k-1}} = \frac{q}{p^{k-1}} \) и, если эта дробь неправильная, выделить из неё целую часть, которая на самом деле должна быть нулевой, потому что исходная дробь правильная. Проведём аккуратное рассуждение.
	
	Так как \( \deg h < \deg(p^k) = k \deg p \), а \( \deg(q p) = \deg q + \deg p \), то из \( h = q p + r \) и \( \deg r < \deg p \) следует, что \( \deg q \leq (k-1)\deg p - 1 \) (иначе степень левой части была бы не меньше \( k \deg p \)). Таким образом, \( \frac{q}{p^{k-1}} \) — правильная дробь (степень числителя меньше степени знаменателя). Действительно, \( \deg q < (k-1)\deg p = \deg(p^{k-1}) \).
	
	Итак, \( \frac{q}{p^{k-1}} \) — правильная дробь того же типа, но с меньшим показателем \( k-1 \). По предположению индукции она раскладывается в сумму простейших дробей со знаменателями \( p, \dots, p^{k-1} \).
	
	Дробь \( \frac{r}{p^k} \) уже является простейшей (или суммой простейших, если \( r \) не взаимно просто с \( p \)? Но \( \operatorname{НОД}(r, p) = 1 \), так как иначе исходная дробь была бы сократима, что противоречит каноничности записи). Следовательно, она сама есть простейшая.
	
	Таким образом, \( \frac{h}{p^k} \) представлена как сумма простейших дробей. Индукция завершена. \hfill \(\square\)
	
	\subsection{Явные формулы для случая простых корней (над \( \mathbb{C} \))}
	
	Если поле \( K \) таково, что все неприводимые множители знаменателя линейны (например, \( K = \mathbb{C} \) или мы рассматриваем разложение над алгебраическим замыканием), то разложение на простейшие принимает особенно простой вид.
	
	Пусть \( g(t) = (t - d_1)(t - d_2)\dots(t - d_n) \), где все \( d_i \) различны (простые корни), и \( \frac{f}{g} \) — правильная дробь. Тогда существует единственное представление
	\[
	\frac{f(t)}{g(t)} = \sum_{i=1}^n \frac{A_i}{t - d_i},
	\]
	где \( A_i \in K \).
	
	Для нахождения коэффициентов \( A_i \) умножим обе части на \( g(t) \):
	\[
	f(t) = \sum_{i=1}^n A_i \prod_{j \neq i} (t - d_j).
	\]
	Подставим \( t = d_k \). Все слагаемые в сумме, кроме \( i = k \), обратятся в ноль, так как содержат множитель \( (d_k - d_i) \). Получим:
	\[
	f(d_k) = A_k \prod_{j \neq k} (d_k - d_j).
	\]
	Заметим, что произведение \( \prod_{j \neq k} (d_k - d_j) \) равно \( g'(d_k) \), где \( g'(t) \) — производная \( g(t) \). Действительно, из \( g(t) = \prod_{i=1}^n (t - d_i) \) имеем \( g'(t) = \sum_{i=1}^n \prod_{j \neq i} (t - d_j) \), и при \( t = d_k \) все члены, кроме \( i = k \), равны нулю, так что \( g'(d_k) = \prod_{j \neq k} (d_k - d_j) \).
	
	Таким образом,
	\[
	A_k = \frac{f(d_k)}{g'(d_k)}.
	\]
	
	\begin{tcolorbox}[colback=yellow!5, colframe=orange!50!black, title=Формула для простых корней]
		Если \( g(t) = \prod_{i=1}^n (t - d_i) \) с попарно различными \( d_i \), то
		\[
		\frac{f(t)}{g(t)} = \sum_{i=1}^n \frac{f(d_i)}{g'(d_i)(t - d_i)}.
		\]
	\end{tcolorbox}
	
	Эта формула особенно полезна при интегрировании рациональных функций и в других приложениях.
	
	\subsection{Пример}
	
	Разложить на простейшие дроби над \( \mathbb{C} \):
	\[
	\frac{3t + 5}{(t-1)(t+2)}.
	\]
	
	Здесь \( g(t) = (t-1)(t+2) \), \( g'(t) = (t+2) + (t-1) = 2t + 1 \). Корни: \( d_1 = 1 \), \( d_2 = -2 \).
	
	Вычисляем:
	\[
	A_1 = \frac{f(1)}{g'(1)} = \frac{3 \cdot 1 + 5}{2 \cdot 1 + 1} = \frac{8}{3},
	\]
	\[
	A_2 = \frac{f(-2)}{g'(-2)} = \frac{3 \cdot (-2) + 5}{2 \cdot (-2) + 1} = \frac{-6 + 5}{-4 + 1} = \frac{-1}{-3} = \frac{1}{3}.
	\]
	
	Таким образом,
	\[
	\frac{3t + 5}{(t-1)(t+2)} = \frac{8/3}{t-1} + \frac{1/3}{t+2}.
	\]
	
	\subsection{Единственность разложения}
	
	Единственность разложения на простейшие следует из единственности разложения многочлена на неприводимые множители и из линейной независимости дробей вида \( \frac{t^j}{p(t)^k} \) (где \( 0 \leq j < \deg p \)) над полем \( K \). В частности, для случая простых корней коэффициенты \( A_i \) определяются однозначно по формуле выше
	
	
	
	
	
	
	% ==================== РАЗДЕЛ 45: МАТРИЦЫ И ДЕЙСТВИЯ НАД НИМИ ====================
	\newpage
	\section{Матрицы и действия над ними (45 вопрос)}
	
	\subsection{Определение матрицы}
	
	Пусть \( \mathbb{R} \) — некоторое кольцо (в частности, может быть полем действительных чисел, но в общем случае — любое коммутативное кольцо с единицей). Рассмотрим два множества индексов \( I \) и \( J \) (обычно конечные множества, например \( I = \{1, \dots, m\} \), \( J = \{1, \dots, n\} \)).
	
	\begin{tcolorbox}[colback=blue!5, colframe=blue!40!black, title=Определение]
		\textit{Матрицей} размера \( I \times J \) над кольцом \( \mathbb{R} \) называется семейство элементов \( (a_{ij})_{i \in I, j \in J} \), где \( a_{ij} \in \mathbb{R} \). Множество всех таких матриц обозначается \( \mathbb{R}^{I \times J} \).
	\end{tcolorbox}
	
	Если \( |I| = m \), \( |J| = n \), то говорят о матрице размера \( m \times n \). При \( m = n \) матрица называется \textit{квадратной}, число \( n \) называется её \textit{порядком}.
	
	Матрицу принято записывать в виде прямоугольной таблицы:
	\[
	A = 
	\begin{pmatrix}
		a_{11} & a_{12} & \dots & a_{1n} \\
		a_{21} & a_{22} & \dots & a_{2n} \\
		\vdots & \vdots & \ddots & \vdots \\
		a_{m1} & a_{m2} & \dots & a_{mn}
	\end{pmatrix}.
	\]
	
	\subsection{Подматрицы}
	
	Если \( I' \subset I \) и \( J' \subset J \), то матрица \( (a_{ij})_{i \in I', j \in J'} \) называется \textit{подматрицей} матрицы \( A \). В частности, при \( |I'| = 1 \) получаем \textit{строку}, при \( |J'| = 1 \) — \textit{столбец}.
	
	\subsection{Сложение матриц}
	
	\begin{tcolorbox}[colback=green!5, colframe=green!50!black, title=Определение сложения]
		Пусть \( A, B \in \mathbb{R}^{I \times J} \). Их \textit{суммой} называется матрица \( A + B \in \mathbb{R}^{I \times J} \), определяемая поэлементно:
		\[
		(A + B)_{ij} = A_{ij} + B_{ij} \quad \forall i \in I, \forall j \in J.
		\]
	\end{tcolorbox}
	
	\medskip
	\noindent \textbf{Свойства сложения:}
	\begin{enumerate}
		\item Если сложение в \( \mathbb{R} \) ассоциативно, то сложение матриц ассоциативно: \( (A + B) + C = A + (B + C) \).
		\item Если сложение в \( \mathbb{R} \) коммутативно, то сложение матриц коммутативно: \( A + B = B + A \).
		\item Существует \textit{нулевая матрица} \( O \in \mathbb{R}^{I \times J} \), у которой все элементы равны нулю. Для любой матрицы \( A \):
		\[
		A + O = O + A = A.
		\]
		\item Для любой матрицы \( A \) существует \textit{противоположная матрица} \( -A \), определяемая как \( (-A)_{ij} = -A_{ij} \). Тогда
		\[
		A + (-A) = (-A) + A = O.
		\]
	\end{enumerate}
	
	\subsection{Умножение матрицы на скаляр}
	
	\begin{tcolorbox}[colback=green!5, colframe=green!50!black, title=Определение умножения на скаляр]
		Пусть \( A \in \mathbb{R}^{I \times J} \) и \( d \in \mathbb{R} \). Произведение \( dA \) (и \( Ad \)) определяется поэлементно:
		\[
		(dA)_{ij} = d \cdot A_{ij}, \quad (Ad)_{ij} = A_{ij} \cdot d \quad \forall i \in I, \forall j \in J.
		\]
		Если кольцо \( \mathbb{R} \) коммутативно, то \( dA = Ad \).
	\end{tcolorbox}
	
	\medskip
	\noindent \textbf{Свойства умножения на скаляр:}
	\begin{enumerate}
		\item Если умножение в \( \mathbb{R} \) ассоциативно, то для любых \( \alpha, \beta \in \mathbb{R} \) и матрицы \( A \):
		\[
		(\alpha \beta) A = \alpha (\beta A), \quad (\alpha A) \beta = \alpha (A \beta), \quad A (\alpha \beta) = (A \alpha) \beta.
		\]
		\item Если в \( \mathbb{R} \) выполнены дистрибутивные законы, то для любых \( \alpha, \beta \in \mathbb{R} \) и матриц \( A, B \):
		\[
		(\alpha + \beta) A = \alpha A + \beta A, \quad A (\alpha + \beta) = A \alpha + A \beta,
		\]
		\[
		\alpha (A + B) = \alpha A + \alpha B, \quad (A + B) \alpha = A \alpha + B \alpha.
		\]
		\item Если в \( \mathbb{R} \) есть единица \( 1 \), то \( 1 \cdot A = A \cdot 1 = A \).
	\end{enumerate}
	
	\subsection{Умножение матриц}
	
	\begin{tcolorbox}[colback=red!5, colframe=red!50!black, title=Определение умножения матриц]
		Пусть \( A \in \mathbb{R}^{I \times J} \) и \( B \in \mathbb{R}^{J \times K} \). Их \textit{произведением} называется матрица \( AB \in \mathbb{R}^{I \times K} \), определяемая формулой:
		\[
		(AB)_{ik} = \sum_{j \in J} A_{ij} B_{jk} \quad \forall i \in I, \forall k \in K.
		\]
	\end{tcolorbox}
	
	\myremark
	Умножение определено только тогда, когда число столбцов первой матрицы равно числу строк второй (множество индексов \( J \) совпадает).
	
	\medskip
	\noindent \textbf{Свойства умножения матриц:}
	
	\begin{enumerate}
		\item \textbf{Ассоциативность.} Если \( A \in \mathbb{R}^{I \times J} \), \( B \in \mathbb{R}^{J \times K} \), \( C \in \mathbb{R}^{K \times L} \), то
		\[
		(AB)C = A(BC).
		\]
		
		\textit{Доказательство.} Обе части определены и являются матрицами размера \( I \times L \). Вычислим элемент с индексами \( (i, l) \):
		\[
		((AB)C)_{il} = \sum_{k \in K} (AB)_{ik} C_{kl} = \sum_{k \in K} \left( \sum_{j \in J} A_{ij} B_{jk} \right) C_{kl} = \sum_{j \in J} A_{ij} \left( \sum_{k \in K} B_{jk} C_{kl} \right) = \sum_{j \in J} A_{ij} (BC)_{jl} = (A(BC))_{il}.
		\]
		
		\item \textbf{Дистрибутивность.} Если определены соответствующие суммы и произведения, то
		\[
		(A + B)C = AC + BC, \quad C(A + B) = CA + CB.
		\]
		
		\item \textbf{Связь с умножением на скаляр.} Для любых \( \alpha \in \mathbb{R} \) и матриц \( A, B \) (подходящих размеров):
		\[
		(\alpha A)B = \alpha (AB) = A(\alpha B).
		\]
	\end{enumerate}
	
	\subsection{Единичная матрица}
	
	Пусть \( I \) — конечное множество. Определим символ Кронекера:
	\[
	\delta_{ij} = 
	\begin{cases}
		1, & i = j, \\
		0, & i \neq j.
	\end{cases}
	\]
	
	\begin{tcolorbox}[colback=yellow!5, colframe=orange!50!black, title=Определение]
		\textit{Единичной матрицей} порядка \( |I| \) называется матрица \( E_I \in \mathbb{R}^{I \times I} \) с элементами \( (E_I)_{ij} = \delta_{ij} \).
	\end{tcolorbox}
	
	Для квадратной матрицы порядка \( n \) единичная матрица имеет вид:
	\[
	E_n = 
	\begin{pmatrix}
		1 & 0 & \dots & 0 \\
		0 & 1 & \dots & 0 \\
		\vdots & \vdots & \ddots & \vdots \\
		0 & 0 & \dots & 1
	\end{pmatrix}.
	\]
	
	Основное свойство: для любой матрицы \( A \in \mathbb{R}^{I \times J} \)
	\[
	E_I A = A, \quad A E_J = A.
	\]
	
	\subsection{Обратная матрица}
	
	\begin{tcolorbox}[colback=blue!5, colframe=blue!40!black, title=Определение]
		Пусть \( A \in \mathbb{R}^{I \times I} \) — квадратная матрица. Матрица \( X \in \mathbb{R}^{I \times I} \) называется:
		\begin{itemize}
			\item \textit{левой обратной} к \( A \), если \( XA = E_I \);
			\item \textit{правой обратной} к \( A \), если \( AX = E_I \);
			\item \textit{обратной} к \( A \), если она является одновременно левой и правой обратной. Обозначение: \( A^{-1} \).
		\end{itemize}
	\end{tcolorbox}
	
	\begin{tcolorbox}[colback=green!5, colframe=green!50!black, title=Свойство]
		Если для матрицы \( A \) существуют левая обратная \( X \) и правая обратная \( Y \), то \( X = Y \), и эта матрица является единственной обратной к \( A \).
	\end{tcolorbox}
	
	\medskip
	\textit{Доказательство.}
	\[
	X = X E = X (A Y) = (X A) Y = E Y = Y.
	\]
	Таким образом, левая и правая обратные совпадают. Если бы существовали две разные обратные матрицы \( A_1^{-1} \) и \( A_2^{-1} \), то, взяв \( X = A_1^{-1} \) как левую, а \( Y = A_2^{-1} \) как правую, получили бы \( A_1^{-1} = A_2^{-1} \). \hfill \(\square\)
	
	\subsection{Транспонирование матриц}
	
	\begin{tcolorbox}[colback=red!5, colframe=red!50!black, title=Определение]
		Пусть \( A \in \mathbb{R}^{I \times J} \). \textit{Транспонированной матрицей} к \( A \) называется матрица \( A^T \in \mathbb{R}^{J \times I} \), определяемая условием:
		\[
		(A^T)_{ji} = A_{ij} \quad \forall i \in I, \forall j \in J.
		\]
	\end{tcolorbox}
	
	\medskip
	\noindent \textbf{Свойства транспонирования:}
	\begin{enumerate}
		\item \((A^T)^T = A\).
		\item \((A + B)^T = A^T + B^T\).
		\item \((\alpha A)^T = \alpha A^T\) (если кольцо коммутативно, то \(\alpha\) можно выносить).
		\item \((AB)^T = B^T A^T\).
	\end{enumerate}
	
	\textit{Доказательство свойства 4.}
	Пусть \( A \in \mathbb{R}^{I \times J} \), \( B \in \mathbb{R}^{J \times K} \). Тогда \( AB \in \mathbb{R}^{I \times K} \), и \( (AB)^T \in \mathbb{R}^{K \times I} \). С другой стороны, \( B^T \in \mathbb{R}^{K \times J} \), \( A^T \in \mathbb{R}^{J \times I} \), так что \( B^T A^T \in \mathbb{R}^{K \times I} \). Для любых \( k \in K \), \( i \in I \):
	\[
	((AB)^T)_{ki} = (AB)_{ik} = \sum_{j \in J} A_{ij} B_{jk} = \sum_{j \in J} (A^T)_{ji} (B^T)_{kj} = \sum_{j \in J} (B^T)_{kj} (A^T)_{ji} = (B^T A^T)_{ki}.
	\]
	(В предпоследнем равенстве использована коммутативность умножения в кольце \( \mathbb{R} \); если кольцо некоммутативно, формула остаётся верной, но порядок сомножителей важен: \((AB)^T = B^T A^T\) всегда, независимо от коммутативности). \hfill \(\square\)
	
	\subsection{Кольцо квадратных матриц}
	
	Множество квадратных матриц \( \mathbb{R}^{I \times I} \) с операциями сложения и умножения образует кольцо (вообще говоря, некоммутативное). Единицей в этом кольце служит единичная матрица \( E_I \), нулём — нулевая матрица. Отображение \( d \mapsto d E_I \) вкладывает кольцо \( \mathbb{R} \) в центр кольца матриц (если \( \mathbb{R} \) коммутативно)
	
	
	
	
	% ==================== РАЗДЕЛ 46: ИНВЕРСИИ. ЗНАК ПЕРЕСТАНОВКИ ====================
	\newpage
	\section{Инверсии. Определение знака перестановки через число инверсий (46 вопрос)}
	
	\subsection{Определение перестановки}
	
	Пусть \( I = \{1, 2, \dots, n\} \) — конечное множество (обычно рассматривают первые \( n \) натуральных чисел).
	
	\begin{tcolorbox}[colback=blue!5, colframe=blue!40!black, title=Определение]
		\textit{Перестановкой} множества \( I \) называется биективное отображение \( \sigma: I \to I \). Множество всех перестановок обозначается \( S_n \).
	\end{tcolorbox}
	
	Перестановку удобно записывать в виде двух строк:
	\[
	\sigma = \begin{pmatrix}
		1 & 2 & \dots & n \\
		\sigma(1) & \sigma(2) & \dots & \sigma(n)
	\end{pmatrix}.
	\]
	
	\subsection{Инверсии в перестановке}
	
	\begin{tcolorbox}[colback=green!5, colframe=green!50!black, title=Определение инверсии]
		Пусть \( \sigma \in S_n \). Пара \( (i, j) \) с \( 1 \leq i < j \leq n \) называется \textit{инверсией} в перестановке \( \sigma \), если \( \sigma(i) > \sigma(j) \).
	\end{tcolorbox}
	
	Иными словами, инверсия — это беспорядок: больший элемент стоит раньше меньшего.
	
	Число инверсий перестановки \( \sigma \) обозначается \( \operatorname{inv}(\sigma) \) или \( I(\sigma) \).
	
	\begin{tcolorbox}[colback=yellow!5, colframe=orange!50!black, title=Определение чётности]
		Перестановка называется \textit{чётной}, если число её инверсий чётно, и \textit{нечётной}, если число инверсий нечётно.
	\end{tcolorbox}
	
	\subsection{Транспозиции}
	
	\begin{tcolorbox}[colback=blue!5, colframe=blue!40!black, title=Определение]
		\textit{Транспозицией} называется перестановка, которая меняет местами два элемента, оставляя все остальные на месте. Если меняются элементы \( i \) и \( j \) (\( i \neq j \)), транспозиция обозначается \( (i \; j) \).
	\end{tcolorbox}
	
	Транспозиция является частным случаем цикла длины 2.
	
	\subsection{Изменение числа инверсий при элементарных преобразованиях}
	
	\begin{tcolorbox}[colback=red!5, colframe=red!50!black, title=Лемма 1 (Транспозиция соседних элементов)]
		При перестановке местами двух соседних элементов в перестановке число инверсий изменяется ровно на 1.
	\end{tcolorbox}
	
	\medskip
	\textit{Доказательство.}
	Рассмотрим перестановку \( (\dots, a, b, \dots) \), где \( a \) и \( b \) — два соседних элемента. Поменяем их местами: получим \( (\dots, b, a, \dots) \).
	
	Разобьём все пары элементов, участвующие в подсчёте инверсий, на три категории:
	\begin{itemize}
		\item Пары, не содержащие ни \( a \), ни \( b \). Число инверсий среди них не меняется.
		\item Пары, содержащие ровно один из элементов \( a \) или \( b \), и другой элемент, отличный от \( a, b \). Для каждого такого элемента \( x \) (стоящего слева или справа от пары) относительный порядок с \( a \) и \( b \) может измениться, но вклад в изменение числа инверсий от этих пар в сумме даёт ноль, так как если \( x \) образовывал инверсию с \( a \), то после перестановки он может образовывать инверсию с \( b \), и наоборот. Более аккуратное рассуждение: рассмотрим все элементы, расположенные левее пары; для них порядок с \( a \) и \( b \) не меняется, так как они остаются перед обоими. Для элементов правее пары аналогично — они остаются после обоих. Таким образом, единственная пара, которая может изменить свой статус, — это сама пара \( (a, b) \).
		\item Пара \( (a, b) \). До перестановки она могла быть инверсией (если \( a > b \)) или нет (если \( a < b \)). После перестановки статус меняется на противоположный.
	\end{itemize}
	
	Следовательно, общее число инверсий изменяется ровно на \( \pm 1 \). \hfill \(\square\)
	
	\begin{tcolorbox}[colback=red!5, colframe=red!50!black, title=Лемма 2 (Произвольная транспозиция)]
		Любая транспозиция может быть представлена в виде произведения нечётного числа транспозиций соседних элементов.
	\end{tcolorbox}
	
	\medskip
	\textit{Доказательство.}
	Рассмотрим транспозицию \( (c \; h) \), где \( c \) и \( h \) — два элемента, между которыми находится \( m \) других элементов: \( d, e, \dots, f, g \) (всего \( m \) элементов). Тогда
	\[
	(c \; h) = (c \; d)(c \; e)\dots(c \; f)(c \; g)(c \; h)(g \; h)(f \; h)\dots(e \; h)(d \; h).
	\]
	В этом произведении каждый множитель — транспозиция соседних элементов (после выполнения предыдущих шагов элементы сближаются). Число множителей равно \( (m+1) + m = 2m + 1 \), т.е. нечётно. \hfill \(\square\)
	
	\subsection{Знак перестановки}
	
	\begin{tcolorbox}[colback=purple!5, colframe=purple!50!black, title=Определение знака]
		\textit{Знаком} (или сигнатурой) перестановки \( \sigma \in S_n \) называется число
		\[
		\operatorname{sgn}(\sigma) = (-1)^{\operatorname{inv}(\sigma)}.
		\]
		Таким образом, \( \operatorname{sgn}(\sigma) = 1 \), если перестановка чётная, и \( \operatorname{sgn}(\sigma) = -1 \), если нечётная.
	\end{tcolorbox}
	
	\subsection{Свойства знака перестановки}
	
	\begin{enumerate}
		\item Знак произведения перестановок равен произведению знаков:
		\[
		\operatorname{sgn}(\sigma \tau) = \operatorname{sgn}(\sigma) \operatorname{sgn}(\tau).
		\]
		
		\item Знак транспозиции всегда равен \( -1 \).
		
		\item Знак обратной перестановки равен знаку исходной:
		\[
		\operatorname{sgn}(\sigma^{-1}) = \operatorname{sgn}(\sigma).
		\]
	\end{enumerate}
	
	\textit{Доказательство свойства 1 (идея).}
	Из леммы 2 любая перестановка может быть разложена в произведение транспозиций. При умножении на транспозицию чётность меняется, что соответствует умножению знака на \( -1 \). Следовательно, знак произведения совпадает с произведением знаков. \hfill \(\square\)
	
	\subsection{Пример}
	
	Рассмотрим перестановку
	\[
	\sigma = \begin{pmatrix}
		1 & 2 & 3 & 4 \\
		3 & 1 & 4 & 2
	\end{pmatrix}.
	\]
	Найдём все инверсии:
	\begin{itemize}
		\item \( i=1, j=2 \): \( \sigma(1)=3 > \sigma(2)=1 \) — инверсия.
		\item \( i=1, j=3 \): \( 3 < 4 \) — не инверсия.
		\item \( i=1, j=4 \): \( 3 > 2 \) — инверсия.
		\item \( i=2, j=3 \): \( 1 < 4 \) — не инверсия.
		\item \( i=2, j=4 \): \( 1 < 2 \) — не инверсия.
		\item \( i=3, j=4 \): \( 4 > 2 \) — инверсия.
	\end{itemize}
	Всего инверсий: 3 (пары (1,2), (1,4), (3,4)). Число нечётное, значит, перестановка нечётная, \( \operatorname{sgn}(\sigma) = -1 \).
	
	\subsection{Связь с определителем}
	
	Знак перестановки используется при определении определителя матрицы:
	\[
	\det A = \sum_{\sigma \in S_n} \operatorname{sgn}(\sigma) a_{1,\sigma(1)} a_{2,\sigma(2)} \dots a_{n,\sigma(n)}.
	\]
	
	\subsection{Замечание}
	
	Определение чётности через инверсии эквивалентно определению через разложение в произведение транспозиций. Оба подхода широко используются в комбинаторике и линейной алгебре
	
	
	
	
	% ==================== РАЗДЕЛ 47: ПОНЯТИЕ ОПРЕДЕЛИТЕЛЯ ====================
	\newpage
	\section{Понятие определителя. Определитель транспонированной матрицы (47 вопрос)}
	
	\subsection{Определение определителя}
	
	Пусть \( R \) — коммутативное кольцо (в частности, может быть полем \( \mathbb{R} \) или \( \mathbb{C} \)). Рассмотрим квадратную матрицу \( A = (a_{ij}) \in R^{n \times n} \) порядка \( n \).
	
	\begin{tcolorbox}[colback=blue!5, colframe=blue!40!black, title=Определение]
		\textit{Определителем} (или детерминантом) матрицы \( A \) называется элемент кольца \( R \), обозначаемый \( \det A \) или \( |A| \), определяемый формулой:
		\[
		\det A = \sum_{\sigma \in S_n} \operatorname{sgn}(\sigma) \, a_{1,\sigma(1)} a_{2,\sigma(2)} \dots a_{n,\sigma(n)},
		\]
		где \( S_n \) — группа всех перестановок множества \( \{1, 2, \dots, n\} \), а \( \operatorname{sgn}(\sigma) \) — знак перестановки (\( +1 \) для чётных, \( -1 \) для нечётных).
	\end{tcolorbox}
	
	Для матрицы, записанной в виде таблицы, определитель часто обозначают так:
	\[
	\det A = \begin{vmatrix}
		a_{11} & a_{12} & \dots & a_{1n} \\
		a_{21} & a_{22} & \dots & a_{2n} \\
		\vdots & \vdots & \ddots & \vdots \\
		a_{n1} & a_{n2} & \dots & a_{nn}
	\end{vmatrix}.
	\]
	
	\myremark
	Сумма в определении содержит \( n! \) слагаемых. Каждое слагаемое — это произведение \( n \) элементов матрицы, взятых по одному из каждой строки и каждого столбца. Знак слагаемого определяется чётностью перестановки \( \sigma \).
	
	\subsection{Эквивалентная форма определения}
	
	В приведённом определении множители в произведении расположены в порядке следования номеров строк: \( a_{1,\sigma(1)} a_{2,\sigma(2)} \dots a_{n,\sigma(n)} \). Однако, поскольку кольцо \( R \) коммутативно, порядок сомножителей не важен. Можно переставить их так, чтобы вторые индексы шли по порядку, а первые задавались перестановкой.
	
	Рассмотрим перестановку \( \sigma^{-1} \). Тогда
	\[
	a_{1,\sigma(1)} a_{2,\sigma(2)} \dots a_{n,\sigma(n)} = a_{\sigma^{-1}(1),1} a_{\sigma^{-1}(2),2} \dots a_{\sigma^{-1}(n),n}.
	\]
	Так как \( \operatorname{sgn}(\sigma) = \operatorname{sgn}(\sigma^{-1}) \), а отображение \( \sigma \mapsto \sigma^{-1} \) является биекцией на \( S_n \), получаем эквивалентное определение:
	
	\begin{tcolorbox}[colback=green!5, colframe=green!50!black, title=Эквивалентная форма]
		\[
		\det A = \sum_{\sigma \in S_n} \operatorname{sgn}(\sigma) \, a_{\sigma(1),1} a_{\sigma(2),2} \dots a_{\sigma(n),n}.
		\]
	\end{tcolorbox}
	
	В этой форме множители соответствуют выбору по одному элементу из каждого столбца, а перестановка указывает, из какой строки берётся элемент для данного столбца.
	
	\subsection{Определитель транспонированной матрицы}
	
	Напомним, что транспонированная матрица \( A^T \) получается из \( A \) заменой строк на столбцы: \( (A^T)_{ij} = A_{ji} \).
	
	\begin{tcolorbox}[colback=red!5, colframe=red!50!black, title=Теорема]
		Для любой квадратной матрицы \( A \) над коммутативным кольцом
		\[
		\det(A^T) = \det A.
		\]
	\end{tcolorbox}
	
	\medskip
	\textit{Доказательство.}
	Пусть \( A = (a_{ij}) \), тогда \( A^T = (a_{ji}) \). По определению,
	\[
	\det(A^T) = \sum_{\sigma \in S_n} \operatorname{sgn}(\sigma) \, (A^T)_{1,\sigma(1)} (A^T)_{2,\sigma(2)} \dots (A^T)_{n,\sigma(n)}.
	\]
	Но \( (A^T)_{i,\sigma(i)} = a_{\sigma(i), i} \). Следовательно,
	\[
	\det(A^T) = \sum_{\sigma \in S_n} \operatorname{sgn}(\sigma) \, a_{\sigma(1),1} a_{\sigma(2),2} \dots a_{\sigma(n),n}.
	\]
	Это в точности эквивалентная форма определителя, доказанная выше. Таким образом, \( \det(A^T) = \det A \). \hfill \(\square\)
	
	\myremark
	Это свойство означает, что определитель не меняется при транспонировании, т.е. строки и столбцы равноправны.
	
	\subsection{Примеры определителей малых порядков}
	
	\textbf{Порядок \( n = 1 \):}
	Матрица \( A = (a) \). Группа \( S_1 \) состоит из одной тождественной перестановки (знак \( +1 \)). Поэтому
	\[
	\det A = a.
	\]
	
	\textbf{Порядок \( n = 2 \):}
	Матрица \( A = \begin{pmatrix} a & b \\ c & d \end{pmatrix} \). Группа \( S_2 \) содержит две перестановки:
	\begin{itemize}
		\item Тождественная \( \sigma = (1)(2) \), \( \operatorname{sgn}=+1 \): произведение \( a_{11} a_{22} = ad \).
		\item Транспозиция \( \sigma = (1\;2) \), \( \operatorname{sgn}=-1 \): произведение \( a_{12} a_{21} = bc \).
	\end{itemize}
	Таким образом,
	\[
	\det A = ad - bc.
	\]
	
	\textbf{Порядок \( n = 3 \):}
	Группа \( S_3 \) содержит 6 перестановок. Можно вычислить определитель по правилу Саррюса (правило треугольников) или непосредственно по формуле. Результат:
	\[
	\begin{vmatrix}
		a_{11} & a_{12} & a_{13} \\
		a_{21} & a_{22} & a_{23} \\
		a_{31} & a_{32} & a_{33}
	\end{vmatrix}
	= a_{11}a_{22}a_{33} + a_{12}a_{23}a_{31} + a_{13}a_{21}a_{32}
	- a_{11}a_{23}a_{32} - a_{12}a_{21}a_{33} - a_{13}a_{22}a_{31}.
	\]
	
	\subsection{Замечания}
	
	\begin{itemize}
		\item Определитель является важнейшей характеристикой квадратной матрицы. Он используется при решении систем линейных уравнений (правило Крамера), при вычислении обратной матрицы, при исследовании линейной зависимости строк/столбцов и во многих других вопросах.
		\item Из свойства \( \det(A^T) = \det A \) следует, что все утверждения об определителе, сформулированные для строк, автоматически переносятся на столбцы.
	\end{itemize}
	
	
	
	
	% ==================== РАЗДЕЛ 48: СВОЙСТВА ОПРЕДЕЛИТЕЛЯ ====================
	\newpage
	\section{Свойства определителя (48 вопрос)}
	
	\subsection{Предварительные замечания}
	
	Пусть \( R \) — коммутативное кольцо (в частности, поле \( \mathbb{R} \) или \( \mathbb{C} \)). Для матрицы \( A = (a_{ij}) \in R^{n \times n} \) определитель \( \det A \) задаётся формулой
	\[
	\det A = \sum_{\sigma \in S_n} \operatorname{sgn}(\sigma) a_{1,\sigma(1)} a_{2,\sigma(2)} \dots a_{n,\sigma(n)}.
	\]
	
	В данном разделе мы рассмотрим основные свойства определителя, кроме свойства линейности по столбцу (свойство I), которое обычно выделяют отдельно. Свойства формулируются для столбцов; для строк они выполняются аналогично в силу равенства \( \det A^T = \det A \).
	
	\subsection{Свойство II (Однородность по столбцу)}
	
	\begin{tcolorbox}[colback=blue!5, colframe=blue!40!black, title=Свойство II]
		При умножении любого столбца матрицы на скаляр \( c \in R \) определитель умножается на этот скаляр. Схематически:
		\[
		\det(\dots, cX, \dots) = c \det(\dots, X, \dots),
		\]
		где \( X \) — исходный столбец.
	\end{tcolorbox}
	
	\medskip
	\textit{Доказательство.}
	Пусть матрица \( A' \) получена из \( A \) умножением \( k \)-го столбца на \( c \), т.е. \( a'_{ik} = c a_{ik} \) для всех \( i \), а остальные столбцы не изменены. Тогда
	\[
	\det A' = \sum_{\sigma \in S_n} \operatorname{sgn}(\sigma) a_{\sigma(1)1} \dots (c a_{\sigma(k)k}) \dots a_{\sigma(n)n}.
	\]
	Вынося общий множитель \( c \) за знак суммы (поскольку сумма конечна и умножение в кольце дистрибутивно), получаем
	\[
	\det A' = c \sum_{\sigma \in S_n} \operatorname{sgn}(\sigma) a_{\sigma(1)1} \dots a_{\sigma(k)k} \dots a_{\sigma(n)n} = c \det A.
	\]
	Для строки доказательство аналогично с использованием эквивалентной формы определителя (через перестановки строк). \hfill \(\square\)
	
	\myremark
	Из этого свойства следует, что общий множитель столбца можно выносить за знак определителя.
	
	\subsection{Свойство III (Определитель с равными столбцами)}
	
	\begin{tcolorbox}[colback=red!5, colframe=red!50!black, title=Свойство III]
		Если матрица содержит два одинаковых столбца (или две одинаковые строки), то её определитель равен нулю.
	\end{tcolorbox}
	
	\medskip
	\textit{Доказательство.}
	Пусть в матрице \( A \) столбцы с номерами \( k \) и \( l \) (\( k \neq l \)) совпадают: \( a_{ik} = a_{il} \) для всех \( i \). Рассмотрим транспозицию \( \tau \in S_n \), которая меняет местами \( k \) и \( l \):
	\[
	\tau(k) = l,\quad \tau(l) = k,\quad \tau(j) = j \text{ при } j \neq k,l.
	\]
	Разобьём группу \( S_n \) на два непересекающихся подмножества: множество чётных перестановок \( A_n \) и множество нечётных \( A_n \tau \) (так как умножение на транспозицию меняет чётность). Тогда
	\[
	\det A = \sum_{\sigma \in A_n} \operatorname{sgn}(\sigma) a_{\sigma(1)1} \dots a_{\sigma(n)n} + \sum_{\sigma \in A_n} \operatorname{sgn}(\sigma \tau) a_{\sigma\tau(1)1} \dots a_{\sigma\tau(n)n}.
	\]
	Во второй сумме \( \operatorname{sgn}(\sigma \tau) = -\operatorname{sgn}(\sigma) \). Заметим, что для любого \( \sigma \in A_n \) произведение \( a_{\sigma\tau(1)1} \dots a_{\sigma\tau(n)n} \) отличается от произведения \( a_{\sigma(1)1} \dots a_{\sigma(n)n} \) лишь тем, что множители, соответствующие столбцам \( k \) и \( l \), переставлены. Но поскольку эти столбцы равны, \( a_{\sigma\tau(k)k} = a_{\sigma(l)k} = a_{\sigma(l)l} \) и \( a_{\sigma\tau(l)l} = a_{\sigma(k)l} = a_{\sigma(k)k} \), то произведение не меняется. Таким образом, для каждого \( \sigma \in A_n \) соответствующие слагаемые из первой и второй сумм равны по модулю, но противоположны по знаку. Следовательно, вся сумма равна нулю. \hfill \(\square\)
	
	\myremark
	Для строк доказательство аналогично с использованием эквивалентной формы определителя.
	
	\subsection{Свойство IV (Определитель единичной матрицы)}
	
	\begin{tcolorbox}[colback=green!5, colframe=green!50!black, title=Свойство IV]
		Определитель единичной матрицы \( E_n \) равен 1:
		\[
		\det E_n = 1.
		\]
	\end{tcolorbox}
	
	\medskip
	\textit{Доказательство.}
	Единичная матрица \( E_n = (\delta_{ij}) \), где \( \delta_{ij} \) — символ Кронекера (\( \delta_{ij} = 1 \) при \( i = j \), и \( 0 \) при \( i \neq j \)). В сумме, определяющей определитель,
	\[
	\det E_n = \sum_{\sigma \in S_n} \operatorname{sgn}(\sigma) \delta_{1,\sigma(1)} \delta_{2,\sigma(2)} \dots \delta_{n,\sigma(n)},
	\]
	каждое слагаемое, соответствующее не тождественной перестановке \( \sigma \), содержит хотя бы один множитель \( \delta_{i,\sigma(i)} = 0 \) (поскольку при \( \sigma(i) \neq i \) этот множитель равен нулю). Единственное ненулевое слагаемое — при \( \sigma = \operatorname{id} \), для которого все множители равны 1 и \( \operatorname{sgn}(\operatorname{id}) = 1 \). Следовательно, сумма равна 1. \hfill \(\square\)
	
	\myremark
	Обобщение: определитель скалярной матрицы \( cE_n \) (т.е. матрицы с \( c \) на диагонали и нулями вне) равен \( c^n \). Доказательство аналогично, так как каждый диагональный элемент равен \( c \), и только тождественная перестановка даёт ненулевой вклад \( c^n \).
	
	\subsection{Следствия из основных свойств}
	
	\begin{tcolorbox}[colback=yellow!5, colframe=orange!50!black, title=Следствие 1]
		Определитель матрицы, содержащей нулевой столбец (нулевую строку), равен нулю.
	\end{tcolorbox}
	\textit{Доказательство.} Нулевой столбец можно рассматривать как произведение любого столбца на скаляр \( 0 \). По свойству II определитель такой матрицы равен \( 0 \) умноженный на определитель некоторой матрицы, т.е. \( 0 \). \hfill \(\square\)
	
	\begin{tcolorbox}[colback=yellow!5, colframe=orange!50!black, title=Следствие 2]
		Определитель матрицы, содержащей два пропорциональных столбца (строки), равен нулю.
	\end{tcolorbox}
	\textit{Доказательство.} Если один столбец равен другому, умноженному на скаляр \( c \), то по свойству II
	\[
	\det(\dots, X, \dots, cX, \dots) = c \det(\dots, X, \dots, X, \dots).
	\]
	Но определитель с двумя равными столбцами равен нулю по свойству III. Следовательно, исходный определитель равен нулю. \hfill \(\square\)
	
	\begin{tcolorbox}[colback=yellow!5, colframe=orange!50!black, title=Следствие 3]
		Определитель не меняется, если к одному столбцу прибавить другой столбец, умноженный на произвольный скаляр (аналогично для строк).
	\end{tcolorbox}
	\textit{Доказательство.}
	\[
	\det(\dots, X + cY, \dots, Y, \dots) = \det(\dots, X, \dots, Y, \dots) + \det(\dots, cY, \dots, Y, \dots)
	\]
	(по свойству I, линейности по столбцу). Второе слагаемое равно \( c \det(\dots, Y, \dots, Y, \dots) = c \cdot 0 = 0 \) по свойству III. Таким образом, определитель не изменился. \hfill \(\square\)
	
	\begin{tcolorbox}[colback=yellow!5, colframe=orange!50!black, title=Следствие 4]
		При перестановке двух столбцов (строк) определитель меняет знак.
	\end{tcolorbox}
	\textit{Доказательство.}
	Рассмотрим определитель \( \Delta = \det(\dots, X, \dots, Y, \dots) \) с двумя столбцами \( X \) и \( Y \). Тогда, используя свойство I и следствие 3, можно показать, что \( \det(\dots, Y, \dots, X, \dots) = -\Delta \). Стандартное доказательство:
	\[
	0 = \det(\dots, X+Y, \dots, X+Y, \dots) = \det(\dots, X, \dots, X+Y, \dots) + \det(\dots, Y, \dots, X+Y, \dots) = \det(\dots, X, \dots, X, \dots) + \det(\dots, X, \dots, Y, \dots) + \det(\dots, Y, \dots, X, \dots) + \det(\dots, Y, \dots, Y, \dots).
	\]
	Первое и последнее слагаемые равны нулю (равные столбцы). Остаётся \( \det(\dots, X, \dots, Y, \dots) + \det(\dots, Y, \dots, X, \dots) = 0 \), откуда и следует смена знака. \hfill \(\square\)
	
	\subsection{Замечание об определителе неквадратной матрицы}
	
	Иногда рассматривают прямоугольные матрицы размера \( m \times n \). Если \( m \neq n \), то определитель не определён. Однако можно рассматривать максимальные квадратные подматрицы и говорить об определителях таких подматриц (минорах). При разных способах выбора подматрицы определители могут различаться знаком, что иногда полезно (например, при изучении ранга матрицы)
	
	
	
	
	
	
	% ==================== РАЗДЕЛ 49: РАЗЛОЖЕНИЕ ОПРЕДЕЛИТЕЛЯ ПО СТРОКЕ (СТОЛБЦУ) ====================
	\newpage
	\section{Разложение определителя по строке (столбцу) (49 вопрос)}
	
	\subsection{Миноры и алгебраические дополнения}
	
	Пусть \( R \) — коммутативное кольцо (в частности, поле) и \( A = (a_{ij}) \in R^{n \times n} \) — квадратная матрица порядка \( n \).
	
	\begin{tcolorbox}[colback=blue!5, colframe=blue!40!black, title=Определения]
		\begin{itemize}
			\item \textit{Минором} \( M_{ij} \) элемента \( a_{ij} \) называется определитель матрицы порядка \( n-1 \), полученной из \( A \) вычёркиванием \( i \)-й строки и \( j \)-го столбца.
			\item \textit{Алгебраическим дополнением} элемента \( a_{ij} \) называется число
			\[
			A_{ij} = (-1)^{i+j} M_{ij}.
			\]
		\end{itemize}
	\end{tcolorbox}
	
	\myremark
	Знак \( (-1)^{i+j} \) часто называют \textit{знаковой функцией} или \textit{четностью позиции} \( (i,j) \).
	
	\subsection{Теорема о разложении определителя по строке}
	
	\begin{tcolorbox}[colback=green!5, colframe=green!50!black, title=Теорема (Разложение по i-й строке)]
		Для любой фиксированной строки \( i \) (\( 1 \leq i \leq n \)) определитель матрицы \( A \) равен сумме произведений элементов этой строки на их алгебраические дополнения:
		\[
		\det A = \sum_{j=1}^{n} a_{ij} A_{ij}.
		\]
	\end{tcolorbox}
	
	\medskip
	\textit{Доказательство.}
	Воспользуемся определением определителя через перестановки:
	\[
	\det A = \sum_{\sigma \in S_n} \operatorname{sgn}(\sigma) a_{1,\sigma(1)} a_{2,\sigma(2)} \dots a_{n,\sigma(n)}.
	\]
	
	Разобьём все перестановки на группы в зависимости от значения \( \sigma(i) \), т.е. от того, какой элемент \( i \)-й строки участвует в произведении. Для фиксированного \( j \) рассмотрим множество перестановок \( \sigma \), для которых \( \sigma(i) = j \). Каждая такая перестановка определяет произведение, содержащее множитель \( a_{ij} \). Оставшиеся \( n-1 \) множителей соответствуют выбору по одному элементу из остальных строк и столбцов (кроме \( i \)-й строки и \( j \)-го столбца).
	
	Суммируя по всем \( j \), получаем:
	\[
	\det A = \sum_{j=1}^{n} a_{ij} \left( \sum_{\substack{\sigma \in S_n \\ \sigma(i)=j}} \operatorname{sgn}(\sigma) \prod_{k \neq i} a_{k,\sigma(k)} \right).
	\]
	
	Покажем, что внутренняя сумма равна \( A_{ij} = (-1)^{i+j} M_{ij} \). Для этого заметим, что существует биекция между перестановками \( \sigma \) с \( \sigma(i)=j \) и перестановками \( \tau \) множества \( \{1,\dots,n\}\setminus\{i\} \) и \( \{1,\dots,n\}\setminus\{j\} \). При этой биекции знак \( \sigma \) связан со знаком \( \tau \) множителем \( (-1)^{i+j} \) (из-за необходимости переставить строки и столбцы, чтобы привести к стандартному порядку). Действительно, если вычеркнуть \( i \)-ю строку и \( j \)-й столбец и рассмотреть матрицу порядка \( n-1 \), то её определитель \( M_{ij} \) равен сумме по перестановкам \( \tau \) с соответствующими знаками. После учёта перестановок, возвращающих \( i \)-ю строку и \( j \)-й столбец на место, появляется множитель \( (-1)^{i+j} \). Более формально: каждая перестановка \( \sigma \) с \( \sigma(i)=j \) может быть получена из перестановки \( \tau \) (на \( n-1 \) элементе) путём вставки \( i \) на \( i \)-е место и \( j \) на \( j \)-е место с соответствующим циклическим сдвигом; знак при этом меняется в зависимости от чётности числа транспозиций, что даёт \( (-1)^{i+j} \).
	
	Таким образом,
	\[
	\sum_{\substack{\sigma \in S_n \\ \sigma(i)=j}} \operatorname{sgn}(\sigma) \prod_{k \neq i} a_{k,\sigma(k)} = (-1)^{i+j} M_{ij} = A_{ij}.
	\]
	
	Подставляя в исходную сумму, получаем требуемое. \hfill \(\square\)
	
	\subsection{Разложение по столбцу}
	
	Аналогичная теорема верна для разложения по столбцу. Её можно получить либо повторением рассуждения с перестановками, либо применить теорему о разложении по строке к транспонированной матрице и использовать \( \det A^T = \det A \).
	
	\begin{tcolorbox}[colback=green!5, colframe=green!50!black, title=Теорема (Разложение по j-му столбцу)]
		Для любой фиксированной строки \( j \) (\( 1 \leq j \leq n \))
		\[
		\det A = \sum_{i=1}^{n} a_{ij} A_{ij}.
		\]
	\end{tcolorbox}
	
	\subsection{Примеры}
	
	\textbf{Пример 1 (матрица 2×2).} 
	\[
	A = \begin{pmatrix} a & b \\ c & d \end{pmatrix}.
	\]
	Разложим по первой строке:
	\[
	\det A = a \cdot A_{11} + b \cdot A_{12}.
	\]
	Вычислим алгебраические дополнения:
	\[
	A_{11} = (-1)^{1+1} M_{11} = 1 \cdot d = d,
	\]
	\[
	A_{12} = (-1)^{1+2} M_{12} = -1 \cdot c = -c.
	\]
	Таким образом, \( \det A = a d + b (-c) = ad - bc \), что совпадает с известной формулой.
	
	\textbf{Пример 2 (матрица 3×3).}
	\[
	A = \begin{pmatrix}
		1 & 2 & 3 \\
		0 & 4 & 5 \\
		1 & 0 & 6
	\end{pmatrix}.
	\]
	Разложим по второй строке, так как в ней есть нуль (это упрощает вычисления):
	\[
	\det A = a_{21} A_{21} + a_{22} A_{22} + a_{23} A_{23} = 0 \cdot A_{21} + 4 \cdot A_{22} + 5 \cdot A_{23}.
	\]
	Вычисляем миноры:
	\[
	M_{22} = \begin{vmatrix} 1 & 3 \\ 1 & 6 \end{vmatrix} = 1\cdot6 - 3\cdot1 = 3,
	\quad A_{22} = (-1)^{2+2} \cdot 3 = 3.
	\]
	\[
	M_{23} = \begin{vmatrix} 1 & 2 \\ 1 & 0 \end{vmatrix} = 1\cdot0 - 2\cdot1 = -2,
	\quad A_{23} = (-1)^{2+3} \cdot (-2) = -1 \cdot (-2) = 2.
	\]
	Таким образом,
	\[
	\det A = 4 \cdot 3 + 5 \cdot 2 = 12 + 10 = 22.
	\]
	
	Можно проверить прямым вычислением по правилу Саррюса: \( 1\cdot4\cdot6 + 2\cdot5\cdot1 + 3\cdot0\cdot0 - 3\cdot4\cdot1 - 2\cdot0\cdot6 - 1\cdot5\cdot0 = 24 + 10 + 0 - 12 - 0 - 0 = 22 \).
	
	\subsection{Следствия из теоремы о разложении}
	
	\begin{enumerate}
		\item \textbf{Определитель треугольной матрицы} равен произведению диагональных элементов. Действительно, разлагая по первому столбцу (для нижнетреугольной) или по первой строке (для верхнетреугольной), получаем, что все слагаемые, кроме диагонального, содержат нулевой множитель.
		
		\item \textbf{Определитель блочно-треугольной матрицы}. Если матрица имеет вид
		\[
		\begin{pmatrix}
			A & B \\
			0 & D
		\end{pmatrix},
		\]
		где \( A \) и \( D \) — квадратные блоки, то \( \det = \det A \cdot \det D \). Это можно доказать индукцией, разлагая по строкам, содержащим нули.
		
		\item \textbf{Формула для обратной матрицы}. Если \( \det A \neq 0 \), то
		\[
		A^{-1} = \frac{1}{\det A} (A_{ji})^T,
		\]
		т.е. элемент обратной матрицы \( (A^{-1})_{ij} = \frac{A_{ji}}{\det A} \). Это следует из того, что сумма произведений элементов строки на алгебраические дополнения другой строки равна нулю.
	\end{enumerate}
	
	\subsection{Замечание о вычислительной эффективности}
	
	Разложение определителя по строке или столбцу даёт рекуррентный способ вычисления, но его непосредственное применение требует \( O(n!) \) операций, что неприемлемо для больших \( n \). Однако для матриц с большим количеством нулей (разреженных) этот метод может быть эффективен. На практике обычно используют метод Гаусса (приведение к треугольному виду), который требует \( O(n^3) \) операций.
	
	
	
	% ==================== РАЗДЕЛ 50: ОПРЕДЕЛИТЕЛЬ СТУПЕНЧАТОЙ МАТРИЦЫ. ОПРЕДЕЛИТЕЛЬ ПРОИЗВЕДЕНИЯ ====================
	\newpage
	\section{Определитель ступенчатой матрицы. Определитель произведения матриц (50 вопрос)}
	
	\subsection{Определитель треугольной (ступенчатой) матрицы}
	
	\begin{tcolorbox}[colback=blue!5, colframe=blue!40!black, title=Теорема]
		Определитель верхней треугольной матрицы (т.е. матрицы, у которой все элементы ниже главной диагонали равны нулю) равен произведению её диагональных элементов. То же верно для нижней треугольной матрицы.
	\end{tcolorbox}
	
	\medskip
	\textit{Доказательство.}
	Пусть \( A = (a_{ij}) \) — верхняя треугольная матрица порядка \( n \), т.е. \( a_{ij} = 0 \) при \( i > j \). Вычислим её определитель по определению:
	\[
	\det A = \sum_{\sigma \in S_n} \operatorname{sgn}(\sigma) a_{1,\sigma(1)} a_{2,\sigma(2)} \dots a_{n,\sigma(n)}.
	\]
	Рассмотрим произвольную перестановку \( \sigma \), отличную от тождественной. Для неё существует хотя бы один индекс \( k \) такой, что \( \sigma(k) < k \) (иначе, если \( \sigma(i) \geq i \) для всех \( i \), то в силу биективности \( \sigma \) должна быть тождественной). Но если \( \sigma(k) < k \), то в произведении присутствует множитель \( a_{k,\sigma(k)} \), который равен нулю, так как \( k > \sigma(k) \) (элемент ниже главной диагонали). Следовательно, все слагаемые, кроме соответствующего тождественной перестановке, равны нулю.
	
	Для тождественной перестановки \( \sigma(i) = i \), знак \( \operatorname{sgn}(\operatorname{id}) = 1 \), и произведение равно \( a_{11} a_{22} \dots a_{nn} \). Таким образом,
	\[
	\det A = a_{11} a_{22} \dots a_{nn}.
	\]
	
	Для нижней треугольной матрицы доказательство аналогично (рассматриваем транспонированную матрицу или проводим симметричное рассуждение). \hfill \(\square\)
	
	\myremark
	Это свойство часто используется при вычислении определителей методом приведения к треугольному виду (метод Гаусса).
	
	\subsection{Поведение определителя при элементарных преобразованиях}
	
	Напомним, что существуют три типа элементарных преобразований строк (столбцов) матрицы:
	\begin{enumerate}
		\item Перестановка двух строк (столбцов).
		\item Умножение строки (столбца) на ненулевой скаляр.
		\item Прибавление к одной строке (столбцу) другой строки (столбца), умноженной на скаляр.
	\end{enumerate}
	
	Им соответствуют \textit{элементарные матрицы} — матрицы, полученные из единичной применением соответствующего преобразования.
	
	Из свойств определителя (см. вопрос 48) следует:
	\begin{itemize}
		\item При перестановке двух строк (столбцов) определитель меняет знак (свойство 4).
		\item При умножении строки (столбца) на скаляр \( c \) определитель умножается на \( c \) (свойство II).
		\item При прибавлении к строке (столбцу) другой строки (столбца), умноженной на скаляр, определитель не меняется (следствие 3).
	\end{itemize}
	
	Таким образом, для любой элементарной матрицы \( E \) и любой матрицы \( A \) подходящего размера
	\[
	\det(EA) = \det E \cdot \det A, \quad \det(AE) = \det A \cdot \det E,
	\]
	где
	\[
	\det E = 
	\begin{cases}
		-1, & \text{если } E \text{ — матрица перестановки двух строк}, \\
		c, & \text{если } E \text{ получена умножением строки на } c, \\
		1, & \text{если } E \text{ получена прибавлением строки}.
	\end{cases}
	\]
	
	\subsection{Определитель произведения матриц над полем}
	
	\begin{tcolorbox}[colback=green!5, colframe=green!50!black, title=Теорема]
		Для любых квадратных матриц \( A, B \in K^{n \times n} \) над полем \( K \) выполняется
		\[
		\det(AB) = \det A \cdot \det B.
		\]
	\end{tcolorbox}
	
	\medskip
	\textit{Доказательство.}
	Используем приведение матрицы \( A \) к каноническому виду с помощью элементарных преобразований. Из линейной алгебры известно, что любую матрицу можно представить в виде
	\[
	A = U_1 U_2 \dots U_r \begin{pmatrix} E_k & 0 \\ 0 & 0 \end{pmatrix} V_1 V_2 \dots V_s,
	\]
	где \( U_i, V_j \) — элементарные матрицы, а \( E_k \) — единичная матрица порядка \( k = \operatorname{rank} A \). Если матрица невырождена (ранг \( n \)), то средний блок равен \( E_n \), и \( A \) есть произведение элементарных матриц. Если матрица вырождена, то средний блок содержит нулевую строку.
	
	Рассмотрим два случая.
	
	\textit{Случай 1:} \( \det A = 0 \) (матрица вырождена). Тогда в каноническом виде \( \begin{pmatrix} E_k & 0 \\ 0 & 0 \end{pmatrix} \) имеет нулевую строку. При умножении на любую матрицу справа нулевая строка останется нулевой, поэтому матрица \( A' = \begin{pmatrix} E_k & 0 \\ 0 & 0 \end{pmatrix} V_1 \dots V_s B \) также содержит нулевую строку, и её определитель равен 0. Тогда \( \det(AB) = 0 \), так как умножение на элементарные матрицы слева не меняет нулевого определителя. Следовательно, \( \det(AB) = 0 = 0 \cdot \det B \).
	
	\textit{Случай 2:} \( \det A \neq 0 \) (матрица невырождена). Тогда \( A \) представима в виде произведения элементарных матриц:
	\[
	A = E_1 E_2 \dots E_m.
	\]
	Тогда
	\[
	AB = E_1 E_2 \dots E_m B.
	\]
	Применяя последовательно свойство мультипликативности для элементарных матриц (доказанное выше), получаем
	\[
	\det(AB) = \det E_1 \det E_2 \dots \det E_m \det B = \det A \cdot \det B.
	\]
	\hfill \(\square\)
	
	\subsection{Обобщение на матрицы над коммутативным кольцом}
	
	\begin{tcolorbox}[colback=yellow!5, colframe=orange!50!black, title=Теорема (Общий случай)]
		Для любых квадратных матриц \( A, B \) порядка \( n \) над коммутативным кольцом \( R \) с единицей
		\[
		\det(AB) = \det A \cdot \det B.
		\]
	\end{tcolorbox}
	
	\medskip
	\textit{Замечание.}
	Прямое доказательство этой теоремы над кольцом довольно громоздко, так как элементарные преобразования над кольцом не всегда позволяют привести матрицу к диагональному виду (из-за отсутствия деления). Однако результат остаётся верным. Один из способов доказательства — использовать принцип продолжения алгебраических тождеств: рассмотреть матрицы с неопределёнными элементами \( x_{ij}, y_{ij} \) как элементы кольца многочленов \( \mathbb{Z}[x_{ij}, y_{ij}] \), доказать тождество для матриц над этим кольцом (где можно делить на ненулевые многочлены, работая в поле частных), а затем применить специализацию, подставляя конкретные элементы из \( R \). \hfill \(\square\)
	
	\subsection{Следствие для обратимых матриц}
	
	\begin{tcolorbox}[colback=red!5, colframe=red!50!black, title=Критерий обратимости]
		Матрица \( A \in R^{n \times n} \) над коммутативным кольцом с единицей обратима (т.е. существует \( A^{-1} \)) тогда и только тогда, когда её определитель \( \det A \) обратим в \( R \).
	\end{tcolorbox}
	
	\medskip
	\textit{Доказательство.}
	Если \( A \) обратима, то \( AA^{-1} = E \). Применяя теорему о произведении, получаем
	\[
	\det A \cdot \det A^{-1} = \det E = 1,
	\]
	следовательно, \( \det A \) обратим (обратным служит \( \det A^{-1} \)).
	
	Обратно, если \( \det A \) обратим, то можно построить обратную матрицу по формуле
	\[
	A^{-1} = \frac{1}{\det A} \operatorname{adj}(A),
	\]
	где \( \operatorname{adj}(A) \) — присоединённая матрица (транспонированная матрица алгебраических дополнений). Проверка показывает, что \( A \cdot \operatorname{adj}(A) = \operatorname{adj}(A) \cdot A = (\det A) E \). Деление на обратимый элемент \( \det A \) даёт обратную матрицу. \hfill \(\square\)
	
	\myremark
	В частности, над полем условие обратимости равносильно \( \det A \neq 0 \), так как в поле любой ненулевой элемент обратим.
	
	\subsection{Пример}
	
	Рассмотрим матрицы над полем:
	\[
	A = \begin{pmatrix} 1 & 2 \\ 3 & 4 \end{pmatrix}, \quad B = \begin{pmatrix} 0 & 1 \\ 1 & 0 \end{pmatrix}.
	\]
	\[
	\det A = 1\cdot4 - 2\cdot3 = 4 - 6 = -2, \quad \det B = 0\cdot0 - 1\cdot1 = -1.
	\]
	\[
	AB = \begin{pmatrix} 1 & 2 \\ 3 & 4 \end{pmatrix} \begin{pmatrix} 0 & 1 \\ 1 & 0 \end{pmatrix} = \begin{pmatrix} 2 & 1 \\ 4 & 3 \end{pmatrix}, \quad \det(AB) = 2\cdot3 - 1\cdot4 = 6 - 4 = 2.
	\]
	Действительно, \( \det A \cdot \det B = (-2) \cdot (-1) = 2 \), что совпадает с \( \det(AB) \).
	
	
	
	
	% ==================== РАЗДЕЛ 51: ЭЛЕМЕНТАРНЫЕ ПРЕОБРАЗОВАНИЯ. ОПРЕДЕЛИТЕЛЬ ВАНДЕРМОНДА ====================
	\newpage
	\section{Элементарные преобразования. Определитель Вандермонда (51 вопрос)}
	
	\subsection{Элементарные преобразования и определитель}
	
	Напомним основные типы элементарных преобразований строк (столбцов) матрицы и их влияние на определитель.
	
	\begin{tcolorbox}[colback=blue!5, colframe=blue!40!black, title=Элементарные преобразования]
		\begin{enumerate}
			\item \textbf{Перестановка двух строк (столбцов).} При такой операции определитель меняет знак.
			\item \textbf{Умножение строки (столбца) на ненулевой скаляр \(c\).} Определитель умножается на \(c\).
			\item \textbf{Прибавление к одной строке (столбцу) другой строки (столбца), умноженной на произвольный скаляр.} Определитель не меняется.
		\end{enumerate}
	\end{tcolorbox}
	
	Этим преобразованиям соответствуют \textit{элементарные матрицы} — матрицы, полученные из единичной применением одного такого преобразования. Определители элементарных матриц равны соответственно \(-1\), \(c\) и \(1\).
	
	Благодаря этим свойствам, элементарные преобразования широко используются для упрощения вычисления определителей, в частности, для приведения матрицы к треугольному виду.
	
	\subsection{Определитель Вандермонда}
	
	Особый интерес представляет определитель матрицы Вандермонда, который часто возникает в задачах интерполяции, алгебры и комбинаторики.
	
	\begin{tcolorbox}[colback=green!5, colframe=green!50!black, title=Определение]
		\textit{Определителем Вандермонда} порядка \(n\) называется определитель матрицы
		\[
		V(x_1, x_2, \dots, x_n) = 
		\begin{vmatrix}
			1 & x_1 & x_1^2 & \dots & x_1^{n-1} \\
			1 & x_2 & x_2^2 & \dots & x_2^{n-1} \\
			1 & x_3 & x_3^2 & \dots & x_3^{n-1} \\
			\vdots & \vdots & \vdots & \ddots & \vdots \\
			1 & x_n & x_n^2 & \dots & x_n^{n-1}
		\end{vmatrix},
		\]
		где \(x_1, x_2, \dots, x_n\) — элементы некоторого поля (или коммутативного кольца).
	\end{tcolorbox}
	
	\subsection{Вычисление определителя Вандермонда}
	
	\begin{tcolorbox}[colback=red!5, colframe=red!50!black, title=Теорема (Формула Вандермонда)]
		\[
		V(x_1, x_2, \dots, x_n) = \prod_{1 \leq i < j \leq n} (x_j - x_i).
		\]
	\end{tcolorbox}
	
	\medskip
	\textit{Доказательство.}
	Докажем формулу индукцией по \(n\).
	
	\textit{База:} \(n = 1\). Определитель матрицы \(1 \times 1\) равен 1, и произведение по пустому множеству считается равным 1. Формула верна.
	
	\textit{Индукционный переход:} Предположим, что формула верна для порядка \(n-1\). Рассмотрим определитель \(V(x_1, \dots, x_n)\). Выполним следующие элементарные преобразования столбцов, не меняющие определитель (тип 3): для каждого \(k\) от \(n-1\) до 1 вычтем из \(k+1\)-го столбца \(k\)-й столбец, умноженный на \(x_1\). Точнее, произведём последовательно:
	\[
	C_{n} \leftarrow C_{n} - x_1 C_{n-1},\quad C_{n-1} \leftarrow C_{n-1} - x_1 C_{n-2},\quad \dots,\quad C_2 \leftarrow C_2 - x_1 C_1.
	\]
	После этих преобразований первая строка примет вид \((1, 0, 0, \dots, 0)\). Действительно, элемент \(j\)-го столбца (\(j \geq 2\)) в первой строке после вычитаний станет \(x_1^{j-1} - x_1 \cdot x_1^{j-2} = 0\).
	
	Разложим полученный определитель по первой строке. Единственное ненулевое слагаемое соответствует первому столбцу, и мы получаем:
	\[
	V(x_1, \dots, x_n) = 1 \cdot (-1)^{1+1} \cdot \Delta,
	\]
	где \(\Delta\) — определитель матрицы порядка \(n-1\), полученной вычёркиванием первой строки и первого столбца. Эта матрица имеет вид:
	\[
	\Delta = 
	\begin{vmatrix}
		1 & x_2 - x_1 & x_2^2 - x_1 x_2 & \dots & x_2^{n-1} - x_1 x_2^{n-2} \\
		1 & x_3 - x_1 & x_3^2 - x_1 x_3 & \dots & x_3^{n-1} - x_1 x_3^{n-2} \\
		\vdots & \vdots & \vdots & \ddots & \vdots \\
		1 & x_n - x_1 & x_n^2 - x_1 x_n & \dots & x_n^{n-1} - x_1 x_n^{n-2}
	\end{vmatrix}.
	\]
	
	Заметим, что в каждой строке, начиная со второй, можно вынести общий множитель \((x_i - x_1)\):
	\[
	x_i^{k} - x_1 x_i^{k-1} = x_i^{k-1}(x_i - x_1).
	\]
	Таким образом, для \(i \geq 2\) и \(j \geq 2\) элемент \(a_{ij} = x_i^{j-2}(x_i - x_1)\). Вынесем из каждой строки множитель \((x_i - x_1)\) (это преобразование типа 2, умножающее определитель на произведение этих множителей):
	\[
	\Delta = \left( \prod_{i=2}^n (x_i - x_1) \right) \cdot 
	\begin{vmatrix}
		1 & x_2 & x_2^2 & \dots & x_2^{n-2} \\
		1 & x_3 & x_3^2 & \dots & x_3^{n-2} \\
		\vdots & \vdots & \vdots & \ddots & \vdots \\
		1 & x_n & x_n^2 & \dots & x_n^{n-2}
	\end{vmatrix}.
	\]
	
	Полученный определитель — это в точности определитель Вандермонда порядка \(n-1\) от переменных \(x_2, x_3, \dots, x_n\). По предположению индукции он равен
	\[
	\prod_{2 \leq i < j \leq n} (x_j - x_i).
	\]
	
	Следовательно,
	\[
	V(x_1, \dots, x_n) = \left( \prod_{i=2}^n (x_i - x_1) \right) \cdot \prod_{2 \leq i < j \leq n} (x_j - x_i) = \prod_{1 \leq i < j \leq n} (x_j - x_i).
	\]
	\hfill \(\square\)
	
	\subsection{Следствия и применения}
	
	\begin{enumerate}
		\item \textbf{Критерий равенства элементов.} Определитель Вандермонда равен нулю тогда и только тогда, когда среди \(x_1, \dots, x_n\) есть равные. Это следует из формулы: произведение содержит множитель \((x_j - x_i)\), который обращается в ноль при \(x_i = x_j\).
		
		\item \textbf{Интерполяция.} При построении интерполяционного многочлена Лагранжа система линейных уравнений для коэффициентов имеет матрицу Вандермонда. Её невырожденность при различных узлах гарантирует существование и единственность решения.
		
		\item \textbf{Линейная независимость функций.} Функции \(1, t, t^2, \dots, t^{n-1}\) линейно независимы над любым полем, что также следует из невырожденности матрицы Вандермонда при различных значениях аргумента.
	\end{enumerate}
	
	\subsection{Пример}
	
	Вычислим определитель Вандермонда для \(x_1 = 1, x_2 = 2, x_3 = 3\):
	\[
	V(1,2,3) = 
	\begin{vmatrix}
		1 & 1 & 1 \\
		1 & 2 & 4 \\
		1 & 3 & 9
	\end{vmatrix}
	= (2-1)(3-1)(3-2) = 1 \cdot 2 \cdot 1 = 2.
	\]
	
	Проверим прямым вычислением:
	\[
	\begin{vmatrix}
		1 & 1 & 1 \\
		1 & 2 & 4 \\
		1 & 3 & 9
	\end{vmatrix}
	= 1\cdot2\cdot9 + 1\cdot4\cdot1 + 1\cdot1\cdot3 - 1\cdot2\cdot1 - 1\cdot4\cdot3 - 1\cdot1\cdot9 = 18 + 4 + 3 - 2 - 12 - 9 = 2.
	\]
	
	\subsection{Замечание}
	
	Определитель Вандермонда является важным примером того, как элементарные преобразования позволяют эффективно вычислять определители специального вида. Существуют также обобщения определителя Вандермонда (например, для случая, когда строки содержат не обязательно последовательные степени, но разные функции).
	
	
	
	
	% ==================== РАЗДЕЛ 52: ОБРАТНАЯ МАТРИЦА ====================
	\newpage
	\section{Обратная матрица (52 вопрос)}
	
	\subsection{Определение взаимной (присоединённой) матрицы}
	
	Пусть \( R \) — коммутативное кольцо с единицей (в частности, поле). Рассмотрим квадратную матрицу \( A = (a_{ij}) \in R^{n \times n} \). Напомним, что алгебраическим дополнением элемента \( a_{ij} \) называется
	\[
	A_{ij} = (-1)^{i+j} M_{ij},
	\]
	где \( M_{ij} \) — минор (определитель матрицы порядка \( n-1 \), полученной вычёркиванием \( i \)-й строки и \( j \)-го столбца).
	
	\begin{tcolorbox}[colback=blue!5, colframe=blue!40!black, title=Определение]
		\textit{Взаимной} (или \textit{присоединённой}) матрицей к матрице \( A \) называется матрица
		\[
		\tilde{A} = (\tilde{a}_{ij}) \in R^{n \times n}, \quad \text{где } \tilde{a}_{ij} = A_{ji}.
		\]
		Иными словами, \( \tilde{A} \) есть транспонированная матрица алгебраических дополнений:
		\[
		\tilde{A} = 
		\begin{pmatrix}
			A_{11} & A_{21} & \dots & A_{n1} \\
			A_{12} & A_{22} & \dots & A_{n2} \\
			\vdots & \vdots & \ddots & \vdots \\
			A_{1n} & A_{2n} & \dots & A_{nn}
		\end{pmatrix}.
		\]
	\end{tcolorbox}
	
	\myremark
	В литературе присоединённую матрицу часто обозначают \( \operatorname{adj}(A) \) (от англ. \textit{adjugate matrix}).
	
	\subsection{Основное свойство взаимной матрицы}
	
	\begin{tcolorbox}[colback=green!5, colframe=green!50!black, title=Теорема]
		Для любой квадратной матрицы \( A \in R^{n \times n} \) выполняется
		\[
		A \tilde{A} = \tilde{A} A = (\det A) \cdot E_n,
		\]
		где \( E_n \) — единичная матрица порядка \( n \).
	\end{tcolorbox}
	
	\medskip
	\textit{Доказательство.}
	Вычислим элемент произведения \( A \tilde{A} \) с индексами \( (i, j) \):
	\[
	(A \tilde{A})_{ij} = \sum_{k=1}^{n} a_{ik} \tilde{a}_{kj} = \sum_{k=1}^{n} a_{ik} A_{jk}.
	\]
	Рассмотрим два случая.
	
	\textit{Случай 1:} \( i = j \). Тогда
	\[
	(A \tilde{A})_{ii} = \sum_{k=1}^{n} a_{ik} A_{ik}.
	\]
	Это в точности разложение определителя матрицы \( A \) по \( i \)-й строке (см. теорему Лапласа, частный случай \( s = 1 \)). Следовательно, \( (A \tilde{A})_{ii} = \det A \).
	
	\textit{Случай 2:} \( i \neq j \). Тогда сумма
	\[
	\sum_{k=1}^{n} a_{ik} A_{jk}
	\]
	представляет собой сумму произведений элементов \( i \)-й строки на алгебраические дополнения элементов \( j \)-й строки. Такая сумма равна нулю, так как это определитель матрицы, у которой \( i \)-я и \( j \)-я строки совпадают (а именно, матрицы, полученной из \( A \) заменой \( j \)-й строки на \( i \)-ю). По свойству III определителя (определитель с двумя одинаковыми строками равен нулю) эта сумма равна нулю.
	
	Таким образом, все диагональные элементы произведения равны \( \det A \), а недиагональные равны нулю, т.е. \( A \tilde{A} = (\det A) E_n \).
	
	Равенство \( \tilde{A} A = (\det A) E_n \) доказывается аналогично (разложение по столбцам или применение к транспонированной матрице). \hfill \(\square\)
	
	\subsection{Обратная матрица}
	
	Из основного свойства немедленно следует конструкция обратной матрицы, если определитель обратим в кольце.
	
	\begin{tcolorbox}[colback=red!5, colframe=red!50!black, title=Следствие (Формула обратной матрицы)]
		Если \( \det A \) обратим в \( R \) (т.е. существует элемент \( (\det A)^{-1} \in R \)), то матрица \( A \) обратима, и
		\[
		A^{-1} = \frac{1}{\det A} \, \tilde{A}.
		\]
	\end{tcolorbox}
	
	\medskip
	\textit{Доказательство.}
	По условию существует \( (\det A)^{-1} \). Умножим равенство \( A \tilde{A} = (\det A) E \) на \( (\det A)^{-1} \) слева:
	\[
	\left( \frac{1}{\det A} \tilde{A} \right) A = E.
	\]
	Аналогично, из \( \tilde{A} A = (\det A) E \) получаем
	\[
	A \left( \frac{1}{\det A} \tilde{A} \right) = E.
	\]
	Таким образом, матрица \( \frac{1}{\det A} \tilde{A} \) является обратной к \( A \). \hfill \(\square\)
	
	\subsection{Критерий обратимости над коммутативным кольцом}
	
	\begin{tcolorbox}[colback=yellow!5, colframe=orange!50!black, title=Теорема (Критерий обратимости)]
		Квадратная матрица \( A \in R^{n \times n} \) над коммутативным кольцом с единицей обратима (слева, справа или двусторонне) тогда и только тогда, когда её определитель \( \det A \) обратим в \( R \).
	\end{tcolorbox}
	
	\medskip
	\textit{Доказательство.}
	Если \( A \) обратима, то существует \( A^{-1} \) с \( A A^{-1} = E \). Тогда
	\[
	\det A \cdot \det A^{-1} = \det E = 1,
	\]
	следовательно, \( \det A \) обратим (обратным служит \( \det A^{-1} \)).
	
	Обратно, если \( \det A \) обратим, то по следствию выше матрица \( \frac{1}{\det A} \tilde{A} \) является обратной к \( A \). \hfill \(\square\)
	
	\myremark
	Из этого критерия, в частности, следует, что для матриц над полем обратимость равносильна \( \det A \neq 0 \) (так как в поле любой ненулевой элемент обратим). Матрицы с ненулевым определителем над полем называют \textit{невырожденными}.
	
	\subsection{Пример}
	
	Найдём обратную матрицу для
	\[
	A = \begin{pmatrix} 1 & 2 \\ 3 & 4 \end{pmatrix}
	\]
	над полем \( \mathbb{R} \).
	
	Вычислим определитель: \( \det A = 1\cdot4 - 2\cdot3 = 4 - 6 = -2 \neq 0 \), значит, матрица обратима.
	
	Найдём алгебраические дополнения:
	\[
	A_{11} = (-1)^{1+1} \cdot 4 = 4, \quad A_{12} = (-1)^{1+2} \cdot 3 = -3,
	\]
	\[
	A_{21} = (-1)^{2+1} \cdot 2 = -2, \quad A_{22} = (-1)^{2+2} \cdot 1 = 1.
	\]
	
	Присоединённая матрица:
	\[
	\tilde{A} = \begin{pmatrix} A_{11} & A_{21} \\ A_{12} & A_{22} \end{pmatrix} = \begin{pmatrix} 4 & -2 \\ -3 & 1 \end{pmatrix}.
	\]
	
	Тогда
	\[
	A^{-1} = \frac{1}{\det A} \tilde{A} = \frac{1}{-2} \begin{pmatrix} 4 & -2 \\ -3 & 1 \end{pmatrix} = \begin{pmatrix} -2 & 1 \\ 1.5 & -0.5 \end{pmatrix}.
	\]
	
	Проверка:
	\[
	A A^{-1} = \begin{pmatrix} 1 & 2 \\ 3 & 4 \end{pmatrix} \begin{pmatrix} -2 & 1 \\ 1.5 & -0.5 \end{pmatrix} = \begin{pmatrix} -2+3 & 1-1 \\ -6+6 & 3-2 \end{pmatrix} = \begin{pmatrix} 1 & 0 \\ 0 & 1 \end{pmatrix}.
	\]
	
	\subsection{Замечание о существовании обратной матрицы}
	
	Для матриц над кольцом, не являющимся полем, определитель может быть отличен от нуля, но не иметь обратного элемента. Например, над кольцом целых чисел \( \mathbb{Z} \) матрица \( \begin{pmatrix} 2 & 0 \\ 0 & 2 \end{pmatrix} \) имеет определитель 4, который не обратим в \( \mathbb{Z} \), поэтому матрица не обратима над \( \mathbb{Z} \) (хотя она обратима над \( \mathbb{Q} \)). Обратимость над \( \mathbb{Z} \) означает, что матрица имеет целочисленную обратную, что возможно только если \( \det A = \pm 1 \).
	
	
	
	
	
	% ==================== РАЗДЕЛ 53: ФОРМУЛЫ КРАМЕРА ====================
	\newpage
	\section{Формулы Крамера (53 вопрос)}
	
	\subsection{Постановка задачи}
	
	Рассмотрим систему линейных уравнений с квадратной матрицей:
	\[
	\begin{cases}
		a_{11}x_1 + a_{12}x_2 + \dots + a_{1n}x_n = b_1, \\
		a_{21}x_1 + a_{22}x_2 + \dots + a_{2n}x_n = b_2, \\
		\vdots \\
		a_{n1}x_1 + a_{n2}x_2 + \dots + a_{nn}x_n = b_n,
	\end{cases}
	\]
	где \(a_{ij}, b_i\) — элементы некоторого поля \(K\) (или коммутативного кольца с единицей, в котором определитель матрицы системы обратим). В матричном виде:
	\[
	A X = B,
	\]
	где \(A = (a_{ij}) \in K^{n \times n}\), \(X = (x_1, \dots, x_n)^T\), \(B = (b_1, \dots, b_n)^T\).
	
	\subsection{Теорема Крамера}
	
	\begin{tcolorbox}[colback=blue!5, colframe=blue!40!black, title=Теорема (Формулы Крамера)]
		Если матрица \(A\) обратима (т.е. \(\det A \neq 0\) над полем или \(\det A\) обратим над кольцом), то система \(AX = B\) имеет единственное решение, которое выражается формулами:
		\[
		x_j = \frac{\det A_j}{\det A}, \quad j = 1, 2, \dots, n,
		\]
		где \(A_j\) — матрица, полученная из \(A\) заменой \(j\)-го столбца на столбец свободных членов \(B\).
	\end{tcolorbox}
	
	\medskip
	\textit{Доказательство.}
	Так как \(A\) обратима, решение существует и единственно: \(X = A^{-1}B\). Воспользуемся выражением обратной матрицы через присоединённую (см. вопрос 52):
	\[
	A^{-1} = \frac{1}{\det A} \tilde{A},
	\]
	где \(\tilde{A} = (A_{ji})\) — матрица алгебраических дополнений (транспонированная). Тогда
	\[
	x_j = \sum_{k=1}^n (A^{-1})_{jk} b_k = \frac{1}{\det A} \sum_{k=1}^n A_{kj} b_k.
	\]
	
	Рассмотрим определитель матрицы \(A_j\). Разложим его по \(j\)-му столбцу. В матрице \(A_j\) на месте \(j\)-го столбца стоит столбец \(B\), а остальные столбцы такие же, как у \(A\). При разложении по \(j\)-му столбцу:
	\[
	\det A_j = \sum_{k=1}^n b_k \cdot (\text{алгебраическое дополнение элемента } a_{kj} \text{ в матрице } A_j).
	\]
	Но алгебраическое дополнение элемента, стоящего в \(k\)-й строке и \(j\)-м столбце матрицы \(A_j\), совпадает с алгебраическим дополнением элемента \(a_{kj}\) в исходной матрице \(A\) (поскольку вычёркивание \(k\)-й строки и \(j\)-го столбца даёт одну и ту же подматрицу). Это дополнение есть \((-1)^{k+j} M_{kj} = A_{kj}\). Таким образом,
	\[
	\det A_j = \sum_{k=1}^n b_k A_{kj}.
	\]
	
	Сравнивая с выражением для \(x_j\), получаем:
	\[
	x_j = \frac{1}{\det A} \det A_j.
	\]
	\hfill \(\square\)
	
	\subsection{Пример}
	
	Решим систему над полем \(\mathbb{R}\):
	\[
	\begin{cases}
		2x + 3y = 8, \\
		x - 2y = -3.
	\end{cases}
	\]
	
	Здесь
	\[
	A = \begin{pmatrix} 2 & 3 \\ 1 & -2 \end{pmatrix}, \quad B = \begin{pmatrix} 8 \\ -3 \end{pmatrix}.
	\]
	
	Вычислим определитель основной матрицы:
	\[
	\det A = 2 \cdot (-2) - 3 \cdot 1 = -4 - 3 = -7 \neq 0.
	\]
	
	Матрица \(A_1\) (заменяем первый столбец на \(B\)):
	\[
	A_1 = \begin{pmatrix} 8 & 3 \\ -3 & -2 \end{pmatrix}, \quad \det A_1 = 8 \cdot (-2) - 3 \cdot (-3) = -16 + 9 = -7.
	\]
	
	Матрица \(A_2\) (заменяем второй столбец на \(B\)):
	\[
	A_2 = \begin{pmatrix} 2 & 8 \\ 1 & -3 \end{pmatrix}, \quad \det A_2 = 2 \cdot (-3) - 8 \cdot 1 = -6 - 8 = -14.
	\]
	
	По формулам Крамера:
	\[
	x = \frac{\det A_1}{\det A} = \frac{-7}{-7} = 1, \quad y = \frac{\det A_2}{\det A} = \frac{-14}{-7} = 2.
	\]
	
	Ответ: \(x = 1, y = 2\).
	
	\subsection{Случай, когда определитель системы равен нулю}
	
	Если \(\det A = 0\) (над полем), то формулы Крамера неприменимы. Возможны два случая:
	\begin{itemize}
		\item Система несовместна (решений нет).
		\item Система имеет бесконечно много решений (если столбец \(B\) линейно выражается через столбцы \(A\), но ранг меньше \(n\)).
	\end{itemize}
	В этих случаях требуется исследование методом Гаусса или с помощью ранга матрицы.
	
	\subsection{Обобщение на системы с квадратной матрицей над кольцом}
	
	Если коэффициенты \(a_{ij}\) и свободные члены \(b_i\) лежат в коммутативном кольце \(R\) с единицей, и \(\det A\) обратим в \(R\), то формулы Крамера дают решение, причём \(x_j = (\det A)^{-1} \det A_j\) также лежат в \(R\). Это следует из того, что \(\det A_j\) выражается как линейная комбинация \(b_k\) с коэффициентами — алгебраическими дополнениями, которые лежат в \(R\). Таким образом, решение существует и единственно в \(R\).
	
	\myremark
	Если \(\det A\) не обратим, то система может не иметь решений в \(R\) или иметь несколько решений (в зависимости от конкретного вида).
	
	\subsection{Вывод формул Крамера через обратную матрицу}
	
	Другой способ вывода - обратную матрицу. Поскольку \(X = A^{-1}B\), а \(A^{-1} = \frac{1}{\det A} \tilde{A}\), то
	\[
	x_j = \frac{1}{\det A} \sum_{k=1}^n \tilde{A}_{jk} b_k = \frac{1}{\det A} \sum_{k=1}^n A_{kj} b_k,
	\]
	что совпадает с полученным выше выражением.
	
	
	
	
	% ==================== РАЗДЕЛ 54: ВЕКТОРНОЕ ПРОСТРАНСТВО ====================
	\newpage
	\section{Векторное пространство. Следствия из аксиом. Примеры (54 вопрос)}
	
	\subsection{Определение векторного пространства}
	
	Пусть \( K \) — поле (элементы поля будем называть \textit{скалярами}).
	
	\begin{tcolorbox}[colback=blue!5, colframe=blue!40!black, title=Определение]
		\textit{Векторным} (или \textit{линейным}) \textit{пространством} над полем \( K \) называется множество \( L \), на котором заданы две операции:
		\begin{itemize}
			\item \textit{Сложение векторов:} каждой паре \( x, y \in L \) сопоставляется элемент \( x + y \in L \);
			\item \textit{Умножение вектора на скаляр:} каждому \( \lambda \in K \) и \( x \in L \) сопоставляется элемент \( \lambda x \in L \).
		\end{itemize}
		Эти операции удовлетворяют следующим аксиомам (для любых \( x, y, z \in L \) и любых \( \lambda, \mu \in K \)):
		\begin{enumerate}
			\item \( (x + y) + z = x + (y + z) \) (ассоциативность сложения);
			\item \( x + y = y + x \) (коммутативность сложения);
			\item существует элемент \( 0 \in L \) такой, что \( x + 0 = x \) для всех \( x \in L \) (существование нулевого вектора);
			\item для каждого \( x \in L \) существует элемент \( -x \in L \) такой, что \( x + (-x) = 0 \) (существование противоположного вектора);
			\item \( \lambda (\mu x) = (\lambda \mu) x \) (ассоциативность умножения на скаляр);
			\item \( \lambda (x + y) = \lambda x + \lambda y \) (дистрибутивность умножения относительно сложения векторов);
			\item \( (\lambda + \mu) x = \lambda x + \mu x \) (дистрибутивность умножения относительно сложения скаляров);
			\item \( 1 \cdot x = x \), где \( 1 \) — единица поля \( K \).
		\end{enumerate}
	\end{tcolorbox}
	
	\myremark
	Аксиомы 1–4 означают, что \( L \) является абелевой группой по сложению. Эту группу называют \textit{аддитивной группой} пространства \( L \).
	
	\subsection{Простейшие следствия из аксиом}
	
	Из аксиом векторного пространства вытекают следующие свойства, которыми мы будем пользоваться в дальнейшем.
	
	\begin{tcolorbox}[colback=green!5, colframe=green!50!black, title=Следствие 1 (Умножение на нулевой скаляр)]
		Для любого вектора \( x \in L \) выполнено \( 0 \cdot x = 0 \).
	\end{tcolorbox}
	
	\medskip
	\textit{Доказательство.}
	Используем дистрибутивность (аксиома 7):
	\[
	0 \cdot x = (0 + 0) \cdot x = 0 \cdot x + 0 \cdot x.
	\]
	Прибавим к обеим частям вектор \( -(0 \cdot x) \):
	\[
	0 \cdot x = 0.
	\]
	\hfill \(\square\)
	
	\begin{tcolorbox}[colback=green!5, colframe=green!50!black, title=Следствие 2 (Умножение скаляра на нулевой вектор)]
		Для любого скаляра \( \lambda \in K \) выполнено \( \lambda \cdot 0 = 0 \).
	\end{tcolorbox}
	
	\medskip
	\textit{Доказательство.}
	Используем дистрибутивность (аксиома 6):
	\[
	\lambda \cdot 0 = \lambda (0 + 0) = \lambda \cdot 0 + \lambda \cdot 0.
	\]
	Прибавим \( -(\lambda \cdot 0) \) к обеим частям:
	\[
	\lambda \cdot 0 = 0.
	\]
	\hfill \(\square\)
	
	\begin{tcolorbox}[colback=green!5, colframe=green!50!black, title=Следствие 3 (Свойство противоположного вектора)]
		Для любого вектора \( x \in L \) выполнено \( (-1) \cdot x = -x \).
	\end{tcolorbox}
	
	\medskip
	\textit{Доказательство.}
	Используем дистрибутивность и следствие 1:
	\[
	x + (-1) \cdot x = 1 \cdot x + (-1) \cdot x = (1 + (-1)) \cdot x = 0 \cdot x = 0.
	\]
	Таким образом, вектор \( (-1) \cdot x \) является противоположным к \( x \), а по аксиоме 4 противоположный вектор единствен. Следовательно, \( (-1) \cdot x = -x \). \hfill \(\square\)
	
	\begin{tcolorbox}[colback=green!5, colframe=green!50!black, title=Следствие 4 (Умножение на скаляр равно нулю)]
		Если \( \lambda x = 0 \), то либо \( \lambda = 0 \), либо \( x = 0 \).
	\end{tcolorbox}
	
	\medskip
	\textit{Доказательство.}
	Предположим, что \( \lambda \neq 0 \). Тогда в поле \( K \) существует обратный элемент \( \lambda^{-1} \). Умножив обе части равенства \( \lambda x = 0 \) на \( \lambda^{-1} \), получим:
	\[
	\lambda^{-1} (\lambda x) = (\lambda^{-1} \lambda) x = 1 \cdot x = x = \lambda^{-1} \cdot 0 = 0,
	\]
	где последнее равенство следует из следствия 2. Таким образом, \( x = 0 \). \hfill \(\square\)
	
	\subsection{Примеры векторных пространств}
	
	\begin{enumerate}
		\item \textbf{Геометрические векторы.} Множество всех векторов плоскости (или трёхмерного пространства) с обычными операциями сложения по правилу параллелограмма и умножения на вещественное число образует векторное пространство над полем \( \mathbb{R} \). Именно эти пространства послужили прототипом общего понятия.
		
		\item \textbf{Арифметическое пространство \( K^n \).} Множество всех упорядоченных наборов \( (x_1, x_2, \dots, x_n) \) элементов поля \( K \) с покомпонентным сложением и умножением на скаляр:
		\[
		(x_1, \dots, x_n) + (y_1, \dots, y_n) = (x_1 + y_1, \dots, x_n + y_n),
		\]
		\[
		\lambda (x_1, \dots, x_n) = (\lambda x_1, \dots, \lambda x_n).
		\]
		Проверка аксиом сводится к свойствам операций в поле.
		
		\item \textbf{Пространство многочленов.} Множество \( K[t]_{\leq n} \) всех многочленов степени не выше \( n \) с коэффициентами из поля \( K \) (включая нулевой многочлен) с обычным сложением многочленов и умножением на скаляр. Аксиомы выполняются в силу свойств многочленов.
		
		\item \textbf{Пространство функций.} Множество всех функций из произвольного множества \( X \) в поле \( K \) с поточечными операциями:
		\[
		(f + g)(x) = f(x) + g(x), \quad (\lambda f)(x) = \lambda f(x).
		\]
		Проверка аксиом тривиальна.
		
		\item \textbf{Нулевое пространство.} Множество \( \{0\} \), состоящее из одного нулевого вектора, с операциями \( 0+0 = 0 \) и \( \lambda \cdot 0 = 0 \) является векторным пространством над любым полем \( K \).
		
		\item \textbf{Поле как пространство над своим подполем.} Если \( K \) — поле, а \( F \) — его подполе, то \( K \) можно рассматривать как векторное пространство над \( F \) (умножение элементов \( K \) на скаляры из \( F \) — это просто умножение в поле). Например, \( \mathbb{C} \) — векторное пространство над \( \mathbb{R} \), а \( \mathbb{R} \) — векторное пространство над \( \mathbb{Q} \).
	\end{enumerate}
	
	\subsection{Замечания о записи скаляров}
	
	В определении 19.1 скаляры записаны слева от вектора. Иногда удобна правая запись: \( x \lambda \). При этом аксиомы 5–7 принимают вид:
	\[
	(x \mu) \lambda = x (\mu \lambda), \quad (x + y) \lambda = x \lambda + y \lambda, \quad x (\lambda + \mu) = x \lambda + x \mu,
	\]
	а аксиома 8: \( x \cdot 1 = x \). В случае поля это эквивалентно левой записи благодаря коммутативности умножения в поле. Для некоммутативных тел различают левые и правые векторные пространства.
	
	
	\subsection{Дополнительные примеры векторных пространств}
	
	Продолжим список примеров, начатый в предыдущем разделе.
	
	\begin{enumerate}
		\setcounter{enumi}{6}
		\item \textbf{Пространство семейств \(K^I\).} Пусть \(K\) — поле, \(I\) — произвольное множество. Рассмотрим множество \(K^I\) всех семейств \((\alpha_i)_{i \in I}\) элементов поля \(K\) (т.е. функций \(I \to K\)). Определим сложение и умножение на скаляр покомпонентно:
		\[
		(\alpha_i) + (\beta_i) = (\gamma_i), \quad \text{где } \gamma_i = \alpha_i + \beta_i,
		\]
		\[
		\lambda (\alpha_i) = (\delta_i), \quad \text{где } \delta_i = \lambda \alpha_i.
		\]
		Аксиомы линейного пространства легко проверяются. Нулевым элементом служит семейство, все члены которого равны нулю.
		
		В частности, при \(I = \emptyset\) получаем нулевое пространство: \(K^\emptyset = \{\emptyset\} = \{0\}\).
		
		\item \textbf{Пространство формальных линейных комбинаций \(K^{(I)}\).} Рассмотрим подмножество \(K^{(I)} \subset K^I\), состоящее из семейств, у которых лишь конечное число компонент отлично от нуля (такие семейства называют \textit{финитными}). Те же операции, что и в предыдущем примере, корректно определены на этом подмножестве (сумма двух финитных семейств финитна, произведение финитного семейства на скаляр финитно). Полученное пространство называется \textit{пространством формальных линейных комбинаций} множества \(I\) над полем \(K\). Оно будет обозначаться \(\langle I \rangle_K\) или просто \(\langle I \rangle\).
		
		Если множество \(I\) конечно, то \(K^{(I)} = K^I\), и разница между примерами 6 и 7 исчезает. В частности, арифметическое пространство \(K^n\) есть частный случай как пространства \(K^I\), так и пространства формальных линейных комбинаций.
		
		\item \textbf{Пространство матриц \(K^{I \times J}\).} Пусть \(I, J\) — произвольные множества. Множество \(K^{I \times J}\) всех матриц размера \(I \times J\) (т.е. функций \(I \times J \to K\)) с обычными операциями сложения матриц и умножения на скаляр (покомпонентными) образует линейное пространство над \(K\). В частности, множество \(K^{m \times n}\) всех прямоугольных матриц фиксированного размера является линейным пространством.
		
		Можно рассматривать также подпространства:
		\begin{itemize}
			\item матрицы, у которых в каждой строке лишь конечное число ненулевых элементов (строки финитны);
			\item матрицы, у которых в каждом столбце лишь конечное число ненулевых элементов;
			\item матрицы, у которых лишь конечное число элементов отлично от нуля (такие матрицы образуют пространство формальных линейных комбинаций множества \(I \times J\)).
		\end{itemize}
	\end{enumerate}
	
	\subsection{Свойства действий в векторном пространстве}
	
	Из аксиом векторного пространства вытекают следующие важные свойства. В дальнейшем \(L\) обозначает линейное пространство над полем \(K\), символ \(0\) используется как для нулевого вектора, так и для нуля поля \(K\) (смысл ясен из контекста), строчные латинские буквы обозначают векторы, греческие — скаляры.
	
	\begin{tcolorbox}[colback=blue!5, colframe=blue!40!black, title=Свойство 1]
		\(\alpha a = 0\) тогда и только тогда, когда \(\alpha = 0\) или \(a = 0\).
	\end{tcolorbox}
	
	\medskip
	\textit{Доказательство.}
	\begin{itemize}
		\item[\(\Leftarrow\)] Если \(\alpha = 0\), то \(0 \cdot a = (0+0)a = 0a + 0a\), откуда \(0a = 0\). Если \(a = 0\), то \(\alpha \cdot 0 = \alpha(0+0) = \alpha 0 + \alpha 0\), откуда \(\alpha 0 = 0\).
		\item[\(\Rightarrow\)] Пусть \(\alpha a = 0\) и \(\alpha \neq 0\). Тогда в поле \(K\) существует \(\alpha^{-1}\). Умножая обе части на \(\alpha^{-1}\), получаем:
		\[
		a = 1 \cdot a = (\alpha^{-1} \alpha) a = \alpha^{-1} (\alpha a) = \alpha^{-1} \cdot 0 = 0.
		\]
		Следовательно, \(a = 0\). \hfill \(\square\)
	\end{itemize}
	
	\begin{tcolorbox}[colback=blue!5, colframe=blue!40!black, title=Свойство 2]
		\(\forall a \in L, \forall \alpha \in K\) выполнено:
		\[
		\alpha(-a) = -(\alpha a) = (-\alpha)a.
		\]
		В частности, \((-1)a = -a\).
	\end{tcolorbox}
	
	\medskip
	\textit{Доказательство.}
	\[
	0 = 0 \cdot a = (\alpha + (-\alpha))a = \alpha a + (-\alpha)a \quad \Rightarrow \quad (-\alpha)a = -(\alpha a).
	\]
	Аналогично,
	\[
	0 = \alpha \cdot 0 = \alpha(a + (-a)) = \alpha a + \alpha(-a) \quad \Rightarrow \quad \alpha(-a) = -(\alpha a).
	\]
	Таким образом, оба выражения равны \(-(\alpha a)\). Подставляя \(\alpha = 1\), получаем \((-1)a = -a\). \hfill \(\square\)
	
	\begin{tcolorbox}[colback=blue!5, colframe=blue!40!black, title=Свойство 3 (Дистрибутивность относительно вычитания)]
		Для любых \(a, b \in L\) и любых \(\alpha, \beta \in K\):
		\[
		\alpha(a - b) = \alpha a - \alpha b, \quad (\alpha - \beta)a = \alpha a - \beta a.
		\]
	\end{tcolorbox}
	
	\medskip
	\textit{Доказательство.}
	\[
	\alpha(a - b) = \alpha(a + (-b)) = \alpha a + \alpha(-b) = \alpha a + (-(\alpha b)) = \alpha a - \alpha b.
	\]
	Второе равенство:
	\[
	(\alpha - \beta)a = (\alpha + (-\beta))a = \alpha a + (-\beta)a = \alpha a + (-(\beta a)) = \alpha a - \beta a.
	\]
	\hfill \(\square\)
	
	\begin{tcolorbox}[colback=blue!5, colframe=blue!40!black, title=Свойство 4 (Линейность конечных сумм)]
		Для любого конечного семейства векторов \((a_i)_{i=1}^n\) и любого скаляра \(\alpha\):
		\[
		\alpha \sum_{i=1}^n a_i = \sum_{i=1}^n \alpha a_i.
		\]
		Для любого конечного семейства скаляров \((\alpha_i)_{i=1}^n\) и любого вектора \(a\):
		\[
		\left( \sum_{i=1}^n \alpha_i \right) a = \sum_{i=1}^n \alpha_i a.
		\]
	\end{tcolorbox}
	
	\medskip
	\textit{Доказательство.}
	Оба утверждения легко доказываются индукцией по \(n\) с использованием аксиом дистрибутивности. \hfill \(\square\)
	
	\myremark
	Свойство 4 обобщается на любые финитные семейства (т.е. семейства, в которых лишь конечное число членов отлично от нуля): сумма по такому семейству определена, и линейность сохраняется.
	
	\subsection{Единственность нейтральных и противоположных элементов}
	
	Как уже отмечалось, из аксиом абелевой группы следует, что нулевой вектор единствен, и для каждого вектора \(a\) противоположный вектор \(-a\) также единствен. Более того, для любых \(a, b \in L\) определена разность \(a - b = a + (-b)\), и операция вычитания всюду определена.
	
	
	
	
	% ==================== РАЗДЕЛ 55: ЛИНЕЙНАЯ ЗАВИСИМОСТЬ И НЕЗАВИСИМОСТЬ ====================
	\newpage
	\section{Линейная зависимость и независимость векторов. Свойства (55 вопрос)}
	
	\subsection{Линейные комбинации}
	
	Пусть \( L \) — линейное пространство над полем \( K \). Рассмотрим семейство векторов \( a = (a_i)_{i \in I} \), где \( I \) — некоторое множество индексов.
	
	\begin{tcolorbox}[colback=blue!5, colframe=blue!40!black, title=Определение]
		\textit{Линейной комбинацией} семейства \( a \) называется выражение вида
		\[
		\sum_{i \in I} \alpha_i a_i,
		\]
		где \( (\alpha_i)_{i \in I} \) — семейство скаляров из \( K \), причём лишь конечное число коэффициентов \( \alpha_i \) отлично от нуля (такие семейства называют \textit{финитными}).
	\end{tcolorbox}
	
	Линейная комбинация называется \textit{тривиальной}, если все коэффициенты \( \alpha_i \) равны нулю. В противном случае она называется \textit{нетривиальной}.
	
	Если вектор \( b \in L \) может быть представлен как линейная комбинация семейства \( a \), то говорят, что \( b \) \textit{линейно выражается} через векторы \( a_i \).
	
	Множество всех линейных комбинаций семейства \( a \) обозначается \( \langle a \rangle \) (или \( \operatorname{span}\{a_i\} \)).
	
	\subsection{Свойства линейных комбинаций}
	
	\begin{tcolorbox}[colback=green!5, colframe=green!50!black, title=Свойство 1]
		Линейная комбинация подсемейства семейства \( a \) является линейной комбинацией всего семейства \( a \).
	\end{tcolorbox}
	
	\medskip
	\textit{Доказательство.}
	Пусть \( J \subset I \) и дана линейная комбинация \( \sum_{j \in J} \alpha_j a_j \). Доопределим коэффициенты для всех \( i \in I \setminus J \) как нулевые. Тогда полученная сумма \( \sum_{i \in I} \beta_i a_i \) (где \( \beta_i = \alpha_i \) при \( i \in J \) и \( \beta_i = 0 \) при \( i \notin J \)) равна исходной, но теперь это линейная комбинация всего семейства. \hfill \(\square\)
	
	\begin{tcolorbox}[colback=green!5, colframe=green!50!black, title=Свойство 2]
		Линейная комбинация семейства линейных комбинаций семейства \( a \) есть линейная комбинация семейства \( a \).
	\end{tcolorbox}
	
	\medskip
	\textit{Доказательство.}
	Пусть \( b_1, \dots, b_m \) — векторы, каждый из которых линейно выражается через \( a \):
	\[
	b_j = \sum_{i \in I} \alpha_{ji} a_i, \quad j = 1, \dots, m,
	\]
	где для каждого \( j \) лишь конечное число коэффициентов \( \alpha_{ji} \) отлично от нуля. Рассмотрим произвольную линейную комбинацию векторов \( b_j \):
	\[
	c = \sum_{j=1}^m \beta_j b_j = \sum_{j=1}^m \beta_j \left( \sum_{i \in I} \alpha_{ji} a_i \right) = \sum_{i \in I} \left( \sum_{j=1}^m \beta_j \alpha_{ji} \right) a_i.
	\]
	Внутренняя сумма \( \gamma_i = \sum_{j=1}^m \beta_j \alpha_{ji} \) является конечной суммой для каждого \( i \), и лишь для конечного числа индексов \( i \) она может быть ненулевой (так как каждый \( b_j \) выражается через конечное число \( a_i \), а \( m \) конечно). Таким образом, \( c \) есть линейная комбинация семейства \( a \). \hfill \(\square\)
	
	\subsection{Матричная запись линейных комбинаций}
	
	Рассмотрим семейство \( a = (a_i)_{i \in I} \) как строку, т.е. матрицу размера \( 1 \times I \) с элементами из \( L \):
	\[
	a = \begin{pmatrix} a_i \end{pmatrix}_{i \in I}.
	\]
	Пусть \( X = (\xi_i)_{i \in I} \) — столбец скаляров (матрица размера \( I \times 1 \)), причём лишь конечное число его элементов отлично от нуля. Тогда произведение \( aX \) определено и равно одночленной матрице (размера \( 1 \times 1 \)), единственный элемент которой есть
	\[
	\sum_{i \in I} a_i \xi_i.
	\]
	Одночленную матрицу можно отождествить с её единственным элементом — вектором из \( L \). Таким образом, линейная комбинация семейства \( a \) записывается как произведение строки \( a \) на столбец коэффициентов \( X \).
	
	Если теперь \( b = (b_j)_{j \in J} \) — семейство (строка) векторов, каждый из которых линейно выражается через \( a \), то для каждого \( j \) существует столбец коэффициентов \( X_j \) такой, что \( b_j = a X_j \). Составим из этих столбцов матрицу \( B \) размера \( I \times J \): \( B = (X_1 | X_2 | \dots | X_J) \). Тогда
	\[
	b = a B,
	\]
	т.е. строка \( b \) получается умножением строки \( a \) на матрицу \( B \) справа. Это представление часто используется в линейной алгебре.
	
	\subsection{Линейная зависимость и независимость}
	
	\begin{tcolorbox}[colback=red!5, colframe=red!50!black, title=Определение]
		Семейство векторов \( (a_i)_{i \in I} \) называется \textit{линейно зависимым}, если существует его нетривиальная линейная комбинация, равная нулевому вектору. В противном случае семейство называется \textit{линейно независимым}.
	\end{tcolorbox}
	
	Иными словами, линейная независимость означает, что из равенства
	\[
	\sum_{i \in I} \alpha_i a_i = 0
	\]
	следует, что все коэффициенты \( \alpha_i = 0 \).
	
	\myremark
	Пустое семейство по определению считается линейно независимым. Свойство линейной зависимости (независимости) сохраняется при переиндексации семейства.
	
	В матричной форме: семейство (строка) \( a \) линейно зависимо, если существует ненулевой столбец \( X \) над \( K \) такой, что \( aX = 0 \). Линейная независимость означает, что из \( aX = 0 \) следует \( X = 0 \).
	
	\subsection{Однозначность разложения по линейно независимому семейству}
	
	Если семейство \( a \) линейно независимо и \( aX = aY \), то \( X = Y \). Действительно,
	\[
	aX - aY = a(X - Y) = 0 \Rightarrow X - Y = 0.
	\]
	Таким образом, линейная комбинация линейно независимого семейства однозначно определяет свои коэффициенты.
	
	\subsection{Простейшие свойства линейной зависимости и независимости}
	
	\begin{tcolorbox}[colback=yellow!5, colframe=orange!50!black, title=Свойство 1]
		Семейство \( (a_i)_{i \in I} \) линейно зависимо тогда и только тогда, когда один из его векторов линейно выражается через остальные.
	\end{tcolorbox}
	
	\medskip
	\textit{Доказательство.}
	\begin{itemize}
		\item[\(\Rightarrow\)] Пусть семейство линейно зависимо. Тогда существует нетривиальная линейная комбинация \( \sum_{i \in I} \alpha_i a_i = 0 \). Выберем индекс \( j \) такой, что \( \alpha_j \neq 0 \). Тогда
		\[
		a_j = \sum_{i \neq j} \left( -\frac{\alpha_i}{\alpha_j} \right) a_i,
		\]
		т.е. \( a_j \) выражается через остальные.
		\item[\(\Leftarrow\)] Если для некоторого \( j \) имеем \( a_j = \sum_{i \neq j} \beta_i a_i \), то
		\[
		\sum_{i \neq j} \beta_i a_i - a_j = 0,
		\]
		и это нетривиальная линейная комбинация (коэффициент при \( a_j \) равен \(-1 \neq 0\)). \hfill \(\square\)
	\end{itemize}
	
	\begin{tcolorbox}[colback=yellow!5, colframe=orange!50!black, title=Свойство 2]
		Одноэлементное семейство \( (a) \) линейно зависимо тогда и только тогда, когда \( a = 0 \).
	\end{tcolorbox}
	
	\medskip
	\textit{Доказательство.}
	Если \( a = 0 \), то \( 1 \cdot a = 0 \) — нетривиальная комбинация. Если \( a \neq 0 \) и \( \alpha a = 0 \), то из свойства 19.4(1) следует \( \alpha = 0 \). \hfill \(\square\)
	
	\begin{tcolorbox}[colback=yellow!5, colframe=orange!50!black, title=Свойство 3]
		Всякое надсемейство линейно зависимого семейства линейно зависимо.
	\end{tcolorbox}
	
	\medskip
	\textit{Доказательство.}
	Нетривиальная линейная комбинация исходного семейства, равная нулю, является также линейной комбинацией любого надсемейства (добавляем нулевые коэффициенты). \hfill \(\square\)
	
	\begin{tcolorbox}[colback=yellow!5, colframe=orange!50!black, title=Свойство 3']
		Каждое подсемейство линейно независимого семейства линейно независимо.
	\end{tcolorbox}
	
	\begin{tcolorbox}[colback=yellow!5, colframe=orange!50!black, title=Свойство 4]
		Семейство, содержащее нулевой вектор, линейно зависимо.
	\end{tcolorbox}
	
	\medskip
	\textit{Доказательство.}
	Следует из свойств 2 и 3, так как одноэлементное семейство из нулевого вектора линейно зависимо. \hfill \(\square\)
	
	\begin{tcolorbox}[colback=yellow!5, colframe=orange!50!black, title=Свойство 5]
		Семейство, содержащее пропорциональные (в частности, равные) векторы, линейно зависимо.
	\end{tcolorbox}
	
	\medskip
	\textit{Доказательство.}
	Пусть в семействе есть два вектора \( a_i \) и \( a_j \) такие, что \( a_j = \lambda a_i \) (пропорциональны). Тогда подсемейство \( (a_i, a_j) \) линейно зависимо, так как \( a_j - \lambda a_i = 0 \). По свойству 3 исходное семейство также линейно зависимо. \hfill \(\square\)
	
	\begin{tcolorbox}[colback=yellow!5, colframe=orange!50!black, title=Свойство 6]
		Если семейство \( a = (a_i)_{i \in I} \) линейно зависимо, а его подсемейство \( a' = (a_i)_{i \neq j} \) линейно независимо, то вектор \( a_j \) линейно выражается через \( a' \).
	\end{tcolorbox}
	
	\medskip
	\textit{Доказательство.}
	В силу линейной зависимости \( a \) существует нетривиальная комбинация \( \sum_{i \in I} \alpha_i a_i = 0 \). Если бы \( \alpha_j = 0 \), то получилась бы нетривиальная комбинация подсемейства \( a' \), равная нулю, что противоречит его линейной независимости. Следовательно, \( \alpha_j \neq 0 \), и тогда
	\[
	a_j = \sum_{i \neq j} \left( -\frac{\alpha_i}{\alpha_j} \right) a_i.
	\]
	\hfill \(\square\)
	
	
	
	
	
	
	
	
	\newpage	
	\section{Лемма о линейной зависимости линейных комбинаций (56 вопрос)}
	
	\begin{tcolorbox}[colback=blue!5, colframe=blue!40!black, title=Лемма (о линейной зависимости линейных комбинаций)]
		Пусть \( a = (a_i)_{i \in I} \) — семейство векторов в линейном пространстве \( L \), и \( b = (b_j)_{j \in J} \) — семейство линейных комбинаций семейства \( a \). Если мощность множества \( J \) больше мощности множества \( I \) (в частности, если \( I \) конечно и \( |J| > |I| \)), то семейство \( b \) линейно зависимо.
	\end{tcolorbox}
	
	\medskip
	\textit{Доказательство.}
	Рассмотрим два случая.
	
	\textit{Случай 1:} \( I \) конечно. Пусть \( |I| = n \). Без ограничения общности можно считать индексы натуральными: \( a = (a_1, \dots, a_n) \). Так как \( |J| > n \), выберем в \( J \) подмножество из \( n+1 \) элемента; соответствующее подсемейство \( b' \) семейства \( b \) также состоит из линейных комбинаций \( a \). Тогда существует матрица \( B \in K^{n \times (n+1)} \) такая, что \( b' = aB \). По лемме 17.3 (о существовании ненулевого решения однородной системы) существует ненулевой столбец \( X \in K^{(n+1) \times 1} \) такой, что \( BX = 0 \). Тогда
	\[
	b'X = (aB)X = a(BX) = a \cdot 0 = 0,
	\]
	и это нетривиальная линейная комбинация векторов \( b' \) (так как \( X \neq 0 \)). Следовательно, \( b' \) линейно зависимо, а тогда и \( b \) линейно зависимо (по свойству 3).
	
	\textit{Случай 2:} \( I \) бесконечно. Обозначим через \( \mathcal{J} \) множество всех конечных подмножеств множества \( I \). Известно, что \( |\mathcal{J}| = |I| \). Для каждого \( H \in \mathcal{J} \) определим \( J_H = \{ j \in J \mid b_j \text{ линейно выражается через } (a_i)_{i \in H} \} \). Поскольку любая линейная комбинация использует лишь конечное число векторов из \( a \), имеем \( \bigcup_{H \in \mathcal{J}} J_H = J \).
	
	Если бы для всех \( H \) выполнялось \( |J_H| \leq |H| \), то из теоремы о мощности объединения (см. 7.7) следовало бы \( |J| \leq |I| \), что противоречит условию. Следовательно, существует такое \( H \in \mathcal{J} \), что \( |J_H| > |H| \). По доказанному в первом случае, семейство \( (b_j)_{j \in J_H} \) линейно зависимо, а значит, и \( b \) линейно зависимо. \hfill \(\square\)
	
	\myremark
	Эта лемма является ключевым техническим результатом, используемым при доказательстве теоремы о размерности.
	
	
	\newpage	
	\section{Теорема о связи количества линейно независимых векторов и порождающих (57 вопрос)}
	
	\begin{tcolorbox}[colback=red!5, colframe=red!50!black, title=Теорема]
		Пусть в линейном пространстве \( L \) заданы два семейства: \( a = (a_i)_{i \in I} \) — линейно независимое, и \( b = (b_j)_{j \in J} \) — порождающее (т.е. каждый вектор из \( L \) линейно выражается через \( b \), в частности, все \( a_i \) выражаются через \( b \)). Тогда мощность множества \( I \) не превосходит мощности множества \( J \): \( |I| \leq |J| \).
	\end{tcolorbox}
	
	\medskip
	\textit{Доказательство.}
	Так как семейство \( b \) порождающее, каждый вектор \( a_i \) линейно выражается через \( b \). Рассмотрим семейство \( a \) как семейство линейных комбинаций семейства \( b \). Если бы \( |I| > |J| \), то по лемме о линейной зависимости линейных комбинаций семейство \( a \) было бы линейно зависимым, что противоречит условию. Следовательно, \( |I| \leq |J| \). \hfill \(\square\)
	
	\myremark
	Эта теорема является основой для определения размерности линейного пространства: все базисы (максимальные линейно независимые семейства) имеют одинаковую мощность, которая и называется размерностью.
	
	\subsection{Следствия}
	
	\begin{enumerate}
		\item В конечномерном пространстве любое линейно независимое семейство можно дополнить до базиса.
		\item Любое порождающее семейство содержит базис (как подсемейство).
		\item Размерность пространства однозначно определена.
	\end{enumerate}
	
	
	
	
	
	% ==================== РАЗДЕЛ 58: ПОНЯТИЕ БАЗИСА И РАЗМЕРНОСТИ ====================
	\newpage
	\section{Понятие базиса и размерности. Примеры. Единственность разложения по базису (58 вопрос)}
	
	\subsection{Порождающие семейства}
	
	Пусть \( L \) — линейное пространство над полем \( K \).
	
	\begin{tcolorbox}[colback=blue!5, colframe=blue!40!black, title=Определение]
		Семейство векторов \( a = (a_i)_{i \in I} \) в пространстве \( L \) называется \textit{порождающим} (или \textit{системой образующих}), если каждый вектор пространства \( L \) является линейной комбинацией этого семейства, т.е.
		\[
		L = \langle a \rangle.
		\]
	\end{tcolorbox}
	
	\begin{tcolorbox}[colback=green!5, colframe=green!50!black, title=Свойства порождающих семейств]
		\begin{enumerate}
			\item Всякое надсемейство порождающего семейства также порождает \( L \).
			\item Всякое нетривиальное надсемейство порождающего семейства линейно зависимо.
			\item Всякое семейство, длина которого больше длины порождающего семейства, линейно зависимо (следует из леммы 20.4).
			\item Если семейство \( (a_i)_{i \in I} \) порождает \( L \) и для некоторого подмножества \( J \subset I \) каждый вектор \( a_i \) при \( i \in I \setminus J \) линейно выражается через \( (a_j)_{j \in J} \), то \( (a_j)_{j \in J} \) также порождает \( L \).
		\end{enumerate}
	\end{tcolorbox}
	
	\subsection{Определение базиса}
	
	\begin{tcolorbox}[colback=red!5, colframe=red!50!black, title=Определение]
		Семейство векторов в линейном пространстве называется \textit{базисом} этого пространства, если оно:
		\begin{enumerate}
			\item линейно независимо;
			\item порождает это пространство.
		\end{enumerate}
	\end{tcolorbox}
	
	Ясно, что базис остаётся базисом при любой его переиндексации. Базис можно рассматривать как множество (со стандартной индексацией), поскольку векторы в базисе различны (из линейной независимости следует, что они попарно различны).
	
	\subsection{Единственность разложения по базису}
	
	Пусть \( e = (e_i)_{i \in I} \) — базис пространства \( L \), и \( x \in L \). Так как базис порождает \( L \), вектор \( x \) можно представить в виде линейной комбинации:
	\[
	x = \sum_{i \in I} \xi_i e_i = eX,
	\]
	где \( X = (\xi_i)_{i \in I} \) — столбец коэффициентов (лишь конечное число коэффициентов отлично от нуля). Такое представление называется \textit{разложением вектора \( x \) по базису \( e \)}.
	
	\begin{tcolorbox}[colback=yellow!5, colframe=orange!50!black, title=Теорема (Единственность разложения)]
		Разложение вектора по базису единственно. Иными словами, коэффициенты \( \xi_i \) однозначно определяются вектором \( x \) и базисом \( e \).
	\end{tcolorbox}
	
	\medskip
	\textit{Доказательство.}
	Предположим, что существуют два разложения:
	\[
	x = \sum_{i \in I} \xi_i e_i = \sum_{i \in I} \eta_i e_i.
	\]
	Тогда
	\[
	0 = x - x = \sum_{i \in I} (\xi_i - \eta_i) e_i.
	\]
	В силу линейной независимости базиса все коэффициенты \( \xi_i - \eta_i \) должны быть равны нулю, следовательно, \( \xi_i = \eta_i \) для всех \( i \in I \). \hfill \(\square\)
	
	\myremark
	Коэффициенты \( \xi_i \) называются \textit{координатами} вектора \( x \) в базисе \( e \). Столбец, составленный из этих координат, обозначается \( [x]_e \). Таким образом,
	\[
	x = e [x]_e.
	\]
	
	Базис \( e \) устанавливает взаимно однозначное соответствие между векторами пространства \( L \) и пространством финитных семейств \( K^{(I)} \). Если базис конечен и натурально индексирован: \( e = (e_1, \dots, e_n) \), то это соответствие превращается в изоморфизм между \( L \) и арифметическим пространством \( K^n \).
	
	\subsection{Примеры базисов}
	
	\begin{enumerate}
		\item \textbf{Нулевое пространство.} Единственным базисом нулевого пространства является пустое семейство.
		
		\item \textbf{Поле как одномерное пространство.} В поле \( K \), рассматриваемом как линейное пространство над самим собой, любой ненулевой элемент образует базис (одноэлементное семейство).
		
		\item \textbf{Стандартный базис арифметического пространства \( K^n \).} Векторы
		\[
		e_1 = (1,0,\dots,0),\; e_2 = (0,1,\dots,0),\; \dots,\; e_n = (0,0,\dots,1)
		\]
		образуют базис пространства \( K^n \). Действительно, любой вектор \( (\alpha_1,\dots,\alpha_n) \) представляется как \( \sum_{i=1}^n \alpha_i e_i \), и если такая сумма равна нулю, то все \( \alpha_i = 0 \).
		
		\item \textbf{Стандартный базис пространства формальных линейных комбинаций \( \langle I \rangle_K \).} Для каждого \( i \in I \) определим вектор \( e_i = (\delta_{ij})_{j \in I} \) (единица на \( i \)-м месте, нули на остальных). Тогда семейство \( (e_i)_{i \in I} \) является базисом пространства \( \langle I \rangle_K \). Это позволяет отождествить само множество \( I \) с базисными векторами.
		
		\item \textbf{Матричные единицы в пространстве матриц.} В пространстве \( K^{I \times J} \) (или \( K^{m \times n} \)) базис образуют матрицы \( e_{ij} \), у которых элемент с индексами \( (i,j) \) равен 1, а все остальные — 0. Эти матрицы называются \textit{матричными единицами}.
		
		\item \textbf{Геометрические векторы.} На плоскости любые два неколлинеарных вектора образуют базис. В трёхмерном пространстве любые три некомпланарных вектора образуют базис.
	\end{enumerate}
	
	\subsection{Размерность пространства}
	
	Из теоремы о связи количеств линейно независимых и порождающих векторов (вопрос 57) следует, что все базисы данного пространства имеют одинаковую мощность.
	
	\begin{tcolorbox}[colback=blue!5, colframe=blue!40!black, title=Определение]
		\textit{Размерностью} линейного пространства \( L \) называется мощность любого его базиса. Обозначение: \( \dim_K L \) или просто \( \dim L \).
	\end{tcolorbox}
	
	Если пространство имеет конечный базис, оно называется \textit{конечномерным}. В противном случае — \textit{бесконечномерным}.
	
	\myremark
	Размерность нулевого пространства равна 0 (базис пуст). Размерность пространства \( K^n \) равна \( n \). Размерность пространства матриц \( K^{m \times n} \) равна \( mn \).
	
	
	
	
	
	
	
	% ==================== РАЗДЕЛ 59: ТЕОРЕМА О СУЩЕСТВОВАНИИ БАЗИСА ====================
	\newpage
	\section{Теорема о существовании базиса (59 вопрос)}
	
	\begin{tcolorbox}[colback=red!5, colframe=red!50!black, title=Теорема (Существование базиса)]
		В любом линейном пространстве существует базис.
	\end{tcolorbox}
	
	\medskip
	\textit{Доказательство.}
	Достаточно заметить, что в любом линейном пространстве есть порождающее семейство (например, множество всех векторов пространства). По предложению 21.4 из любого порождающего семейства можно выделить базис. Следовательно, базис существует. \hfill \(\square\)
	
	Приведём два варианта доказательства предложения 21.4 (о выделении базиса из порождающего семейства).
	
	\begin{tcolorbox}[colback=green!5, colframe=green!50!black, title=Предложение 21.4]
		У каждого порождающего пространство семейства существует подсемейство, являющееся базисом этого пространства.
	\end{tcolorbox}
	
	\medskip
	\textit{Доказательство (вариант A — для конечного случая).}
	Пусть \( a = (a_1, \dots, a_n) \) — конечное порождающее семейство. Применим индукцию по \( n \).
	
	\textit{База:} \( n = 0 \) (пустое семейство). Тогда \( L = \{0\} \), и пустое семейство само является базисом (и подсемейством самого себя).
	
	\textit{Индукционный переход.} Если семейство \( a \) линейно независимо, то оно само и есть искомый базис. Если же оно линейно зависимо, то по свойству 20.3(1) один из его векторов (скажем, \( a_j \)) линейно выражается через остальные. Рассмотрим семейство \( a' = (a_i)_{i \neq j} \). По свойству 21.2(4) оно также порождает \( L \). Длина \( a' \) равна \( n-1 < n \), поэтому по предположению индукции у \( a' \) есть подсемейство, являющееся базисом. Это подсемейство является подсемейством и исходного семейства \( a \). \hfill \(\square\)
	
	\medskip
	\textit{Доказательство (вариант B — общий случай, с использованием леммы Цорна).}
	Пусть \( a = (a_i)_{i \in I} \) — порождающее семейство. Рассмотрим множество
	\[
	\mathcal{J} = \{ J \subset I \mid \text{семейство } (a_i)_{i \in J} \text{ линейно независимо} \}.
	\]
	Множество \( \mathcal{J} \) непусто, так как \( \emptyset \in \mathcal{J} \). Оно частично упорядочено по включению. Покажем, что \( \mathcal{J} \) индуктивно.
	
	Пусть \( \mathcal{M} \) — цепь в \( \mathcal{J} \). Положим \( H = \bigcup_{J \in \mathcal{M}} J \). Покажем, что семейство \( (a_i)_{i \in H} \) линейно независимо. Если бы существовала нетривиальная линейная комбинация
	\[
	\sum_{i \in H} \alpha_i a_i = 0,
	\]
	то множество носителей \( \{ i \in H \mid \alpha_i \neq 0 \} \) конечно и содержится в некотором \( J \in \mathcal{M} \) (так как \( \mathcal{M} \) — цепь, можно взять \( J \), содержащее все эти индексы). Тогда мы получили бы нетривиальную линейную комбинацию семейства \( (a_i)_{i \in J} \), равную нулю, что противоречит его линейной независимости. Следовательно, \( (a_i)_{i \in H} \) линейно независимо, и \( H \in \mathcal{J} \). Очевидно, \( H \) является верхней гранью цепи \( \mathcal{M} \). Таким образом, \( \mathcal{J} \) индуктивно.
	
	По лемме Цорна в \( \mathcal{J} \) существует максимальный элемент \( I^0 \). Пусть \( a' = (a_i)_{i \in I^0} \). По построению \( a' \) линейно независимо. Покажем, что \( a' \) порождает \( L \). Для любого \( j \in I \setminus I^0 \) семейство \( (a_i)_{i \in I^0 \cup \{j\}} \) линейно зависимо (иначе \( I^0 \) не было бы максимальным). Тогда по свойству 20.3(6) вектор \( a_j \) линейно выражается через \( a' \). Таким образом, все векторы исходного порождающего семейства выражаются через \( a' \), а значит, \( a' \) порождает \( L \). Следовательно, \( a' \) — базис. \hfill \(\square\)
	
	
	
	
	
	
	
	% ==================== РАЗДЕЛ 60: ТРИ ОПРЕДЕЛЕНИЯ БАЗИСА ====================
	\newpage
	\section{Три определения базиса (60 вопрос)}
	
	Базис линейного пространства можно охарактеризовать тремя эквивалентными способами.
	
	\begin{tcolorbox}[colback=blue!5, colframe=blue!40!black, title=Определение 1 (Максимальное линейно независимое семейство)]
		Линейно независимое семейство векторов пространства \( L \) называется \textit{максимальным линейно независимым}, если любое его нетривиальное надсемейство в \( L \) линейно зависимо.
	\end{tcolorbox}
	
	\begin{tcolorbox}[colback=blue!5, colframe=blue!40!black, title=Определение 2 (Минимальное порождающее семейство)]
		Порождающее семейство пространства \( L \) называется \textit{минимальным порождающим}, если ни одно из его нетривиальных подсемейств не порождает \( L \).
	\end{tcolorbox}
	
	\begin{tcolorbox}[colback=blue!5, colframe=blue!40!black, title=Определение 3 (Базис в обычном смысле)]
		Семейство векторов называется \textit{базисом}, если оно линейно независимо и порождает \( L \).
	\end{tcolorbox}
	
	\begin{tcolorbox}[colback=red!5, colframe=red!50!black, title=Предложение 21.9]
		Для семейства \( a \) в пространстве \( L \) следующие утверждения равносильны:
		\begin{enumerate}
			\item \( a \) — максимальное линейно независимое семейство;
			\item \( a \) — минимальное порождающее семейство;
			\item \( a \) — базис.
		\end{enumerate}
	\end{tcolorbox}
	
	\medskip
	\textit{Доказательство.}
	
	\textbf{1 \(\Rightarrow\) 3.} Пусть \( a \) максимально линейно независимо. По определению оно линейно независимо. Докажем, что оно порождает \( L \). Возьмём произвольный вектор \( x \in L \). Рассмотрим семейство \( a' \), полученное добавлением \( x \) к \( a \). Это надсемейство, нетривиальное, поэтому по максимальности оно линейно зависимо. Тогда по свойству 20.3(6) вектор \( x \) линейно выражается через \( a \). Ввиду произвольности \( x \) это означает, что \( a \) порождает \( L \).
	
	\textbf{3 \(\Rightarrow\) 1.} Если \( a \) — базис, то оно линейно независимо. Его максимальность следует из свойства 21.2(2): любое нетривиальное надсемейство порождающего семейства линейно зависимо.
	
	\textbf{2 \(\Rightarrow\) 3.} Пусть \( a \) — минимальное порождающее семейство. По предложению 21.4 у него существует подсемейство, являющееся базисом. В силу минимальности это подсемейство не может быть собственным, следовательно, совпадает с \( a \). Таким образом, \( a \) само является базисом.
	
	\textbf{3 \(\Rightarrow\) 2.} Если \( a \) — базис, то оно порождает \( L \). Его минимальность следует из свойства 21.2(2): любое нетривиальное подсемейство порождающего семейства, будучи надсемейством чего-то? На самом деле надо рассуждать иначе: если бы существовало собственное подсемейство, порождающее \( L \), то оно было бы порождающим семейством меньшей длины, что противоречило бы теореме о связи количеств (вопрос 57), применённой к базису и этому подсемейству. Более прямо: пусть \( a' \) — собственное подсемейство \( a \). Так как \( a \) линейно независимо, в \( a \setminus a' \) есть вектор, не выражающийся через \( a' \) (иначе \( a \) было бы линейно зависимо). Следовательно, \( a' \) не порождает \( L \). Значит, \( a \) минимально. \hfill \(\square\)
	
	\myremark
	Таким образом, базис можно определить как максимальное линейно независимое семейство или как минимальное порождающее семейство. Эти определения часто используются в различных контекстах
	
	
	
	
	
	% ==================== РАЗДЕЛ 61: СВОЙСТВА ПРОСТРАНСТВ ДАННОЙ РАЗМЕРНОСТИ ====================
	\newpage
	\section{Свойства пространств данной размерности. Размерность подпространства (61 вопрос)}
	
	\subsection{Корректность определения размерности}
	
	\begin{tcolorbox}[colback=blue!5, colframe=blue!40!black, title=Предложение 22.1]
		Все базисы линейного пространства \( L \) имеют одну и ту же длину.
	\end{tcolorbox}
	
	\medskip
	\textit{Доказательство.}
	Пусть \( e \) и \( e' \) — два базиса пространства \( L \). Так как \( e \) порождает \( L \), а \( e' \) линейно независимо, то по теореме о связи количеств (вопрос 57) длина \( e' \) не превосходит длины \( e \). Меняя роли базисов, получаем, что длина \( e \) не превосходит длины \( e' \). Следовательно, длины равны. \hfill \(\square\)
	
	\begin{tcolorbox}[colback=blue!5, colframe=blue!40!black, title=Определение 22.2]
		\textit{Размерностью} линейного пространства \( L \) называется длина его базиса. Обозначение: \( \dim L \).
	\end{tcolorbox}
	
	Корректность определения следует из предложения 22.1. Пространства, имеющие конечный базис, называются \textit{конечномерными}; их размерность — конечное число. Пространства, не имеющие конечного базиса, называются \textit{бесконечномерными}; их размерность — бесконечная мощность.
	
	\subsection{Основные свойства конечномерных пространств}
	
	\begin{tcolorbox}[colback=green!5, colframe=green!50!black, title=Предложение 22.3]
		Любое семейство векторов, длина которого больше размерности пространства, линейно зависимо.
	\end{tcolorbox}
	
	\medskip
	\textit{Доказательство.}
	Пусть \( \dim L = n \) и \( a \) — семейство векторов длины \( > n \). Возьмём любой базис \( e \) пространства \( L \). Так как базис порождает \( L \), каждый вектор из \( a \) линейно выражается через \( e \). По лемме о линейной зависимости линейных комбинаций (вопрос 56) семейство \( a \) линейно зависимо. \hfill \(\square\)
	
	\begin{tcolorbox}[colback=green!5, colframe=green!50!black, title=Предложение 22.4]
		Пусть \( L \) конечномерно и \( a \) — семейство векторов в \( L \), длина которого равна \( \dim L \). Тогда следующие условия равносильны:
		\begin{enumerate}
			\item \( a \) линейно независимо;
			\item \( a \) порождает \( L \);
			\item \( a \) является базисом.
		\end{enumerate}
	\end{tcolorbox}
	
	\medskip
	\textit{Доказательство.}
	\begin{itemize}
		\item[1 \(\Rightarrow\) 2.] Семейство \( a \) конечно (так как его длина конечна). Любое его нетривиальное надсемейство имеет длину \( > \dim L \) и по предложению 22.3 линейно зависимо. Таким образом, \( a \) — максимальное линейно независимое семейство, а значит, по предложению 21.9 оно порождает \( L \).
		\item[2 \(\Rightarrow\) 1.] Поскольку \( a \) порождает \( L \), у него есть подсемейство \( a' \), являющееся базисом (предложение 21.4). Длина \( a' \) равна \( \dim L \), а так как \( a' \subset a \) и оба семейства конечны и имеют одинаковую длину, то \( a' = a \). Следовательно, \( a \) линейно независимо.
		\item Эквивалентность с пунктом 3 очевидна из определения базиса. \hfill \(\square\)
	\end{itemize}
	
	\subsection{Размерность подпространства}
	
	\begin{tcolorbox}[colback=red!5, colframe=red!50!black, title=Теорема]
		Пусть \( L \) — конечномерное линейное пространство, и \( M \subset L \) — линейное подпространство. Тогда \( M \) также конечномерно, и \( \dim M \leq \dim L \). При этом \( \dim M = \dim L \) тогда и только тогда, когда \( M = L \).
	\end{tcolorbox}
	
	\medskip
	\textit{Доказательство.}
	Рассмотрим множество всех линейно независимых семейств векторов в \( M \). Их длины ограничены сверху числом \( \dim L \) (по предложению 22.3, применённому к \( L \)). Возьмём максимальное линейно независимое семейство \( e \) в \( M \) (оно существует, например, можно выбрать семейство максимальной длины, которая конечна). Это семейство будет базисом \( M \): оно линейно независимо по построению, а его максимальность в \( M \) означает, что любой вектор из \( M \) линейно выражается через \( e \). Следовательно, \( M \) конечномерно и \( \dim M = |e| \leq \dim L \).
	
	Если \( \dim M = \dim L \), то базис \( M \) является линейно независимым семейством в \( L \) длины \( \dim L \), а значит, по предложению 22.4 он порождает \( L \). Таким образом, \( M = L \). Обратное очевидно. \hfill \(\square\)
	
	\myremark
	Для бесконечномерных пространств утверждение также верно: любое подпространство имеет размерность, не превосходящую размерности всего пространства (в смысле сравнения мощностей), но доказательство требует аксиомы выбора.
	
	
	
	
	
	
	
	% ==================== РАЗДЕЛ 62: РАЗМЕРНОСТИ СУММЫ И ПЕРЕСЕЧЕНИЯ ====================
	\newpage
	\section{Размерности суммы и пересечения подпространств (62 вопрос)}
	
	\subsection{Сумма и пересечение подпространств}
	
	Напомним, что для подпространств \( M \) и \( N \) линейного пространства \( L \) определены:
	\begin{itemize}
		\item \textit{Пересечение} \( M \cap N \) — множество векторов, принадлежащих обоим подпространствам. Это также подпространство.
		\item \textit{Сумма} \( M + N = \{ x + y \mid x \in M, y \in N \} \) — также подпространство.
	\end{itemize}
	
	\begin{tcolorbox}[colback=blue!5, colframe=blue!40!black, title=Теорема (Формула Грассмана)]
		Для конечномерных подпространств \( M \) и \( N \) линейного пространства \( L \) выполняется равенство:
		\[
		\dim(M + N) + \dim(M \cap N) = \dim M + \dim N.
		\]
	\end{tcolorbox}
	
	\medskip
	\textit{Доказательство.}
	Пусть \( \dim(M \cap N) = k \). Выберем базис \( e_1, \dots, e_k \) в \( M \cap N \). Дополним его до базиса \( M \): векторами \( f_1, \dots, f_m \) (так что \( \dim M = k + m \)). Аналогично, дополним базис пересечения до базиса \( N \): векторами \( g_1, \dots, g_n \) (так что \( \dim N = k + n \)).
	
	Покажем, что семейство
	\[
	e_1, \dots, e_k, f_1, \dots, f_m, g_1, \dots, g_n
	\]
	является базисом суммы \( M + N \).
	
	\begin{itemize}
		\item \textit{Порождающее свойство.} Любой вектор из \( M + N \) есть сумма вектора из \( M \) и вектора из \( N \). Вектор из \( M \) выражается через \( e_i \) и \( f_j \), вектор из \( N \) — через \( e_i \) и \( g_l \). Следовательно, их сумма выражается через объединённое семейство.
		\item \textit{Линейная независимость.} Предположим, что
		\[
		\sum_{i=1}^k \alpha_i e_i + \sum_{j=1}^m \beta_j f_j + \sum_{l=1}^n \gamma_l g_l = 0.
		\]
		Перенесём часть с \( g_l \) в правую часть:
		\[
		\sum_{i=1}^k \alpha_i e_i + \sum_{j=1}^m \beta_j f_j = -\sum_{l=1}^n \gamma_l g_l.
		\]
		Левая часть принадлежит \( M \), правая — \( N \). Следовательно, обе части принадлежат \( M \cap N \). Поэтому левая часть выражается через базис \( e_i \) пересечения:
		\[
		\sum_{i=1}^k \alpha_i e_i + \sum_{j=1}^m \beta_j f_j = \sum_{i=1}^k \delta_i e_i.
		\]
		Тогда
		\[
		\sum_{i=1}^k (\alpha_i - \delta_i) e_i + \sum_{j=1}^m \beta_j f_j = 0.
		\]
		Векторы \( e_i, f_j \) линейно независимы (они образуют базис \( M \)), поэтому все коэффициенты \( \beta_j = 0 \) и \( \alpha_i = \delta_i \). Подставляя \( \beta_j = 0 \) в исходное равенство, получаем
		\[
		\sum_{i=1}^k \alpha_i e_i + \sum_{l=1}^n \gamma_l g_l = 0.
		\]
		Аналогично, рассматривая это как комбинацию базиса \( N \), заключаем, что все \( \alpha_i = 0 \) и \( \gamma_l = 0 \). Таким образом, все коэффициенты равны нулю.
	\end{itemize}
	
	Итак, объединённое семейство — базис суммы. Число векторов в нём равно \( k + m + n = (k + m) + (k + n) - k = \dim M + \dim N - \dim(M \cap N) \). Следовательно,
	\[
	\dim(M + N) = \dim M + \dim N - \dim(M \cap N).
	\]
	\hfill \(\square\)
	
	\subsection{Следствия}
	
	\begin{enumerate}
		\item Если \( M \cap N = \{0\} \), то \( \dim(M + N) = \dim M + \dim N \). В этом случае сумма называется \textit{прямой} (обозначение \( M \oplus N \)).
		\item Для любого подпространства \( M \subset L \) существует такое подпространство \( N \), что \( L = M \oplus N \) (такое \( N \) называется \textit{дополнением} к \( M \)).
	\end{enumerate}
	
	
	
	
	% ==================== РАЗДЕЛ 63: ПРЯМАЯ СУММА ПОДПРОСТРАНСТВ ====================
	\newpage
	\section{Прямая сумма подпространств: эквивалентные определения (63 вопрос)}
	
	\subsection{Определение прямой суммы}
	
	Пусть \( L \) — линейное пространство, и \( M_1, \dots, M_k \) — его подпространства.
	
	\begin{tcolorbox}[colback=blue!5, colframe=blue!40!black, title=Определение]
		Говорят, что \( L \) является \textit{прямой суммой} подпространств \( M_1, \dots, M_k \), и пишут
		\[
		L = M_1 \oplus M_2 \oplus \dots \oplus M_k,
		\]
		если каждый вектор \( x \in L \) единственным образом представляется в виде
		\[
		x = x_1 + x_2 + \dots + x_k, \quad x_i \in M_i.
		\]
	\end{tcolorbox}
	
	\subsection{Эквивалентные условия}
	
	\begin{tcolorbox}[colback=green!5, colframe=green!50!black, title=Теорема]
		Для подпространств \( M_1, \dots, M_k \) конечномерного пространства \( L \) следующие условия равносильны:
		\begin{enumerate}
			\item \( L = M_1 \oplus \dots \oplus M_k \) (т.е. каждый вектор представляется единственным образом).
			\item \( L = M_1 + \dots + M_k \) и \( M_i \cap \left( \sum_{j \neq i} M_j \right) = \{0\} \) для каждого \( i \).
			\item Объединение базисов подпространств \( M_i \) даёт базис всего пространства \( L \).
			\item \( \dim L = \sum_{i=1}^k \dim M_i \) и \( L = M_1 + \dots + M_k \).
		\end{enumerate}
	\end{tcolorbox}
	
	\medskip
	\textit{Доказательство.}
	\begin{itemize}
		\item[1 \(\Rightarrow\) 2.] Из единственности разложения следует, что сумма всех подпространств равна \( L \). Если бы существовал ненулевой вектор \( x \in M_i \cap \sum_{j \neq i} M_j \), то он имел бы два различных представления: как элемент \( M_i \) (с нулевыми остальными слагаемыми) и как сумма векторов из остальных подпространств. Это противоречит единственности.
		
		\item[2 \(\Rightarrow\) 3.] Выберем в каждом \( M_i \) базис \( e_{i1}, \dots, e_{i, \dim M_i} \). Покажем, что объединённое семейство линейно независимо. Пусть
		\[
		\sum_{i,j} \alpha_{ij} e_{ij} = 0.
		\]
		Сгруппируем слагаемые по \( i \): для каждого \( i \) положим \( x_i = \sum_j \alpha_{ij} e_{ij} \in M_i \). Тогда \( x_1 + \dots + x_k = 0 \). Для любого \( i \) имеем
		\[
		x_i = -\sum_{j \neq i} x_j \in M_i \cap \left( \sum_{j \neq i} M_j \right) = \{0\},
		\]
		следовательно, \( x_i = 0 \). Так как \( x_i \) — линейная комбинация базисных векторов \( M_i \), все коэффициенты \( \alpha_{ij} \) равны нулю. Объединённое семейство линейно независимо. Его длина равна сумме размерностей, и оно порождает сумму всех подпространств, которая по условию равна \( L \). Таким образом, это базис.
		
		\item[3 \(\Rightarrow\) 4.] Очевидно, что объединение базисов порождает сумму подпространств, а так как оно базис, то эта сумма равна \( L \). Размерность равна количеству векторов в объединённом базисе, т.е. сумме размерностей.
		
		\item[4 \(\Rightarrow\) 1.] Из условия \( L = \sum M_i \) следует, что любой вектор представим в виде суммы. Пусть вектор имеет два представления:
		\[
		x_1 + \dots + x_k = y_1 + \dots + y_k, \quad x_i, y_i \in M_i.
		\]
		Тогда \( (x_1 - y_1) + \dots + (x_k - y_k) = 0 \). Выберем в каждом \( M_i \) базис; пусть общее число базисных векторов равно сумме размерностей. Расписывая каждую разность \( x_i - y_i \) по базису \( M_i \), получим линейную комбинацию объединённого семейства, равную нулю. Поскольку объединённое семейство имеет длину, равную размерности \( L \), оно линейно независимо (по предложению 22.4, так как порождает \( L \)). Следовательно, все коэффициенты равны нулю, и \( x_i = y_i \) для всех \( i \). Единственность доказана. \hfill \(\square\)
	\end{itemize}
	
	% ==================== РАЗДЕЛ 64: БАЗИС ПРЯМОЙ СУММЫ ====================
	\newpage
	\section{Базис прямой суммы (64 вопрос)}
	
	\subsection{Построение базиса прямой суммы}
	
	Из предыдущего раздела непосредственно вытекает описание базиса прямой суммы.
	
	\begin{tcolorbox}[colback=red!5, colframe=red!50!black, title=Теорема]
		Пусть \( L = M_1 \oplus M_2 \oplus \dots \oplus M_k \) — прямая сумма подпространств. Если в каждом \( M_i \) выбран базис \( e_{i1}, e_{i2}, \dots, e_{i, \dim M_i} \), то объединение этих базисов
		\[
		e_{11}, \dots, e_{1, \dim M_1}, e_{21}, \dots, e_{2, \dim M_2}, \dots, e_{k1}, \dots, e_{k, \dim M_k}
		\]
		является базисом всего пространства \( L \).
	\end{tcolorbox}
	
	\medskip
	\textit{Доказательство.}
	Это было доказано при установлении эквивалентности условий в предыдущей теореме (пункт 2 \(\Rightarrow\) 3). Действительно, из определения прямой суммы следует условие 2, которое влечёт линейную независимость объединённого семейства, а порождающее свойство очевидно, так как сумма всех \( M_i \) равна \( L \). \hfill \(\square\)
	
	\subsection{Обратное утверждение}
	
	Если объединение базисов подпространств \( M_i \) даёт базис \( L \), то сумма этих подпространств является прямой. Это также было доказано (условие 3 \(\Rightarrow\) 1).
	
	\subsection{Примеры}
	
	\begin{enumerate}
		\item В пространстве \( \mathbb{R}^3 \) плоскость \( M = \{ (x,y,0) \} \) и прямая \( N = \{ (0,0,z) \} \) находятся в прямой сумме: \( \mathbb{R}^3 = M \oplus N \). Базисом служит объединение базиса плоскости (например, \( (1,0,0), (0,1,0) \)) и базиса прямой (\( (0,0,1) \)).
		
		\item В пространстве многочленов степени не выше \( n \) можно представить прямую сумму подпространства чётных и нечётных многочленов.
		
		\item В пространстве матриц \( K^{n \times n} \) можно рассматривать прямую сумму подпространств симметрических и кососимметрических матриц (при \( \operatorname{char} K \neq 2 \)).
	\end{enumerate}
	
	\subsection{Размерность прямой суммы}
	
	Из теоремы о базисе непосредственно следует, что
	\[
	\dim (M_1 \oplus \dots \oplus M_k) = \dim M_1 + \dots + \dim M_k.
	\]
	Это согласуется с формулой Грассмана для случая попарно пересекающихся только по нулю подпространств
	
	
	
	
	% ==================== РАЗДЕЛ 65: МАТРИЦЫ ПЕРЕХОДА МЕЖДУ БАЗИСАМИ ====================
	\newpage
	\section{Матрицы перехода между базисами (65 вопрос)}
	
	\subsection{Необходимость преобразования координат}
	
	Соответствие между векторами линейного пространства и их координатами зависит от выбора базиса. Часто приходится сначала выбрать некоторый базис, провести вычисления, а затем перейти к другому, более удобному базису. Для этого нужен способ пересчитывать координаты вектора при замене базиса — так называемое \textit{преобразование координат}. Оказывается, что такое преобразование удобно описывается некоторой матрицей.
	
	\subsection{Определение матрицы перехода}
	
	Пусть \( L \) — линейное пространство над полем \( K \). Рассмотрим два базиса пространства \( L \): старый базис \( e = (e_i)_{i \in I} \) и новый базис \( e' = (e'_j)_{j \ \in J} \) (для простоты будем считать индексацию натуральной, если не оговорено иное).
	
	Так как \( e \) — базис, он порождает \( L \), и каждый вектор нового базиса \( e'_j \) может быть разложен по старому базису:
	\[
	e'_j = \sum_{i \in I} c_{ij} e_i, \quad j \in J.
	\]
	
	\begin{tcolorbox}[colback=blue!5, colframe=blue!40!black, title=Определение]
		\textit{Матрицей перехода} от старого базиса \( e \) к новому базису \( e' \) называется матрица \( C = (c_{ij}) \) размера \( I \times J \), столбцами которой являются координатные столбцы векторов нового базиса в старом базисе:
		\[
		e' = e C.
		\]
	\end{tcolorbox}
	
	Если пространство конечномерно и \( \dim L = n \), а базисы натурально индексированы: \( e = (e_1, \dots, e_n) \), \( e' = (e'_1, \dots, e'_n) \), то матрица перехода \( C \) является квадратной порядка \( n \):
	\[
	C = \begin{pmatrix}
		c_{11} & c_{12} & \dots & c_{1n} \\
		c_{21} & c_{22} & \dots & c_{2n} \\
		\vdots & \vdots & \ddots & \vdots \\
		c_{n1} & c_{n2} & \dots & c_{nn}
	\end{pmatrix},
	\quad \text{где } e'_j = \sum_{i=1}^n c_{ij} e_i.
	\]
	
	\subsection{Свойства матрицы перехода}
	
	\begin{tcolorbox}[colback=green!5, colframe=green!50!black, title=Предложение]
		Матрица перехода от одного базиса к другому обратима. Обратной к ней является матрица перехода от нового базиса к старому.
	\end{tcolorbox}
	
	\medskip
	\textit{Доказательство.}
	Пусть \( C \) — матрица перехода от \( e \) к \( e' \), т.е. \( e' = eC \). Пусть \( C' \) — матрица перехода от \( e' \) обратно к \( e \), т.е. \( e = e' C' \). Тогда
	\[
	e = e' C' = (eC) C' = e (C C').
	\]
	В силу линейной независимости строки \( e \) (как семейства векторов) получаем \( C C' = E \). Аналогично, подставляя \( e' = e C \) в \( e' = e' C' C \), получаем \( C' C = E \). Таким образом, \( C' = C^{-1} \). \hfill \(\square\)
	
	\myremark
	В случае конечномерного пространства обратимость матрицы перехода означает, что её определитель отличен от нуля (матрица неособенная).
	
	
	
	
	
	
	% ==================== РАЗДЕЛ 66: ОБРАТИМАЯ МАТРИЦА КАК МАТРИЦА ПЕРЕХОДА ====================
	\newpage
	\section{Обратимая матрица как матрица перехода между базисами (66 вопрос)}
	
	\subsection{Критерий того, что матрица является матрицей перехода}
	
	Не любая квадратная матрица может служить матрицей перехода между базисами данного пространства. Однако над полем любая обратимая матрица может быть реализована как матрица перехода.
	
	\begin{tcolorbox}[colback=red!5, colframe=red!50!black, title=Теорема]
		Пусть \( L \) — конечномерное линейное пространство над полем \( K \), \( \dim L = n \), и \( e = (e_1, \dots, e_n) \) — некоторый базис \( L \). Тогда для любой обратимой матрицы \( C \in K^{n \times n} \) существует единственный базис \( e' = (e'_1, \dots, e'_n) \) такой, что \( C \) является матрицей перехода от \( e \) к \( e' \).
	\end{tcolorbox}
	
	\medskip
	\textit{Доказательство.}
	Определим векторы \( e'_j \) формулой
	\[
	e'_j = \sum_{i=1}^n c_{ij} e_i, \quad j = 1, \dots, n,
	\]
	т.е. положим \( e' = e C \). Покажем, что \( e' \) — базис.
	
	Так как \( C \) обратима, существует обратная матрица \( C^{-1} \). Тогда \( e = e' C^{-1} \), следовательно, каждый вектор старого базиса (а значит, и любой вектор пространства) линейно выражается через \( e' \). Таким образом, \( e' \) порождает \( L \). Осталось проверить линейную независимость. Пусть
	\[
	\sum_{j=1}^n \alpha_j e'_j = 0.
	\]
	В матричной записи это означает \( e' \alpha = 0 \), где \( \alpha = (\alpha_1, \dots, \alpha_n)^T \). Но \( e' = e C \), поэтому \( e (C \alpha) = 0 \). В силу линейной независимости \( e \) получаем \( C \alpha = 0 \). Умножая слева на \( C^{-1} \), находим \( \alpha = 0 \). Следовательно, \( e' \) линейно независимо и является базисом. \hfill \(\square\)
	
	\myremark
	Таким образом, множество всех обратимых матриц порядка \( n \) находится во взаимно однозначном соответствии с множеством всех базисов пространства \( L \) при фиксированном исходном базисе. Это соответствие задаётся формулой \( e' = e C \).
	
	\subsection{Следствие}
	
	Любые два базиса конечномерного пространства связаны обратимой матрицей перехода. Обратно, любая обратимая матрица задаёт переход от данного базиса к некоторому новому базису.
	
	
	
	
	
	% ==================== РАЗДЕЛ 67: ИЗМЕНЕНИЕ КООРДИНАТ ПРИ ЗАМЕНЕ БАЗИСА ====================
	\newpage
	\section{Изменение координат при замене базиса (67 вопрос)}
	
	\subsection{Формула преобразования координат}
	
	Пусть в пространстве \( L \) заданы два базиса: старый \( e = (e_1, \dots, e_n) \) и новый \( e' = (e'_1, \dots, e'_n) \), связанные соотношением \( e' = e C \), где \( C \) — матрица перехода от \( e \) к \( e' \).
	
	Для произвольного вектора \( x \in L \) рассмотрим его координатные столбцы в этих базисах:
	\[
	[x]_e = \begin{pmatrix} \xi_1 \\ \vdots \\ \xi_n \end{pmatrix}, \quad [x]_{e'} = \begin{pmatrix} \xi'_1 \\ \vdots \\ \xi'_n \end{pmatrix},
	\]
	так что \( x = e [x]_e = e' [x]_{e'} \).
	
	\begin{tcolorbox}[colback=blue!5, colframe=blue!40!black, title=Теорема (Формула преобразования координат)]
		Координаты вектора в старом и новом базисах связаны соотношением:
		\[
		[x]_e = C [x]_{e'}.
		\]
		Или, что то же самое,
		\[
		[x]_{e'} = C^{-1} [x]_e.
		\]
	\end{tcolorbox}
	
	\medskip
	\textit{Доказательство.}
	Подставим выражение нового базиса через старый в разложение вектора:
	\[
	x = e' [x]_{e'} = (e C) [x]_{e'} = e (C [x]_{e'}).
	\]
	С другой стороны, \( x = e [x]_e \). В силу единственности разложения по базису \( e \) получаем \( [x]_e = C [x]_{e'} \). \hfill \(\square\)
	
	\myremark
	Обратите внимание: матрица перехода \( C \) связывает \textit{старые координаты с новыми} по формуле \( [x]_e = C [x]_{e'} \). Это означает, что если мы хотим выразить новые координаты через старые, нужно использовать обратную матрицу: \( [x]_{e'} = C^{-1} [x]_e \). Принятое определение матрицы перехода (новый базис через старый) именно таково, что она даёт выражение старых координат через новые. Это удобно, так как часто приходится преобразовывать выражения, содержащие координаты, заменяя в них старые координаты их выражением через новые.
	
	\subsection{Пример}
	
	В пространстве \( \mathbb{R}^2 \) задан стандартный базис \( e_1 = (1,0), e_2 = (0,1) \). Новый базис: \( e'_1 = (1,1), e'_2 = (1,-1) \). Найдём матрицу перехода от \( e \) к \( e' \).
	
	Координаты \( e'_1 \) в старом базисе: \( (1,1)^T \), координаты \( e'_2 \): \( (1,-1)^T \). Следовательно,
	\[
	C = \begin{pmatrix} 1 & 1 \\ 1 & -1 \end{pmatrix}.
	\]
	
	Пусть вектор \( x \) имеет в старом базисе координаты \( [x]_e = (3,2)^T \). Найдём его координаты в новом базисе:
	\[
	[x]_{e'} = C^{-1} [x]_e = \frac{1}{-2} \begin{pmatrix} -1 & -1 \\ -1 & 1 \end{pmatrix} \begin{pmatrix} 3 \\ 2 \end{pmatrix} = -\frac{1}{2} \begin{pmatrix} -3-2 \\ -3+2 \end{pmatrix} = -\frac{1}{2} \begin{pmatrix} -5 \\ -1 \end{pmatrix} = \begin{pmatrix} 2.5 \\ 0.5 \end{pmatrix}.
	\]
	
	Проверим: \( x = 3e_1 + 2e_2 = (3,2) \). В новом базисе: \( 2.5 e'_1 + 0.5 e'_2 = 2.5(1,1) + 0.5(1,-1) = (2.5+0.5, 2.5-0.5) = (3,2) \).
	
	\subsection{Геометрический смысл}
	
	Формула преобразования координат показывает, что при замене базиса координаты вектора преобразуются с помощью матрицы, обратной к матрице перехода. Это согласуется с тем, что векторы базиса преобразуются с помощью матрицы \( C \), а координаты — контравариантно (противоположным образом).
	
	\subsection{Общий случай бесконечномерных пространств}
	
	В бесконечномерном пространстве матрица перехода также существует, но она может быть бесконечной. При этом она по-прежнему обратима в том смысле, что существует матрица обратного перехода. Однако понятие обратимости для бесконечных матриц требует более тонкого рассмотрения (матрица должна иметь конечное число ненулевых элементов в каждом столбце и строке, чтобы умножение было определено).
	
	
	
	
	
\end{document}
